\documentclass[11pt]{article}
\usepackage[T1]{fontenc}
\usepackage[utf8]{inputenc}
\usepackage{lmodern,amsmath,amssymb,amsthm,mathtools}
\usepackage[a4paper,margin=29mm]{geometry}
\usepackage{microtype,enumitem,booktabs}
\usepackage[colorlinks=true,allcolors=blue,breaklinks=true]{hyperref}
\setlength{\emergencystretch}{2em}
\newtheorem{theorem}{Theorem}[section]
\newtheorem{proposition}[theorem]{Proposition}
\newtheorem{lemma}[theorem]{Lemma}
\newtheorem{corollary}[theorem]{Corollary}
\theoremstyle{remark}\newtheorem{remark}[theorem]{Remark}
\newcommand{\OO}{\mathcal O}
\newcommand{\F}{\mathbf F}
\newcommand{\Z}{\mathbf Z}
\newcommand{\Pic}{\operatorname{Pic}}
\newcommand{\Spec}{\operatorname{Spec}}
\newcommand{\Sym}{\operatorname{Sym}}
\newcommand{\PP}{\mathbf P}
\newcommand{\workbenchurl}{https://github.com/KokunoYumeto/schemes-and-moduli-workbench}
\newcommand{\recordurl}{https://github.com/KokunoYumeto/mathematics-workbenches}
\title{Picard schemes, very ample surfaces\\and metabelian local groups\\[0.3em]\large After Grothendieck, Raynaud and Suzuki}
\author{}
\date{}
\begin{document}
\maketitle
\vspace{-2em}
\begin{abstract}
Grothendieck's Picard and local fundamental group questions ask for canonical divisor classes and finite descriptions of covers near a singularity. We prove that a universal polarization class exists exactly when at most one elementary divisor is odd. Using Raynaud's descent criterion, we prove that changing the abelian structure group preserves schematicity of a torsor, and obtain schematicity of its full Picard functor. For abelian-surface torsors, a projective scheme modification with flat, geometrically integral fibres also forces the torsor to be a scheme. Over any base, a schematic projective bundle forces its abelian torsor base to be a scheme, in every relative dimension. After Raynaud, we construct a smooth projective surface in characteristic two with a very ample line bundle whose inverse has one-dimensional first cohomology; its space of sections has dimension 562 and its irregularity is 565. For geometric Picard kernels, we prove a strict genus bound for a nontrivial order-$p$ class on a smooth very ample curve in every positive characteristic, together with a big adjoint divisor and a canonical map with surface image. At exact conductor equality, we determine the canonical class, exclude conductors one through four and conductor six, and show that the normalized torsor has ample canonical bundle and at most one characteristic-two $A_1$ point. For algebraically trivial classes we give an explicit positive lower bound for the excess in the conductor inequality, even on a normal projective surface when the curve lies in its smooth locus and has positive self-intersection. The same bound holds when the N\'eron--Severi class has the specified finite power divisibility. The first-valuation bound is strict in odd characteristic. In characteristic two, equality forces nontrivial two-torsion on the curve; at logarithmic order two its actual two-torsion rank is at least two. The canonical discrepancy line has exact order two at logarithmic order two and is trivial at every larger odd half-order. At every odd half-order the suitable resolved cover admits a fibration over the projective line. We also determine further necessary restrictions in the smallest elliptic and inseparable Albanese cases. A torsion class pulled back from a curve satisfies an additional bound involving the genus of that curve. In the single-singularity case, the class cannot be a double. After Suzuki's study of nilpotent local fundamental groups, we prove a uniform cohomology bound for tame coefficients, an explicit finite bound on abelian inertia, and finite generation of the maximal metabelian quotient. The conductor argument uses Kato's published integrality theorem. The mathematics and text were produced by GPT-6 Astra (ChatGPT Pro mode). The separate checks and their scope are described at the end; no human has checked the paper.
\end{abstract}

\section{The questions and the results}

The first question is about choosing a divisor class in a family of polarized abelian varieties. A polarization determines a connected component of the relative Picard scheme, but does not in general specify a point of that component. Grothendieck studies canonical trivializations of such torsors in his Montpellier manuscript, cote 59, sheet 13, where he writes ``torseur trivial canonique'' in the comparison of polarization torsors \cite{Groth59}. We determine the exact obstruction at the prime two.

Let $D=(d_1\mid\cdots\mid d_g)$ be a polarization type. On the moduli stack of abelian schemes with polarization of type $D$, write $\Pic^\delta$ for the component of the relative fppf Picard sheaf corresponding to the polarization. It is a torsor under the dual abelian scheme. A universal section means one compatible with all base changes and all isomorphisms of polarized abelian schemes.

\begin{theorem}\label{thm:polarization}
Over $\Z[1/2]$, the torsor $\Pic^\delta$ has a universal section if and only if at most one of the $d_i$ is odd. More intrinsically, the construction exists when the alternating pairing on $A[2]$ induced by the polarization has rank at most two, and is canonical on that locus. No odd prime dividing the polarization degree needs to be inverted. If at least two elementary divisors are odd, no universal section exists even over $\mathbf C$.
\end{theorem}

We formulate the section as a relative Picard class; the zero section of an abelian scheme gives its normalized rigidified line-bundle representative. For an elliptic curve the distinguished degree-one class is that of the origin. For type $(1,2)$ there is again a distinguished class; for a principally polarized abelian surface there is none. The distinction depends on the rank of the pairing on two-torsion, not just on the parity of the polarization degree.

The proof uses the classical correspondence between symmetric line bundles and quadratic refinements, in the scheme-theoretic form of Poonen and Rains \cite[Proposition 3.6 and diagram (3.7), pp.~1312--1313]{PR}. The canonical construction in rank at most two is already the proof of their Proposition~3.12(f), pp.~1314--1315; it works in families. For principal type over $\mathbf C$ in dimension at least two, the obstruction also follows from Fringuelli--Pirisi \cite[Theorem~C, author version p.~3]{FP}. We combine these known cases with the integral-transvection argument for arbitrary type. No priority claim is made for this synthesis.

Grothendieck also asks whether the full Picard functor is representable when the family is ``propre plat et à fibres géom. normales, voire seulement intègres'' \cite[cote~57, sheet~314; PDF p.~315]{Groth57}. Fppf triviality of a torsor alone does not answer this question, because descent of schemes can fail. Raynaud's criterion gives the following positive case.

\begin{theorem}\label{thm:schematic-abelian-pushout}
Let $S$ be any scheme, let $A$ and $B$ be smooth proper group schemes of finite presentation over $S$ with geometrically connected fibres, and let $h:A\to B$ be a homomorphism. If an fppf $A$-torsor $T$ is represented by a scheme, then its pushout $h_*T$ is represented by a scheme. No condition on the kernel of $h$ is imposed.
\end{theorem}

\begin{corollary}\label{cor:picard-schematic-abelian-torsor}
For every scheme $S$, every abelian $S$-scheme $A$ in this sense, and every fppf $A$-torsor $T$ represented by a scheme, the full fppf Picard functor $\Pic_{T/S}$ is a separated $S$-scheme locally of finite presentation.
\end{corollary}

The Picard conclusion over a connected locally noetherian base is stated by Achter and Casalaina-Martin \cite[\S1.7.2, first paragraph, p.~13]{ACMPrym}. We give a proof over an arbitrary base that also handles nonprojective abelian schemes and changing N\'eron--Severi groups. Raynaud's affine-neighborhood and descent results carry the proof \cite[Corollary~3.8 and Theorem~4.2]{RaynaudDescent}; the connected-group action theorem in \cite[\S6.6, Theorem~2]{BLR} allows a finite cover on which a projective multisection can be constructed. The theorem then handles the potentially degenerate homomorphism associated to a Picard class. This settles the abelian-torsor case of schematicity for proper families with geometrically normal fibres. The question for general such families remains open here.

There is also an obstruction to constructing such a family by modifying a nonschematic torsor. Theorem~\ref{thm:surface-torsor-modification} proves the following. Let $S$ be any scheme, $A/S$ an abelian scheme of relative dimension two, and $Q$ an fppf $A$-torsor algebraic space. If a scheme $X$ is proper, flat and finitely presented over $S$ with geometrically integral fibres, and a projective $S$-morphism $X\to Q$ is an isomorphism over a dense open of $Q$, then $Q$ is a scheme. The same conclusion holds if the morphism is birational on every geometric fibre. Thus a projective modification of a nonschematic abelian-surface torsor cannot preserve all these family hypotheses. The proof turns a rank-one torsion-free pushforward into a finite flat subscheme that splits the torsor; Raynaud's criterion then descends schematicity. It does not settle the modification question for higher-dimensional torsors.

Projective bundles give another construction that can be excluded over every base. Let $S$ be any scheme, $A/S$ a smooth proper group scheme of finite presentation with geometrically connected fibres, $Q$ an fppf $A$-torsor algebraic space, and $E$ a finite locally free sheaf on $Q$ whose rank is everywhere positive. Theorem~\ref{thm:torsor-projective-bundle} states that $\PP_Q(E)$ is a scheme if and only if $Q$ is a scheme. The base need not be reduced or noetherian, and the bundle need not be homogeneous. Over a finite normalization of a local base, the Ramanujam--Samuel theorem turns affine boundaries into Cartier divisors. The projective bundle formula makes their positivity incompatible with a common positive-dimensional kernel of the line-bundle homomorphisms of the torsor. The resulting finite map to projective Picard-class torsors supplies the affine neighborhoods required by Raynaud's descent criterion. Thus projective bundles over nonschematic abelian torsors cannot give counterexamples to Picard schematicity. Brion classifies homogeneous projective bundles over abelian varieties over an algebraically closed field \cite[pp.~2475--2479]{BrionProjective}; the relative assertion here imposes no homogeneity hypothesis.

The second question concerns Kodaira vanishing and its relation to restriction on Picard schemes. Raynaud constructed smooth surfaces in positive characteristic with an ample line bundle $Q$ for which $H^1(Q^{-1})\ne0$ \cite[pp.~273--278]{Raynaud}. He explicitly asks whether the line bundle can be globally generated or very ample \cite[\S3, Remark~3, p.~278]{Raynaud}. Lauritzen and Rao constructed very ample counterexamples in dimension at least six; their opening discussion distinguishes that question from the surface case \cite[pp.~1--2 of the author manuscript]{LR}. Zheng's later constructions refine the Raynaud method and study further failures of positivity \cite{Zheng}. Here we require very ampleness of the line bundle producing the negative first cohomology. We give the following surface construction.

\begin{theorem}\label{thm:surface}
There exist an algebraically closed field $k$ of characteristic two, a smooth geometrically integral projective surface $Y/k$, and a very ample line bundle $L$ on $Y$ such that
\[
 h^0(Y,L)=562,\qquad h^1(Y,L^{-1})=1,\qquad h^1(Y,\OO_Y)=565.
\]
For any smooth connected hyperplane section $H\in|L|$, the tangent space at the identity of the kernel of $\Pic_{Y/k}\to\Pic_{H/k}$ has dimension one.
\end{theorem}

The assertion about tangent spaces follows from the theorem on cohomology; it does not assert the existence of a nontrivial geometric point in that kernel. This distinction matters when the characteristic is positive and Picard schemes need not be reduced.

A general-basis form of this surface construction over every algebraically
closed field of characteristic two is given in
\cite[Theorem~1 and Sections~3--6, pp.~2, 4--8]{HexagonSurface}.
The construction below uses the explicitly specified ground field in
Theorem~\ref{thm:surface}.

Kweon separates these two questions in \cite[Question~7.7, arXiv v2 p.~19]{KweonNS}: whether restriction of $\Pic^\tau$ to a smooth hyperplane section is injective, and whether it is a closed immersion. His Theorem~7.2, pp.~17--18, makes the kernel finite with connected Cartier dual; this still permits nontrivial geometric $p$-torsion. Theorem~\ref{thm:surface} gives a negative answer to the closed-immersion question, already for a smooth surface and a very ample curve, because the tangent map has a nonzero kernel. It leaves the geometric-point question open here.

The construction begins with a concrete Artin--Schreier curve and the degree-three cyclic surface used in the Raynaud construction. The decisive extra step is to take ninth roots of every member of a generic basis of a very ample power of $Q$. Their roots become sections of the pullback of $Q$. Together with the original sections they generate the coordinate rings of affine charts, which proves very ampleness. At the same time the cover algebra is a direct sum of negative powers of $Q$, so the cohomology can be read off exactly. Taking roots of a basis is a standard cyclic-cover operation: Bhatt uses it for global generation in the proof of \cite[Proposition~7.2, p.~1771]{Bhatt}. The affine-chart argument here proves very ampleness while retaining the original sections; no historical priority is claimed for the operation.

The surface theorem gives a very ample failure of Kodaira vanishing in the smallest possible dimension: on a smooth projective curve an anti-ample line bundle has no global sections. It does not classify such surfaces, minimize the degree of the cover, or give an example over a specified finite field. Its ground field will be an algebraic closure of an explicitly specified purely transcendental extension of $\F_2$.

A recent draft by ChatGPT and Claude in collaboration with Ajit and Lee gives positive-characteristic counterexamples concerning splinters and positivity: its Theorem~A constructs globally $+$-regular rational surfaces whose anticanonical divisor is nef but neither big nor semiample \cite[Theorem~A, p.~2]{AjitLee}. This is a different failure of positivity from the very ample negative cohomology here. The draft is contextual and is not an input to our construction. Proposition~\ref{prop:tangent-ample} below also excludes a smooth finite flat torsor approach to passing from an infinitesimal Picard kernel to a nontrivial point of that kernel.

The geometric-point question gives a separate genus obstruction. A nontrivial torsion line bundle can disappear on a curve only if the curve is sufficiently large relative to its self-intersection. Very ampleness rules out the first equality permitted by the conductor calculation in every positive characteristic.

\begin{theorem}\label{thm:strict-picard-all}
Let $X$ be a smooth projective integral surface over an algebraically
closed field $k$ of characteristic $p>0$, and let $Y\subset X$ be a
smooth integral very ample divisor. If $M\ne\OO_X$ is a line bundle
of order $p$ with $M|_Y\simeq\OO_Y$, then
\[
 2g(Y)-2>pY^2.
\]
Moreover $|K_X-(p-2)Y|$ contains a linear subsystem with
two-dimensional image. In particular $K_X-(p-2)Y$ is big,
$X$ is of general type, and its canonical map has two-dimensional image.
\end{theorem}

The exact conductor gives a stronger structural conclusion.
The logarithmic form of $M$ is obtained from local units defining a
trivialization of $M^p$; its exact order on $Y$ is a multiple of $p$.
The logarithmic construction is classical; see the proof of
Lang's Lemma~3 \cite[pp.~198--199]{LangPTorsion}.
We prove the needed normality and vanishing-order statements below.
Our notation distinguishes three related integers:
\[
 \begin{array}{c|l}
 m=pa&\text{logarithmic order}\\
 a&\text{conductor integer}\\
 a(p-1)&\text{conductor-ideal coefficient}
 \end{array}
\]
The last coefficient is measured along the inverse-image Cartier curve
in the normalized torsor.

\begin{theorem}\label{thm:conductor-equality-structure}
Let $X$ be a smooth projective integral surface over an algebraically
closed field $k$ of characteristic $p>0$, let $Y\subset X$ be a
smooth integral very ample divisor, and let $M\ne\OO_X$ have
order $p$ with $M|_Y\simeq\OO_Y$.
Let $pa$ be the exact vanishing order along $Y$ of its logarithmic
form, and let $\pi:N\to X$ be the normalization of its
$\mu_p$-torsor. Suppose that $2g(Y)-2=paY^2$.
Then
\[
 a\geq5,\quad a\ne6,\qquad K_X\sim(pa-1)Y,\qquad
 K_N\sim(a-1)\pi^*Y,
\]
and $K_N$ is ample. In odd characteristic $N$ is smooth.
In characteristic two it is smooth or has exactly one $A_1$
singularity, with completed local equation $uv=w^2$.
Write $\epsilon=0$ in the smooth case and $\epsilon=1$ otherwise,
and put $d=Y^2$. Then $p\mid(d-\epsilon)$.
The minimal resolution $T$ of $N$ is a smooth minimal surface of
general type, with
\[
 K_T^2=p(a-1)^2d,\qquad c_2(T)=-pa\,d+2\epsilon,
 \qquad 12\chi(\OO_T)=p(a^2-3a+1)d+2\epsilon.
\]
\end{theorem}

Positive self-intersection alone gives a quantitative restriction when
the torsion class has enough tensor roots. Here $\operatorname{NS}(X)$
is the group of line bundles modulo algebraic equivalence.

\begin{theorem}\label{thm:positive-excess-picard}
Let $X$ be a normal integral projective surface over an algebraically
closed field $k$ of characteristic $p>0$. Let $Y\subset X$ be a smooth
integral effective Cartier divisor in its smooth locus, with $d=Y^2>0$.
Let $M\ne\OO_X$ have order $p$ and $M|_Y\simeq\OO_Y$, and let $m$
be the exact order along $Y$ of its logarithmic form. Then $m>0$,
$p\mid m$, and
\[
 G=2g(Y)-2,\qquad \delta=G-md\geq0.
\]
Put $r=v_p(m)\geq1$ and
\[
 Q_r(m)=\prod_{j=1}^{r}\left(1+2(p-1)\frac{m}{p^j}\right).
\]
If $[M]\in p^r\operatorname{NS}(X)$, then
\begin{equation}\label{eq:positive-excess-main-bound}
 (Q_r(m)-1)\delta\geq md.
\end{equation}
In particular this bound holds whenever $M$ is algebraically trivial.
For $r=1$ it reads
\[
 \delta\geq\frac{p}{2(p-1)}d.
\]
If $r=1$ and $p$ is odd, this last inequality is strict.
\end{theorem}

Thus an algebraically trivial class requires a definite gap above the
logarithmic genus bound. In characteristic two with $v_2(m)=1$, the
gap is at least $Y^2$. More generally, a smaller gap forces the class
outside $p^r\operatorname{NS}(X)$. No ampleness assumption is needed.
The equality obstruction is stated as
Corollary~\ref{thm:positive-square-picard};
Theorem~\ref{thm:odd-positive-excess-strict} proves strictness
of the first-valuation excess bound in odd characteristic.

At the first characteristic-two boundary, the curve must itself
carry nontrivial two-torsion. More precisely, the discrepancy
between its canonical line and the indicated power of its normal
line has the following restricted form.

\begin{corollary}\label{cor:positive-excess-characteristic-two-boundary}
Under the hypotheses of Theorem~\ref{thm:positive-excess-picard},
suppose $p=2$, $m=2a$ with $a$ odd, $[M]\in2\operatorname{NS}(X)$,
and $\delta=d$. Put
\[
 N=\OO_Y(Y),\qquad
 \mathcal T=\omega_Y\otimes N^{-(2a+1)}.
\]
If $a=1$, the line $\mathcal T$ has exact order two and
$\dim_{\F_2}\Pic(Y)[2](k)\geq2$.
If $a\geq3$, the line $\mathcal T$ is trivial.
In both cases $\Pic(Y)[2]\ne0$, where this denotes actual
two-torsion line-bundle classes over $k$.
Under Theorem~\ref{thm:positive-excess-picard} with $p=2$,
$m=2a$ for odd $a$, and $[M]\in2\operatorname{NS}(X)$,
the condition $\Pic(Y)[2]=0$ therefore implies $\delta>d$.
When $m=2$, the same strict inequality holds whenever
$\dim_{\F_2}\Pic(Y)[2](k)\leq1$.
Since $\delta$ is even, this gives
$\delta\geq d+2$ for even $d$ and $\delta\geq d+1$ for odd $d$.
\end{corollary}

At $a=1$, Proposition~\ref{prop:smallest-endpoint-rank} also
identifies the zero divisor of the logarithmic form of $\mathcal T$:
it is three times the tangential leading divisor. The logarithmic form
of every restricted square root of $M$ is nonzero on that divisor.
For every odd $a$, Proposition~\ref{prop:boundary-resolved-fibration}
gives a fibration over $\PP^1$ on a suitable resolution of the normalized
cover, with the smooth inverse curve as the support of a fibre of
power-of-two multiplicity.
These are necessary conditions on an equality configuration.
They do not construct one or assert that the remaining alternatives
occur. In particular the general Picard point-kernel question
remains open.

There is also a quantitative restriction in every characteristic $p>0$
when a class with logarithmic order $pa$, $p\nmid a$, comes from a curve.
Under the surface, curve and torsion hypotheses above, suppose
$f:X\to B$ is a nonconstant morphism to a smooth projective curve
of genus $b\geq1$, and $M=f^*P$ for a nontrivial order-$p$ line
bundle $P$ on $B$. Then $Y$ dominates $B$ and, writing
$n=\deg(f|_Y)$, Proposition~\ref{prop:boundary-curve-pullback-excess}
gives
\[
 \delta\geq\frac{pd}{2(p-1)}+
                  \frac{n(b-1)}{(p-1)a}.
\]
No separate N\'eron--Severi divisibility hypothesis is needed here.
Thus a base of genus at least two excludes the first-valuation endpoint.
In characteristic two the bound is $\delta\geq d+n(b-1)/a$.
In characteristic two, for an algebraically trivial class on a smooth surface whose
Albanese image is a curve, that curve must be elliptic at the endpoint.
Still in characteristic two, when $a=1$ and $Y$ is ample, the Albanese map cannot be
generically separable onto a surface image; the precise alternatives
are in Corollary~\ref{cor:boundary-smallest-albanese-gate}.
In the elliptic case, the restricted curve map has exactly one
inseparable step of degree two, followed by a separable map of degree
at least two. The first normalized base change has smooth generic
fibre. In the surface-image case, the Albanese differential has rank
exactly one, and the restricted curve maps separably to its normalized
image. The common global conductor and cotangent conclusions below
require only smoothness of $X$, ampleness of $Y$, and the stated square
root condition.

At the same smallest endpoint, very ampleness also gives
\[
 d\ge12,\qquad
 h^0(X,\OO_X(Y))\le\left\lfloor d/4\right\rfloor+2
\]
by Proposition~\ref{prop:smallest-endpoint-very-ample-degree}.
The degree-two Albanese case has a stronger relation between the
surface and the singularities of its image.  The following statement
requires only ampleness of the distinguished curve.

\begin{theorem}[The smallest endpoint in Albanese degree two]
\label{thm:smallest-endpoint-albanese-summary}
Let $X$ be a smooth integral projective surface over an algebraically
closed field of characteristic two, and let $Y\subset X$ be a smooth
integral ample divisor.  Put $d=Y^2$.
Let $M\ne\OO_X$ be algebraically trivial of exact order two,
with $M|_Y\simeq\OO_Y$.  Suppose that its logarithmic form has
exact order two along $Y$ and that $2g(Y)-2=3d$.
Assume that $A=\operatorname{Alb}(X)$ is a surface and that
$\alpha:X\to A$ has generic degree two.
Then
\[
 d=4\chi(\OO_X),\qquad 2d<K_X^2<3d,\qquad
 1\le c_2(X)=3d-K_X^2\le d-1.
\]
In the Stein factorization $X\xrightarrow{r}X_0\xrightarrow{a_0}A$,
the scheme $X_0$ is a normal Gorenstein finite flat double cover of $A$.
Set
\[
 t=\operatorname{length}(R^1r_*\OO_X),\qquad
 \mathcal L=(a_{0*}\OO_{X_0}/\OO_A)^{-1}.
\]
Let $D_A=\alpha(Y)$ with its reduced integral structure.
Its normalization is $Y$.
Write $M=\alpha^*P$, and let $D'$ be the image of the lift of
$Y\to D_A\subset A$ to the degree-two $\mu_2$-isogeny defined by $P$.
For $\epsilon=p_a(D')-g(Y)\ge0$, one has
\[
 D_A\equiv2\mathcal L,\qquad
 2t=d+\epsilon,\qquad
 \mathcal L^2=\frac{3d}{2}+\epsilon.
\]
In particular $\epsilon$ is even, and
\[
 t\ge\max\left\{\frac d2,
       \left\lceil\frac{d^2}{4c_2(X)}\right\rceil\right\}.
\]
If $Y$ is very ample, then $4\mid d$, $d\ge16$, and
\[
 6\le h^0(X,\OO_X(Y))\le d/4+2.
\]
\end{theorem}

Proposition~\ref{prop:smallest-endpoint-albanese-degree-two},
Corollary~\ref{cor:smallest-endpoint-albanese-very-ample-degree},
and Theorem~\ref{thm:smallest-endpoint-strict-canonical}
prove these assertions.
The latter also determines the exceptional contribution to
$\alpha^*D_A$ and the equality condition in $4c_2(X)t\ge d^2$.
The finite Stein factor must have a nonrational singularity.
Neither the normalization defect nor the resolution defect is
bounded above by this result, and no endpoint surface is constructed.

In the single-singularity case of
Theorem~\ref{thm:conductor-equality-structure}, there is a stronger
obstruction. Proposition~\ref{prop:internal-cartier-descent-obstruction}
shows that $M$ cannot descend as an invertible sheaf to the radial
quotient and that $[M]\notin2\operatorname{NS}(X)$, regardless of
the $2$-adic valuation of its logarithmic order. Its proof treats the
exceptional curve on the smooth quotient resolution explicitly.

Under Theorem~\ref{thm:strict-picard-all}, the genus bound in
characteristic two says $g(Y)\geq Y^2+2$.
The adjoint conclusion restricts the ambient surface as well as the
curve. Proposition~\ref{prop:residual-adjoint} gives the exact fixed
divisor and a nef residual adjoint, and bounds the rank of exterior
differentiation on global one-forms. These are necessary conditions
for a nontrivial geometric point of the Picard kernel. They do not
construct such a point and do not follow from the tangent-space
calculation in Theorem~\ref{thm:surface}.

The general-type consequence alone follows from known vanishing
results. Ekedahl proves vanishing for numerically positive line
bundles on surfaces not of general type, except possibly for
quasi-elliptic surfaces of Kodaira dimension one
\cite[II, Theorem~1.6, p.~125]{EkedahlCanonical}.
In that exceptional case, the theorem of Langer proved by Zhang
\cite[Theorem~3.1, author pp.~4--5]{ZhangQuasiElliptic}
gives vanishing for a nef and big bundle of fibre degree greater
than two in characteristic two, or greater than one in characteristic
three. Apply these results to $\OO_X(Y)\otimes M^{-1}$:
the restriction sequence gives nonzero first cohomology of its
inverse, whereas its numerical class is that of $Y$.
On an integral genus-one fibre, a very ample line bundle has degree
at least three, so the quasi-elliptic exception is excluded as well.
The result here specifies a big adjoint and its canonical subsystem;
it makes no priority claim for the general-type consequence.

Langer's vanishing theorem supplies a neighbouring slope obstruction. Writing $\mu_{\max,Y}(\Omega_X^1)$ for the largest slope of a nonzero subsheaf, measured by intersection with $Y$, \cite[Theorem~9.1, pp.~2105--2106]{LangerS} gives
\[
 pY^2\leq\mu_{\max,Y}(\Omega_X^1)
\]
under the hypotheses of Theorem~\ref{thm:strict-picard-all}. Indeed, $M$ is numerically flat and the restriction sequence injects $H^0(Y,\OO_Y)$ into $H^1(X,M(-Y))$, since a nontrivial torsion line bundle has no nonzero section. The vanishing criterion therefore cannot hold. Langer's restriction theorem also gives faithful flatness of the map on $S$-fundamental groups when $D\cdot H>\mu_{\max,H}(\Omega_X^1)$ for an ample $H$ \cite[Theorem~10.2 and proof, pp.~2109--2111]{LangerS}. A nonzero logarithmic form of order $m$ along $Y$ gives an injection $\OO_X(mY)\hookrightarrow\Omega_X^1$, hence $\mu_{\max,H}(\Omega_X^1)\geq mY\cdot H$. For $D=Y$ this criterion is therefore unavailable. The strict conductor theorem uses the roots of the torsion class and its normalized covers.

Cotner and Zavyalov prove a Lefschetz theorem for flat cohomology under a Hodge-equivalence hypothesis \cite[Corollary~3.1.2]{CotnerZavyalov}. For a smooth surface and a smooth curve whose inclusion is a Hodge $1$-equivalence, their result makes $H^1_{\mathrm{fppf}}(X,\mu_p)\to H^1_{\mathrm{fppf}}(Y,\mu_p)$ injective, hence also $\Pic(X)[p]\to\Pic(Y)[p]$. Their strongly ample condition includes this extra cohomological hypothesis \cite[Definition~2.3.5]{CotnerZavyalov}. Very ampleness alone does not imply it: the tangent kernel in Theorem~\ref{thm:surface} already violates its coherent-cohomology requirement. A nonzero logarithmic form vanishing on the curve also violates the required injection on global one-forms \cite[Definition~2.2.1]{CotnerZavyalov}. Their Picard-scheme isomorphism for a strongly ample divisor assumes ambient dimension at least three \cite[Theorem~3.4.4]{CotnerZavyalov}.

Biswas and Holla prove surjectivity of Nori's fundamental group for smooth divisors in sufficiently high powers of an ample bundle, and an isomorphism in ambient dimension at least three \cite[Theorem~1.1, author p.~2]{BiswasHolla}. Their discussion of ample divisors also gives a failure of faithful flatness using an $F$-trivial rank-two extension \cite[Section~2, especially Remark~2.4, author pp.~3--7]{BiswasHolla}. That extension does not itself exhibit an order-$p$ line bundle trivial on the divisor.

The proof begins with the closed logarithmic differential of $M$.
Its exact vanishing order along $Y$ is a positive multiple of $p$.
Multiplying by ratios of hyperplane sections and differentiating gives
a system of canonical forms. After its fixed divisor is removed, it
restricts rationally to the original embedding of $Y$.
Even when base points lie on $Y$, this system cannot be composed with
a pencil: the pencil base would be $Y$, every exceptional rational
curve would be vertical, and the generic fibre would carry both a
globally generated line bundle of degree one and a canonical section
of negative degree. The residual adjoint therefore has positive
square. A second proof of bigness interprets wedges with coordinate
differentials as a family of tangent lines; a one-dimensional image
would cover the surface by lines on which the original cotangent
section cannot exist.

At conductor equality, the positive square makes a general
determinant divisor connected. Its intersection with $Y$ makes it
reduced, so its only regular functions are constants. This splits
the Euler extension and gives $K_X\sim(pa-1)Y$.
The compatible splitting supplies a fixed ambient point contained
in every tangent plane. Such an embedding is called strange.
Furukawa identifies the relation between this point and the Euler
sequence; his converse from a twisted tangent section assumes
$H^1(X,\OO_X(-Y))=0$
\cite[Propositions~2.5 and~2.10, author pp.~4--6]{FurukawaStrange}.
The determinant argument here supplies the lift under conductor
equality without that vanishing assumption.

Fujino and Sato prove that a smooth projective toric variety with
an ample locally free subsheaf of its tangent bundle is projective
space, in arbitrary characteristic
\cite[Theorem~1.1, p.~1, and proof in \S5, pp.~8--10]{FujinoSatoToric}.
Their proof uses the toric Euler sequence and cohomological lifting.
The internal-center argument here treats a nef and big line bundle
on a resolution whose exceptional curve is contracted; no toric
structure is assumed. Fujino and Sato credit Rethlas and ChatGPT
with assistance in developing the proof ideas and assign independent
formalization and final verification to the authors
\cite[Acknowledgments, p.~3]{FujinoSatoToric}.

The resulting multiplicative quotient is the Frobenius twist of
the distinguished normalized torsor. Seshadri's regular-foliation
quotient theorem supplies smoothness away from the possible radial
point \cite[Proposition~6 and Theorem~2, pp.~6-11--6-16]{SeshadriCartier}.
Closedness excludes that point in odd characteristic; in
characteristic two its invariant ring gives one $A_1$ singularity.
The classification of surfaces with numerically trivial canonical
divisor excludes conductor one, and the surface inequality of
Gu--Sun--Zhou excludes conductors two and three. Its genus-two
clause, combined with integral canonical divisibility, excludes
conductor four. The same fiber-degree argument excludes conductor
six in Proposition~\ref{prop:conductor-six-excluded}. Thus the strict
genus bound is a consequence of the structure theorem.
Furukawa's classification of smooth strange complete intersections
also shows that the embedding cannot be a complete intersection.

The broader ample setting still needs separate arguments.
The global degree filtration gives the cover's Euler characteristic
without very ampleness. Ekedahl's foliation theorem and the ruled-cover and canonical-defect
calculations give conditional results for the merely ample cases
that the Euler splitting does not cover. The separate obstruction
using Kadets's sectional-monodromy theorem requires very ampleness. In odd characteristic,
Proposition~\ref{prop:odd-higher-conductor-equality} refines the
structure theorem by canonical divisibility and the Gu--Sun--Zhou
and Tate bounds. At most three Albanese fiber genera survive.
At $a=5$ only $p=3,7$ remain, while for $p\geq17$ the
necessary conductor is at least $(p-1)/2$.
No remaining equality case is constructed here.

For Theorem~\ref{thm:positive-excess-picard}, normalization of the
torsor converts genus excess into the sum of a conductor correction
and twice the genus defect of its inverse-image curve. Resolving that
curve loses at most four times its genus defect in self-intersection.
A $p$th root of the bundle then gives a new logarithmic form of order
$m/p$. Its curve still has positive square under the contrary numerical
inequality. The excess grows by at most the factor
$1+2(p-1)m/p$. Iterating through the specified finite chain of roots
would give a positive logarithmic order prime to $p$, contrary to
closedness. This proves the product bound while allowing singular
intermediate surfaces and singular inverse-image curves.

Nagamachi and Takamatsu relate conductors and genus changes under
purely inseparable extensions of the constant field
\cite[Theorem~0.2 and discussion, pp.~2--3]{NagamachiTakamatsu}.
Their ramification-index-one formula is neighbouring local background.
Here the cover comes from a torsion line bundle on a surface;
the argument combines its conductor with resolution loss and tensor
roots to bound restriction to a positive-square curve.
Section~\ref{subsec:positive-square-picard} gives the complete argument.

The classical logarithmic form attached to a torsion line bundle
is closed and Cartier invariant; see Esnault's formula~(2.4),
author p.~3, in \cite{EsnaultPartialConnection}.
The endpoint argument uses the vanishing order and conormal
coefficient of this form. In odd characteristic, equality forces
cusps $y^2=x^p$ and a zero in a degree-zero conormal section,
which is impossible.

At the characteristic-two endpoint, the resolved curve has square
zero. Equality in the genus and resolution estimates forces the
normalized cover to be smooth along a cuspidal inverse-image curve.
Resolving its cusps identifies the normal line of the square-zero
curve and forces the new form to have conormal order zero.
A leading two-form and, at the smallest order, a conormal
coefficient that does not descend then give
Corollary~\ref{cor:positive-excess-characteristic-two-boundary}.
Castryck, Streng and Testa obtain nontrivial rational two-torsion
for a nonempty open family of curves with specified Newton polygons
\cite[Theorem~2, author p.~2; introduction, pp.~1--3]{CastryckStrengTesta}.
Their construction compares characteristic-two and geometric
theta characteristics. Here two-torsion is forced by a surface
restriction class at an extremal genus equality, and the particular
line $\mathcal T$ is constrained as well.
Section~\ref{subsec:positive-excess-boundary} proves this endpoint
criterion. At logarithmic order two, the conormal trivialization
also gives the exact logarithmic divisor of $\mathcal T$, distinguishing
it from every restricted square root of $M$.
The same actual normal trivialization holds for all odd half-orders.
Totaro's moving-divisor lemma \cite[Lemma~4.1, pp.~7--8]{TotaroMovingDivisors}
then makes a power-of-two multiple of the resolved curve move.
The last cusp exceptional curve forces the base of the resulting
fibration to be $\PP^1$.
Section~\ref{subsec:conductor-picard-descent}
identifies the exact thickening depth and the obstruction to
lifting the radial action for general torsion classes.
Section~\ref{subsec:internal-cartier-descent} treats invertible
descent at the internal characteristic-two center. An exact
exceptional multiplicity makes a determinant divisor reduced on
the quotient resolution. Its connectedness then splits the Euler
extension, whose radial lift contradicts the saturated tangent
direction on the exceptional curve.

The third question is Grothendieck's finite-generation conjecture for the local fundamental group. In its surface case, let $A$ be a complete noetherian normal local domain of dimension two and characteristic $p>0$, with algebraically closed residue field $k$, and put $U=\Spec A\setminus\{\mathfrak m\}$. Fix a proper SNC resolution $X\to\Spec A$, an isomorphism over $U$, with reduced exceptional divisor $D=\bigcup_iC_i$. Such resolutions exist for excellent surfaces; if $A$ is regular, one may first blow up its closed point. A locally constant $\F_\ell$-sheaf here has finite-dimensional fibres and finite continuous monodromy. It is \emph{tame along $D$} when its monodromy has prime-to-$p$ ramification index and separable residue extension at every generic point of $D$.

\begin{theorem}\label{thm:local-coefficients}
For the fixed resolution above there is a finite constant $C_X$ such that
\[
 \dim_{\F_p}H^1(U,M)\leq C_X\operatorname{rank}M
\]
for every finite locally constant $\F_p$-sheaf $M$ tame along $D$. The same constant works for every rank, every finite monodromy group, and every tame inertia order. The monodromy group may be nonabelian and its order may be divisible by $p$.
\end{theorem}

\begin{theorem}\label{thm:abelian-inertia}
Let $B=(C_i\cdot C_j)$, let $g_i$ be the genus of $C_i$, and let $\nu_i$ be its number of crossing points. Put
\[
 k_i=\max(0,2g_i-2+\nu_i),\qquad t=(-B)^{-1}k.
\]
Every character $\chi\in H^1(U,\mathbf Q_p/\mathbf Z_p)$ has logarithmic Swan conductor $s_i(\chi)\leq t_i$ at $C_i$. For any integer $e\geq0$ with $p^e>\max_i t_i$, the kernel
\[
 I=\ker\bigl(\pi_1(U)^{\mathrm{ab}}\to\pi_1(X)^{\mathrm{ab}}\bigr)
\]
is finite and is annihilated by $p^e|\det B|$.
\end{theorem}

\begin{theorem}\label{thm:metabelian}
For every complete noetherian normal local domain $A$ of dimension two and characteristic $p>0$ with algebraically closed residue field, the maximal metabelian quotient of $\pi_1(\Spec A\setminus\{\mathfrak m\})$ is topologically finitely generated. Equivalently, there is one finite bound, depending on $A$, for the numbers of generators of all its finite metabelian quotients.
\end{theorem}

A group is metabelian when its commutator subgroup is abelian. The final theorem controls all covers with such a monodromy group, uniformly, rather than one chosen finite cover. It is a part of Grothendieck's conjecture on the entire local fundamental group, recorded in SGA~2, Expos\'e XIII, Conjecture 3.1; Suzuki states this conjecture precisely and proves finite generation of the maximal pro-nilpotent quotient \cite[Conjecture 1.1 and Theorem 1.2, author pp.~1--2]{Suzuki}. The prime-to-characteristic theory of Grothendieck and Murre supplies another part of that setting. The metabelian assertion addresses a different quotient; it makes no assertion about arbitrary solvable derived length or the entire group.

Suzuki's equal-characteristic proof uses Artin--Schreier nearby cycles and the negative intersection matrix of a resolution. We extend that calculation to varying tame coefficient sheaves. Roots of the exceptional divisors remove their inertia; the main issue is to make the resulting estimates independent of root order. The degrees of the roots and the distribution of their stabilizer characters provide that cancellation. Sun's Frobenius filtration controls the pole cutoff, and the tame curve bound of SGA~1 controls the prime-to-$p$ coefficient part. Kato's integrality theorem, in the precise published restatement of Bright and Newton \cite[Theorem 4.1, pp.~848--849]{BN}, bounds every abelian conductor by the intersection matrix. These two estimates give a uniform first-cohomology bound for abelian-monodromy modules. Counting complements to a minimal normal subgroup then proves the metabelian assertion.

\section{Quadratic refinements and the universal class}

This section proves Theorem~\ref{thm:polarization} by reducing it to finite symplectic linear algebra. Let $A/S$ be an abelian scheme with $2$ invertible on $S$, dual $A^\vee$, and polarization $\lambda:A\to A^\vee$. Use the convention
\[
 \lambda_L(a)=[t_a^*L\otimes L^{-1}].
\]
The symmetric classes in $\Pic^\delta$ form an \'etale torsor under $A^\vee[2]$. Normalize a symmetry $[-1]^*L\simeq L$ to be the identity at the origin. Its values on $A[2]$ give a quadratic refinement of
\[
 b(x,y)=e_2(x,\lambda(y)).
\]
We identify $\mu_2$ with the constant group $\F_2$ additively. The identity is
\begin{equation}\label{eq:quadratic}
 q(x+y)=q(x)+q(y)+b(x,y).
\end{equation}
Tensoring $L$ with a two-torsion degree-zero class adds the corresponding linear functional on $A[2]$. The perfect Weil pairing and the symmetric-line-bundle sequence identify the entire torsor with the sheaf of quadratic refinements \cite[Proposition 3.6]{PR}. This identification is compatible with arbitrary base change, so a choice of refinement made without auxiliary data produces a section of the relative Picard sheaf.

For the positive construction we follow \cite[Proposition~3.12(f)]{PR}. Let $R$ be the radical of $b$. Its formation commutes with base change on any locus of constant rank, since it is the kernel of a homomorphism of finite \'etale group schemes. When $b=0$, choose $q=0$. When $b$ has rank two, put
\[
 q(x)=0\quad(x\in R),\qquad q(x)=1\quad(x\notin R).
\]
The two loci are open and closed, so this formula defines a morphism of finite \'etale schemes. On the two-dimensional quotient, any two distinct nonzero vectors have pairing one and sum to the third nonzero vector. This verifies \eqref{eq:quadratic}; the cases involving zero or equal vectors give the same identity directly. The definition is invariant under every isomorphism. Since the rank is twice the number of odd elementary divisors, it proves existence in precisely the stated cases.

For necessity, work over $\mathbf C$. The automorphism $[-1]$ of a polarized abelian variety fixes any universal Picard class; hence a universal section would be symmetric. Pull it back to the space of marked polarized complex tori. This space is the connected Siegel upper half-space. The two-torsion local system is constant there, so a section of its finite quadratic-refinement torsor is constant. It must therefore be fixed by the integral automorphisms of the polarization lattice.

Choose a lattice basis $e_i,f_i$ with alternating form $E(e_i,f_j)=d_i\delta_{ij}$ and zero pairings between the $e_i$ and between the $f_i$. For every integral vector $v$ the map
\[
 T_v(x)=x+E(x,v)v
\]
is integral and preserves $E$; its inverse replaces the plus sign by a minus sign. Let $W$ be the subspace modulo two spanned by the pairs with odd $d_i$. Its form is nondegenerate, and the reductions of these integral transformations include the transvection along every nonzero $v\in W$. If $q$ were invariant, choose $x$ with $b(x,v)=1$. Then
\[
 q(x)=q(T_vx)=q(x+v)=q(x)+q(v)+1,
\]
so $q(v)=1$ for every nonzero $v\in W$. In dimension at least four, two independent orthogonal vectors $v,w$ exist. Equation~\eqref{eq:quadratic} gives $q(v+w)=0$, contradicting the preceding conclusion. This proves the negative assertion and Theorem~\ref{thm:polarization}.

\section{Schematicity of abelian torsors and their Picard functors}
\label{sec:abelian-torsor-schemes}

An abelian scheme in this section is a smooth proper group scheme of
finite presentation with geometrically connected fibres. Projectivity is
not included in the definition. All torsors and contracted products are
for the fppf topology.

The distinction between an algebraic space and a scheme matters for
torsors under abelian schemes. We prove that changing the structure group
preserves schematicity when the original torsor is a scheme.

This section proves Theorem~\ref{thm:schematic-abelian-pushout} and Corollary~\ref{cor:picard-schematic-abelian-torsor}.

There is no assumption on the kernel of $h$. In particular, the theorem
applies when the kernel has positive dimension. The proof combines
Raynaud's criterion for effective descent with the affine-neighborhood
and ampleness results for connected group actions in
\cite[\S6.6, Theorem~2(c),(d), pp.~164--166]{BLR}.
We give the finite-cover and approximation steps explicitly.

Call a scheme \emph{AF} if every finite set of its points is contained
in an affine open subscheme. We use the following form of Raynaud's
criterion \cite[Theorem~4.2, pp.~5-16--5-18]{RaynaudDescent}.
Let $F$ be an fpqc sheaf over $S$. Suppose that an fpqc cover $V\to S$
makes $F_V$ a locally noetherian scheme, and that a finite surjective
map $W\to S$ makes $F_W$ a scheme. Set
\[
 Y=F_W,\qquad Z=F_{W\times_S W},
\]
and let $v_1,v_2:Z\rightrightarrows Y$ be the projections.
Then $F$ is a scheme if and only if
$v_2(v_1^{-1}(y))$ is contained in an affine open of $Y$ for every
$y\in Y$. Thus an AF scheme $Y$ suffices. The finite cover $W\to S$
is not required to be flat, and the criterion does not require a
descended line bundle.

\begin{lemma}\label{lem:semilocal-multisection}
Let $R$ be a noetherian semilocal ring. Let $Y\to\Spec R$ be projective,
smooth and surjective, of pure relative dimension $d$.
There is a closed subscheme $Z\subset Y$ that is finite locally free
and surjective over $\Spec R$.
\end{lemma}

\begin{proof}
For $d=0$ take $Z=Y$. Suppose $d>0$, and choose a very ample line
bundle $L$ on $Y$. Write $\mathfrak m_1,\ldots,\mathfrak m_r$ for the
maximal ideals of $R$ and $Y_i$ for the corresponding fibres.
Graded Noether normalization on each $Y_i$ supplies $d$ homogeneous
sections of positive powers of $L_i$ whose common zero scheme is
nonempty and zero-dimensional: take $d$ of the $d+1$ homogeneous
parameters of a finite map to $\PP^d_{\kappa(\mathfrak m_i)}$.
This construction works over finite residue fields as well.
Taking powers puts all the sections in a common degree without
changing their common zero sets.

Set $I=\bigcap_i\mathfrak m_i$. For sufficiently large $n$, Serre
vanishing gives
\[
 H^1(Y,I\OO_Y\otimes L^n)=0.
\]
After taking further common powers, the chosen sections on the disjoint
union of the $Y_i$ therefore lift to $d$ sections of $L^n$ on $Y$.
Let $Z$ be their common zero scheme. It is projective, and its closed
fibres are nonempty and zero-dimensional.

The locus where a fibre of $Z\to\Spec R$ has positive local dimension
is closed in $Z$ by \cite[IV$_3$, 13.1.3]{EGA}. Its image is closed
by properness and contains no closed point of $\Spec R$, so is empty.
Thus $Z\to\Spec R$ is proper and quasi-finite, hence finite.

At each point above a maximal ideal of $R$, the chosen fibre sections
form a system of parameters in a regular local ring of dimension $d$.
They form a regular sequence. Since $Y$ is $R$-flat, the fibrewise
regular-sequence criterion \cite[IV$_3$, 11.3.8]{EGA} makes $Z$ flat
at those points. The flat locus is open in $Z$. The image of its
closed complement under the finite map to $\Spec R$ avoids every
closed point, so that complement is empty. Therefore $Z$ is finite
locally free over $R$. Its closed fibres are nonempty, and its rank
is locally constant; the map is consequently surjective.
\end{proof}

\begin{proof}[Proof of Theorem~\ref{thm:schematic-abelian-pushout}]
Write $Q=h_*T$. It is a $B$-torsor algebraic space, and
$Q_T\simeq B_T$. We first work over an excellent noetherian local
scheme $S=\Spec R$.

Choose an affine open $U\subset T$ meeting the closed fibre.
The map $T\to S$ is smooth, so the image of $U$ is open and contains
the closed point of $S$; it is all of $S$. The $A$-action is transitive
on geometric fibres, whence $A\cdot U=T$.
The affine-neighborhood theorem \cite[\S6.6, Theorem~2(c)]{BLR}
makes $T$ an AF scheme. The same argument, starting with an affine
neighborhood of the identity in $B$, makes $B$ AF.

Let $W\to S$ be the normalization of $S_{\mathrm{red}}$.
It is finite and surjective because $S$ is excellent; $W$ is normal,
affine and semilocal. A finite morphism preserves AF on its source:
put the finitely many image points in an affine open and take its
affine inverse image. Thus $T_W$ and $B_W$ are AF.

Choose one point of $T_W$ above each closed point of $W$.
A single affine open contains all these points. Its image in $W$
is open by smoothness, contains every closed point and hence equals
$W$. Its saturation under $A_W$ is all of $T_W$.
Normality of $W$ and \cite[\S6.6, Theorem~2(d)]{BLR} now supply
a relatively ample line bundle on $T_W$; therefore $T_W$ is
projective over $W$. The same argument makes $B_W$ projective.
The use of one affine open meeting all closed fibres is essential.

By Lemma~\ref{lem:semilocal-multisection}, there is a finite locally
free surjective closed subscheme $Z\subset T_W$. Its inclusion gives
a section of $T_Z$, and thus $Q_Z\simeq B_Z$.
Since $B_W$ is projective and $Z$ is affine, $B_Z$ is AF.
Apply Raynaud's criterion over $W$, using $Z\to W$ both as the fpqc
cover and as the finite cover. It follows first that $Q_W$ is a
scheme.

We next prove projectivity of this scheme. The map
$q:Q_Z\to Q_W$ is finite locally free and surjective, and $Q_Z$ is
projective over the affine scheme $W$. Choose an ample line bundle
$L$ on $Q_Z$. Its norm
\[
 \operatorname{Nm}_{q}(L)
 =\det(q_*L)\otimes\det(q_*\OO_{Q_Z})^{-1}
\]
is an invertible sheaf on $Q_W$, and is ample by
\cite[II, Proposition~6.6.1 and Corollary~6.6.4]{EGA}.
Properness of $Q_W\to W$ follows by fpqc descent from the $B$-torsor
structure. Hence $Q_W$ is projective over $W$, and in particular AF.
The norm is used only after the target has been proved to be a scheme.

Apply Raynaud's criterion a second time, over $S$.
The fpqc cover is $T\to S$, on which $Q_T\simeq B_T$ is a locally
noetherian scheme. The separate finite surjective cover is $W\to S$,
on which $Q_W$ is AF. The criterion proves that $Q$ is a scheme.
No flatness of $W\to S$ has been used.

We justify passage from local rings to open neighborhoods.
Start with an excellent noetherian affine $S$ and $s\in S$, and
form its finite normalization $W\to S$ before localizing.
Over $W_s=W\times_S\Spec\OO_{S,s}$ the preceding argument constructs
projective embeddings of $T_{W_s}$ and $B_{W_s}$, and a
finite locally free surjective closed subscheme $Z_s\subset T_{W_s}$.
As $W_s$ is the inverse limit of the $W_U$ for affine neighborhoods
$U$ of $s$, these finitely presented objects and maps descend to
some $W_U$. After shrinking $U$, the embeddings remain closed
immersions and $Z_U\to W_U$ is finite locally free and surjective.
These are finite-presentation limit properties
\cite[IV$_3$, 8.8--8.10 and 11.2.6(ii)]{EGA}.
The two applications of Raynaud's criterion then take place over
the open neighborhood $U$. Thus $Q_U$ is a scheme. Such
neighborhoods cover $S$.

Finally let $S$ be arbitrary, and work on an affine part of it.
The schemes $A,B,T$, their group laws, the action, the homomorphism
$h$, and the inverse of the torsor isomorphism
$A\times_S T\simeq T\times_S T$ are finite-presentation data.
They descend to a model over a finitely generated $\Z$-algebra.
After increasing the model, the finitely many identities hold,
and properness, smoothness and surjectivity hold as well
\cite[IV$_3$, 8.10.5; IV$_4$, 17.7.8(ii)]{EGA}.
For a proper smooth morphism the number of geometric connected
components is locally constant \cite[IV$_3$, 15.5.9(ii)]{EGA};
restrict the model to the open and closed locus through which
the original base factors. The fibres are then geometrically
connected. The resulting base is excellent and noetherian.
The pushout is compatible with base change, so the result there
proves schematicity over the original base.
\end{proof}



The conclusion over a connected locally noetherian base is stated
in \cite[\S1.7.2, p.~13]{ACMPrym}. The proof below spells out the varying Néron--Severi
parameter and the possibly nontrivial torsor of line-bundle classes
on $A$, then applies Theorem~\ref{thm:schematic-abelian-pushout}.
The contracted-product identity is classical; see
\cite[Theorem~3.0.3, p.~664]{Alexeev}.

\begin{proof}[Proof of Corollary~\ref{cor:picard-schematic-abelian-torsor}]
The same approximation as in the theorem reduces to a noetherian
base, since relative Picard sheaves commute with arbitrary base
change. Write $P_A=\Pic_{A/S}$ and $B=A^\vee$.
The Picard functor is a separated algebraic space locally of finite
presentation \cite[Section~108.9, Lemmas~3--4]{Stacks}:
the proper smooth morphism has geometrically integral fibres and
its degree-zero direct image is universally the structure sheaf.
Its open neutral subfunctor $B$ is an abelian scheme; see,
with the definition of abelian scheme used here,
\cite[Corollary~4.6, author version pp.~14--15]{LaurentSchroer}.
Since $B$ is proper over $S$ and $P_A$ is separated, $B$ is also
closed in $P_A$.

Form the quotient algebraic space $N=P_A/B$ for the free action
of the smooth group $B$ \cite[Tag~06PH]{Stacks}. The map $P_A\to N$ is a $B_N$-torsor.
Local finite presentation of $N\to S$ descends through this smooth
surjection. The pullback of $\Delta_{N/S}$ to $P_A\times_S P_A$
is the inverse image of $B$ under the difference map, so is closed.
Thus $N$ is separated over $S$. Its geometric fibres are discrete,
being the quotients of the geometric Picard groups by their open
identity components. Hence $N\to S$ is locally quasi-finite, and
the separated locally quasi-finite algebraic-space theorem
\cite[Tag~03XX]{Stacks} makes $N$ a scheme. We do not assert
that $N$ is flat or that its ranks are locally constant.

Regard $Q_A=P_A\to N$ as the $B_N$-torsor of the universal class
$\delta\in N(N)$. Inversion $\iota=[-1]^*$ acts by $-1$ on $B$
and by $+1$ on $N$
\cite[\S\S3.1--3.2, pp.~1311--1312]{PR}.
Its fixed subsheaf $R\subset Q_A$ is a $B[2]_N$-torsor:
if $x$ is a local point of $Q_A$ and $d=\iota(x)-x\in B$, then
\[
 \iota(x+b)-(x+b)=d-2b.
\]
The finite locally free surjection $[2]:B\to B$ therefore gives
fixed points locally, whose differences lie in $B[2]$.
Consequently $R\to N$ is finite locally free and surjective, and
$R$ is a scheme. It supplies a tautological section of $(Q_A)_R$,
so $(Q_A)_R\simeq B_R$.

For each affine open $V\subset N$, the scheme $R_V$ is affine and
noetherian. An abelian scheme over an affine base is AF
\cite[application after Corollary~3.8, p.~5-15]{RaynaudDescent}.
Raynaud's criterion, using $R_V\to V$ for both covers, makes
$(Q_A)_V$ a scheme. Hence $Q_A$ is a scheme over $N$.
This proves schematicity of the component torsor; it does not
replace that torsor by the trivial one.

For a local section $t$ of $T$, let
$j_t:A\to T$ send $a$ to $t+a$.
The induced identifications of $\Pic_{T/S}$ with $P_A$ identify
the neutral parts with $B$. Changing $t$ acts by translation on
$A$, and translation fixes both $B$ and $N$. The quotient parameters
therefore glue to a canonical map
\[
 P_T=\Pic_{T/S}\longrightarrow N,
\]
which is a $B_N$-torsor. The universal class $\delta$ defines
a homomorphism
\[
 \lambda_\delta:A_N\longrightarrow B_N,\qquad
 a\longmapsto \tau_a^*L\otimes L^{-1},
\]
where $L$ is any local representative of $\delta$.
The expression is independent of that representative because
translation fixes $B$, and is a homomorphism by the theorem of
the square.

The scheme $U=T_N\times_N Q_A$ is a torsor under the abelian
$N$-scheme $A_N\times_N B_N$. On relative Picard classes define
\[
 m:U\longrightarrow P_T,\qquad
 (t,L)\longmapsto(j_t^{-1})^*L.
\]
There is no need to choose a universal line bundle in this formula.
Since $j_{t+a}^{-1}=\tau_{-a}\circ j_t^{-1}$, one has
\[
 m(t+a,L+b)=m(t,L)+b-\lambda_\delta(a).
\]
Thus $m$ is equivariant for
$h(a,b)=b-\lambda_\delta(a)$ and induces an isomorphism
$h_*U\simeq P_T$ of $B_N$-torsors. It is an isomorphism because,
after choosing $t$ and $L$ fppf locally, it is the identity of $B_N$.
Theorem~\ref{thm:schematic-abelian-pushout} now makes $P_T$ a
scheme. Its separatedness and local finite presentation were
already properties of the Picard algebraic space.
\end{proof}

\begin{remark}
These results settle the schematic abelian-torsor subcase of the
question whether a proper flat family with geometrically normal
or geometrically integral fibres has a schematic full Picard functor.
They do not settle that question for general higher-dimensional
families. In particular, geometric normality alone does not imply
that the whole neutral Picard space is smooth or flat.
\end{remark}

\section{Projective modifications of abelian-surface torsors}
\label{sec:projective-modifications-torsors}

A projective modification of an algebraic space can have a scheme as
its source. In a family, however, the source need not remain flat or
have geometrically integral fibres. For torsors under abelian surfaces,
these requirements impose the following obstruction.

\begin{theorem}\label{thm:surface-torsor-modification}
Let $S$ be an arbitrary scheme, let $A$ be an abelian $S$-scheme of
relative dimension two, and let $Q$ be an fppf $A$-torsor algebraic space.
Let $X$ be a scheme, proper, flat and finitely presented over $S$, with
geometrically integral fibres. Suppose that $f:X\to Q$ is a projective $S$-morphism
and that one of the following conditions holds:
\begin{enumerate}
\item every geometric fibre $f_{\bar s}:X_{\bar s}\to Q_{\bar s}$ is
birational;
\item there is a dense open subspace $V\subset Q$ such that
$f^{-1}V\to V$ is an isomorphism.
\end{enumerate}
Then $Q$ is a scheme.
\end{theorem}

Here an abelian scheme is a smooth proper finitely presented group
scheme with geometrically connected fibres; projectivity is not
included in the definition. No group action or torsor structure
is assumed on $X$. Its fibres may be singular. The theorem therefore
includes projective modifications whose geometrically integral fibres
are geometrically normal. In the second condition,
birationality on every fibre is a consequence of the argument below;
it is not an additional hypothesis.

The proof constructs a finite locally free subscheme of $Q$ from a
rank-one pushforward. Its inclusion splits the torsor after a finite
flat base change, to which Raynaud's schematic descent criterion
applies. We give the sheaf construction over bases that may be
nonnormal and nonreduced.

The decomposition of a rank-one torsion-free sheaf into its determinant
and an ideal of points is classical. Su and Wang give its relative
Picard--Hilbert formulation for smooth spectral surfaces
\cite[Lemmas~3.22--3.23 and Remark~3.24, pp.~19--21]{SuWang},
crediting the double-dual line-bundle statement to Koll\'ar.
We give the relative sheaf argument in full, then use the resulting
finite subscheme for schematic descent of the torsor.

\begin{lemma}\label{lem:smooth-surface-flat-perfect}
Let $S$ be locally noetherian, let $Y\to S$ be a smooth algebraic space
of relative dimension two, and let $\mathcal F$ be a coherent
$\OO_Y$-module flat over $S$. Then $\mathcal F$ is perfect, locally
with a finite locally free resolution of length at most two.
\end{lemma}

\begin{proof}
Work on an affine \textup{\'etale} chart of $Y$, at a point $y$ over
$s\in S$. Put $R=\OO_{S,s}$ and let $B$ be the local ring of the chart
at $y$. The ring $B$ is flat over $R$, and $B/\mathfrak m_RB$ is a
regular local ring of dimension at most two. Write $M=\mathcal F_y$.
Choose finite free $B$-modules in exact sequences
\[
 0\longrightarrow K_1\longrightarrow P_0\longrightarrow M
 \longrightarrow0,
 \qquad
 0\longrightarrow K_2\longrightarrow P_1\longrightarrow K_1
 \longrightarrow0.
\]
The modules $K_1$ and $K_2$ are $R$-flat, and these sequences remain
exact modulo $\mathfrak m_R$. Since a regular local ring of dimension
at most two has global dimension at most two, the second syzygy
$K_2/\mathfrak m_RK_2$ is free over $B/\mathfrak m_RB$.
The fibre criterion for flatness
\cite[IV$_3$, 11.3.10]{EGA}, or \cite[Tag~00MP]{Stacks}, makes
$K_2$ flat over $B$. It is finite, hence free over the local ring $B$.
The zero-module case causes no exception. Spreading this resolution
to a neighborhood proves the assertion.
\end{proof}

\begin{lemma}\label{lem:relative-surface-rank-one-defect}
Let $S$ be locally noetherian, and let $Y\to S$ be a smooth algebraic space of relative
dimension two with geometrically integral fibres. Let $\mathcal F$ be
a coherent $\OO_Y$-module, flat over $S$, whose restriction to every
geometric fibre is torsion-free of rank one. There are an invertible
sheaf $\mathcal D$ and a closed subspace $Z\subset Y$, flat over $S$
with zero-dimensional fibres, fitting into an exact sequence
\[
 0\longrightarrow\mathcal F\longrightarrow\mathcal D
 \longrightarrow\mathcal D\otimes\OO_Z\longrightarrow0.
\]
The first map remains injective on every fibre. If $Y\to S$ is proper,
then $Z\to S$ is finite locally free.
\end{lemma}

\begin{proof}
Let $U\subset Y$ be the locus where $\mathcal F$ is locally free of
rank one. If $\mathcal F_s$ is free at a point of $Y_s$, the fibre
criterion for flatness makes $\mathcal F$ free at that point as well.
A torsion-free rank-one module over a regular local ring of dimension
zero or one is free. Hence every point of $Y_s\setminus U_s$ has local
dimension two in $Y_s$.

By Lemma~\ref{lem:smooth-surface-flat-perfect}, $\mathcal F$ has finite
locally free resolutions on \textup{\'etale} charts. The alternating
determinant of such a resolution is a line bundle with a canonical
identification with $\mathcal F$ on $U$. We explain both its gluing
and the required extension of $\mathcal F$.

At a point $y\notin U$, choose a regular sequence of length two in
the regular fibre local ring $\OO_{Y_s,y}$ and lift it to the local
ring of an \textup{\'etale} chart of $Y$. Since $Y$ is $S$-flat, the
fibrewise regular-sequence criterion
\cite[IV$_3$, 11.3.8]{EGA} makes the lifts a regular sequence in that
local ring. Thus every invertible sheaf has depth at least two there.
Similarly, choose a nonzero element of the maximal ideal of
$\OO_{Y_s,y}$. It is a nonzerodivisor on the torsion-free module
$(\mathcal F_s)_y$. Flatness of $\mathcal F$ over $S$ and the same
criterion show that a lift is a nonzerodivisor on $\mathcal F_y$.
Consequently $\mathcal F_y$ has depth at least one.

Put $j:U\hookrightarrow Y$. The depth-one assertion means that no
associated point of $\mathcal F$ lies outside $U$, so
\[
 \mathcal F\longrightarrow j_*(\mathcal F|_U)
\]
is injective. Indeed, a nonzero kernel would have an associated point
of $\mathcal F$ in the complement. The depth-two extension lemma
\cite[Tag~0E9I]{Stacks} says that any invertible sheaf $\mathcal L$
on a chart satisfies $\mathcal L=j_*(\mathcal L|_U)$ there.
This lemma requires no normality or reducedness.

On a double overlap of two charts, the determinant lines have an
isomorphism over $U$ through their identifications with $\mathcal F$.
Their Hom sheaf is invertible. The depth-two extension lemma extends
that isomorphism uniquely; it also extends its inverse. Restriction
is injective, so the extended maps are inverse and satisfy the cocycle
identity on triple overlaps. \textup{\'Etale} descent therefore gives
an invertible sheaf $\mathcal D$ on $Y$ with
$\mathcal D|_U=\mathcal F|_U$. This construction does not require a
global locally free resolution on $Y$. The same extension lemma gives
\[
 \mathcal F\lhook\joinrel\longrightarrow
 j_*(\mathcal F|_U)=j_*(\mathcal D|_U)=\mathcal D.
\]

For every geometric $s$, this map restricts to an isomorphism on the
dense open $U_s$ of the integral surface $Y_s$. Its fibre kernel has
rank zero, so torsion-freeness of $\mathcal F_s$ makes that kernel
zero. Let $\mathcal C$ be the cokernel on $Y$. Since $\mathcal D$ is
$S$-flat, the fibrewise injectivity gives
\[
 \operatorname{Tor}_1^{\OO_{S,s}}
       (\mathcal C_y,\kappa(s))=0
 \quad\text{for every }y\text{ over }s.
\]
The local criterion for flatness \cite[Tag~00MK]{Stacks} applies to
the finite module $\mathcal C_y$ over the noetherian local ring at
$y$, and shows that $\mathcal C$ is $S$-flat.

Tensoring the inclusion by $\mathcal D^{-1}$ gives an ideal
$\mathcal I\subset\OO_Y$. Define $Z$ by $\mathcal I$. Then
$\mathcal C=\mathcal D\otimes\OO_Z$, so $Z$ is flat over $S$.
Its fibres are supported in $Y_s\setminus U_s$, which has dimension
zero. If $Y/S$ is proper, $Z/S$ is proper and quasi-finite, hence
finite \cite[Tag~0A4X]{Stacks}. It is then finite locally free.
\end{proof}

\begin{lemma}\label{lem:ample-pushforward-flat-basechange}
Let $S$ be noetherian affine, let $Y$ be a quasi-compact separated
algebraic space of finite type over $S$, and let $f:X\to Y$ be proper
with an $f$-ample invertible sheaf $L$. Suppose that $X$ is a scheme,
flat and of finite type over $S$. For sufficiently large $n$, the
coherent sheaf
\[
 \mathcal F=f_*L^n
\]
is flat over $S$, its formation commutes with arbitrary base change
on $S$, and the evaluation $f^*\mathcal F\to L^n$ is surjective.
\end{lemma}

\begin{proof}
Relative Serre vanishing and generation, applied on a finite affine
\textup{\'etale} cover of $Y$, give a common $n$ for which
$R^if_*L^n=0$ for $i>0$ and the evaluation map is surjective.
Properness makes $\mathcal F$ coherent.

Let $V\to Y$ be an affine \textup{\'etale} chart. The scheme $X_V$
is quasi-compact, separated and $S$-flat. A finite affine cover of
$X_V$ has affine finite intersections and gives a bounded ordered
\v Cech complex $C^\bullet$ computing the cohomology of $L^n$ on
$X_V$. Every term is flat over $\OO(S)$: the coordinate ring of each
affine intersection is flat over $\OO(S)$, and the module associated
to the restricted line bundle is finite projective over that ring.
The positive cohomology of $C^\bullet$ vanishes by the chosen
relative Serre vanishing and flat base change to $V$.

Starting at the final nonzero term, exactness and flatness of the
successive quotients show that every successive kernel is flat,
including $H^0(C^\bullet)$. The corresponding short exact sequences
remain exact after tensoring with any $\OO(S)$-module. Thus tensoring
the complex preserves its vanishing in positive degrees and computes
the tensor product of $H^0(C^\bullet)$ in degree zero.
After any affine base change $S'\to S$, the base-changed affine cover
has exactly this tensor-product \v Cech complex. This proves both
flatness and the base-change isomorphism
\[
 \mathcal F|_{Y_{S'}}\simeq
 (f_{S'})_*\bigl(L^n|_{X_{S'}}\bigr).
\]
The assertions glue across the \textup{\'etale} cover of $Y$ and
across an affine cover of $S'$. In particular $S'\to S$ need not be
flat or noetherian.
\end{proof}

\begin{proof}[Proof of Theorem~\ref{thm:surface-torsor-modification}]
We first assume that $S$ is noetherian affine and condition~(1) holds.
Choose an $f$-ample line bundle $L$ and take $n$ as in
Lemma~\ref{lem:ample-pushforward-flat-basechange}. Write
$\mathcal F=f_*L^n$. On every geometric fibre,
\[
 \mathcal F_{\bar s}=(f_{\bar s})_*L_{\bar s}^n
\]
is torsion-free of rank one. Birationality gives rank one. To check
torsion-freeness, a section killed by a nonzero function on the
integral target vanishes over a dense open. Its pullback is a section
of a line bundle on the integral $X_{\bar s}$, so it is zero.

Apply Lemma~\ref{lem:relative-surface-rank-one-defect} to obtain an
invertible $\mathcal D$ and a finite locally free closed subscheme
$Z\subset Q$ with $\mathcal F=\mathcal D\otimes\mathcal I_Z$.
For every geometric $s$, we claim
\[
 Z_{\bar s}=\varnothing
 \quad\Longleftrightarrow\quad
 f_{\bar s}\text{ is an isomorphism}.
\]
If $f_{\bar s}$ is an isomorphism then $\mathcal F_{\bar s}$ is
invertible. The fibre criterion makes $\mathcal F$ invertible at
all points of that fibre, so $Z_{\bar s}$ is empty.
Conversely, if $Z_{\bar s}$ is empty, the evaluation map is a
surjection of line bundles
\[
 f_{\bar s}^*\mathcal D_{\bar s}\longrightarrow L_{\bar s}^n,
\]
hence an isomorphism. The line bundle $L_{\bar s}^n$ is therefore
trivial on each fibre of $f_{\bar s}$. Since it is relatively ample,
these proper fibres are zero-dimensional. Thus $f_{\bar s}$ is
finite. It is birational onto the normal surface $Q_{\bar s}$ and
has integral source, so it is an isomorphism.

The rank of $Z/S$ is locally constant. Decompose $S$ into the open
and closed parts $S_0$ and $S_+$ on which that rank is respectively
zero and positive. Over $S_0$ every fibre map is an isomorphism, and
so is $f$. Indeed, its non-quasi-finite locus is empty, making $f$
finite. The coherent cokernel of $\OO_Q\to f_*\OO_X$ has zero fibres,
so it is zero by Nakayama's lemma. The source and target of this
surjection are $S$-flat; hence its kernel also has zero fibres and
is zero by Nakayama's lemma. It follows that $Q_{S_0}\simeq X_{S_0}$
is a scheme. This also explains why, over a local base, an isomorphism
on the closed fibre cannot hide a nonisomorphic fibre elsewhere:
the proper image of the closed non-quasi-finite locus would avoid
the closed point, so that locus would be empty. For the resulting
finite map, the cokernel of $\OO_Q\to f_*\OO_X$ vanishes along the
closed fibre by Nakayama. Its support is proper over the local base,
and therefore it is empty. The kernel then vanishes by the same
argument, using $S$-flatness to preserve exactness on the closed fibre.

Over $S_+$, the map $Z\to S_+$ is finite locally free and surjective.
Its inclusion in $Q$ gives a section of $Q_Z$, so $Q_Z\simeq A_Z$.
An abelian scheme over an affine base is AF, meaning that every finite
set of its points lies in an affine open; see the application following
\cite[Corollary~3.8, p.~5-15]{RaynaudDescent}. Finite morphisms preserve AF on
their source. Thus, after restricting to an affine part of $S_+$,
$A_Z$ is AF because it is finite over $A$.

Apply \cite[Theorem~4.2, pp.~5-16--5-18]{RaynaudDescent} with
$Z\to S_+$ as both its fpqc cover and its finite cover. On that cover
$Q_Z=A_Z$ is a noetherian scheme, and its finite descent orbits have
affine neighborhoods because it is AF. Therefore $Q_{S_+}$ is a
scheme. This proves the noetherian case under condition~(1).

We now handle both conditions over an arbitrary base. Work on an
affine open $S=\Spec R$ and write $R$ as a filtered union of finitely
generated $\mathbf Z$-subalgebras. Descend $A$ and $X$ as schemes,
$Q$ as an algebraic space, and the morphism $f$, the group law,
the action, and the torsor isomorphism, together with its inverse.
The descent of finitely presented algebraic spaces and their maps
is \cite[Tag~07SK]{Stacks}; the scheme version is
\cite[IV$_3$, 8.8--8.10]{EGA}. All finitely many identities hold
after increasing the index.

One can retain the torsor property by also descending an actual
trivializing cover. Choose an affine scheme $W$ with an \textup{\'etale}
surjection $W\to Q$ of finite presentation. Then $W\to S$ is smooth
and surjective, and $W\to Q$ gives a section of $Q_W$. Descend $W$,
this section, and the resulting isomorphism $Q_W\simeq A_W$ with
its inverse. Retain smoothness and surjectivity of $W_i\to S_i$.
The descended $Q_i$ is consequently an $A_i$-torsor; in particular
it is proper and smooth over $S_i$.

Retain properness of the scheme morphisms $X_i\to S_i$ and
$A_i\to S_i$ by \cite[IV$_3$, 8.10.5]{EGA}, flatness by
\cite[IV$_3$, 11.2.6(ii)]{EGA}, and smoothness of $A_i$ and $W_i$
by \cite[IV$_4$, 17.7.8(ii)]{EGA}. Geometric integrality of the
$X_i$-fibres and geometric connectedness of the $A_i$-fibres can
also be retained. The bad loci on a noetherian model are constructible
by \cite[IV$_3$, 9.7.7]{EGA}, and their inverse images in $S$ are
empty. Each bad locus is a finite union of locally closed pieces
$\Spec((R_i/I)[1/a])$. The tensor product of the coordinate ring
of such a piece with $R$ is zero; the equation $1=0$ already holds
at a finite stage. Increasing the index removes all the pieces.
The relative dimension of $A_i$ can be fixed to two by removing its
other open-and-closed dimension loci, whose pullback to $S$ is empty.

Finally descend $L$ as an invertible sheaf and retain its relative
ampleness. Choose a finite affine \textup{\'etale} cover of a model
of $Q$ and pull it back through the system. On each such affine
target, relative ampleness is ampleness on the inverse-image scheme.
The limit criterion \cite[IV$_3$, 8.10.5.2]{EGA}, equivalently
\cite[Tag~09MT]{Stacks}, makes $L_i$ ample there at a common finite
stage. Relative ampleness is \textup{\'etale} local on the target,
so $L_i$ is $f_i$-ample. The map $f_i$ is proper because $X_i$ is
proper over $S_i$ and $Q_i$ is separated over $S_i$. At this point
no birationality assertion about $f_i$ has been imposed.

On this noetherian model choose $n$ by
Lemma~\ref{lem:ample-pushforward-flat-basechange} and put
\[
 \mathcal F_i=(f_i)_*L_i^n.
\]
This sheaf is $S_i$-flat and commutes with arbitrary base change.
By Lemma~\ref{lem:smooth-surface-flat-perfect} it is perfect on
$Q_i$. Its virtual rank $r_i$, the alternating sum of the ranks in
a finite locally free resolution, is locally constant and is preserved
by base change. Since the fibres of $Q_i/S_i$ are geometrically
connected, $r_i$ is constant on each fibre. It therefore comes from
a locally constant function $\rho_i$ on $S_i$: the images of its
disjoint rank loci are disjoint opens, because $Q_i\to S_i$ is
smooth, open and surjective.

Under condition~(1), the base change $\mathcal F=f_*L^n$ has generic
rank one on every geometric fibre, hence virtual rank one everywhere.
To see this directly, a finite locally free resolution remains exact
after base change on $S_i$, since $\mathcal F_i$ and its syzygies
are $S_i$-flat. At the generic point of an integral fibre its Euler
rank is the dimension of its sole cohomology module. Under
condition~(2), $\mathcal F$ is invertible on the dense open $V$ where
$f$ is an isomorphism. Its locally constant virtual rank is then one
on a dense subset, hence everywhere. In either case $S\to S_i$
factors through the open-and-closed locus where $\rho_i=1$.
Replace $S_i$ by that locus.

Let $\bar s$ be a geometric point of the resulting model, and let
$\eta$ be the generic point of the integral surface $Q_{i,\bar s}$.
Flatness and base change give
\[
 \dim_{\kappa(\eta)}
 H^0\bigl((X_i)_\eta,L_i^n|_{(X_i)_\eta}\bigr)=1.
\]
Indeed, base-changing a finite locally free resolution to $\bar s$
and then localizing at $\eta$ gives a complex of finite-dimensional
vector spaces with sole cohomology $(\mathcal F_i)_{\bar s,\eta}$,
so that dimension equals its virtual rank one.
The evaluation map is surjective; consequently the single nonzero
global section generates the line bundle on $(X_i)_\eta$, making
it trivial. It is also ample. A proper scheme over a field with
ample trivial line bundle is zero-dimensional and hence finite.
Its function algebra has dimension one, so
\[
 (X_i)_\eta=\Spec\kappa(\eta).
\]
The integral source $X_{i,\bar s}$ therefore maps birationally to
$Q_{i,\bar s}$. This holds for every geometric $\bar s$ of the
model. The noetherian case already proved applies, making $Q_i$ a
scheme. Its base change $Q$ is a scheme. The affine opens of $S$
cover the original base, which proves the theorem.
\end{proof}

\begin{remark}\label{rem:surface-modification-scope}
The theorem excludes projective scheme modifications of nonschematic
abelian-surface torsors that retain flatness and geometrically integral
fibres. It does not assert schematicity of the full Picard functor of
an arbitrary proper flat family. The surface hypothesis enters when
the rank-one defect has zero-dimensional fibres. In higher relative
dimension the defect can be positive-dimensional, so the argument
does not supply the finite splitting cover required by schematic
descent. The proof also shows that projectivity of $f$ may be replaced
by properness together with an $f$-ample invertible sheaf.
\end{remark}

\section{Projective bundles over abelian torsors}
\label{sec:torsor-projective-bundles}

The following theorem excludes scheme projective bundles over
nonschematic abelian torsors over every base.

\begin{theorem}\label{thm:torsor-projective-bundle}
Let $S$ be an arbitrary scheme, let $A$ be an abelian $S$-scheme,
and let $Q$ be an fppf $A$-torsor algebraic space. Let $E$ be a finite
locally free $\OO_Q$-module whose rank is everywhere positive.
Then $\PP_Q(E)$ is a scheme if and only if $Q$ is a scheme.
\end{theorem}

An abelian scheme here is a smooth proper group scheme of finite
presentation with geometrically connected fibres; projectivity is not part of
the definition. We use the convention that $\PP_Q(E)$ parametrizes
invertible quotients, with tautological sheaf $\OO_{\PP_Q(E)}(1)$.
The bundle is locally free on the \'etale site of $Q$ and need not
have constant rank over the whole base. The converse implication
in the theorem is the usual construction of a projective bundle
over a scheme.

The proof begins with the normal local case. Affine boundaries and
the Ramanujam--Samuel theorem supply line bundles whose numerical
classes cannot vanish along a common positive-dimensional
translation kernel. For a general local base we apply this argument
simultaneously over its finite normalization. Selected relative
Picard classes provide projective targets, making that normalized
torsor projective. Raynaud's finite-cover criterion then descends
schematicity. Approximation removes the noetherian hypothesis.

Nonschematic elliptic torsors already occur over normal
noetherian local surfaces \cite[Theorem~2.1 and \S3]{Zomervrucht}.
The theorem applies to vector bundles of every positive rank on
these torsors, in every relative dimension, without a homogeneity
assumption. Brion's classification concerns homogeneous projective
bundles over abelian varieties over an algebraically closed field
\cite[pp.~2475--2479]{BrionProjective}.

\subsection{Affine boundaries and translation kernels}

\begin{lemma}\label{lem:normal-affine-relative-boundary}
Let $R$ be a normal noetherian local ring, with closed point
$s\in S=\Spec R$. Let $X\to S$ be a smooth proper scheme with
geometrically integral fibres. For every affine open $U\subset X$
meeting $X_s$, there is an effective Cartier divisor $D$ on $X$
whose support is $X\setminus U$. Its restriction to every fibre is
an effective Cartier divisor, and $D$ is flat over $S$.
\end{lemma}

\begin{proof}
The ring $R$ is a domain. Smoothness over the normal base makes $X$
normal. It is also irreducible. Indeed, a nonempty open of $X$ has
nonempty open image in the integral scheme $S$, so it meets the
generic fibre. Two such opens meet because that fibre is integral.
Thus $X$ is integral.

The map $U\to S$ is smooth and hence open. Its image contains the
closed point of the local scheme $S$, so its image is all of $S$.
Consequently $U_t$ is nonempty and dense in $X_t$ for every $t\in S$.
Since $X$ is separated over the affine scheme $S$, it is a separated
scheme. Its affine opens have affine intersections, and therefore
the inclusion $U\hookrightarrow X$ is an affine morphism.

By EGA IV$_4$, 21.12.7, every irreducible component of $X\setminus U$
has codimension one \cite[pp.~307--308]{EGA}. There are finitely many
such components. Let $D$ be the Weil cycle which is their sum, with
coefficient one. Its support contains no irreducible component of
any fibre, since $U_t$ is dense in the integral fibre $X_t$.
The second assertion of EGA IV$_4$, 21.14.3(i), applies to the smooth
morphism $X\to S$ and the normal locally noetherian base: precisely
such codimension-one cycles are locally principal
\cite[pp.~325--328]{EGA}.

Thus $D$ is a Cartier divisor. It is effective as a Cartier divisor:
a local rational equation has nonnegative valuation at every
codimension-one point, so normality makes it regular. Since $X$ is
integral, that equation is a non-zero-divisor. Its zero set is exactly
$X\setminus U$. On each integral fibre the equation remains nonzero,
because $U_t$ is nonempty; it is again a non-zero-divisor. Finally,
$\OO_X$ and $\OO_X(-D)$ are flat over $S$, and
\[
 0\longrightarrow\OO_X(-D)\longrightarrow\OO_X
 \longrightarrow\OO_D\longrightarrow0
\]
remains exact on every fibre. The fibre criterion for flatness gives
the flatness of $D\to S$.
\end{proof}

For a line bundle $L$ on an abelian variety $B$ over an algebraically
closed field, write
\[
 \phi_L:B\longrightarrow B^\vee,
 \qquad b\longmapsto[t_b^*L\otimes L^{-1}].
\]
The theorem of the square makes this a homomorphism. Its kernel is
all of $B$ exactly when $L$ lies in $\Pic^0(B)$, equivalently when
$L$ is algebraically trivial
\cite[Proposition~9.2, Remark~9.3 and Proposition~10.1,
pp.~14--17 of the author revision]{MilneAV}.
If $i:B'\hookrightarrow B$ is an abelian subvariety and
$\phi_L\circ i=0$, then
\[
 \phi_{i^*L}=i^\vee\circ\phi_L\circ i=0.
\]
It follows that $i^*L$ is algebraically, and hence numerically, trivial.
We will apply this implication to restrictions of global line
bundles on a torsor.

\begin{lemma}\label{lem:torsor-finite-translation-kernel}
Let $S$ be the spectrum of a noetherian local ring, with closed point
$s$, let $A/S$ be an abelian scheme, and let $Q$ be an fppf
$A$-torsor algebraic space. Choose a geometric point $\bar s$ above
$s$ and a point $q_0\in Q_{\bar s}$. For a line bundle $M$ on $Q$,
put $L_M=\tau_{q_0}^*(M|_{Q_{\bar s}})$, where
$\tau_{q_0}:A_{\bar s}\to Q_{\bar s}$ is the torsor isomorphism.
If there are finitely many line bundles $M_1,\ldots,M_n$ on $Q$ such
that
\[
 \bigcap_{i=1}^n\ker\phi_{L_{M_i}}\subset A_{\bar s}
\]
is finite, then $Q$ is a scheme.
\end{lemma}

\begin{proof}
Corollary~\ref{cor:picard-schematic-abelian-torsor}, applied to $A$
itself, shows that $P=\Pic_{A/S}$ is a separated $S$-scheme locally
of finite presentation. For every line bundle $M$ on $Q$, the action
map $A\times_S Q\to Q$ pulls $M$ back to a family of line bundles
on $A$, and therefore defines a morphism
\[
 f_M:Q\longrightarrow P.
\]
More explicitly, a point $q\in Q(T)$ is sent to the relative Picard
class of $\tau_q^*M$ on $A_T$. On the geometric closed fibre, the
identification given by $q_0$ writes $f_M$ as
\[
 a\longmapsto[L_M]+\phi_{L_M}(a).
\]
In particular the nonempty fibres of
$f=(f_{M_1},\ldots,f_{M_n}):Q\to P^n$ over $\bar s$ are translates
of the indicated finite intersection of kernels.

The morphism $f$ is proper. Its graph in $Q\times_S P^n$ is a closed
immersion because $P^n$ is separated over $S$, and the projection
$Q\times_S P^n\to P^n$ is proper because $Q\to S$ is proper.
It is quasi-finite at every point over $s$, as can be checked after
the residue-field extension to $\bar s$. The non-quasi-finite locus
is closed in $Q$. Its image in $S$ is closed by properness, and does
not contain $s$. A nonempty closed subset of a local scheme contains
its closed point, so this locus is empty. Hence $f$ is proper and
quasi-finite, and is finite \cite[Tag~0A4X]{Stacks}.
A finite algebraic space over a scheme is a scheme, proving the lemma.
\end{proof}

The lemma does not require a polarization, flatness of the finite
kernel, or local constancy of N\'eron--Severi groups. The use of the
full Picard scheme as target avoids choosing components in a family.

\subsection{The normal noetherian local case}

\begin{proposition}\label{prop:normal-local-torsor-projective-bundle}
Theorem~\ref{thm:torsor-projective-bundle} holds when $S$ is
the spectrum of a normal noetherian local ring.
\end{proposition}

\begin{proof}
Write $X=\PP_Q(E)$, let $\pi:X\to Q$ be the projection, and suppose
that $X$ is a scheme. The torsor $Q\to S$ is smooth, proper and
finitely presented, with geometrically integral fibres. It is
connected, since its fibres are connected and the base is local.
Thus $E$ has constant rank $r$. If $r=1$, the morphism $\pi$ is an
isomorphism and the conclusion follows. Assume $r\geq2$.
The scheme $X\to S$ is smooth and proper with geometrically integral
fibres, so Lemma~\ref{lem:normal-affine-relative-boundary} applies.

We first record the line-bundle decomposition on $X$. For every
invertible sheaf $H$ on $X$ there are an integer $m$ and an invertible
sheaf $M$ on $Q$ with
\begin{equation}\label{eq:torsor-bundle-picard-decomposition}
 H\simeq\OO_X(m)\otimes\pi^*M.
\end{equation}
This formula is valid for the algebraic-space base $Q$. On affine
\'etale charts of $Q$, the ordinary projective-bundle formula gives
a locally constant fibre degree. These degrees agree on overlaps,
and connectedness makes their value $m$ constant. After twisting
by $\OO_X(-m)$, every projective-space fibre has trivial restriction.
Cohomology and base change makes $\pi_*(H(-m))$ invertible, and the
evaluation map is an isomorphism on every fibre, hence an isomorphism.
The Hilbert-polynomial and base-change assertions used here are
EGA III$_2$, 7.9.11 and 7.9.8--7.9.9, together with 7.7.5
\cite{EGA}. These statements and the resulting line bundle
descend along the
\'etale charts, giving \eqref{eq:torsor-bundle-picard-decomposition}.
The same calculation proves uniqueness.

Let $U\subset X$ be an affine open meeting the closed fibre, and let
$D$ be its effective Cartier boundary from
Lemma~\ref{lem:normal-affine-relative-boundary}. In
\eqref{eq:torsor-bundle-picard-decomposition} for $H=\OO_X(D)$,
one has $m\geq0$. Indeed, choose a geometric point of $U$ and
restrict the canonical section of $\OO_X(D)$ to its
$\pi$-fibre $\PP^{r-1}$. This is a nonzero section of $\OO(m)$,
which is impossible for $m<0$. Other $\pi$-fibres may be contained
in $D$; this sign argument makes no flatness assertion over $Q$.

Suppose, towards a contradiction, that $Q$ is not a scheme. Put
$k=\overline{\kappa(s)}$, choose $q_0\in Q_k$, and use it to
identify $Q_k$ with $A_k$. For every global line bundle $M$ on $Q$
define $L_M$ as in Lemma~\ref{lem:torsor-finite-translation-kernel}.
Consider the scheme-theoretic intersection
\[
 K=\bigcap_{M\in\Pic(Q)}\ker\phi_{L_M}\subset A_k.
\]
This intersection is already a finite intersection. To see this,
its defining quasi-coherent ideal is the sum of the ideals of the
individual kernels. On the noetherian scheme $A_k$ this ideal is
of finite type. On each member of a finite affine cover, finitely
many sections drawn from finitely many of those ideals generate
the sum. Taking the union of the finite lists proves the assertion.
Thus no finite-generation assertion about $\Pic(Q)$ or its image
in a N\'eron--Severi group is needed.

By Lemma~\ref{lem:torsor-finite-translation-kernel}, $K$ is not
finite. A zero-dimensional closed subscheme of the proper
$k$-scheme $A_k$ is finite, so $K$ has positive dimension. Over the
algebraically closed field $k$, its reduced identity component
\[
 B=(K_{\mathrm{red}})^0
\]
is consequently a positive-dimensional abelian subvariety of $A_k$.
Choose an integral projective curve $C_0\subset B$ and let
$C=\tau_{q_0}(C_0)\subset Q_k$. Such a curve exists because an
abelian variety over a field is projective. Smoothness of the curve
is not needed.

For every global line bundle $M$ on $Q$, the homomorphism
$\phi_{L_M}$ vanishes on $B$. The observation preceding
Lemma~\ref{lem:torsor-finite-translation-kernel} shows that
$L_M|_B$ is algebraically trivial. In particular,
\begin{equation}\label{eq:torsor-curve-degrees-zero}
 \deg(M|_C)=0\quad\text{for every }M\in\Pic(Q),
 \qquad \deg((\det E)|_C)=0.
\end{equation}
Set $Y=\PP_C(E|_C)$, an integral projective $k$-scheme of dimension
$r$, and put $\xi=c_1(\OO_Y(1))$. The inclusion $C\hookrightarrow
Q_k$ identifies $Y$ with a closed subscheme of $X_k$.
The projective-bundle formula and Chern relation give
\begin{equation}\label{eq:torsor-bundle-top-intersection-zero}
 \deg\bigl(\xi^r\cap[Y]\bigr)
   =\deg((\det E)|_C)=0.
\end{equation}
For clarity, with the quotient convention the relation over the
curve is
$\xi^r\cap[Y]=\pi_C^*c_1(E|_C)\,\xi^{r-1}\cap[Y]$;
pushing forward uses
$\pi_{C*}(\xi^{r-1}\cap[Y])=[C]$.
These formulas apply to an integral, possibly singular curve
\cite[Tags~02TX and~02TZ, equation~(42.37.1.1)]{Stacks}.

We now prove, for every positive-dimensional integral closed
subvariety $Z\subset Y$, that
\begin{equation}\label{eq:torsor-bundle-positive-intersections}
 \deg\bigl(\xi^d\cap[Z]\bigr)>0,
 \qquad d=\dim Z.
\end{equation}
Choose a point of $Z$ and an affine open $U\subset X$ containing
its image. This open meets the closed fibre of $X\to S$.
Let $D$ be its boundary from
Lemma~\ref{lem:normal-affine-relative-boundary}, and write
\[
 \OO_X(D)\simeq\OO_X(m)\otimes\pi^*M,
 \qquad m\geq0.
\]
The mixed intersection contributed by $M$ is zero. More explicitly,
the projection formula gives
\[
 \deg\bigl(c_1(\pi_C^*(M|_C))\,\xi^{d-1}\cap[Z]\bigr)
 =\deg\bigl(c_1(M|_C)\cap
       \pi_{C*}(\xi^{d-1}\cap[Z])\bigr)=0.
\]
Indeed, the pushforward on the right is a one-cycle on the integral
curve $C$, hence an integer multiple of $[C]$, and
\eqref{eq:torsor-curve-degrees-zero} applies. This also verifies the
required numerical equality when $C$ is singular or $Z$ lies in
a fibre of $Y\to C$.

The open $U_k\cap Z$ is nonempty and affine: it is a closed
subscheme of the affine scheme $U_k$. It cannot be all of $Z$, since
a positive-dimensional integral projective scheme over a field is
not affine. Thus the restriction $D_Z$ is a nonempty effective
Cartier divisor on $Z$. Its associated cycle has the form
\[
 [D_Z]=\sum_j a_j[Z_j],\qquad a_j>0,
\]
where every $Z_j$ is integral of dimension $d-1$.
If $d=1$, its degree is strictly positive, and the preceding
projection calculation yields
\[
 m\deg\bigl(\xi\cap[Z]\bigr)=\deg[D_Z]>0.
\]
Because $m\geq0$, this proves $m>0$ and
\eqref{eq:torsor-bundle-positive-intersections} in dimension one.

For $d>1$, assume the assertion in all smaller positive dimensions.
Intersecting $[D_Z]$ with $\xi^{d-1}$ gives
\[
 \begin{split}
 m\deg\bigl(\xi^d\cap[Z]\bigr)
  &=\deg\bigl(\xi^{d-1}\cap[D_Z]\bigr)\\
  &=\sum_j a_j\deg\bigl(\xi^{d-1}\cap[Z_j]\bigr)>0.
 \end{split}
\]
Again $m\geq0$ proves the desired strict positivity. The induction
is complete. Taking $Z=Y$ contradicts
\eqref{eq:torsor-bundle-top-intersection-zero}.
Therefore $Q$ is a scheme.
\end{proof}

The argument uses a boundary chosen separately for each integral
subvariety $Z$. A single effective boundary need not be ample.
What matters is that every such boundary has, on $Y$, its numerical
class on the same nonnegative ray, and that an appropriate boundary
cuts every positive-dimensional $Z$ nontrivially. The calculation
after extension to $k$ is entirely an intersection calculation;
schematicity is not descended from an algebraic closure of the
residue field.

\subsection{Projective targets for relative Picard classes}

We will need projectivity over a finite normal cover of a local
base. The following argument verifies the simultaneous affine
neighborhood needed for the group scheme.

\begin{lemma}\label{lem:dual-projective-finite-normal-cover}
Let $S$ be noetherian local, let $W\to S$ be finite and surjective
with $W$ normal, and let $B$ be an abelian $S$-scheme. Then
$B_W$ is projective over $W$.
\end{lemma}

\begin{proof}
Choose an affine open $V\subset B$ meeting the closed fibre.
Since $V\to S$ is smooth and open, its image is all of $S$.
Translation therefore saturates $V$ to $B$, as can be checked on
geometric fibres. The hypotheses of \cite[6.6, Theorem~2(c)]{BLR}
are satisfied: the base is affine and locally noetherian, the group
is smooth with connected fibres, and the saturated open is
quasi-affine. Hence every finite set in $B$ has a common affine
open neighborhood.

The map $B_W\to B$ is finite. Thus the same property holds in $B_W$:
put the finitely many image points in an affine open of $B$ and
take its affine inverse image. The scheme $W$ is affine,
noetherian and semilocal. Choose one point in every closed
$W$-fibre of $B_W$, and put these finitely many points in a single
affine open $V_W\subset B_W$. Its image in $W$ is open by
smoothness and contains all closed points, so is all of $W$.
Indeed a nonempty closed subset of an affine scheme contains a
closed point. Translation saturates $V_W$ to $B_W$.
Since $W$ is normal, \cite[6.6, Theorem~2(d)]{BLR} gives a
$W$-ample invertible sheaf on $B_W$. Properness makes $B_W$
projective.
\end{proof}

\begin{lemma}\label{lem:projective-picard-class-torsor}
Let $T$ be a noetherian affine scheme and let $C/T$ be an abelian
scheme whose dual $B=C^\vee$ is projective over $T$. Put
$P=\Pic_{C/T}$, and let $\delta$ be a section of the fppf quotient
$P/B$ over $T$. Then
\[
 P_\delta=P\times_{P/B,\delta}T
\]
is a projective $T$-scheme.
\end{lemma}

\begin{proof}
The sheaf $P_\delta$ is a $B$-torsor. In particular it is an
algebraic space, smooth and proper of finite presentation over $T$.
No representability or local constancy assertion about the whole
quotient $P/B$ is needed. The Picard identities
\cite[\S\S3.1--3.2]{PR} say that
$\iota=[-1]^*$ acts as $-1$ on $B$ and as $+1$ on $P/B$.
Thus it preserves $P_\delta$. If $x$ is a local section of this
torsor and $b$ is a section of $B$, then
\[
 \iota(x+b)-(x+b)=\iota(x)-x-2b.
\]
Multiplication by $2$ on $B$ is finite locally free and surjective,
including in characteristic $2$
\cite[Corollary~4.6]{LaurentSchroer}.
The fixed subsheaf $R_\delta$ is consequently a torsor under
$B[2]$. By descent of finite locally free affine morphisms,
$R_\delta\to T$ is a finite locally free surjective scheme. On
a relative-dimension-$g$ component its degree is $2^{2g}$.

The inclusion $R_\delta\to P_\delta$ supplies a tautological
section after base change, and hence an isomorphism
\[
 (P_\delta)_{R_\delta}\simeq B_{R_\delta}.
\]
Only a relative Picard class is used here; a universal line bundle
on $C\times_TP_\delta$ is not required.
The right side is projective over the affine scheme $R_\delta$,
and therefore every finite set in it has an affine neighborhood.
Apply \cite[Theorem~4.2]{RaynaudDescent} with both the fpqc cover
and the separate finite surjective cover equal to
$R_\delta\to T$. The pullback is a locally noetherian scheme
with the required affine neighborhoods. It follows that
$P_\delta$ is a scheme.

We can now apply the norm to schemes. Write
$g:(P_\delta)_{R_\delta}\to P_\delta$ for the finite locally free
projection, and pull a $T$-ample invertible sheaf on $B$ back to
$B_{R_\delta}\simeq(P_\delta)_{R_\delta}$. Denote the resulting
sheaf by $L'$. It is ample over $T$, since $B_{R_\delta}\to B$ is
finite. Its norm is the invertible sheaf
\[
 N_g(L')=
 \det(g_*L')\otimes
 \det\bigl(g_*\OO_{(P_\delta)_{R_\delta}}\bigr)^{-1}.
\]
EGA II, 6.6.1, makes this sheaf ample over $T$
\cite{EGA}. Since $P_\delta$ is proper over $T$, it is projective.
\end{proof}

The torsor $P_\delta$ in this lemma need not be trivial.
Its finite flat symmetric-representative cover is essential to the
argument, and need not be \'etale. Likewise, an arbitrary schematic
abelian torsor over a normal semilocal base need not be projective:
\cite[Example~6.1, p.~1574]{BhattAnnihilation} gives a Zariski
locally trivial torsor of infinite order which cannot be split by
a finite cover. A projective smooth torsor over that semilocal base
would have a finite locally free multisection by
Lemma~\ref{lem:semilocal-multisection}. We next obtain the
additional projectivity needed here from the original scheme
projective bundle.

\subsection{Uniform boundaries over a finite normalization}

\begin{proposition}\label{prop:torsor-projective-after-normalization}
Let $S$ be the spectrum of an excellent noetherian local ring,
let $A/S$ be an abelian scheme, and let $Q$ be an fppf $A$-torsor
algebraic space. Suppose that $E$ is finite locally free of
everywhere positive rank and $X=\PP_Q(E)$ is a scheme. If
$W\to S$ is the normalization of $S_{\mathrm{red}}$, then
$Q_W$ is projective over $W$.
\end{proposition}

\begin{proof}
Let $s$ be the closed point of $S$. The morphism $W\to S$ is
finite and surjective by excellence. The scheme $W$ is normal,
noetherian, affine and semilocal, and all its closed points lie
over $s$. Its finitely many integral components are open and
closed. We do not assume that $W\to S$ is flat.
The dual $B=A^\vee$ is an abelian $S$-scheme and commutes with
base change \cite[\S3.1]{PR}\cite[Corollary~4.6]{LaurentSchroer}.
Lemma~\ref{lem:dual-projective-finite-normal-cover} makes $B_W$
projective over $W$.

The rank $r$ of $E$ is constant. Its locally constant rank
function is constant on each connected geometric fibre of
$Q\to S$, hence descends to the connected local base.
If $r=1$, then $X=Q$ is a scheme. An affine open in $Q$ meeting
the closed fibre is saturated by the $A$-action, so
\cite[6.6, Theorem~2(c)]{BLR} makes $Q$ AF. This property passes
through the finite map $Q_W\to Q$. One affine open in $Q_W$
therefore meets all closed $W$-fibres, and part (d) of the same
theorem makes $Q_W$ projective, exactly as in
Lemma~\ref{lem:dual-projective-finite-normal-cover}.
Assume henceforth that $r\geq2$.

\emph{Global boundaries.}
Let $U\subset X$ be an affine open meeting $X_s$. The map
$U\to S$ is smooth and open, so is surjective. Its finite
pullback $U_W$ is affine and meets every fibre over $W$.
Those fibres are geometrically integral, so $U_W$ is dense
in every fibre, including after a residue-field extension.
The scheme $X_W$ is smooth over normal $W$, hence normal;
on every integral component of $W$ it is integral, by the
generic-fibre argument in
Lemma~\ref{lem:normal-affine-relative-boundary}.

The inclusion $U_W\hookrightarrow X_W$ is affine because
$X_W$ is separated over affine $W$. EGA IV$_4$, 21.12.7,
therefore makes its boundary a finite union of prime divisors.
Give each component multiplicity one. The support of this
effective cycle contains no irreducible component of any
fibre, so EGA IV$_4$, 21.14.3(i), makes it locally principal
\cite{EGA}. Normality makes the local equations regular
non-zero-divisors, and equations with the same Weil divisor
differ by a unit. We obtain a global effective Cartier divisor
$D_U$ on $X_W$ with support $X_W\setminus U_W$.
It remains effective Cartier on every geometric fibre,
because the section is nonzero on the nonempty open
$(U_W)_{\bar w}$ of the integral fibre.

On each open-and-closed component of $W$, the chartwise
argument proving \eqref{eq:torsor-bundle-picard-decomposition} gives
\[
 \OO_{X_W}(D_U)\simeq
 \OO_{X_W}(m_U)\otimes\pi_W^*M_U.
\]
Here $M_U$ is a global line bundle on $Q_W$, and $m_U$ is a
locally constant integer function on $W$. The componentwise
line bundles glue over the disjoint union. Every value of
$m_U$ is nonnegative: on a projective-space fibre containing
a point of $U_W$, the canonical section restricts to a nonzero
section of $\OO(m_U)$. These boundaries and their factors
$M_U$ are defined over the whole of $W$.

\emph{Finite common kernels on every closed fibre.}
For every $M\in\Pic(Q_W)$ let
$\phi_M:A_W\to B_W$ be its translation homomorphism.
It is defined after a local choice of torsor point and is
independent of that choice \cite[\S3.6]{PR}.
For a closed $w\in W$, choose
$k=\overline{\kappa(w)}$ and put
\[
 K_w=\bigcap_{M\in\Pic(Q_W)}\ker(\phi_{M,k})\subset A_k.
\]
This is a finite intersection, by the noetherian ideal argument
in the proof of
Proposition~\ref{prop:normal-local-torsor-projective-bundle}.
We claim it is finite.

If not, its reduced identity component is a positive-dimensional
abelian subvariety. Choose a translated integral curve
$C\subset Q_k$ in that component, as in that proposition.
Every global line bundle on $Q_W$ has degree zero on $C$,
including $\det E_W$. Set
\[
 Y=\PP_C(E_W|_C),\qquad \xi=c_1(\OO_Y(1)).
\]
Then $Y$ is integral projective of dimension $r$, is closed
in $X_k$, and the Chern relation gives
\[
 \deg(\xi^r\cap[Y])=0.
\]

For completeness, the boundaries just constructed give the
contradictory positivity uniformly over $W$.
For a positive-dimensional integral $Z\subset Y$, choose
an affine open $U$ in the original scheme $X$ containing the
image of a point of $Z$. That image lies over $s$, so $U$
belongs to the preceding construction. The intersection
$(U_W)_k\cap Z$ is nonempty and affine, but cannot equal the
positive-dimensional projective $Z$.
Thus $D_U$ restricts to a nonempty effective Cartier divisor
$D_Z$ on $Z$. If $d=\dim Z$ and $m=m_U(w)\geq0$, projection
formula gives
\begin{equation}\label{eq:uniform-boundary-intersection}
 m\deg(\xi^d\cap[Z])=
 \deg(\xi^{d-1}\cap[D_Z]).
\end{equation}
Indeed the omitted term is
\[
 \deg\bigl(c_1(M_U|_C)\cap
       \pi_{C*}(\xi^{d-1}\cap[Z])\bigr)=0:
\]
the pushforward is a one-cycle on the integral curve $C$,
and $M_U|_C$ has degree zero.

For $d=1$, the right side of
\eqref{eq:uniform-boundary-intersection} is strictly positive.
For $d>1$, write
$[D_Z]=\sum_j a_j[Z_j]$ with $a_j>0$ and $\dim Z_j=d-1$.
Induction makes every degree
$\deg(\xi^{d-1}\cap[Z_j])$ positive. In either case
\eqref{eq:uniform-boundary-intersection}, together with $m\geq0$,
gives $\deg(\xi^d\cap[Z])>0$.
Taking $Z=Y$ contradicts its zero top intersection.
This proves that every $K_w$ is finite.

\emph{A finite map to projective Picard-class torsors.}
There are finitely many closed points in $W$. At each one,
select finitely many global line bundles defining the
finite intersection $K_w$, and take the union of these finite
selections. Denote the resulting bundles by $M_1,\ldots,M_n$.
Their common kernel is finite on every closed geometric fibre.

For each $M_i$, the action defines
\[
 f_i:Q_W\longrightarrow\Pic_{A_W/W},
 \qquad q\longmapsto[\tau_q^*M_i].
\]
Its class modulo $B_W$ is a single section $\delta_i$ over $W$.
Indeed local choices of $q$ differ by translation, and the
corresponding Picard classes differ by an element of $B_W$.
They therefore agree in the quotient and descend.
The associated homomorphism is $\phi_{M_i}$.
Thus $f_i$ factors through the torsor $P_{\delta_i}$.
Lemma~\ref{lem:projective-picard-class-torsor} makes each of
these torsors projective over $W$.

Consider
\[
 f=(f_1,\ldots,f_n):Q_W\longrightarrow
       \prod_{i=1}^n P_{\delta_i}.
\]
It is proper because its target is separated over $W$ and
$Q_W$ is proper over $W$. On each closed geometric fibre it is,
after a torsor-point choice, a translate of the product
homomorphism $(\phi_{M_i})_i$. Its nonempty fibres are
translates of the finite common kernel.
The non-quasi-finite locus of $f$ is closed in $Q_W$, and its
proper image in $W$ contains no closed point. It is therefore
empty. Hence $f$ is proper and quasi-finite, and so finite
\cite[Tag~0A4X]{Stacks}.
A finite algebraic space over its scheme target is a scheme.
The finite pullback of a $W$-ample invertible sheaf on that
target is $W$-ample on $Q_W$. Thus $Q_W$ is projective.

If $A$ has relative dimension zero, the common kernel is
already finite without choosing line bundles; take $n=0$
and the target to be $W$. The same finiteness argument applies.
\end{proof}

The affine opens in this proof come from $X$ over the original
local base. Their finite pullbacks meet every closed fibre of
$W$ at once. This is what makes their line bundles global over
the semilocal normalization and permits one finite family of
translation homomorphisms.

\subsection{Descent and approximation over arbitrary bases}

\begin{proof}[Proof of Theorem~\ref{thm:torsor-projective-bundle}]
First let $S$ be excellent noetherian local, with closed point
$s$, and suppose $X=\PP_Q(E)$ is a scheme. Let $W\to S$ be the
finite normalization of $S_{\mathrm{red}}$.
Proposition~\ref{prop:torsor-projective-after-normalization}
makes $Q_W$ projective over the affine $W$, so every finite
set of its points has a common affine open neighborhood.

Choose an affine open $U\subset X$ meeting $X_s$.
The smooth map $U\to S$ is surjective, hence is an fpqc cover.
The morphism $U\to X\to Q$ gives a section of $Q_U$, so
$Q_U\simeq A_U$ is a locally noetherian scheme.
Apply \cite[Theorem~4.2]{RaynaudDescent} to the fppf sheaf $Q$,
using $U\to S$ as the fpqc cover and $W\to S$ as the separate
finite surjective cover. All finite orbits associated to
$W\times_S W$ lie in affine opens of $Q_W$. It follows that
$Q$ is a scheme. The normalization was not assumed flat;
the original nilpotent structure of $S$ is present in
$U\to S$.

Now let $S$ be excellent and noetherian. For $s\in S$, the
local result makes $Q\times_S\Spec\OO_{S,s}$ a scheme.
We spread a finite affine open cover of it to a neighborhood
of $s$. Choose an affine neighborhood $S_0$.
The finitely many affine charts and their quasi-compact open
immersions into the localized $Q$ are of finite presentation.
EGA IV$_3$, 8.8.2 and 8.10.5(iii), together with
\cite[Tag~07SK]{Stacks}, descend them to open immersions over
an affine neighborhood $T\subset S_0$ of $s$.
For an algebraic-space target the open-immersion assertion is
checked on a finite affine \'etale presentation and descends
from those charts. The complement of their finite union is
closed in $Q_T$ and has empty fibre over $s$. Its image in
$T$ is closed by properness. Removing that image gives a
neighborhood on which the charts cover, so $Q$ is a scheme
there. These neighborhoods cover $S$.

Finally take an arbitrary affine base $S=\Spec R$ and write
$R=\varinjlim R_i$ as the filtered union of its finitely
generated $\Z$-subalgebras. Each $R_i$ is excellent and
noetherian. We descend all the finite-presentation data needed
for a model of the hypothesis.
Descend the schemes $A$ and $X$, the algebraic space $Q$,
the group and action maps, and the isomorphism
\[
 A\times_R Q\longrightarrow Q\times_R Q,\qquad
 (a,q)\longmapsto(a+q,q),
\]
with its inverse and both composition identities.
Choose a finite affine \'etale presentation $V\to Q$
on which $E$ is trivial on each chart. Descend that
presentation, its overlaps, the vector-bundle transition
matrices and their inverse and cocycle identities.
This produces a finite locally free $E_i$ of positive ranks.
Also descend
\[
 X\simeq\PP_Q(E),
\]
its inverse and its two identity equations.
Thus $X_i$ is a scheme and is isomorphic to
$\PP_{Q_i}(E_i)$.
These are the scheme and algebraic-space limit theorems
EGA IV$_3$, \S\S8.8--8.10, and \cite[Tag~07SK]{Stacks}.

After increasing $i$, properness of $A_i$ follows from
EGA IV$_3$, 8.10.5(xii), and smoothness by EGA IV$_4$,
17.7.8(ii). The group identities hold at the finite stage.
The geometric-integrality failure locus of its fibres is
constructible by EGA IV$_3$, 9.7.7.
Its inverse image in $\Spec R$ is empty and the property
is preserved and reflected by residue-field extension.
Hence that constructible locus disappears at some later
stage. To see this explicitly, decompose it into finitely
many locally closed affine pieces; emptiness after the
filtered colimit means that $1$ belongs to their defining
ideals after localization, a finite relation already true
at a finite stage. Thus $A_i$ is an abelian scheme.

We also arrange \'etaleness and surjectivity of the presentation, and
smoothness and surjectivity of $V_i\to\Spec R_i$, using
EGA IV$_4$, 17.7.8(ii), and IV$_3$, 8.10.5(vi).
The map $V_i\to Q_i$ supplies local sections. Together with
the descended simple-transitivity isomorphism, it makes
$Q_i$ an fppf $A_i$-torsor.
Its smoothness and properness follow by fppf descent from
$A_i$. The construction therefore gives an fppf torsor,
a vector bundle of positive rank, and a scheme projective
bundle over the excellent noetherian model.
The case already proved makes $Q_i$ a scheme, and base
change makes $Q$ a scheme. Covering an arbitrary $S$ by
affines proves the theorem.
\end{proof}

The theorem excludes the projective-bundle construction over
nonschematic abelian torsors for all bases. Schematicity of the
Picard functor of a general proper flat family with geometrically
normal or geometrically integral fibres remains a separate
question: such a family need not have the torsor action and
translation homomorphisms used above.


\section{A generated ample bundle on a Raynaud surface}

We construct the surface and compute the cohomology that the later root cover will preserve. All objects in this section are first defined over $\F_2$.

\subsection{The curve and a ruled surface}

Let $C$ be the smooth projective model of
\[
 y^2+y=x^{15}.
\]
It is geometrically integral: a pole of $g^2+g$ has even order, whereas $x^{15}$ has an odd pole at infinity. The Artin--Schreier cover of the $x$-line has one point $\infty$ above infinity, with
\[
 v_\infty(x)=-2,\qquad v_\infty(y)=-15.
\]
Write $P_0=(0,0)$, $P_1=(0,1)$, and $t=y/x^8$. Directly from the equation,
\begin{equation}\label{eq:curve}
 \operatorname{div}(t)=7P_0+\infty-8P_1,
 \qquad x+t^{-2}=t^{-1}x^{-7},\qquad dt=x^6dx.
\end{equation}
The affine curve is smooth and $dx$ is nowhere zero there. Since $t$ is a uniformizer at infinity, the last identity gives $\operatorname{div}(dx)=12\infty$. Thus $g(C)=7$.

Put $N=\OO_C(2\infty)$. On $U=C\setminus\{\infty\}$ and $V=C\setminus\{P_0,P_1\}$ define a rank-two vector bundle $E$ by frames related on the overlap by
\[
 v_0=e_0,\qquad v_1=t^{-6}(e_1-t^{-1}e_0).
\]
This gives an exact sequence $0\to\OO_C\to E\to N^3\to0$ and $\det E=N^3$. Let $\pi:P=\PP_C(E)\to C$ use the quotient convention, so $\pi_*\OO_P(j)=\Sym^jE$ for $j\geq0$. The section $E_\infty$ defined by $e_0=0$ has $\OO_P(E_\infty)=\OO_P(1)$.

The quadratic expression $e_1^2+xe_0^2$ defines a curve $C''$ in $P$. In the other frame it becomes
\[
 v_1^2+t^{-12}(x+t^{-2})v_0^2.
\]
Consequently
\[
 \OO_P(C'')=\OO_P(2)\otimes\pi^*N^{-6}.
\]
The curves $E_\infty$ and $C''$ are disjoint: when $e_0=0$ the quadratic expression is $e_1^2\ne0$. Away from infinity, in the coordinate $w=e_1/e_0$, the equation is $w^2+x=0$, whose differential $dx$ never vanishes. At infinity the coefficient $t^{-12}(x+t^{-2})$ has valuation one by \eqref{eq:curve}; its differential is a unit. Thus $C''$ is smooth as well.

\subsection{The cyclic cover and direct images}

Set $M=\OO_P(-1)\otimes\pi^*N^2$. The equality
\[
 M^{-3}=\OO_P(E_\infty+C'')
\]
defines a degree-three cyclic cover
\[
 \psi:S=\Spec_P(\OO_P\oplus M\oplus M^2)\longrightarrow P,
\]
where the algebra multiplication uses the section of $M^{-3}$ vanishing on $E_\infty+C''$. The branch divisors are smooth and disjoint, and three is invertible; locally the cover is $z^3=u$ along a branch parameter $u$, and is \'etale elsewhere. Hence $S$ is smooth. The finite map is projective. Since $M$ and $M^2$ have negative degree on the fibres of $P\to C$, its global functions after any algebraic extension are the constants. Smoothness and connectedness therefore imply geometric integrality.

Write $f=\pi\psi$ and let $T$ be the reduced inverse image of $E_\infty$. Then $\psi^*E_\infty=3T$, and $T\simeq C$. On a dense affine chart the equations are
\[
 y^2+y=x^{15},\qquad z^3=w^2+x.
\]
The line bundle to be used is
\[
 Q=\OO_S(2T)\otimes f^*N.
\]
For every integer $a$, the local valuation calculation at $T$ gives
\begin{equation}\label{eq:floor}
 \psi_*\OO_S(aT)
 =\bigoplus_{j=0}^2M^j\bigl(\lfloor(a+j)/3\rfloor E_\infty\bigr).
\end{equation}
Indeed, the eigenfunction $z^j$ has valuation $j$ at $T$, and a pole of order $h$ along $E_\infty$ pulls back to a pole of order $3h$. The largest allowed $h$ is $\lfloor(a+j)/3\rfloor$. This also proves the formula for negative $a$.

Applying \eqref{eq:floor} with $a=2$ gives $f_*Q=N\oplus N^3$. These line bundles have bases $(1,x)$ and $(1,x,x^2,x^3)$ and are globally generated. The first summand generates $Q$ away from $T$. Along $T$ the second has local valuation $1-3+2=0$, so it generates there. Therefore $Q$ is globally generated and
\begin{equation}\label{eq:qsections}
 h^0(S,Q)=6.
\end{equation}
The corresponding sections on the displayed affine chart are $1,x,z,zx,zx^2,zx^3$.

\subsection{The negative powers}

We record all the negative cohomology needed later in one calculation. For $n\geq1$, the $j$th summand of $\psi_*Q^{-n}$ is
\[
 \OO_P(a_j)\otimes\pi^*N^{-n+2j},
 \qquad a_j=\left\lfloor\frac{-2n+j}{3}\right\rfloor-j.
\]
Each $a_j\leq-1$, so all zeroth direct images under $\pi$ vanish. Relative duality gives
\[
 R^1\pi_*\OO_P(a)=\bigl(\Sym^{-a-2}E\bigr)^\vee\otimes N^{-3}
 \quad(a\leq-2),
\]
whereas $\OO_P(-1)$ has no direct images. For $n=1$ this yields
\[
 H^1(S,Q^{-1})=H^0(C,N^{-2}\oplus\OO_C)=\F_2.
\]
For $n\geq2$, dualizing the filtration of $\Sym^{-a_j-2}E$ with line quotients $N^{3h}$ gives a filtration of every nonzero first direct-image summand with quotients
\[
 N^{-n+2j-3-3h},\qquad h\geq0.
\]
Every exponent is at most $1-n<0$. Such a line bundle has no section, and the same is true of a successive extension of them. Leray therefore proves
\begin{equation}\label{eq:negative}
 h^1(S,Q^{-1})=1,\qquad H^1(S,Q^{-n})=0\quad(n\geq2).
\end{equation}
For $\OO_S$ the only nonzero first direct-image summand is $N$, while $f_*\OO_S=\OO_C$. Hence
\begin{equation}\label{eq:irregularity}
 h^1(S,\OO_S)=g(C)+h^0(C,N)=7+2=9.
\end{equation}
This calculation uses the relative projective-line formulas and line-bundle cohomology on $C$; no numerical computation is needed for its validity.

\section{A very ample ninth power}

We next prove that $Q^9$ embeds $S$ and determine its number of sections. Put
\[
 B=\OO_P(6)\otimes\pi^*N^9.
\]
Then $Q^9=\psi^*B$, and its three eigensummands on $P$ are
\[
 BM^i=\OO_P(6-i)\otimes\pi^*N^{9+2i}\quad(0\leq i\leq2).
\]
The vector bundle $\Sym^aE\otimes N^b$ has a filtration with line quotients $N^{b+3j}$ for $0\leq j\leq a$. For the pairs
\[
 (a,b)=(6,9),(5,11),(4,13),
\]
these quotients still have degree at least sixteen after subtracting an arbitrary effective degree-two divisor on $C$. This exceeds $2g(C)-2=12$. Their first cohomology vanishes by Serre duality, and induction on the filtration gives the same vanishing for the vector bundles. Thus their sections restrict surjectively to every base subscheme of length at most two.

It follows that $B$ is very ample on $P$. A length-two subscheme whose image on $C$ has length two is separated by the preceding restriction surjectivity and the relative complete linear series of $\OO_P(6)$. A length-two subscheme contained in a fibre is separated by $\OO_{\PP^1}(6)$, using surjectivity over that base point. The same argument with length one proves that $BM$ is globally generated.

These two facts prove that $\psi^*B$ is very ample. Indeed, $B$ recovers the embedded base $P$. Near each point choose a nonvanishing section of $BM$. In the cyclic-cover algebra it is a generator of the $M$ summand, up to a unit. Its image as a section of $\psi^*B$ therefore supplies the missing algebra generator of the finite cover. On the affine charts of the embedding by $B$, ratios of these sections generate the coordinate algebra of $S$. The projective map is consequently a closed immersion. In particular $Q$ is ample, since its ninth power is very ample.

Riemann--Roch on $C$ and the same vanishing give
\begin{align*}
 h^0(P,B)&=\sum_{j=0}^6(18+6j-6)=210,\\
 h^0(P,BM)&=\sum_{j=0}^5(22+6j-6)=186,\\
 h^0(P,BM^2)&=\sum_{j=0}^4(26+6j-6)=160.
\end{align*}
Therefore
\begin{equation}\label{eq:556}
 h^0(S,Q^9)=556.
\end{equation}
An independent principal-parts computation on the two affine charts of $C$ gives the same three dimensions; the program \emph{Independent principal-parts verifier} and its exact output are available through the workbench \cite{Workbench}.

\section{Taking roots of a basis}

This section turns the generated bundle $Q$ into a very ample bundle without losing its negative first cohomology. Choose a basis $v_1,\ldots,v_{556}$ of $H^0(S,Q^9)$ over $\F_2$. Let $k$ be an algebraic closure of
\[
 \F_2(\lambda_{ij}:1\leq i,j\leq556),
\]
where the $\lambda_{ij}$ are algebraically independent, and put $s_i=\sum_j\lambda_{ij}v_j$ after base change. The determinant is nonzero, so the $s_i$ are a basis. Bertini's theorem for the embedding by $Q^9$ shows that their divisors are smooth, meet pairwise transversally, and have no triple intersection. To justify the last condition, choose the divisors successively: two of them have a finite intersection and a general subsequent section avoids those finitely many points. Each of the finitely many required conditions is open and nonempty in the parameter space, hence holds at its generic point.

Take all ninth-root covers simultaneously:
\[
 \rho:Y=\Spec_S\left(\bigotimes_{i=1}^{556}
      \bigoplus_{r=0}^8Q^{-r}\right)\longrightarrow S,
\]
with multiplication in the $i$th factor determined by $s_i$. This is a finite flat projective map of degree $9^{556}$. Set $L=\rho^*Q$. Its algebra has the direct-sum decomposition
\begin{equation}\label{eq:rootalgebra}
 \rho_*\OO_Y=\bigoplus_{0\leq r_i\leq8}Q^{-\sum_i r_i}.
\end{equation}
At a point of $S$ at most two branch divisors pass. Their local equations, when present, are part of a regular parameter system. Adjoining their ninth roots again gives regular parameters; adjoining roots of the remaining units is \'etale because nine is invertible. Thus $Y$ is smooth. Since $Q$ is ample, only the summand with every $r_i=0$ has global sections in \eqref{eq:rootalgebra}. Consequently $H^0(Y,\OO_Y)=k$. The proper smooth scheme $Y$ is connected, and its irreducible components are disjoint; hence it is integral. As $k$ is algebraically closed it is geometrically integral.

For each $i$ there is a tautological section $t_i$ of $L$ with $t_i^9=\rho^*s_i$. Together with the six pulled-back sections of $Q$, these sections define a projective map. Let $u$ be any of the six sections, and let $S_u$ denote its nonvanishing locus. The ninth power $u^9$ is a section of the very ample bundle $Q^9$, so $S_u$ is affine. Because the $s_i$ form a basis of that complete linear series, the functions $s_i/u^9$ generate its coordinate ring. The algebra of $Y$ over this chart is therefore generated over $k$ by
\[
 z_i=t_i/u,\qquad z_i^9=s_i/u^9.
\]
These are ratios of the proposed sections of $L$. The corresponding affine chart map is a closed immersion. The charts cover $Y$ because $Q$ is globally generated. Thus the projective map is an immersion, and properness makes it a closed immersion. This proves that $L$ is very ample. Keeping the pulled-back sections of $Q$ is essential in this argument: they provide the affine charts and distinguish simultaneous scalar changes of all roots.

Finally apply \eqref{eq:negative}, \eqref{eq:irregularity}, and \eqref{eq:rootalgebra}. For $L^{-1}$, only the multi-index of total sum zero contributes to first cohomology, so $h^1(Y,L^{-1})=1$. For $L$, the zero multi-index contributes six sections and each of the 556 multi-indices of sum one contributes one; all larger sums are negative powers of an ample bundle and have no section. Hence $h^0(Y,L)=562$. For $\OO_Y$, the zero multi-index contributes nine to first cohomology and those of sum one contribute 556, while all larger sums vanish. Thus $h^1(Y,\OO_Y)=565$.

If $H\in|L|$ is smooth and connected, the sequence
\[
 0\longrightarrow L^{-1}\longrightarrow\OO_Y
   \longrightarrow\OO_H\longrightarrow0
\]
and $H^0(Y,\OO_Y)\simeq H^0(H,\OO_H)=k$ identify the kernel of $H^1(Y,\OO_Y)\to H^1(H,\OO_H)$ with $H^1(Y,L^{-1})$. The tangent space of the Picard scheme is $H^1(\OO)$ and restriction induces this map. This proves the final assertion of Theorem~\ref{thm:surface}.

\section{A smoothness obstruction for finite flat torsors}

The infinitesimal Picard kernel in Theorem~\ref{thm:surface} does not by itself give a nontrivial point of that kernel. The following obstruction explains why one natural finite flat torsor construction cannot bridge the difference while preserving smoothness.

\begin{proposition}\label{prop:tangent-ample}
Let $Z$ be a smooth projective surface over an algebraically closed field and let $D\subset Z$ be a smooth ample integral curve. There is no line subbundle
\[
 \OO_Z(-D)\longrightarrow T_Z
\]
whose restriction to $D$ is tangent to $D$. Here a subbundle has locally free quotient. No hypothesis of $p$-closedness is imposed.
\end{proposition}

\begin{proof}
Write $\mathcal F$ for the proposed subbundle. Its restriction maps into $T_D$ and is nonzero on every fibre. Thus
\[
 \OO_D(-D)\simeq T_D,\qquad D^2=2g(D)-2,\qquad K_Z\cdot D=0.
\]
The quotient of $T_Z$ by $\mathcal F$ has first Chern class $D-K_Z$. Multiplicativity of Chern classes gives
\begin{equation}\label{eq:tangent-chern}
 c_2(Z)=(-D)\cdot(D-K_Z)=-D^2<0.
\end{equation}

In characteristic zero one can already contradict ampleness on the curve. A rank-one distribution is involutive, and its bracket induces a regular connection on $N_{D/Z}$: for a tangent field $v$ and a normal class represented by $w$, put $\nabla_v\overline w=\overline{[v,w]}$. Tangency makes this independent of the lifts, and the bracket identities give the connection rules. A line bundle with a regular connection on a smooth projective curve has degree zero in characteristic zero. Indeed, in a meromorphic frame the connection is a meromorphic differential whose residue at each point is the order of that frame there; the residue theorem says that their sum, the degree, is zero. This contradicts $\deg N_{D/Z}=D^2>0$.

Suppose now that the characteristic is positive. If a positive multiple of $K_Z$ has a nonzero section, its effective divisor has zero intersection with the ample divisor $D$. It is therefore zero. The canonical bundle is torsion, $Z$ is minimal, and its Kodaira dimension is zero. The classification gives $c_2(Z)\in\{0,12,24\}$, including the nonclassical cases in characteristics two and three; see \cite[\S7, author pp.~26--30]{Liedtke}. This contradicts \eqref{eq:tangent-chern}.

Consequently the Kodaira dimension is $-\infty$. The surface is birationally ruled by \cite[Theorems~4.4--4.5, author p.~21]{Liedtke}. Rational surfaces have positive second Chern number, so a minimal model is a ruled surface $Z_0\to C$ of genus $g\geq2$. Let $r$ be the number of point blowups from $Z_0$ to $Z$. Choose a section class $C_0$ with $C_0^2=-e$, a fibre class $F$, and total-transform exceptional classes $E_1,\ldots,E_r$. Pulling back the first two classes to $Z$, write
\[
 D=aC_0+bF-\sum_i m_iE_i,\qquad
 K_Z=-2C_0+(2g-2-e)F+\sum_iE_i.
\]
The intersection form has $C_0\cdot F=1$, $F^2=0$, and $E_i\cdot E_j=-\delta_{ij}$; each $E_i$ is orthogonal to the pullback of the minimal model. Moreover
\begin{equation}\label{eq:exceptional-bounds}
 a\geq1,\qquad 0<m_i<a.
\end{equation}
To check the strict upper bound even at infinitely near centres, let $F_c$ be the full fibre containing the centre that creates $E_i$. At that blowup its total transform contains the exceptional divisor with positive multiplicity. Subtracting one copy leaves a nonzero effective divisor, and subsequent pullbacks preserve this property. Thus both $E_i$ and $F_c-E_i$ are nonzero effective divisors. Their intersections with $D$ are respectively $m_i$ and $a-m_i$, and ampleness proves \eqref{eq:exceptional-bounds}.

Now $K_Z\cdot D=0$ and the intersection form give
\[
 2b=a(e+2g-2)+\sum_i m_i,\qquad
 D^2=2a^2(g-1)+\sum_i m_i(a-m_i).
\]
On the other hand $c_2(Z)=4(1-g)+r$. Equation~\eqref{eq:tangent-chern} therefore becomes
\[
 (2a^2-4)(g-1)+\sum_i\bigl(m_i(a-m_i)+1\bigr)=0.
\]
For $a\geq2$ its left side is strictly positive. For $a=1$, the strict bounds on the integers $m_i$ force $r=0$, and the equation forces $g=1$. Both alternatives contradict the preceding conclusions.
\end{proof}

Regular foliations in positive characteristic can behave differently from their characteristic-zero counterparts. Fassarella, Mendson, dos Santos and Touzet classify them on minimal rational and del Pezzo surfaces and prove an arbitrary-characteristic Camacho--Sad formula \cite[Theorems~1.1--1.2 and 2.8]{FMST}. Their results provide the surrounding foliation theory. Proposition~\ref{prop:tangent-ample} uses the particular tangent line bundle and also excludes the ruled surfaces of higher genus forced by its negative Chern number.

\begin{corollary}\label{cor:smooth-torsor}
Let $X$ be a smooth projective surface over an algebraically closed field of characteristic $p>0$, and let $Y\subset X$ be a smooth ample integral curve. Put $L=\OO_X(Y)$ and let $s$ be its canonical section. Define the finite flat group scheme $G_s$ locally by
\[
 \OO_X[u]/(u^p),\qquad
 \Delta(u)=u\otimes1+1\otimes u+s\,u\otimes u.
\]
If $s_j=g_{ij}^{-1}s_i$ in two trivializations of $L$, use $u_j=g_{ij}u_i$ to glue these algebras. There is no $G_s$-torsor $\pi:Z\to X$ with $Z$ smooth and the full inverse image $D=\pi^*Y$ smooth and integral.
\end{corollary}

\begin{proof}
The formula defines a rank-$p$ group scheme with $\operatorname{Lie}(G_s)=L^{-1}$. The homomorphism to $\mu_p$ given by $1+su$ is an isomorphism off $Y$; over $Y$ the group is the additive height-one group with the indicated line twist.

A torsor is finite flat over $X$, so $Z$ is projective and $D$ is ample. Differentiating the action gives
\[
 \pi^*\operatorname{Lie}(G_s)=\OO_Z(-D)\longrightarrow T_Z.
\]
At every geometric point the orbit map identifies the fibre group with the fibre of the torsor through that point, a closed subscheme of $Z$. Its tangent map is injective. Hence the displayed map has rank one on every fibre and has locally free quotient. Its derivations annihilate pullbacks of functions from $X$, so the image is tangent to $D$. Proposition~\ref{prop:tangent-ample} is the required contradiction.
\end{proof}

Resolving a singular torsor does not preserve these hypotheses: the pullback of an ample curve can become only nef and big, and the tangent subsheaf can acquire zeros. The corollary therefore leaves the general kernel on geometric Picard points undecided.

\section{Restrictions on Picard kernels}
\label{sec:picard-kernel-constraints}
We examine two ways in which a torsion divisor class might disappear on a
smooth ample curve. A purely inseparable multisection of a quasi-elliptic
fibration has the desired effect on torsion, but the intersection calculation
below prevents its ampleness over a base of positive genus. For an arbitrary
order-$p$ class, normalization of its torsor gives a different constraint:
the lifted divisor is reduced, and the conductor along it is determined
exactly. In the equality case, a group action extends to the normalized
cover, and a zero scheme on the smooth base counts its isolated
singularities. These constraints leave the general point-kernel
question open.

For the direct route to the genus results, first read the normalized
torsor calculation in Proposition~\ref{prop:normalized-mup-conductor}.
The adjoint argument is in Section~\ref{subsec:canonical-subsystem-strictness},
and the very ample equality structure is in
Section~\ref{subsec:conductor-equality-structure}; these prove
Theorems~\ref{thm:strict-picard-all} and
\ref{thm:conductor-equality-structure}.
Section~\ref{subsec:positive-square-picard} gives a self-contained
route to the quantitative positive-square bound of
Theorem~\ref{thm:positive-excess-picard}.
The intervening ruled, K3 and rational-cover arguments treat the
merely ample cases, whose equality configurations have separate
hypotheses and remain relevant without very ampleness.

\begin{proposition}\label{prop:quasielliptic-multisection}
Let $k$ be algebraically closed of characteristic $p=2$ or $3$, and let
$f:S\to C$ be a quasi-elliptic fibration from a smooth projective integral
surface to a smooth projective integral curve. Thus $f_*\OO_S=\OO_C$ and
the generic fibre is a regular geometrically integral curve of arithmetic
genus one that is not smooth. Let $R\subset S$ be a smooth integral
horizontal curve for which $R\to C$ is purely inseparable of degree
$q=p^e\geq2$. If $g=g(C)$, then
\[
 R^2\leq\left(2-\frac{5q}{3}\right)(g-1).
\]
In particular, $R$ is not ample when $g\geq1$. The fibration need not be
relatively minimal.
\end{proposition}
\begin{proof}
First $g(R)=g$. Indeed, put $L=k(R)$ and $K=k(C)$. A one-variable
function field over a perfect field satisfies $[L:L^{p^e}]=p^e$.
The purely inseparable extension $L/K$ has exponent at most $e$, so
$L^{p^e}\subset K$; equality of degrees gives $K=L^{p^e}$.
The map of smooth curves is therefore an iterated relative Frobenius
followed by an isomorphism, and Frobenius twisting preserves genus.

Contract the vertical exceptional curves to obtain a relatively minimal
smooth model $f_0:S_0\to C$, with contraction $\beta:S\to S_0$.
Write $\mathcal E$ for the line-bundle quotient of
$R^1(f_0)_*\OO_{S_0}$ by its torsion subsheaf $\mathcal T$, and put
$\mathcal L=\mathcal E^\vee$ and $t=\operatorname{length}(\mathcal T)$.
The canonical-bundle formula, including wild multiple fibres, is
\[
 K_{S_0}\sim f_0^*(K_C+\mathcal L)+\sum_i a_iP_i,
 \qquad 0\leq a_i<m_i,\qquad f_0^*(c_i)=m_iP_i;
\]
see \cite[Theorem~5.6, author pp.~23--24]{Liedtke} and
\cite[Theorem~4.1.6, pp.~364--365]{CDL}.
Leray and Riemann--Roch on $C$ give
\[
 \chi(\OO_{S_0})
 =\chi(\OO_C)-\chi(\mathcal E)-t
 =-\deg\mathcal E-t,
 \qquad \deg\mathcal L=\chi(\OO_{S_0})+t.
\]
Every term in the canonical formula is numerically a rational multiple
of a fibre, so $K_{S_0}^2=0$.

The quasi-elliptic Euler formula is
\[
 c_2(S_0)=4(1-g)+\sum_{c\in C}\bigl(b_2((S_0)_c)-1\bigr)
 \geq4(1-g);
\]
here $b_2$ is the second $\ell$-adic Betti number for $\ell\ne p$,
and only finitely many summands are nonzero
\cite[Corollary~4.1.12 and the preceding quasi-elliptic discussion,
p.~369]{CDL}. Each fibre has at least one irreducible component,
so each summand is nonnegative. Noether's formula now yields
$\chi(\OO_{S_0})=c_2(S_0)/12\geq(1-g)/3$.

Let $R_0=\beta(R)$ with its reduced structure. It is horizontal of degree
$q$ over $C$. Its arithmetic genus need not equal $g$, so we apply
adjunction on $S$, where $R$ is smooth. The divisor
$K_S-\beta^*K_{S_0}$ is effective and exceptional, with no $R$ component.
Its intersection with $R$ is nonnegative. The effective vertical terms
in the canonical formula also have nonnegative intersection with $R_0$.
Consequently
\[
 K_S\cdot R\geq K_{S_0}\cdot R_0
 \geq q\bigl(2g-2+\chi(\OO_{S_0})+t\bigr).
\]
Adjunction gives
\[
 \begin{split}
 R^2
 &=2g-2-K_S\cdot R\\
 &\leq(1-q)(2g-2)-q\chi(\OO_{S_0})-qt\\
 &\leq\left(2-\frac{5q}{3}\right)(g-1).
 \end{split}
\]
This is nonpositive for $g\geq1$ and $q\geq2$, whereas an ample curve
on a projective surface has positive self-intersection.
\end{proof}

A nontrivial $p$-torsion line bundle on $C$ requires $g(C)\geq1$.
Pullback by an iterated relative Frobenius kills such a class, after the
corresponding Frobenius twist of the curve. The proposition explains why
this mechanism on a quasi-elliptic fibration cannot supply a smooth ample
multisection, even after point blowups. It does not exclude maps with an
additional separable part or curves on other kinds of surfaces.

\begin{proposition}\label{prop:normalized-mup-conductor}
Let $X$ be a smooth projective integral surface over an algebraically
closed field $k$ of characteristic $p>0$. Let $Y\subset X$ be a smooth
ample integral curve, and let $M$ be a nontrivial line bundle of order
$p$ with $M|_Y\simeq\OO_Y$. Choose a trivialization of $M^p$ and let
$T\to X$ be its $\mu_p$-torsor. Let $N$ be the normalization of $T$ and
$\pi:N\to X$ the resulting finite map.

The logarithmic differential of $M$ has order $pa$ along $Y$ for an
integer $a\geq1$, with nonzero tangential leading coefficient. The full
inverse image $D=\pi^*Y$ is a reduced integral ample Cartier divisor.
At its generic point the ramification index is one, the residue degree
is $p$, and the conductor of $T$ in $N$ has coefficient $a(p-1)$.
The normalization $\widetilde D$ of $D$ has genus $g(Y)$.

Let $\mathfrak c\subset\OO_N$ be the conductor ideal, put $c=a(p-1)$,
and set $\mathcal J=\mathfrak c\otimes\OO_N(cD)$.
Then $\mathcal J$ is a reflexive ideal in $\OO_N$, and its restriction
embeds in $\OO_D$ with finite-length quotient. Write $\ell$ for that
length and put $\delta=p_a(D)-g(Y)$. Without any smoothness or
Gorenstein assumption on $N$, one has
\[
 2g(Y)-2=ap\,Y^2+\frac{\ell+2\delta}{p-1},
 \qquad \ell\geq0,\quad\delta\geq0.
\]
If $2g(Y)-2=ap\,Y^2$, then $D$ is smooth, $N$ is regular near $D$,
$\mathcal J=\OO_N$, and $N$ is Gorenstein. If $N$ is singular at a
point of $D$, then $2g(Y)-2-apY^2\geq2/(p-1)$.
\end{proposition}
\begin{proof}
On a trivialization of $M$, its chosen $p$th-power trivialization is a
unit $u$, and the torsor has equation $z^p=u$. The forms $du/u$ agree
on overlaps because the transition factors are $p$th powers. They give
a global closed form $\omega$ on $X$. It is nonzero: if a rational
trivialization gave $du=0$, then $u$ would be a $p$th power in $k(X)$,
since $k$ is perfect. The corresponding rational section of $M$ would
have divisor zero, because its $p$th power does. Normality of $X$ would
then trivialize $M$. The same nonvanishing makes $z^p-u$ irreducible
over $k(X)$. Since the torsor algebra is locally free over $X$, it
embeds in its generic field and $T$ is integral.

The image of $\omega|_Y$ in $\Omega^1_{Y/k}$ is zero. To see this,
use a trivialization of $M|_Y$; its $p$th power differs from the chosen
one by a global unit on the proper integral curve $Y$, hence by an
element of $k^*$. The conormal sequence and
$H^0(Y,\OO_Y(-Y))=0$ now give $\omega|_Y=0$.
Let $m\geq1$ be its exact order along $Y$. The nonzero leading section
lies in $\Omega^1_{X/k}|_Y\otimes\OO_Y(-mY)$. Its tangential image
is nonzero, since the kernel on global sections is
$H^0(Y,\OO_Y(-(m+1)Y))=0$. Near the generic point, write the leading
term as $t^m(A\,dt+\eta)$, with $t=0$ defining $Y$ and $\eta$ tangential.
The coefficient of $t^{m-1}$ in $d\omega=0$ is
$m\,dt\wedge\eta$. Since $\eta$ is nonzero modulo $t$, we have
$p\mid m$. Thus $m=pa$.

Put $K=k(Y)$. Complete the local ring at its generic point, choosing a
coefficient field compatible with $k$, to obtain $R=K[[t]]$. In the
completed differential module, write
\[
 u=\sum_{i\geq0}u_it^i.
\]
The equality $\operatorname{ord}_t(du/u)=pa$ implies, for $i<pa$,
that $u_i=0$ if $p\nmid i$ and that $du_i=0$ if $p\mid i$.
The kernel of $d:K\to\Omega^1_{K/k}$ is $K^p$. Choose
$v_j^p=u_{pj}$ for $j<a$, and set $v(t)=\sum_{j<a}v_jt^j$. Then
\[
 u-v(t)^p=t^{pa}b(t),\qquad db(0)\ne0.
\]
In particular $b(0)\notin K^p$. The element
$w=(z-v(t))/t^a$ satisfies $w^p=b(t)$, and
\[
 B=R[w]/(w^p-b(t))
\]
is a one-dimensional local ring with maximal ideal $(t)$ and residue
field $K(b(0)^{1/p})$. It is therefore a discrete valuation ring, and
is the normalization of $A=R[z]/(z^p-u)$ in their common fraction field.
This proves ramification index one and residue degree $p$. Normalization
commutes with this completion by excellence, so these are the invariants
of $N$ along $D$ as well.

As $R$-modules the two rings are
\[
 B=\bigoplus_{i=0}^{p-1}Rw^i,
 \qquad A=\bigoplus_{i=0}^{p-1}Rt^{ai}w^i.
\]
The ideal $t^{a(p-1)}B$ is contained in $A$, hence in its conductor.
Conversely, let $c=\sum c_iw^i$ belong to the conductor. For each $i$,
the product $cw^{p-1-i}$ belongs to $A$, and the coefficient of
$w^{p-1}$ in that product is exactly $c_i$. No other exponent reduces
to $p-1$ using $w^p=b(t)$. Thus each $c_i$ is divisible by
$t^{a(p-1)}$, which proves that the conductor is exactly
$t^{a(p-1)}B$.

The extension of function fields is purely inseparable of degree $p$,
so $\pi$ is radicial. More explicitly, for an integral element $b$ of
the normalization over a normal affine open of $X$, its $p$th power
belongs to that open's fraction field and is integral over its ring;
normality places it in the ring. Thus every prime has a unique prime
above it. The local equation of $Y$ defines a Cartier divisor on $N$;
the calculation shows that this divisor has one component of generic
multiplicity one. Normality gives Serre's condition $S_2$, so its
Cartier quotient is $S_1$. An $S_1$ scheme that is generically reduced
is reduced. Hence the entire inverse image is the integral divisor
$D$. It is ample because it is the finite pullback of an ample divisor.
The smooth projective normalization $\widetilde D$ maps purely
inseparably with degree $p$ to $Y$, so the function-field argument in
the preceding proposition gives $g(\widetilde D)=g(Y)$.

We now compute the correction without resolving $N$. Let $\tau:T\to X$
and $\nu:N\to T$ be the two finite maps. Projection to degree zero
in the cyclic algebra $\mathcal A=\bigoplus_{i\in\mathbf Z/p}M^i$
gives a perfect multiplication pairing: each invertible summand pairs
with its opposite. Thus
$\operatorname{Hom}_{\OO_X}(\mathcal A,\OO_X)\simeq\mathcal A$
as $\mathcal A$-modules. Finite duality gives
$\omega_T\simeq\tau^*\omega_X$. This pairing is not the field trace,
which vanishes for the inseparable function-field extension.

The normal surface $N$ is Cohen--Macaulay. Finite duality for $\nu$
gives
\[
 \omega_N\simeq
 \operatorname{Hom}_{\OO_T}(\nu_*\OO_N,\omega_T)
 \simeq\mathfrak c\otimes\pi^*\omega_X;
\]
see \cite[Tags 0AX2 and 0FVZ]{Stacks}. The second identification
can be checked in a finite birational inclusion of domains $A\subset B$:
every $A$-linear map $B\to A$ is multiplication by an element of their
common fraction field, and its allowed multipliers are exactly the
conductor. The module $\omega_N$ is Cohen--Macaulay of rank one,
hence torsion-free and $S_2$. It is reflexive on the normal surface
\cite[Tag 0AY6]{Stacks}. The same holds for $\mathfrak c$ and
$\mathcal J$.

At the generic point of $D$, the conductor valuation makes
$\mathcal J=\OO_N$; at every other prime divisor its valuation is
nonnegative. Normality therefore gives
$\mathcal J\subset\OO_N$; equivalently, apply the codimension-one
extension criterion \cite[Tag 0AY7]{Stacks}.
At a closed point of $D$, a local equation $t$ is a nonzerodivisor
on the depth-two module $\mathcal J$. The quotient $\mathcal J/t\mathcal J$
has depth one and is torsion-free on the integral curve $D$.
Its map to $\OO_D$ is generically an isomorphism. Any kernel would
have finite support and be torsion, so the kernel is zero. Hence
\[
 0\longrightarrow\mathcal J|_D\longrightarrow\OO_D
 \longrightarrow\mathcal Q\longrightarrow0,
 \qquad\ell=\operatorname{length}_k\mathcal Q\geq0.
\]
This depth argument is what justifies restriction of the possibly
noninvertible conductor ideal.

Cartier-divisor duality \cite[Tag 0AA4]{Stacks}, applied to
$\omega_N[2]$, gives
$\omega_D\simeq(\omega_N\otimes\OO_N(D))|_D$.
There is no higher derived restriction, because $t$ is regular on
$\omega_N$. Define the invertible sheaf
\[
 H=(\pi^*\omega_X\otimes\OO_N((1-c)D))|_D.
\]
Then $\omega_D=H\otimes\mathcal J|_D$, and tensoring the preceding
sequence gives $0\to\omega_D\to H\to H\otimes\mathcal Q\to0$.
Serre duality on the proper Cohen--Macaulay curve gives
$\chi(\omega_D)=-\chi(\OO_D)$
\cite[Tag 0FW0]{Stacks}. The Euler-characteristic definition of degree
therefore yields
\[
 2p_a(D)-2=\deg H-\ell.
\]
Degree multiplies by the degree $p$ of $D\to Y$, even when $D$ is
singular \cite[Tag 0AYZ]{Stacks}. Thus
\[
 D^2=pY^2,\qquad
 \deg H=pK_X\cdot Y+(1-c)pY^2
       =p(2g(Y)-2)-cpY^2.
\]
The normalization sequence gives $p_a(D)=g(Y)+\delta$ with
$\delta\geq0$. Substitution proves the identity.

In the equality case, $\ell=\delta=0$. Vanishing of $\delta$ makes
$D$ normal, hence smooth. At a closed point of $D$, the regular
one-dimensional quotient $\OO_{N,x}/(t)$ has a principal maximal
ideal. Lifting its generator and adjoining $t$ gives two generators
for the maximal ideal of the two-dimensional ring $\OO_{N,x}$;
that ring is regular. The regular locus is open, so $N$ is regular
near $D$.

Vanishing of $\ell$ and Nakayama's lemma make
$\mathcal J=\OO_N$ near $D$. If these reflexive rank-one sheaves
differed anywhere, they would differ along a prime divisor. That
projective curve would be disjoint from $D$, so the canonical section
of $\OO_N(D)$ would trivialize its restriction there. This contradicts
positivity of the degree of an ample line bundle on a curve. Hence
$\mathcal J=\OO_N$ globally. The displayed description of $\omega_N$
now makes it invertible; the Cohen--Macaulay surface is Gorenstein
\cite[Tag 0BFQ]{Stacks}. Finally, a singular point of $N$ on $D$
forces $D$ to be singular by the same local argument, so $\delta\geq1$.
The strict defect bound follows from the identity.
\end{proof}

For the smallest local example, take $p=2$, $K=k(v)$,
$R=K[[t]]$, and $z^2=1+t^2v$. Then $w=(z-1)/t$ satisfies $w^2=v$,
the normalized ring is $B=R[w]$, and $A=R\oplus tRw$ has conductor
$tB$. The uniformizer is still $t$, while the residue field has doubled
its purely inseparable degree. This is a local example of the
normalization calculation, not a global Picard-kernel construction.

If $N$ is smooth near $D$ and its conductor divisor there is
$a(p-1)D+E$, the length $\ell$ equals $E\cdot D$. The module formula
proves the same nonnegative correction without that smoothness
hypothesis. Equality still permits isolated Gorenstein singularities
away from $D$. It does not justify replacing the normal surface by a
smooth one while preserving the ample divisor. On a resolution the
total pullback has zero intersection with exceptional curves, so the
smooth ample-curve obstruction of Proposition~\ref{prop:tangent-ample}
does not apply to that pullback.

\subsection{The group action in the equality case}
Equality in the conductor identity does more than control the canonical
sheaf. It makes the normalized cover a torsor under the same deformation
of $\mu_p$ that appeared in Corollary~\ref{cor:smooth-torsor}.
The group scheme is classical: locally it is the characteristic-$p$
Kummer kernel for the deformed multiplicative group
$\mathcal G^\lambda$ of M\'ezard, Romagny and Tossici
\cite[\S2.2, p.~167, and Lemma~6.1, p.~188]{MRT}.
The assertion here is that the numerical equality forces this action
to extend to the normalization and to be free.

\begin{corollary}\label{cor:normalized-equality-torsor}
In the notation and hypotheses of
Proposition~\ref{prop:normalized-mup-conductor}, suppose
$2g(Y)-2=apY^2$. Let $L=\OO_X(aY)$ and let $s_a$ be its canonical
section. Define $G_{s_a}$ by the formula in
Corollary~\ref{cor:smooth-torsor}, with $L$ and $s_a$ in place of
$\OO_X(Y)$ and $s$. Then $\pi:N\to X$ is a $G_{s_a}$-torsor.
Its infinitesimal action is a homomorphism
\[
 \OO_N(-aD)\longrightarrow T_N
\]
that is nonzero on every fibre, is tangent to $D$, and restricts to
an isomorphism $\OO_D(-aD)\simeq T_D$.
If $a=1$, the surface $N$ has at least one singular point outside $D$.
\end{corollary}
\begin{proof}
The preceding proposition gives the equality of ideal sheaves
$\mathfrak c=\OO_N(-a(p-1)D)$. In particular, $N=T$ away from $D$.
The normal surface $N$ is Cohen--Macaulay and $X$ is regular.
The dimension formula for their finite dominant map and miracle
flatness show that $\pi$ is finite flat of rank $p$
\cite[Tag 00R4]{Stacks}.

The original $\mu_p$-action gives an $X$-derivation $\delta$ of the
function field of $N$. In a local presentation $z^p=u$, it is
characterized by $\delta(z)=z$. It is independent of the chosen
frame of $M$, since changing that frame multiplies $z$ by a function
from $X$. It satisfies $\delta^p=\delta$.

Choose a local equation $t$ for $Y$ and put $\lambda=t^a$ and
$\partial=\lambda\delta$. At the generic point of $D$, the completed
normalization calculation in the proposition writes
$w=(z-v(t))/t^a$. Since $v(t)$ belongs to the completed base ring,
\[
 \partial(t)=0,\qquad \partial(w)=z.
\]
The right side is a unit. Thus $\partial$ preserves the completed
discrete valuation ring. It also preserves the uncompleted one:
its value on an element belongs to the original fraction field and
has nonnegative valuation after completion. The residue derivation is
nonzero as well: a lift to the uncompleted ring of the residue class of
$w$ differs from $w$ by $t$ times the completed ring, and its derivative
therefore has the same nonzero residue as $z$. At every other prime
divisor, $N=T$ and the original action is regular. On an affine open
of the normal integral scheme $N$, its ring is the intersection of
its local rings at height-one primes inside its fraction field.
It follows that $\partial$ is a regular derivation everywhere.

If $\lambda_j=g_{ij}^{-1}\lambda_i$ for two frames of $L$, then
$\partial_j=g_{ij}^{-1}\partial_i$. These derivations therefore
give the displayed homomorphism from $\OO_N(-aD)$. They annihilate
pullbacks from $X$, so their restriction to the smooth curve $D$
is tangent to $D$. The resulting map to $T_D$ is nonzero at its
generic point, by $\partial(w)=z$. Its source and target have the
same degree:
\[
 -aD^2=-apY^2=2-2g(D).
\]
A nonzero map between line bundles of the same degree on a smooth
proper curve is an isomorphism. Hence the infinitesimal action has
no zero on $D$. Away from $D$, its value on the unit $z$ is the unit
$\lambda z$, so it has no zero there either.

We now construct the full group action, rather than only its
derivation. On the chart under consideration, set
\[
 P_0=1,\qquad
 P_i=\frac{\partial(\partial-\lambda)\cdots
 (\partial-(i-1)\lambda)}{i!}\quad(1\leq i<p),
 \qquad
 \rho(b)=\sum_{i=0}^{p-1}P_i(b)u^i.
\]
Here $1$ in $P_0$ denotes the identity operator, and the integers
$i!$ are units in $k$. All these operators are regular, and
$\partial(\lambda)=0$. Where $\lambda$ is invertible, this is the
original $\mu_p$-coaction written through $1+\lambda u$.
Indeed, on an element of weight $j$, for $0\leq j<p$, it is
$b\mapsto b(1+\lambda u)^j$, since
$P_i(b)=\lambda^i\binom ji b$.

Multiplicativity, the identity rule, and coassociativity therefore
hold on this dense open. They hold everywhere because their target
modules, obtained by adjoining one or two variables with $p$th
power zero, are free over the integral local rings of $N$.
On overlaps, each operator $P_i$ scales by $g_{hj}^{-i}$ when the
frame changes by $g_{hj}$; this follows from the fact that the
derivations annihilate $g_{hj}$. Together with $u_j=g_{hj}u_h$, it proves that
the coactions glue. This gives an action of $G_{s_a}$ on $N$ over $X$.

It remains to prove the torsor property. Let $n$ be a closed point
of $N$, lying over $x\in X$. All residue fields are $k$. Since
$\pi$ is radicial and finite flat of rank $p$, its fibre at $x$
is a local scheme of length $p$. The group fibre has coordinate
ring $k[u]/(u^p)$. Nonvanishing of $\partial$ at $n$ gives a local
function whose pullback by the orbit map, after subtracting its
constant term, is $cu+O(u^2)$ with $c\ne0$. Such an element
generates $k[u]/(u^p)$ as an algebra: a truncated compositional
inverse recovers $u$. The orbit map is therefore a closed immersion
between schemes of equal length, hence an isomorphism.

The canonical map $N\times_XG_{s_a}\to N\times_XN$ is a map
between finite locally free $N$-schemes of the same rank $p$.
It is an isomorphism on every closed fibre by the preceding argument.
Its map of finite locally free algebras is consequently an
isomorphism, by its determinant and Nakayama's lemma. Checking
closed points suffices because $N$ is of finite type over $k$.
The finite flat map $\pi$ is surjective, so it is faithfully flat;
the action thus makes it a torsor.

Finally, if $a=1$ and $N$ were smooth, the everywhere nonzero
infinitesimal action would be a line subbundle of $T_N$ with
locally free quotient, tangent to the smooth ample curve $D$.
This contradicts Proposition~\ref{prop:tangent-ample}.
The preceding proposition already gives regularity near $D$,
so a singular point must lie outside $D$.
\end{proof}

The possibility of a singular torsor over a smooth surface is already
visible locally. On the open surface where $1+xy$ is invertible,
the equation $z^p=1+xy$ defines a $\mu_p$-torsor, because $z$ is a
unit. At $z=1$, putting $w=z-1$ gives $xy=w^p$. This is an integral
two-dimensional hypersurface, regular away from its origin and hence
normal by Serre's criterion. It is Gorenstein but singular at the
origin. Thus the torsor property does not eliminate the isolated
singularities allowed by the equality case.

\subsection{The isolated singularities and the tangent bundle}
The remaining singularities can be measured on the smooth base. Dividing
the logarithmic form by its exact vanishing along $Y$ gives a regular
section of a rank-two vector bundle. Its zero scheme detects precisely
the singular locus of the normalized cover and has a computable length.

\begin{proposition}\label{prop:normalized-mup-singularity-defect}
Keep the notation of Proposition~\ref{prop:normalized-mup-conductor},
and suppose that
\[
 2g(Y)-2=apY^2.
\]
Let $\omega$ be the logarithmic differential of $M$, put $m=pa$,
and let $s_Y$ be the canonical section of $\OO_X(Y)$.
The section
\[
 \theta=\omega/s_Y^m\in
 H^0\bigl(X,\Omega^1_{X/k}(-mY)\bigr)
\]
has a finite zero scheme $Z$ disjoint from $Y$, and
\begin{equation}\label{eq:normalized-mup-zero-length}
 \operatorname{length}_k Z=c_2(X)+apY^2=c_2(X)+aD^2.
\end{equation}
The singular subscheme of $N$ defined by
$\operatorname{Fitt}_2(\Omega^1_{N/k})$ is $N\times_X Z$.
Its length is $p\operatorname{length}_k Z$, and $N$ is smooth if and
only if $Z$ is empty.

The tangent sheaf $T_N$ is a vector bundle of rank two, even when $N$
is singular. The infinitesimal action of
Corollary~\ref{cor:normalized-equality-torsor} gives an exact sequence of vector bundles
\begin{equation}\label{eq:normalized-mup-tangent-sequence}
 0\longrightarrow\OO_N(-aD)\longrightarrow T_N
 \longrightarrow\omega_N^{-1}(aD)\longrightarrow0,
\end{equation}
and
\begin{equation}\label{eq:normalized-mup-tangent-chern}
 \deg c_2(T_N)=-aD^2.
\end{equation}
If $a=1$, then $Z$ is nonempty, so $c_2(X)+pY^2\geq1$.
\end{proposition}

\begin{proof}
The equality conclusions of
Proposition~\ref{prop:normalized-mup-conductor} give that $D$ is
smooth, $N$ is regular near $D$, and
\[
 \mathfrak c=\OO_N\bigl(-a(p-1)D\bigr),\qquad
 \omega_N\simeq\pi^*\omega_X\otimes
                \OO_N\bigl(-a(p-1)D\bigr).
\]
In particular the normalization is an isomorphism away from $D$:
$N\setminus D=T\setminus\tau^{-1}Y$. The map $\pi$ is finite flat
of rank $p$. Indeed, its source is a Cohen--Macaulay normal surface,
its target is regular, and its local fibres have dimension zero;
the local dimensions agree because the map is finite and radicial.
Miracle flatness applies \cite[Tag 00R4]{Stacks}.

The exact order $m$ of $\omega$ makes $\theta$ regular. Its
tangential component along $Y$ is a nonzero section of
\[
 \Omega^1_{Y/k}\otimes\OO_Y(-mY).
\]
This line bundle has degree zero by the equality hypothesis, so the
section is nowhere zero. Thus $\theta$ has no zero on $Y$.

Away from $Y$, a trivialization of $M$ gives
$T=\Spec R[z]/(z^p-u)$ with $u\in R^*$, and $\theta$ is a unit
multiple of $du$. Choose an \'etale local basis $dx,dy$ of
$\Omega^1_{R/k}$ and write $du=\alpha\,dx+\beta\,dy$.
The differential module of this hypersurface has three generators
$dx,dy,dz$ and the single relation
$\alpha\,dx+\beta\,dy=0$. Its second Fitting ideal is therefore
$(\alpha,\beta)$. The hypersurface is smooth precisely where this
ideal is the unit ideal.

On this open set $T=N$ is normal, so its singular locus has no
codimension-one component. Finiteness of $\pi$ implies that the
zero scheme of $\theta$ has no curve away from $Y$. Since it also
avoids $Y$, the scheme $Z$ is zero-dimensional, hence finite. At a
point of $Z$, its two local equations form a regular sequence in
the regular local ring of the smooth surface $X$. The top Chern
class of a regular section now gives
\[
 \operatorname{length}_k Z
 =\deg c_2\bigl(\Omega^1_{X/k}(-mY)\bigr);
\]
all the dimension and regular-sequence hypotheses of
\cite[Tag 0FA9]{Stacks} hold here. The rank-two twist formula
\cite[Tag 02UD]{Stacks} gives
\[
 \deg c_2\bigl(\Omega^1_{X/k}(-mY)\bigr)
 =c_2(X)-mK_X\cdot Y+m^2Y^2.
\]
Adjunction and equality imply $K_X\cdot Y=(m-1)Y^2$, proving
\eqref{eq:normalized-mup-zero-length}.

The local Fitting-ideal computation also proves
$\operatorname{Fitt}_2(\Omega^1_{N/k})=\mathcal I_Z\OO_N$
off $D$. Near $D$ both ideals are the unit ideal, because $N$ is
smooth there and $Z$ avoids $Y$. The equality therefore holds
globally. Flatness of rank $p$ gives the asserted length. The same
local hypersurface calculation, together with smoothness near $D$,
shows that $N$ is smooth exactly when $Z$ is empty.

We next check local freeness of the tangent sheaf at a singular
point. In the above presentation put $B=R[z]/(z^p-u)$.
The elements $\alpha,\beta$ form a regular sequence in $R$,
because their zero scheme is finite, and remain a regular sequence
in $B$ by flatness. Thus
\[
 \operatorname{Der}_k(B,B)
 =\ker\bigl(B^3\longrightarrow B,\ (r,s,t)\longmapsto
             \alpha r+\beta s\bigr)
\]
is free with basis $(-\beta,\alpha,0)$ and $(0,0,1)$: the first
generator is the Koszul syzygy of the regular sequence. At every
other point $N$ is smooth. Hence $T_N$ is a rank-two vector bundle
on all of $N$.

Corollary~\ref{cor:normalized-equality-torsor} gives the map
$\OO_N(-aD)\to T_N$, nonzero in every tangent fibre and tangent
to $D$. Since $T_N$ is now known to be locally free, this is a
line subbundle with invertible quotient.

On the regular locus, $\det T_N=\omega_N^{-1}$. Both sheaves are
invertible on all of $N$, and normality extends this identification
across the finite complement. The quotient is therefore
$\omega_N^{-1}(aD)$, proving
\eqref{eq:normalized-mup-tangent-sequence}. Adjunction along $D$
gives $K_N\cdot D=(a-1)D^2$. The Whitney formula then yields
\[
 \deg c_2(T_N)=(-aD)\cdot(-K_N+aD)=-aD^2,
\]
as claimed.

Finally, if $a=1$ and $Z$ were empty, $N$ would be a smooth
projective surface with the line subbundle $\OO_N(-D)\subset T_N$
tangent to its smooth ample curve $D$. This contradicts
Proposition~\ref{prop:tangent-ample}. Thus $Z$ has positive integral
length in this case.
\end{proof}

The length in this proposition is a Fitting-scheme length. It is not
a claimed Milnor number or a discrepancy correction on a resolution.
In particular, the negative second Chern number of the vector bundle
$T_N$ cannot be substituted for the second Chern number of a smooth
resolution. The local model $xy=w^p$ already has a locally free
tangent sheaf and a singular subscheme of length $p$.

\begin{corollary}\label{cor:normalized-mup-plurigenera}
In Proposition~\ref{prop:normalized-mup-singularity-defect}, suppose
$a=1$. Let $f:S\to N$ be a smooth projective resolution that is an
isomorphism over the regular locus of $N$. For every integer $r>0$
there is an injection
\[
 H^0(S,\omega_S^{\otimes r})\longrightarrow
 H^0(N,\omega_N^{\otimes r}),
 \qquad h^0(N,\omega_N^{\otimes r})\leq1.
\]
Every nonzero section of $\omega_N^{\otimes r}$ is nowhere
vanishing. Thus $\kappa(S)\leq0$, and $K_N$ is torsion if some
positive plurigenus of $S$ is nonzero.
\end{corollary}

\begin{proof}
Here $K_N\cdot D=0$, and $D$ is ample. A nonzero section of the
invertible sheaf $\omega_N^{\otimes r}$ has an effective Cartier
zero divisor $E$. If $E$ were nonempty, ampleness would give
$E\cdot D>0$, whereas $E\cdot D=rK_N\cdot D=0$.
Nor can the section have only an isolated zero: a nonzero nonunit
in a noetherian normal local domain has a height-one zero.
The section therefore trivializes $\omega_N^{\otimes r}$.
Since $N$ is proper and integral over the algebraically closed
field, $H^0(N,\OO_N)=k$. It follows that the indicated space of
sections has dimension at most one.

Restrict a pluricanonical section on $S$ to the common regular
open set. The complement of that open set in $N$ is finite.
Normality extends the section uniquely to the invertible sheaf
$\omega_N^{\otimes r}$ on $N$ \cite[Tag 0AY7]{Stacks}.
This gives the injection, which is injective because $S$ is
integral. All positive plurigenera of $S$ are consequently at
most one. If one is nonzero, the corresponding power of
$\omega_N$ is trivial by the first paragraph. The claimed
Kodaira-dimension bound follows.
\end{proof}

For the smallest equality case, Ekedahl's foliation theorem
restricts the surface further. Following his terminology, a surface
is \emph{almost ruled} if it admits a dominant purely inseparable
rational map of degree $p$ from a ruled surface
\cite[p.~108]{EkedahlCanonical}.

\begin{proposition}[After Ekedahl]\label{prop:normalized-mup-ekedahl}
In Proposition~\ref{prop:normalized-mup-conductor}, suppose $a=1$
and $2g(Y)-2=pY^2$. Then either $X$ is almost ruled, or $p=2$ and
\[
 \omega_X\text{ is ample},\qquad
 \chi(\OO_X)=1,\qquad
 \omega_X^{\otimes2}\simeq\OO_X(2Y).
\]
In the latter case $N$ has only rational double points, its minimal
resolution is a K3 surface, and
\[
 1\leq Y^2=K_X^2\leq9,\qquad g(Y)=Y^2+1.
\]
In this latter case the zero scheme $Z$ of
Proposition~\ref{prop:normalized-mup-singularity-defect} has length
$12+Y^2$, and the Fitting singular subscheme of $N$ has length
$2(12+Y^2)$. In particular, when $p>2$ the stated equality forces
$X$ to be almost ruled.
\end{proposition}

\begin{proof}
The wedge pairing for the rank-two cotangent bundle identifies
$\Omega^1_{X/k}$ with $T_X\otimes\omega_X$. Thus the nonzero
section $\theta\in H^0(X,\Omega^1_{X/k}(-pY))$ gives an injection
\[
 \mathcal L:=\OO_X(pY)\otimes\omega_X^{-1}\longrightarrow T_X.
\]
Its image is saturated. Indeed, in smooth local coordinates write
$\theta=A\,dx+B\,dy$. Its zero scheme is finite, so at a zero the
pair $A,B$ is a regular sequence. The kernel of
$(r,s)\mapsto Ar+Bs$ is generated by $(-B,A)$, and its image is
the torsion-free ideal $(A,B)$. Globally this gives
\[
 0\longrightarrow\mathcal L\longrightarrow T_X
 \longrightarrow\mathcal I_Z\otimes\OO_X(-pY)
 \longrightarrow0.
\]
On a Kummer trivialization the form is a nonzero rational multiple
of $du$, where the cover has equation $z^p=u$. Its annihilator
consists of derivations killing $u$, hence is closed under brackets
and $p$th powers. Saturation extends these properties to the whole
sheaf. This is a possibly singular rank-one foliation in the sense
of \cite[pp.~104--105]{EkedahlCanonical}.

Adjunction and the equality hypothesis give
\[
 K_X\cdot Y=(p-1)Y^2,\qquad
 \mathcal L\cdot Y=Y^2>0,
 \qquad (K_X-(p-1)\mathcal L)\cdot Y=0.
\]
Assume that $X$ is not almost ruled. It is then not birationally
ruled: a ruled function field $k(C)(t)$ has the purely inseparable
degree-$p$ extension $k(C)^{1/p}(t)$, which is again ruled.
We may therefore apply
\cite[I, Theorem~2.4(ii), pp.~109--115]{EkedahlCanonical}, with
the injection just constructed and the ample test line bundle
$\OO_X(Y)$.

In alternative~(a) of that theorem, $X$ is minimal of Kodaira
dimension zero or one and
\[
 \omega_X^{\otimes12p}\simeq
 \mathcal L^{\otimes12p(p-1)}.
\]
Substitution gives
$\omega_X^{\otimes12p^2}\simeq
\OO_X(12p^2(p-1)Y)$, which makes $\omega_X$ ample and is
impossible in those Kodaira dimensions. Alternative~(b) makes
$\mathcal L$ trivial, contrary to $\mathcal L\cdot Y>0$.
Alternative~(c) has $p=2$, Kodaira dimension one, and
$\omega_X^{\otimes2}\simeq\mathcal L^{\otimes2}$.
It would give $\omega_X^{\otimes4}\simeq\OO_X(4Y)$,
again a contradiction. Alternative~(d) consequently holds:
$p=2$, $\chi(\OO_X)=1$ and $\mathcal L\simeq\omega_X$.
This proves the asserted square identity and ampleness. It also
shows that $X$ is minimal, since a $(-1)$-curve would have negative
intersection with $K_X$.

We identify the foliation cover in that theorem with $N$.
Put $K=k(X)$ and $E=K(z)$, with $z^p=u$. The nonzero form gives
$u\notin K^p$. Complete $u$ to a $p$-basis $u,v$ of $K$ over
$K^p$. Expansion in the basis $u^iv^j$, for $0\leq i,j<p$,
shows that the invariant field of the derivations killing $u$ is
$F=K^p(u)$. On the other hand,
\[
 E^p=K^p(u)=F,\qquad E=F^{1/p}.
\]
Ekedahl's convention has relative Frobenius $X^{(1)}\to X$;
his foliation cover $(X/\mathcal L)^{(1)}\to X$ has field
$F^{1/p}$ \cite[p.~104]{EkedahlCanonical}.
The identification is over $X$, not merely birational. For a
normal affine ring $R$ on $X$, the invariant ring is $C=R\cap F$.
It is normal and contains $R^p$. The ring $C^{1/p}$ contains $R$,
is a submodule of the finite $R$-module $R^{1/p}$, and is normal
with fraction field $E$. Hence it is the integral closure of $R$
in $E$. These descriptions glue and identify that cover with $N$.

For completeness we spell out the part of the source's proof
that gives the assertion about this cover. Let $\tau:S\to N$
be the minimal resolution and put $f=\pi\tau$. The canonical
formula in \cite[I, Proposition~2.1, p.~108]{EkedahlCanonical}
and the zero intersection above give $K_S\cdot f^*Y=0$.
Thus $\kappa(S)\leq0$, also by
Corollary~\ref{cor:normalized-mup-plurigenera}. If
$\kappa(S)=-\infty$, the surface $S$ is birationally ruled and
its degree-$p$ purely inseparable map to $X$ makes $X$ almost
ruled. Consequently $\kappa(S)=0$.

The surface $S$ is minimal. Otherwise its canonical class is the
pullback of the numerically trivial canonical class of its minimal
model plus an effective exceptional divisor. The zero intersection
with $f^*Y$ forces every component of that divisor to be contracted
by $f$, hence by $\tau$, because $N\to X$ is finite. This
contradicts minimality of the resolution. The surface $N$ is
singular by Proposition~\ref{prop:normalized-mup-singularity-defect}.
Thus \cite[I, Lemma~2.5, p.~111]{EkedahlCanonical} applies:
$N$ has rational double points, and $S$ is K3 or Enriques.

The Enriques possibility is excluded by Euler characteristic.
A finite flat purely inseparable double cover has zero trace and is an
$\alpha_{\mathcal A}$-torsor for an invertible sheaf $\mathcal A$,
and its algebra fits into
\[
 0\longrightarrow\OO_X\longrightarrow\pi_*\OO_N
 \longrightarrow\mathcal A^{-1}\longrightarrow0
 \quad\text{\cite[I, Proposition~1.11, p.~107]{EkedahlCanonical}}.
\]
The associated foliation has
$\mathcal L\simeq\mathcal A^{-2}\otimes\omega_X^{-1}$.
Since $\mathcal L\simeq\omega_X$, the bundle
$\mathcal A^{-1}$ is numerically equivalent to $\omega_X$.
Riemann--Roch gives
$\chi(\mathcal A^{-1})=\chi(\omega_X)=\chi(\OO_X)=1$.
Rational double points are rational singularities, so
$\chi(\OO_S)=\chi(\OO_N)=2$. An Enriques surface has Euler
characteristic one. Therefore $S$ is K3.

The purely inseparable morphism $S\to X$ preserves the first
$\ell$-adic Betti number
\cite[Lemma~0.1.5, p.~103]{EkedahlCanonical}, so $b_1(X)=0$.
Poincar\'e duality gives $c_2(X)=2+b_2(X)\geq3$: the ample
divisor $Y$ has nonzero cycle class because its square is positive.
Noether's formula gives $c_2(X)=12-K_X^2$, while the square
identity gives $K_X^2=Y^2>0$. This proves $1\leq Y^2\leq9$.
The genus formula is the assumed equality for $p=2$.
Finally, \eqref{eq:normalized-mup-zero-length} gives
$\operatorname{length}_k Z=12+Y^2$, and the rank-two flat cover
doubles that length on $N$.
\end{proof}

This proposition applies Ekedahl's published theorem. The
almost-ruled alternative and the characteristic-two K3 case
remain open for the Picard restriction problem. In particular,
a purely inseparable map from a ruled surface does not assert
that $X$ itself is birationally ruled. The strict inequality
$2g(Y)-2>apY^2$ and the cases $a>1$ also remain outside this
reduction.

\subsection{The characteristic-two K3 possibility}

We now describe the K3 case permitted for an ample divisor. The splitting
below follows from duality on the singular surface itself. It does
not require a positive-characteristic vanishing theorem.

\begin{lemma}\label{lem:k3-ample-curve-vanishing}
Let $V$ be a projective integral Gorenstein surface over an
algebraically closed field, with $\omega_V\simeq\OO_V$ and
$H^1(V,\OO_V)=0$. If $C\subset V$ is a smooth ample integral
Cartier curve, then
\[
 H^1(V,\OO_V(nC))=0\quad(n\in\Z\setminus\{0\}),
 \qquad h^0(V,\OO_V(C))=1+g(C)=2+C^2/2.
\]
The line bundle $\OO_V(C)$ is globally generated.
\end{lemma}

\begin{proof}
Adjunction identifies $\OO_C(C)$ with $\omega_C$. The sequence
\[
 0\longrightarrow\OO_V\longrightarrow\OO_V(C)
 \longrightarrow\omega_C\longrightarrow0
\]
therefore gives the asserted dimension of sections. By Serre
duality, $H^2(\OO_V(C))=H^0(\OO_V(-C))^*=0$, whereas
$H^2(\OO_V)$ and $H^1(\omega_C)$ both have dimension one.
The surjection $H^1(\omega_C)\to H^2(\OO_V)$ is consequently
an isomorphism. Its kernel is $H^1(\OO_V(C))$, which vanishes.

For $n\geq2$, use the divisor sequence with successive twists
$(n-1)C$ and $nC$. Curve duality gives
\[
 H^1(C,\OO_C(nC))=H^0(C,\OO_C((1-n)C))^*=0.
\]
Induction proves the positive twists; surface duality proves the
negative ones. These applications of adjunction and duality are
valid for a projective Gorenstein surface
\cite[Tags 0AA4, 0FVZ and 0FW0]{Stacks}.

Finally $C^2=2g(C)-2>0$, so $g(C)\geq2$. The canonical system
of $C$ has no base point. Indeed, for a point $x\in C$ one has
$h^0(\OO_C(x))=1$, since a degree-one map to $\PP^1$ would
make $C$ rational; Riemann--Roch gives
$h^0(\omega_C(-x))=g(C)-1$. The surjection of sections onto
$H^0(\omega_C)$ proves generation along $C$, and the defining
section of $C$ proves it away from $C$.
\end{proof}

\begin{proposition}\label{prop:k3-picard-derivations}
In the non-almost-ruled alternative of
Proposition~\ref{prop:normalized-mup-ekedahl}, put $s=Y^2$.
The tangent sequence has a unique splitting, and
\[
 T_N\simeq\OO_N(-D)\oplus\OO_N(D),\qquad
 h^0(N,T_N)=2+s,\qquad H^1(N,T_N)=0.
\]
All global derivations lie in the positive summand and generate
it. Every nonzero global derivation $v$ satisfies $v^2=c_vv$ for
some $c_v\in k$. Its normal quotient is rational. There are
dominant purely inseparable rational maps from rational surfaces
to $N$ of degree two and to $X$ of degree four.
\end{proposition}

\begin{proof}
Rational double points are rational and crepant singularities.
The K3 resolution therefore gives $\omega_N\simeq\OO_N$ and
$H^1(\OO_N)=0$. The extension class of
\eqref{eq:normalized-mup-tangent-sequence} belongs to
\[
 \operatorname{Ext}^1(\OO_N(D),\OO_N(-D))
 =H^1(\OO_N(-2D))=0
\]
by Lemma~\ref{lem:k3-ample-curve-vanishing}. The splitting is
unique because $H^0(\OO_N(-2D))=0$. The lemma also gives
the cohomology dimensions and global generation of the positive
summand $P=\OO_N(D)$, since $g(D)=s+1$.

For a local generator $e$ of $P$, the class of $e^2$ modulo $P$
transforms quadratically: in characteristic two,
$(fe)^2=f^2e^2+f e(f)e$. Thus it defines a homomorphism
\[
 P^{\otimes2}\longrightarrow T_N/P=\OO_N(-D).
\]
It is zero because $H^0(\OO_N(-3D))=0$. Write $e^2=he$
with $h$ regular. If a global derivation is locally $v=fe$, then
\[
 v^2=(fh+e(f))v.
\]
The coefficient is regular even at the zeros of $f$ and agrees
in different trivializations on the dense locus where $v\ne0$.
It therefore defines a global regular function on $N$, hence a
scalar $c_v$. If $c_v\ne0$, rescale $v$ so that $v^2=v$;
otherwise $v^2=0$. These are respectively a $\mu_2$-action and
an $\alpha_2$-action
\cite[Proposition~2.10, author version p.~5]{MatsumotoActions}.

At every codimension-one point, $N$ is smooth and $P$ is a line
subbundle of its tangent bundle. The divisorial fixed locus of
$v$ is therefore the zero divisor of the corresponding section
of $P$. It is linearly equivalent to $D$, and is nonzero.
The quotient $R=N^v$ is normal and rational by
\cite[Proposition~4.1 and proof, author version pp.~16--17]{MatsumotoActions}:
the nonzero divisorial fixed part places it in the rational
alternative there.

For the degree assertion write $E=k(N)$ and $J=k(R)$.
The quotient has $[E:J]=2$, while $E^2\subset J$ and
$[E:E^2]=4$. Hence $E\subset J^{1/2}$ and
$[J^{1/2}:E]=2$. The field $J^{1/2}$ is rational over the
perfect ground field. It gives a degree-two rational domination
of $N$, and composition with $N\to X$ gives degree four.
\end{proof}

The degree-four conclusion does not imply that $X$ is almost
ruled: that term here requires degree two. In particular the
large space of global derivations does not exclude this case.

\begin{proposition}\label{prop:k3-picard-paracanonical}
In characteristic two, in any equality case
$2g(Y)-2=2aY^2$ of
Proposition~\ref{prop:normalized-mup-conductor}, one has
\[
 \pi_*\OO_N/\OO_X\simeq M^{-1}(aY).
\]
In the non-almost-ruled alternative of
Proposition~\ref{prop:normalized-mup-ekedahl}, the cover
$N\to X$ is an $\alpha_{\omega_X^{-1}}$-torsor and
\[
 0\longrightarrow\OO_X\longrightarrow\pi_*\OO_N
 \longrightarrow\omega_X\longrightarrow0,
 \qquad M\simeq\OO_X(Y)\otimes\omega_X^{-1}.
\]
Thus $Y$ belongs to $|\omega_X\otimes M|$, with $M$ a
nontrivial two-torsion twist, and
$\OO_N(D)\simeq\pi^*\omega_X$.
\end{proposition}

\begin{proof}
The unit in the locally free rank-two algebra $\pi_*\OO_N$
is nonzero on every fibre. Its quotient $Q$ is therefore an
invertible sheaf. The inclusion of the Kummer algebra
$\OO_X\oplus M^{-1}$ gives a nonzero map $M^{-1}\to Q$,
an isomorphism away from $Y$. At the generic point of $Y$,
the normalized basis in the conductor proof is $1,w$, with
$w=(z-v)/t^a$. Modulo the base summand, $[z]=t^a[w]$.
The divisor of the map of line bundles is exactly $aY$;
there is no additional isolated-zero divisor. Hence
$Q=M^{-1}(aY)$.

Now assume the non-almost-ruled case. Since $K_X$ is ample,
$X$ is its own minimal model. Ekedahl's theorem gives a normal
$\alpha_{\omega_X^{-1}}$-torsor $C\to X$ with rational double
points and K3 resolution
\cite[I, Theorem~2.4(ii)(d), p.~109]{EkedahlCanonical}.
We identify this cover with $N$. For an $\alpha_L$-torsor,
the differential of its equation gives $L^{-2}\to\Omega^1_X$.
Normality of $C$ makes its zero scheme finite, so its saturated
annihilator is $L^{-2}\otimes\omega_X^{-1}\subset T_X$.
For $L=\omega_X^{-1}$ this is an embedding $\omega_X\to T_X$.
The foliation belonging to $N$ has the same source because
$\OO_X(2Y)=\omega_X^2$.

Any two such embeddings have zero wedge: the wedge is a section
of $\operatorname{Hom}(\omega_X^2,\det T_X)=\omega_X^{-3}$,
which has no nonzero section. The two generic lines agree and,
by saturation, define the same foliation. Its invariant field
determines the square of the covering field. Uniqueness of square
roots in characteristic two therefore identifies the two field
extensions over $k(X)$, and normality identifies $C$ and $N$
over $X$, as in the proof of
Proposition~\ref{prop:normalized-mup-ekedahl}.

The algebra filtration for this additive torsor has quotient
$\omega_X$ \cite[I, Proposition~1.7, p.~107]{EkedahlCanonical}.
Comparison with $Q=M^{-1}(Y)$ gives the asserted formula for
$M$. The pullback of $M$ to its Kummer torsor is trivial;
pulling this trivialization back along $N\to T$ gives
$\pi^*M\simeq\OO_N$. The final identity follows.
\end{proof}

\begin{proposition}[After Matsumoto]\label{prop:k3-picard-singularities}
In the non-almost-ruled alternative of
Proposition~\ref{prop:normalized-mup-ekedahl}, every singularity
of $N$ has type
\[
 A_1,\qquad D_{2m}^{0}\ (m\geq2),\qquad E_7^{0},\qquad E_8^{0}.
\]
The sum of the Dynkin indices is $12+Y^2$, between thirteen
and twenty-one. This is a necessary condition; it does not assert
that every configuration with this sum occurs.
\end{proposition}

\begin{proof}
All singularities lie away from $D$, where $N=T$ is the original
$\mu_2$-torsor over smooth $X$. Its derivation has no fixed
point there. The smooth-quotient rows in
\cite[Theorem~3.3 and Table~2, author version pp.~10--11]{MatsumotoActions}
give exactly the listed types. The local multiplicative
classification is proved in
\cite[Proposition~4.8 and Table~3, pp.~13--15]{MatsumotoMuN}.
The superscript zero specifies the characteristic-two local type.

The index sum can also be read directly from their equations.
Write the local equation as $z^2+h(x,y)=0$. The Fitting
singular ideal is $(h_x,h_y)$ in a free rank-two algebra over
$k[[x,y]]$. The relevant base lengths are
\[
\begin{array}{c|c|c}
 \text{type}&h&\dim_k k[[x,y]]/(h_x,h_y)\\ \hline
 A_1&xy&1\\
 D_{2m}^{0}&x^2y+xy^m&2m\\
 E_7^{0}&x^3+xy^3&7\\
 E_8^{0}&x^3+y^5&8.
\end{array}
\]
For the $D$ row, the quotient is by
$(y^m,x^2+\epsilon xy^{m-1})$, where $\epsilon=m\bmod2$;
it is free of rank two over $k[[y]]/(y^m)$. For the $E_7$
row, a basis is $1,y,y^2,y^3,y^4,x,xy$. For the $E_8$
row, the ideal is $(x^2,y^4)$. Each Fitting length upstairs is
twice the indicated Dynkin index. Their sum is
$2(12+Y^2)$ by
Proposition~\ref{prop:normalized-mup-ekedahl}, which proves
the assertion.
\end{proof}

For a merely ample divisor these constraints leave a specific
existence question: can such
a characteristic-two surface have a smooth paracanonical curve
on which its nontrivial two-torsion twist becomes trivial? The
splitting, rational domination and local singularity list do not
answer that question.


Gounelas and Liedtke prove that a smooth K3 surface has no nonzero
global symmetric differential of positive degree except in
characteristic two with Artin invariant one
\cite[Theorem~1.1, p.~2]{GounelasLiedtkeSymmetric}.
In the exceptional case the dimension is one in positive even
degrees and zero in odd degrees. Their theorem concerns the
ordinary cotangent bundle of the smooth resolution. Here the
splitting gives the reflexive cotangent bundle on the singular
surface $N$. Applying that theorem would require proving that a
specified nonzero symmetric tensor pulls back without poles along
every exceptional component. The splitting and the singularity
list alone do not provide this extension, so they do not reduce
the remaining alternative to Artin invariant one.

\subsection{Very ampleness excludes the K3 alternative}

Very ampleness gives a stronger conclusion than ampleness in the
smallest equality case. A hyperplane section would have genus too
large for its degree unless the surface lay in projective
four-space. The normal bundle then gives a numerical contradiction.
We first give the curve estimate in arbitrary characteristic.

\begin{lemma}\label{lem:small-degree-curve-bound}
Let $B\subset\PP^n$, $n\geq4$, be a smooth integral nondegenerate
curve over an algebraically closed field. If its degree is
$d\leq9$ and its genus is at least two, then
\[
 g(B)\leq\max(d-4,0)+\max(d-7,0)<d+1.
\]
\end{lemma}

\begin{proof}
We use Kadets's geometric monodromy theorem: the sectional
monodromy of a nondegenerate integral nonstrange curve in
$\PP^n$ is $n$-transitive
\cite[Proposition~2.8 and proof, pp.~4--5]{Kadets}.
A curve is strange if all its tangent lines pass through one
point. The theorem holds in arbitrary characteristic.

The curve $B$ is nonstrange. Indeed, suppose all tangent lines
passed through $q\in\PP^n$. The twisted Euler sequence is
\[
 0\longrightarrow\OO_B(-1)\longrightarrow\OO_B^{n+1}
 \longrightarrow T_{\PP^n}|_B(-1)\longrightarrow0.
\]
A vector representing $q$ gives a section of the last bundle
which is nonzero generically. Its image in
$N_{B/\PP^n}(-1)$ is zero: away from $q$ the section is the
direction of the joining line, which is tangent to $B$.
The tangent--normal sequence consists of vector bundles, so it
factors globally through $T_B(-1)$, also when $q$ lies on $B$.
This would be a nonzero section of a line bundle of degree
$2-2g(B)-d<0$, a contradiction.

Let $\Gamma$ be a geometric generic hyperplane section, a set
of $d$ distinct points in $\PP^{n-1}$. Distinctness follows
also directly: hyperplanes through a fixed tangent line have
dimension $n-2$, so the incidence of tangent hyperplanes has
dimension at most $n-1$ in the $n$-dimensional dual space.
The points of $\Gamma$ span their hyperplane. To see this, use
\[
 0\longrightarrow\mathcal I_B\longrightarrow\mathcal I_B(1)
 \longrightarrow\mathcal I_{\Gamma/\PP^{n-1}}(1)
 \longrightarrow0.
\]
The hyperplane equation is a nonzerodivisor on $B$, which gives
exactness. Nondegeneracy gives $H^0(\mathcal I_B(1))=0$.
Moreover $H^1(\mathcal I_B)=0$, from
$H^0(\OO_{\PP^n})\simeq H^0(\OO_B)$ and
$H^1(\OO_{\PP^n})=0$. Hence
$H^0(\mathcal I_{\Gamma/\PP^{n-1}}(1))=0$.

In particular $d\geq n$ and some $n$ points of $\Gamma$ are
linearly independent. Linear dependence is invariant under the
sectional monodromy group. Its $n$-transitivity therefore makes
every $n$ distinct points independent; thus $\Gamma$ is in
linear general position. This uses no assertion of uniform
position in higher degree.

Write $h_\Gamma(m)$ for the homogeneous-coordinate Hilbert
function. Linear general position implies
\[
 h_\Gamma(m)\geq\min\{d,m(n-1)+1\}\qquad(m\geq1).
\]
Here is a direct interpolation argument. Choose any
$t=\min\{d,m(n-1)+1\}$ of the points and distinguish one.
Partition the other $t-1$ points into $m$ groups of at most
$n-1$ points. Each group lies on a hyperplane avoiding the
distinguished point, by linear general position. The product
of these hyperplane equations vanishes at the other selected
points and is nonzero at the distinguished point. Repeating
this for each selected point proves that evaluation has rank
$t$.

Let $R$ be the homogeneous coordinate ring of $B$, and let
$\ell$ define the generic hyperplane. Multiplication by
$\ell$ is injective. The quotient $R/\ell R$ surjects onto
the homogeneous coordinate ring of $\Gamma$, so
\[
 h_B(m)-h_B(m-1)\geq h_\Gamma(m).
\]
Since $h_B(0)=1$ and $h_B(m)=dm+1-g(B)$ for sufficiently
large $m$, summation gives
\[
 g(B)\leq\sum_{m\geq1}(d-h_\Gamma(m))
 \leq\sum_{m\geq1}\max\{d-m(n-1)-1,0\}.
\]
The possible saturation kernel in $R/\ell R$ has the required
sign; projective normality is not assumed. For $n\geq4$ and
$d\leq9$, the terms with $m\geq3$ vanish. The two remaining
terms are bounded by $\max(d-4,0)+\max(d-7,0)$.
For $d=4,\ldots,9$ these values are respectively
$0,1,2,3,5,7$, each smaller than $d+1$; degrees below four
were already excluded by $d\geq n$.
\end{proof}

\begin{proposition}\label{prop:very-ample-numerical-obstruction}
There is no smooth projective integral surface $X$ over an
algebraically closed field, of any characteristic, with a very
ample divisor $H$ having a smooth integral member, such that
\[
 K_X\equiv H,\qquad \chi(\OO_X)=1,
 \qquad 1\leq H^2\leq9.
\]
Here $\equiv$ denotes numerical equivalence.
\end{proposition}

\begin{proof}
Embed $X$ by the complete system $|H|$ in $\PP^r$ and put
$d=H^2$. A smooth integral member $C\in|H|$ is a hyperplane
section with
\[
 2g(C)-2=(K_X+H)\cdot H=2d,
 \qquad g(C)=d+1.
\]
It is nondegenerate in its hyperplane: the kernel of
$H^0(X,H)\to H^0(C,H|_C)$ is the one-dimensional space
generated by its defining section. If $r\geq5$,
Lemma~\ref{lem:small-degree-curve-bound} contradicts this
genus. Thus $r\leq4$.

An embedded surface requires $r\geq2$. For $r=2$, the image
is all of $\PP^2$, where $K_X=-3H$, a contradiction. For
$r=3$, the image is a degree-$d$ hypersurface and adjunction
gives $K_X=(d-4)H$. Numerical equivalence forces $d=5$.
The sequence
\[
 0\longrightarrow\OO_{\PP^3}(-5)
 \longrightarrow\OO_{\PP^3}\longrightarrow\OO_X
 \longrightarrow0
\]
then gives $h^0(\OO_X)=1$, $h^1(\OO_X)=0$ and
$h^2(\OO_X)=4$, by the cohomology of projective space
\cite[Tag 01XT]{Stacks}. Thus $\chi(\OO_X)=5$, again a
contradiction.

It remains to exclude $r=4$. Write $i:X\hookrightarrow\PP^4$
and $E=N_{X/\PP^4}$. This is a regular immersion of smooth
schemes, with tangent--normal sequence
\[
 0\longrightarrow T_X\longrightarrow T_{\PP^4}|_X
 \longrightarrow E\longrightarrow0.
\]
The Euler sequence gives $c_1(T_{\PP^4}|_X)=5H$ and
$c_2(T_{\PP^4}|_X)=10H^2$. As $c_1(T_X)=-K_X$, the first
two Chern terms give
\[
 c_1(E)=5H+K_X,\qquad
 \deg c_2(E)=10d+5K_X\cdot H+K_X^2-c_2(X).
\]
Noether's formula $K_X^2+c_2(X)=12\chi(\OO_X)$ follows
in every characteristic from Grothendieck--Riemann--Roch for
the structure morphism \cite[Tag 02UO]{Stacks}; it is an
equality of numerical Chern numbers. Consequently
\[
 c_2(X)=12-d,\qquad \deg c_2(E)=17d-12.
\]

Self-intersection gives the other value $\deg c_2(E)=d^2$.
For precision, base-change $i$ along itself. The fibre product
is $X$, its map to $X$ is the identity, and its excess
conormal bundle is $E^\vee$. The excess-intersection formula
\cite[Tag 0FV9]{Stacks} gives
\[
 i^!i_*[X]=c_2(E)\cap[X].
\]
Let $h=c_1(\OO_{\PP^4}(1))$. Since
$i_*[X]=d h^2\cap[\PP^4]$, commutation of $i^!$ with
Chern classes and $i^![\PP^4]=[X]$ give
\[
 i^!i_*[X]=d H^2\cap[X].
\]
These identities are Lemmas~42.59.7--8 of
\cite[Tag 0FEZ]{Stacks}; smoothness supplies the dimension
and Cohen--Macaulay hypotheses. Taking degrees gives $d^2$.
Thus $d^2-17d+12=0$, whose discriminant $241$ lies strictly
between $15^2$ and $16^2$. It has no integer root.
\end{proof}

\subsection{The distinguished cover in the almost-ruled case}

The remaining almost-ruled alternative determines the particular
cover defined by the torsion line bundle, rather than only some
purely inseparable cover of the surface.

\begin{proposition}\label{prop:distinguished-ruled-cover}
Let $X$ be a smooth projective integral surface over an
algebraically closed field $k$ of characteristic $p>0$.
Let $Y\subset X$ be a smooth integral ample divisor, and let
$M\ne\OO_X$ have order $p$ with $M|_Y\simeq\OO_Y$.
Let $N$ be the normalization of the $\mu_p$-torsor defined by
a trivialization of $M^{\otimes p}$.
If a ruled surface $R$ admits a dominant purely inseparable
rational map of degree $p$ to $X$, then
\[
 k(R)=k(N)\quad\text{as extensions of }k(X).
\]
In particular $N$ is birationally ruled, and every smooth
projective model of its function field has Kodaira dimension
$-\infty$. A resolution of the rational map from a smooth
projective ruled model gives a proper birational morphism
to $N$ over $X$. The surface $X$ itself is not birationally ruled.
Moreover, if $k(N)=A(t)$ is a ruling presentation, then
\[
 k(X)\cap A=A^p,\qquad k(X)A=k(N).
\]
\end{proposition}

\begin{proof}
Choose local frames $e_i$ of $M$ and write $u_i\in\OO_X^*$
for their $p$th powers under the chosen trivialization.
On an overlap, $u_i=g_{ij}^p u_j$, so the forms $du_i/u_i$
glue to a regular global form $\omega$ on $X$.
It is nonzero. Indeed, for a rational frame $e$, let
$u\in K^*$ be its $p$th power, where $K=k(X)$.
If $du=0$, the perfect-field differential criterion
\cite[Tag~031W]{Stacks} gives $u=v^p$ with $v\in K$.
The rational section $e/v$ has $p$th power one.
Its divisor has every coefficient zero, so normality makes
it a nowhere-vanishing global section of $M$, a contradiction.
Thus $u\notin K^p$, and the torsor has the degree-$p$ field
\[
 E=K(u^{1/p}).
\]
Its locally free algebra injects into that generic field,
so the torsor is integral. Its normalization equals the
normalization of $X$ in $E$, since integrality is transitive.
This normalization is finite: apply the Nagata finiteness
criterion to $\Spec E\to X$
\cite[Lemma~29.54.14, Tag~0BAK]{Stacks}.

The image of $\omega|_Y$ in $\Omega^1_{Y/k}$ is zero.
A trivialization of $M|_Y$ has $p$th power differing from
the chosen trivialization by a global unit on the proper
integral curve $Y$, hence by an element of $k^*$.
Its logarithmic differential is zero.
The conormal sequence is
\[
 0\longrightarrow\OO_Y(-Y)\longrightarrow
 \Omega^1_{X/k}|_Y\longrightarrow\Omega^1_{Y/k}
 \longrightarrow0.
\]
Since $Y^2>0$, the left-hand line bundle has no nonzero
global section. Hence $\omega|_Y$ vanishes in the full
restricted cotangent bundle. There is consequently a
nonzero section
\[
 \theta\in H^0(X,\Omega^1_{X/k}(-Y)),
 \qquad s_Y\theta=\omega,
\]
where $s_Y$ is the canonical section of $\OO_X(Y)$.

Write the ruled function field as $k(C)(t)$, with $C$ a
smooth projective integral curve, and replace $R$ by
$R_0=C\times\PP^1$. Resolve $R_0\dashrightarrow X$ by
point blowups, obtaining $S\to R_0$ and $g:S\to X$
\cite[Tag~0C5H]{Stacks}. Point blowups of a regular surface
remain regular \cite[Lemma~54.3.2, Tag~0AGP]{Stacks}; since
$k$ is perfect, $S$ is smooth. It is projective, and
$[k(S):K]=p$.

Only finitely many ruling fibres contain images of the
blowup centres. A general fibre $F$ of $S\to C$ is therefore
an unmodified projective line, with
\[
 \Omega^1_{S/k}|_F\simeq\OO_{\PP^1}\oplus\OO_{\PP^1}(-2).
\]
The restriction $g|_F$ is nonconstant: a general ruling
fibre meets the dense quasi-finite locus of $g$ and cannot
be contracted. Ampleness of $Y$ gives
$e_F=\deg(g^*\OO_X(Y)|_F)>0$, without any separability
assumption on this curve map.

Pullback by the morphism $g$ gives a regular section
\[
 \eta\in H^0(S,\Omega^1_{S/k}(-g^*Y)).
\]
It restricts on each general ruling fibre to a section of
\[
 \OO_{\PP^1}(-e_F)\oplus\OO_{\PP^1}(-e_F-2),
\]
and is therefore zero there. These fibres fill a dense open
of the integral surface $S$, so $\eta=0$ and $g^*\omega=0$.
This uses positivity only on the general fibre; $g^*Y$
need not be ample on $S$.

At the generic point, $du/u=0$ in $\Omega^1_{k(S)/k}$.
The same perfect-field criterion gives a $p$th root of $u$
in $k(S)$, and hence a $K$-embedding $E\hookrightarrow k(S)$.
Both extensions have degree $p$, so they are equal.
The $p$th root is unique, and changing the frame or the
trivialization changes $u$ only by a $p$th power; thus this
is the distinguished extension defined by $M$.

The resulting birational map $S\dashrightarrow N$ over $X$
extends to a morphism. For an affine open $U=\Spec B$ of
$X$, let $B'$ be the integral closure of $B$ in $E$.
Every $b'\in B'$ is integral over every local ring
$\OO_{S,s}$ for $s\in g^{-1}(U)$, and belongs to it by
normality. These regular functions give
$g^{-1}(U)\to\Spec B'$ and glue over $X$.
The resulting map is birational by the field equality and
proper by its closed graph and properness of $g$.
Thus $N$ has a ruled function field.

For any smooth projective model $S'$ of that field, resolve
the birational map $R_0\dashrightarrow S'$ in the same way.
The resulting ruled blowup has no positive pluricanonical
sections: its canonical bundle has degree $-2$ on every
general unmodified ruling fibre. Pullback by the birational
morphism injects the pluricanonical sections of $S'$ into
those of the blowup. Hence every positive plurigenus of
$S'$ is zero.

If $X$ itself were birationally ruled, the argument with
$\theta$ could instead be applied to a birational map
$R_0\dashrightarrow X$. It would give $g^*\omega=0$, although
the map of function fields and hence of their differentials
is an isomorphism. This contradicts $\omega\ne0$.

Finally write $L=k(N)=A(t)$. The purely inseparable
degree-$p$ extension $L/K$ has exponent one, so
$L^p\subset K$ and $A^p\subset A\cap K\subset A$.
Since $A$ is the function field of a curve over the perfect
field $k$, one has $[A:A^p]=p$. For example, a separating
generator $x$ gives $[A:k(x)]=[A^p:k(x^p)]$, and the tower
through $k(x^p)\subset k(x)$ proves the assertion.
If $A\subset K$, then
$A(t^p)\subset K\subset A(t)$ and the degree-$p$ equality
forces $K=A(t^p)$, contradicting non-ruledness of $X$.
Thus $A\cap K=A^p$. The strict intermediate compositum
$KA$ in the prime-degree extension $L/K$ is all of $L$.
\end{proof}

The field intersection does not assert that the rational ruling
map on $X$ extends everywhere, or that its fibres are smooth.
The normalization may be singular away from the inverse image
of $Y$; resolving those singularities changes an ample pullback
into a divisor having degree zero on exceptional curves.


\begin{corollary}\label{cor:very-ample-equality-almost-ruled}
Let $X$ be a smooth projective integral surface over an
algebraically closed field of characteristic $p>0$, and let
$Y\subset X$ be a smooth integral very ample divisor. If a
nontrivial line bundle $M$ of order $p$ satisfies
\[
 M|_Y\simeq\OO_Y,\qquad 2g(Y)-2=pY^2,
\]
then $X$ is almost ruled: a ruled surface admits a dominant
purely inseparable rational map of degree $p$ to $X$.
The normalization of the $\mu_p$-torsor defined by $M$ is
birationally ruled, with the same extension of $k(X)$ as
every such degree-$p$ ruled cover. Its smooth projective
models have Kodaira dimension $-\infty$, whereas $X$ is
not birationally ruled.
\end{corollary}

\begin{proof}
The nonnegative terms in the identity of
Proposition~\ref{prop:normalized-mup-conductor} force $a=1$
and both defects to vanish. If $X$ were not almost ruled,
Proposition~\ref{prop:normalized-mup-ekedahl} would give
$p=2$, $\chi(\OO_X)=1$, $\omega_X^2\simeq\OO_X(2Y)$
and $1\leq Y^2\leq9$. In particular $K_X\equiv Y$.
Apply Proposition~\ref{prop:very-ample-numerical-obstruction}
with $H=Y$ and its smooth member $Y$ to obtain a contradiction.
The conclusions about the distinguished cover and non-ruledness
of $X$ now follow from
Proposition~\ref{prop:distinguished-ruled-cover}.
\end{proof}

This excludes the K3 alternative when the divisor is very ample.
The distinguished normal cover in the almost-ruled case is therefore
birationally ruled. The following subsections analyze its positive-genus
and rational rulings, and Theorem~\ref{thm:strict-picard-all} excludes
the smallest very-ample equality in every positive characteristic. For a merely
ample divisor, the K3 existence question stated above remains.

\subsection{Fibrations with positive-genus ruling base}
\label{subsec:positive-genus-picard-fibration}

When the ruling curve of the distinguished normalized cover has positive
genus, its Frobenius twist is a base of a fibration on the smooth quotient.
The resulting morphism on the normalization controls its singularities
and the arithmetic genus of the quotient fibres. We first record a
consequence of the divided differential that applies to every conductor
equality.

\begin{proposition}\label{prop:equality-rational-curves}
Let $X$ be a smooth projective integral surface over an algebraically
closed field $k$ of characteristic $p>0$. Let $Y\subset X$ be a smooth
integral ample divisor, and let $M\ne\OO_X$ have order $p$ with
$M|_Y\simeq\OO_Y$. Suppose that
\[
 2g(Y)-2=apY^2
\]
for the conductor integer $a\geq1$ of
Proposition~\ref{prop:normalized-mup-conductor}.
Every smooth rational curve $E\subset X$ satisfies
\[
 E^2\leq-pa(Y\cdot E).
\]
In particular $X$ has no $(-1)$-curve.
\end{proposition}

\begin{proof}
Put $m=pa$. Proposition~\ref{prop:normalized-mup-singularity-defect}
gives a nonzero section
\[
 \theta\in H^0\bigl(X,\Omega^1_{X/k}(-mY)\bigr)
\]
whose zero scheme is finite. Let $e=Y\cdot E$, which is a positive
integer by ampleness. The twisted conormal sequence of $E$ is
\[
 0\longrightarrow\OO_E(-E-mY)\longrightarrow
 \Omega^1_{X/k}(-mY)|_E\longrightarrow
 \Omega^1_{E/k}\otimes\OO_E(-mY)\longrightarrow0.
\]
Its last term is $\OO_{\PP^1}(-2-me)$ and has no nonzero section.
Consequently $\theta|_E$ comes from a section of
$\OO_E(-E-mY)$. That section is nonzero, since otherwise the zero
scheme of $\theta$ would contain $E$. Its degree is therefore
nonnegative, giving $-E^2-me\geq0$. If $E^2=-1$, this would imply
$-1\leq-pae\leq-2$, a contradiction.
\end{proof}

We will use a genus bound that holds for singular integral curves as
well as smooth ones.

\begin{lemma}\label{lem:very-ample-curve-genus-bound}
Let $F$ be an integral projective curve over an algebraically closed
field, and let $L$ be a very ample invertible sheaf of degree $b$ on $F$.
Then
\[
 p_a(F)\leq\frac{(b-1)(b-2)}2.
\]
\end{lemma}

\begin{proof}
Embed $F$ by $L$, and let $R$ be its standard graded homogeneous
coordinate ring. Choose a hyperplane equation $\ell$ that does not
vanish identically on $F$. Multiplication by $\ell$ is injective because
$R$ is a domain. Put $Q=R/(\ell)$, and let $T$ be the saturated
homogeneous coordinate ring of the finite hyperplane section
$\Gamma=F\cap V(\ell)$, whose length is $b$. There is a graded
surjection $Q\to T$. Writing $H$ for dimensions of graded pieces, we get
\[
 H_R(n)-H_R(n-1)=H_Q(n)\geq H_T(n)\qquad(n\geq1).
\]

Choose a linear form $z$ not vanishing at any point of $\Gamma$.
It is a nonzerodivisor on $T$: its kernel is zero after sheafification
on $\Gamma$, and saturation removes the irrelevant-ideal torsion.
Hence
\[
 H_T(n)-H_T(n-1)=\dim_k(T/zT)_n\geq0.
\]
The algebra $T/zT$ is standard graded. If one of its graded pieces is
zero, all later pieces are zero, since the algebra is generated in
degree one. Thus $H_T$ increases strictly until it reaches its eventual
value $b$. Since $H_T(0)=1$, this proves
\[
 H_T(n)\geq\min(n+1,b)\qquad(n\geq0).
\]
For $n$ sufficiently large the Hilbert polynomial of $F$ gives
$H_R(n)=bn+1-p_a(F)$, whereas $H_R(0)=1$. Summing the first inequality
yields
\[
 \begin{split}
 bn+1-p_a(F)
 &\geq1+\sum_{j=1}^{n}\min(j+1,b)\\
 &=bn+1-\frac{(b-1)(b-2)}2,
 \end{split}
\]
which is the required bound.
\end{proof}

\begin{proposition}\label{prop:positive-genus-picard-fibration}
Let $X$ be a smooth projective integral surface over an algebraically
closed field $k$ of characteristic $p>0$, let $Y\subset X$ be a smooth
integral ample divisor, and let $M\ne\OO_X$ have order $p$ with
$M|_Y\simeq\OO_Y$. Suppose that
\[
 2g(Y)-2=pY^2.
\]
Let $\pi:N\to X$ be the normalization of the $\mu_p$-torsor defined
by a trivialization of $M^p$, and put $D=\pi^*Y$.
Assume that $N$ is birationally ruled, and fix a presentation
\[
 k(N)=k(C)(t)
\]
with $C$ a smooth projective integral curve of genus at least one.
Put $B=C^{(p)}$, and denote relative Frobenius by
$F_{C/k}:C\to B$.

There are projective flat morphisms $f:X\to B$ and $h:N\to C$ such that
\[
 f\pi=F_{C/k}h,
\]
and $N$ is the normalization of the integral surface $X\times_B C$.
The singularities of $N$ are rational double points. Moreover,
\[
 h_*\OO_N=\OO_C,\qquad R^ih_*\OO_N=0\quad(i>0),
 \qquad \chi(\OO_N)=1-g(C).
\]

There is a nonempty open subset $C_0\subset C$ over which
$N_{C_0}\simeq C_0\times\PP^1$ as a $C_0$-scheme.
For corresponding general closed points
$c\in C$ and $b_0=F_{C/k}(c)\in B$, let
$F_N=h^{-1}(c)$ and $F_X=f^{-1}(b_0)$. Then $F_X$ is integral and
geometrically unibranch, and
$\pi|_{F_N}:F_N\simeq\PP^1\to F_X$ is its finite radicial
normalization. One has
\[
 \pi^*F_X=pF_N,\qquad \pi_*F_N=F_X.
\]
Writing $b=Y\cdot F_X=D\cdot F_N$, one obtains
\[
 2p_a(F_X)=(p-1)b,\qquad b>0.
\]
In characteristic two, $b$ is even. In particular the general fibre
$F_X$ is singular. If $Y$ is very ample, then $b\geq p+2$.
\end{proposition}

\begin{proof}
The equality conclusions of
Propositions~\ref{prop:normalized-mup-conductor} and
\ref{prop:normalized-mup-singularity-defect} give that $N$ is normal
and Gorenstein, $\pi$ is finite flat and radicial of degree $p$, and
\begin{equation}\label{eq:positive-genus-canonical}
 \omega_N\simeq\pi^*\omega_X\otimes\OO_N\bigl(-(p-1)D\bigr).
\end{equation}
Put $A=k(C)$, $K=k(X)$ and $L=k(N)=A(t)$.
A smooth ruled model of this assumed function field maps rationally
and purely inseparably to $X$ with degree $p$.
Thus Proposition~\ref{prop:distinguished-ruled-cover} applies and gives
\[
 A\cap K=A^p,\qquad AK=L,
\]
and a proper birational morphism $r:S\to N$, where $S$ is obtained
from $C\times\PP^1$ by finitely many point blowups. Write $q:S\to C$
for the resulting ruling morphism.

\emph{Extension of the map on the quotient.}
The inclusion $A^p\subset K$ defines a rational map $X\dashrightarrow B$.
Resolve it by point blowups $\rho:X'\to X$, with resulting morphism
$f':X'\to B$ \cite[Tag~0C5H]{Stacks}. These blowups remain regular,
and hence smooth over $k$ \cite[Tags~0AGP and~0B8X]{Stacks}.
Each nontrivial fibre of $\rho$ has connected reduced support whose
components are smooth rational curves. Indeed this holds for one point
blowup, and successive blowups insert projective lines meeting the
strict transforms of the components through the centre.

There is no nonconstant morphism $\PP^1\to B$. To check this also for
inseparable maps, suppose that such a map exists. If its map on
differentials is zero, its function-field image lies in $k(t^p)$;
replace $k(t)$ by that rational function field and repeat.
This process stops: a fixed nonconstant rational function has a zero
or pole of finite nonzero order and cannot belong to $k(t^{p^e})$
for every $e$. We obtain a map $\varphi:\PP^1\to B$ with nonzero
differential. The resulting nonzero map of line bundles
$\varphi^*\omega_B\to\omega_{\PP^1}$ would give an effective divisor
of degree
\[
 -2-\deg(\varphi)\bigl(2g(B)-2\bigr)<0,
\]
which is impossible because $g(B)=g(C)\geq1$.
Thus $f'$ is constant on each component of a $\rho$-fibre;
connectedness makes the values agree on the whole fibre.

Let $\Gamma$ be the reduced closed image of
$(\rho,f'):X'\to X\times B$. It is integral, and its projection to $X$
is proper and birational. Every fibre of this projection has one point
in its underlying set, because it is the image of a $\rho$-fibre.
The projection is therefore finite \cite[Tag~02LS]{Stacks}, and is an
isomorphism because $X$ is normal \cite[Tag~0AB1]{Stacks}.
Its second projection gives the required morphism $f:X\to B$.

\emph{The morphism on the normalization.}
The inclusion $A\subset L$ defines a rational lift to $C$.
Let $V=\Spec R\subset B$ be affine, and write
$F_{C/k}^{-1}(V)=\Spec R'$, with $R'$ finite over $R$.
Every element of $R'$ lies in $L$ and is integral over each local ring
of $N$ above $V$. It belongs to those local rings by normality, hence
is regular there. These maps glue to $h:N\to C$, and
$f\pi=F_{C/k}h$ because the equality holds at the generic point.
Likewise $hr=q$, since both morphisms to the separated curve $C$
agree generically.

Relative Frobenius on the smooth curve is finite flat of degree $p$:
its generic degree is $[A:A^p]=p$, and its finite local modules over
the DVRs of $B$ are torsion free. If $\alpha\in A\setminus A^p$, then
$A=A^p(\alpha)$. The polynomial $T^p-\alpha^p$ has no root in $K$,
since a root would be the unique element $\alpha$ and contradict
$A\cap K=A^p$. Thus
\[
 K\otimes_{A^p}A=K(\alpha)=L.
\]
The finite locally free algebra of $X\times_B C$ injects into this
generic field, so the fibre product is integral. Transitivity of
integrality shows that its normalization is the normalization of
$X$ in $L$, namely $N$.

Both $f$ and $h$ are projective, by the closed-graph construction from
the projective $k$-surfaces. They are flat: the local rings of their
integral dominant sources are torsion free over the local DVRs of
the base curves, and a torsion-free module over a DVR is flat.
Their closed fibres are Cartier divisors of pure dimension one.
Also
\begin{equation}\label{eq:positive-genus-functions}
 f_*\OO_X=\OO_B,\qquad h_*\OO_N=\OO_C.
\end{equation}
Indeed proper-cohomology finiteness makes these pushforwards finite
algebras on the base \cite[III, Theorem~3.2.1]{EGA}.
An element of $A(t)$ algebraic over $A$ belongs to $A$; consequently
the algebraic closure of $A^p$ in $K$ is $A\cap K=A^p$.
Thus the two finite algebras lie in their respective base function
fields and equal the base structure sheaves by normality.

\emph{Vanishing and rational singularities.}
For a point blowup $\beta$ of a smooth surface, one has
$\beta_*\OO=\OO$ and $R^i\beta_*\OO=0$ for $i>0$
\cite[Tag~0AGP, Lemma~54.3.4]{Stacks}.
The vanishing can be seen directly. For a regular local ring $R$
with parameters $x,y$, the two blowup charts have rings
$R[y/x]$ and $R[x/y]$, and their intersection has ring
$R[y/x,x/y]$. Every Laurent polynomial in $y/x$ is a sum of one
with nonnegative powers and one with negative powers. The degree-one
\v{C}ech differential is therefore surjective, and the cover has no
higher \v{C}ech term. The degree-zero pushforward follows from
normality and proper-cohomology finiteness.
The same two-chart calculation on $\PP^1_R$ gives $H^0=R$ and
vanishing higher cohomology. Successive applications of Leray yield
\[
 q_*\OO_S=\OO_C,\qquad R^iq_*\OO_S=0\quad(i>0).
\]

We also need $R^ih_*\OO_N=0$ for $i\geq2$, without any smoothness
assumption on the special fibres. Fix a closed point $c\in C$ and a
projective embedding $N\hookrightarrow\PP^m_C$.
Choose a hyperplane $H_1$ containing no component of $N_c$, and then
a hyperplane $H_2$ avoiding the finite scheme's support
$H_1\cap N_c$. These choices are possible over the infinite field $k$.
Lift their equations to a neighbourhood of $c$. The closed subset
$N\cap H_1\cap H_2$ has proper image in $C$ missing $c$.
Delete that image and shrink to an affine neighbourhood $V$.
Then $N_V\setminus H_1$ and $N_V\setminus H_2$ form an affine cover
of $N_V$, with affine intersection. The \v{C}ech complex gives
$H^i(N_V,\OO_N)=0$ for $i\geq2$. Such neighbourhoods prove the
claimed higher-direct-image vanishing.

Normality gives $r_*\OO_S=\OO_N$. The coherent sheaf
$\mathcal T=R^1r_*\OO_S$ is supported on finitely many points:
the proper birational morphism is an isomorphism over an open
containing every codimension-one point of $N$
\cite[Tag~0BFP]{Stacks}.
The low-degree Leray sequence for $q=hr$ is
\[
 0\longrightarrow R^1h_*\OO_N\longrightarrow R^1q_*\OO_S
 \longrightarrow h_*\mathcal T\longrightarrow R^2h_*\OO_N.
\]
The two established vanishings give $R^1h_*\OO_N=0$ and
$h_*\mathcal T=0$. The latter implies $\mathcal T=0$.
Indeed the closed subscheme cut out by
$\operatorname{Ann}(\mathcal T)$ is a zero-dimensional finite-type
$k$-scheme, hence is finite over $k$ and over $C$.
Pushforward along this finite morphism is faithful, being restriction
of scalars on affine opens.

Thus $N$ has rational singularities. To compare with the definition
requiring vanishing on every normal modification
\cite[Tag~0B4N]{Stacks}, localize at a singular closed point of $N$.
Regular points already have rational singularities.
The vanishing just proved gives $H^1(\OO)=0$ on the localized smooth
model. Any other normal modification is dominated by point blowups
of that model, by resolving the rational map to it with Tag~0C5H.
These blowups preserve the vanishing, and Leray injects the
$H^1(\OO)$ of the normal modification into that of the dominating
smooth model. It too is zero.

The local rings at singular points of $N$ are two-dimensional
normal Gorenstein algebraic local rings. They are Nagata and have
normal completions, by excellence
\cite[IV, 7.8.3(ii), (iii), (v), (vi)]{EGA}.
The classical rational-Gorenstein criterion therefore makes them
rational double points in every characteristic
\cite[Sections~54.9 and~54.12, Tags~0B4V and~0BGB]{Stacks}.
Finally, \eqref{eq:positive-genus-functions} and the vanishing of all
higher direct images of $h$ give
\[
 H^i(N,\OO_N)\simeq H^i(C,\OO_C)\quad(i\geq0),
\]
and hence $\chi(\OO_N)=1-g(C)$.

\emph{The general fibres and their multiplicities.}
Remove from $C$ the images of the finitely many points of $N$ over
which $r$ is not an isomorphism, and the images of all blowup centres
in $S\to C\times\PP^1$. Over the remaining nonempty open $C_0$,
both maps are isomorphisms, giving
$N_{C_0}\simeq C_0\times\PP^1$.

We check that the corresponding general fibre on $X$ is reduced.
Choose $v\in A^p\setminus(A^p)^p$. If $v$ were a $p$th power in $K$,
its unique root would belong to $A\cap K=A^p$, a contradiction.
Thus $dv\ne0$ in $\Omega^1_{K/k}$ \cite[Tag~031W]{Stacks}.
The map $f^*\Omega^1_{B/k}\to\Omega^1_{X/k}$ is consequently a
split injection on a dense open subset of $X$. The differential
criterion for morphisms between smooth $k$-schemes makes $f$ smooth
on this open \cite[IV, Theorem~17.11.1]{EGA}. Its image contains a
nonempty open of $B$.

Choose $b_0$ in that open and under $C_0$, and let $c$ be its unique
inverse image under Frobenius. Since $\pi$ is a universal
homeomorphism and $f\pi=F_{C/k}h$, the support of $F_X$ is identified
with the irreducible support of $F_N=\PP^1$.
Moreover $F_X$ meets the smooth locus of $f$, so it is generically
reduced. It is a Cartier divisor in a regular surface, hence is
Cohen--Macaulay and has no embedded associated points.
Generic reducedness then implies reducedness: a nonzero nilpotent
ideal would have an associated prime, necessarily a minimal prime,
contradicting its zero localization there. Thus $F_X$ is integral,
and it is geometrically integral over the algebraically closed field.

Pulling back the point $b_0$ in the commuting square gives the
Cartier-divisor identity
\[
 \pi^*F_X=h^*F_{C/k}^*[b_0]=pF_N.
\]
The residue degree along this divisor is one. Indeed, localizing at
the generic point of $F_X$ gives a finite free extension of DVRs of
rank $p$, with a unique maximal ideal upstairs and ramification
index $p$. Modulo the downstairs uniformizer, its dimension over
the downstairs residue field is both $p$ and $p$ times the residue
degree. Thus $\pi_*F_N=F_X$, and the finite morphism
$F_N\to F_X$ is birational. As $F_N$ is normal, it is the
normalization of $F_X$. The restriction remains radicial, proving
geometric unibranchness.

\emph{Arithmetic genus and the conditional degree bound.}
The projection formula gives
\[
 b=Y\cdot F_X=D\cdot F_N>0.
\]
This equality also follows directly on the singular curve $F_X$.
For its normalization $\nu$, the cokernel of
$\OO_{F_X}\to\nu_*\OO_{\PP^1}$ has finite length. Tensoring by an
invertible sheaf preserves that length; subtracting the two
Euler-characteristic equalities shows that normalization preserves
the degree of an invertible sheaf.

The curve $F_N$ lies in the regular locus of $N$, and $F_N^2=0$.
Adjunction therefore gives $K_N\cdot F_N=-2$.
Using \eqref{eq:positive-genus-canonical} and the residue degree one,
we obtain
\[
 -2=K_X\cdot F_X-(p-1)b.
\]
The integral Cartier fibre $F_X$ also has square zero. Adjunction
for this possibly singular curve gives
\[
 2p_a(F_X)-2=(K_X+F_X)\cdot F_X=K_X\cdot F_X;
\]
the Cartier adjunction formula applies without a smoothness
assumption on the curve \cite[Tags~0AA4 and~0A9X]{Stacks}.
Consequently $2p_a(F_X)=(p-1)b$. It is positive, whereas the
normalization has genus zero, so $F_X$ is singular.
When $p=2$, this genus formula also forces $b$ to be even.

Now suppose that $Y$ is very ample.
Its restriction $\OO_X(Y)|_{F_X}$ is very ample of degree $b$.
Lemma~\ref{lem:very-ample-curve-genus-bound} implies
\[
 (p-1)b\leq(b-1)(b-2),
 \qquad b^2-(p+2)b+2\geq0.
\]
The last quadratic is negative at both endpoints of the interval
$1\leq b\leq p+1$, and by convexity it is negative throughout that
interval. Hence $b\geq p+2$ under this additional hypothesis.
\end{proof}

\begin{remark}\label{rem:positive-genus-higher-conductor}
The fibration and rational-singularity arguments use the compatible
ruled model and the field equalities of
Proposition~\ref{prop:distinguished-ruled-cover}; they do not use
$a=1$ or very ampleness. Suppose more generally that $X,Y,M$ have
the same smoothness, projectivity and ampleness hypotheses, that
$a\geq1$ is the actual conductor integer of
Proposition~\ref{prop:normalized-mup-conductor}, and that
$2g(Y)-2=apY^2$. Assume that the distinguished normalization $N$
is birationally ruled with a positive-genus ruling base.
Then the same arguments apply, with
\[
 K_N=\pi^*K_X-a(p-1)D,
\]
and give
\[
 2p_a(F_X)=a(p-1)b.
\]
For $p=2$, the product $ab$ is even.
If $Y$ is in addition very ample, the embedded-curve genus bound gives
\[
 b\geq a(p-1)+3.
\]
The ruledness hypothesis is independent of the numerical equality.
It is also supplied if $X$ admits a degree-$p$ purely inseparable
ruled rational cover, by
Proposition~\ref{prop:distinguished-ruled-cover}.
\end{remark}

These conclusions apply to a merely ample divisor under the explicit
ruledness hypothesis. They do not identify $M$ with a pullback from
$B$, or make $h$ a projective-line bundle over all of $C$.
The rational-curve inequality in
Proposition~\ref{prop:equality-rational-curves} does not apply to the
singular general fibre $F_X$, even though its normalization is
$\PP^1$. The strict very ample theorem does not exclude these
merely ample configurations.

\subsection{Euler characteristics and the genus-zero ruling}
\label{subsec:equality-torsor-euler}

The group action in Corollary~\ref{cor:normalized-equality-torsor}
gives a global filtration of the normalized cover's direct-image algebra.
Its quotients determine the Euler characteristic. In the genus-zero
ruling case, this also gives a lower bound for the defect of rational
singularities and an obstruction to extending the distinguished ruling.
The canonical-section argument in the following subsection turns these
defects into exclusions and a characteristic-three residual case.

Degree filtrations from affine coordinate changes are standard for
additive torsors; see the $\alpha_{L,a}$-torsor construction in
\cite[equation~(0.2.17), pp.~38--40]{CDL}. The group used here has a
deformed multiplicative coproduct. We give the descent calculation
explicitly, including its torsion character twist.

\begin{proposition}\label{prop:deformed-torsor-degree-filtration}
Let $X$ be a scheme of characteristic $p>0$, let $L$ be an invertible
$\OO_X$-module, and let $s\in H^0(X,L)$. Define $G_s$ as in
Corollary~\ref{cor:smooth-torsor}: if $s_j=g_{ij}^{-1}s_i$ in two
frames of $L$, its local Hopf algebras are
\[
 \OO_X[u_i]/(u_i^p),\qquad
 \Delta(u_i)=u_i\otimes1+1\otimes u_i+s_i u_i\otimes u_i,
 \qquad u_j=g_{ij}u_i.
\]
For every finite flat $G_s$-torsor $\pi:Q\to X$, there is a filtration
by finite locally free submodules
\[
 0=F_{-1}\subset F_0=\OO_X\subset F_1\subset\cdots
 \subset F_{p-1}=\pi_*\OO_Q
\]
with $\operatorname{rank}F_i=i+1$. There is an invertible module $P$ such
that
\begin{equation}\label{eq:deformed-torsor-graded}
 F_i/F_{i-1}\simeq L^{\otimes i}\otimes P^{\otimes i}
 \quad(0\leq i<p),\qquad
 P^{\otimes p}\simeq\OO_X,\qquad
 P|_{Z(s)}\simeq\OO_{Z(s)}.
\end{equation}
Here $Z(s)$ is the scheme-theoretic zero subscheme.
The filtration is intrinsic to the given torsor action.
\end{proposition}

\begin{proof}
The gluing rule identifies the underlying module of the group algebra
with
\[
 \mathcal A=(G_s\to X)_*\OO_{G_s}
 =\bigoplus_{j=0}^{p-1}L^{\otimes j}.
\]
Thus its degree submodules
$V_i=\bigoplus_{j=0}^iL^{\otimes j}$ are global subbundles, with
$V_i/V_{i-1}\simeq L^{\otimes i}$.
The local operation $x*y=x+y+s_i xy$ has inverse
$-x/(1+s_i x)$. The denominator is a unit because $x^p=0$.
The formula
\[
 \vartheta:G_s\longrightarrow\mu_p,\qquad
 \vartheta(u_i)=1+s_i u_i
\]
defines a global character: it is unchanged on overlaps,
is multiplicative for the group operation, and has $p$th power one.

In a fixed frame of $L$, translation by a section $v$ sends
\[
 u\longmapsto v+(1+sv)u.
\]
This substitution preserves every degree submodule $V_i$.
On its degree-$i$ quotient the substitution is multiplication by
$(1+sv)^i=\vartheta(v)^i$, a unit. The inverse translation gives
the inverse linear map. These assertions hold over arbitrary
base rings, including nonreduced ones.

Take an fpqc cover $U\to X$ on which the torsor has a section;
one may take $U=Q$. The section identifies $Q_U$ with $(G_s)_U$,
and the descent datum on
$(\pi_*\OO_Q)|_U\simeq\mathcal A|_U$ is given by translations.
It therefore preserves all the $V_i|_U$, and the cocycle condition
restricts to these submodules. Effective descent and full faithfulness
for quasi-coherent modules give their descents $F_i$, the inclusions
between them, and their maps to $\pi_*\OO_Q$
\cite[Proposition~35.5.2, Tag~023T]{Stacks}.
Injectivity and exactness can be checked after faithfully flat pullback.
Consequently the quotients pull back to $V_i/V_{i-1}$.
The modules $F_i$ and all these quotients are finite locally free,
because that property descends under an fpqc cover
\cite[Lemma~35.7.6, Tag~05B2]{Stacks}.
In particular the construction has no exceptional finite-length
quotient at a point of $X$.

After factoring out the ordinary $L$-transition, the degree-one
quotient has the scalar descent cocycle $\vartheta(v)$.
Let $P$ be its descended line bundle; equivalently,
$P=(F_1/F_0)\otimes L^{-1}$.
The leading coefficient in degree $i$ is the $i$th power of the
degree-one coefficient. This gives the first isomorphism in
\eqref{eq:deformed-torsor-graded}. The scalar cocycle has $p$th
power one, so $P^{\otimes p}\simeq\OO_X$.
On $Z(s)$ the scalar cocycle is itself one, giving the last
isomorphism in \eqref{eq:deformed-torsor-graded}, including when
$Z(s)$ is nonreduced. Constants are fixed by every translation,
so $F_0$ is the unit submodule $\OO_X$.
Changing the chosen torsor section is another translation,
which preserves the flag. The resulting filtration is therefore
independent of these choices.
\end{proof}

There is also an exact identification of the torsion twist.
Suppose $Q\to T$ is equivariant for $\vartheta:G_s\to\mu_p$,
where $T\to X$ is a $\mu_p$-torsor. Let $W$ be the weight-one
line in $(T\to X)_*\OO_T$, defined locally by the coordinate
$z$ with coaction $z\mapsto z\otimes t$.
A section of $Q$ gives a section of $T$. If two sections of $Q$
differ by $v$, the induced sections of $T$ differ by $\vartheta(v)$.
The weight-one coordinate therefore has the same scalar cocycle
as $P$, and $P\simeq W$ by descent.
For the equality torsor $N\to X$, the normalization map $N\to T$
has this equivariance by the coaction constructed in
Corollary~\ref{cor:normalized-equality-torsor}.
Thus its graded pieces are exactly
$\OO_X(iaY)\otimes W^{\otimes i}$.
Defining $W$ by its weight fixes the convention independently of
whether the original Kummer algebra is written using $M$ or $M^{-1}$.

\begin{corollary}\label{cor:normalized-equality-euler}
Let $X$ be a smooth projective integral surface over an algebraically
closed field of characteristic $p>0$. Let $Y\subset X$ be a smooth
ample integral curve, and let $M\ne\OO_X$ have order $p$ with
$M|_Y\simeq\OO_Y$. Let $\pi:N\to X$ be the normalization of the
$\mu_p$-torsor defined by a chosen trivialization of $M^p$.
Let $a\geq1$ be the conductor integer of
Proposition~\ref{prop:normalized-mup-conductor}, characterized by
the order $pa$ of the logarithmic differential along $Y$.
If $2g(Y)-2=apY^2$, then
\begin{equation}\label{eq:normalized-equality-euler}
 \chi(\OO_N)=p\chi(\OO_X)
 -\frac{p(p-1)a\bigl((p+1)a-3\bigr)}{12}Y^2.
\end{equation}
In particular, if $a=1$, then
\begin{equation}\label{eq:normalized-smallest-equality-euler}
 \chi(\OO_N)=p\chi(\OO_X)
 -\frac{p(p-1)(p-2)}{12}Y^2.
\end{equation}
\end{corollary}

\begin{proof}
Corollary~\ref{cor:normalized-equality-torsor} makes $\pi:N\to X$
a finite flat torsor for $G_{s_a}$, with
$L=\OO_X(aY)$ and $s_a$ its canonical section.
Proposition~\ref{prop:deformed-torsor-degree-filtration}
therefore applies. Its line bundle $P$ has $P^p\simeq\OO_X$ and
$P|_Y\simeq\OO_Y$, since $Y\subset Z(s_a)=aY$.

For the line bundle $P$ on the proper scheme $X$, the function
$n\mapsto\chi(P^{\otimes n})$ is a numerical polynomial over
$\mathbf Q$ \cite[Lemma~33.45.1, Tag~0BEM]{Stacks}.
Because $P^p\simeq\OO_X$, it is periodic with period $p$, so it is
constant: the difference of a
nonconstant polynomial at $n+p$ and $n$ has a nonzero leading
coefficient. Hence
\[
 \chi(P^{\otimes i})=\chi(\OO_X)
 \qquad\text{for every integer }i.
\]
Put $d=Y^2$. For $j\geq1$, the Cartier divisor sequence gives
\[
 0\longrightarrow P^{\otimes i}((j-1)Y)
 \longrightarrow P^{\otimes i}(jY)
 \longrightarrow\OO_Y(jY)\longrightarrow0.
\]
The degree formula on the smooth proper integral curve $Y$ gives
\[
 \chi(\OO_Y(jY))=jd+1-g(Y)
\]
\cite[Definition~33.44.1 and Lemmas~33.44.5 and~33.44.7,
Tag~0AYQ]{Stacks}.
Summing these sequences yields
\[
 \chi\bigl(P^{\otimes i}(nY)\bigr)
 =\chi(\OO_X)+\frac{n(n+1)}2d+n(1-g(Y))
 \qquad(n\geq0).
\]
Set $n=ia$ and use $1-g(Y)=-apd/2$. The $i$th graded quotient
of the filtration has Euler characteristic
\[
 \chi(\OO_X)+\frac d2\bigl(a^2i^2-a(ap-1)i\bigr).
\]
Since $\pi$ is finite, cohomology of $\OO_N$ equals cohomology of
$\pi_*\OO_N$: the affine \v{C}ech complexes on $X$ and their inverse
images on $N$ agree. Additivity along the filtration and
\[
 \sum_{i=0}^{p-1}i=\frac{p(p-1)}2,\qquad
 \sum_{i=0}^{p-1}i^2=\frac{p(p-1)(2p-1)}6
\]
now give
\[
 \begin{split}
 \chi(\OO_N)
 &=p\chi(\OO_X)
 +\frac d2\left(
 a^2\frac{p(p-1)(2p-1)}6
 -a(ap-1)\frac{p(p-1)}2\right)\\
 &=p\chi(\OO_X)
 -\frac{p(p-1)a((p+1)a-3)}{12}d.
 \end{split}
\]
Taking $a=1$ proves the second formula.
\end{proof}

\begin{proposition}\label{prop:genus-zero-equality-defect}
Let $X$ be a smooth projective integral surface over an algebraically
closed field $k$ of characteristic $p>0$.
Let $Y\subset X$ be a smooth integral very ample divisor, and let
$M\ne\OO_X$ have order $p$ with $M|_Y\simeq\OO_Y$.
Suppose $2g(Y)-2=pY^2$, and let $\pi:N\to X$ be the normalization
of its chosen $\mu_p$-torsor.
Choose a ruling presentation
$k(N)=k(C)(t)$ as in
Corollary~\ref{cor:very-ample-equality-almost-ruled}, and suppose
$g(C)=0$.
Let $r:S\to N$ be the proper birational morphism from a point-blowup
model of $C\times\PP^1$ supplied by
Proposition~\ref{prop:distinguished-ruled-cover}. Put
\[
 \mathcal T=R^1r_*\OO_S,\qquad
 \epsilon=\dim_kH^0(N,\mathcal T),\qquad d=Y^2.
\]
Then $\mathcal T$ has finite support and
\begin{equation}\label{eq:genus-zero-equality-defect}
 H^1(N,\OO_N)=0,\qquad
 H^2(N,\OO_N)\simeq H^0(N,\mathcal T),\qquad
 \chi(\OO_N)=1+\epsilon.
\end{equation}
If $p\ne3$, then $\epsilon\equiv-1\pmod p$ and
$\epsilon\geq p-1$. In particular $N$ has a nonrational singularity
outside $D=\pi^*Y$, and the distinguished rational map
$N\dashrightarrow C$ does not extend to a morphism.

If $p=3$, write $d=2e$. Then $\epsilon\equiv-e-1\pmod3$.
Rational singularities, or extension of that rational map, can occur
only if $d\equiv4\pmod6$.
\end{proposition}

\begin{proof}
The field is algebraically closed, so $C\simeq\PP^1$.
A point blowup of a smooth surface has pushforward $\OO$ and
vanishing positive higher direct images
\cite[Lemma~54.3.4, Tag~0AGP]{Stacks}.
The vanishing also follows directly from its two affine charts:
on their intersection, every Laurent monomial in the chart parameter
belongs to one chart ring or the other, so the degree-one
\v{C}ech differential is surjective.
The two-chart computation for $\PP^1_R$ has the same conclusion.
Applying Leray to the point-blowup presentation of $S$ gives
\[
 H^0(S,\OO_S)=k,\qquad
 H^1(S,\OO_S)=H^2(S,\OO_S)=0.
\]

Normality gives $r_*\OO_S=\OO_N$. Indeed proper direct-image
coherence \cite[III, Theorem~3.2.1]{EGA} makes the former a
finite algebra, and its sections over an affine open lie in the
common function field and are integral over the normal base ring.
The proper birational map is an isomorphism near every
codimension-one point of $N$ \cite[Tag~0BFP]{Stacks}.
Thus the coherent sheaf $\mathcal T$ has finite support.

We also have $R^ir_*\OO_S=0$ for $i\geq2$.
To verify this without a smooth-fibre assertion, use a projective
embedding of $S$ over $N$; it exists by the closed graph construction.
Every fibre of $r$ has dimension at most one.
At a closed point of $N$, choose a relative hyperplane containing
no component of its fibre, then another avoiding the finite
intersection of the first with that fibre.
These choices are possible because $k$ is infinite.
Properness allows an affine neighbourhood on which the two
hyperplanes have empty common intersection with $S$.
Their complements give an affine cover with affine intersection,
so its \v{C}ech cohomology vanishes in degrees at least two.
These neighbourhoods prove the higher-direct-image assertion.

The low-degree Leray sequence, using
$H^1(S,\OO_S)=H^2(S,\OO_S)=0$, gives the two cohomology assertions
in \eqref{eq:genus-zero-equality-defect}. Since $\mathcal T$
has finite support, its Euler characteristic is $\epsilon$,
and Leray also gives
$1=\chi(\OO_S)=\chi(\OO_N)-\epsilon$.
A nonzero stalk of $\mathcal T$ witnesses a nonrational singularity:
the local ring is a two-dimensional normal Nagata domain and its
localized smooth model is a normal modification with nonzero
$H^1$ \cite[Definition~54.8.3, Tag~0B4N]{Stacks}.
The equality conclusion of
Proposition~\ref{prop:normalized-mup-conductor}
makes $N$ regular near $D$, so this singularity lies outside $D$.

For $p=2$, \eqref{eq:normalized-smallest-equality-euler} gives
$\chi(\OO_N)=2\chi(\OO_X)$.
For a prime $p\geq5$, the equality $2g(Y)-2=pd$ makes $d=2e$
even, and $(p-1)(p-2)$ is divisible by six.
The correction term in
\eqref{eq:normalized-smallest-equality-euler} is therefore
$p(p-1)(p-2)e/6$, an integer divisible by $p$.
For every $p\ne3$, it follows that
$1+\epsilon=\chi(\OO_N)\equiv0\pmod p$.
Nonnegativity of $\epsilon$ yields $\epsilon\geq p-1$.
For $p=3$, the same Euler formula is
$\chi(\OO_N)=3\chi(\OO_X)-e$, giving
$\epsilon\equiv-e-1\pmod3$.
Thus $\epsilon=0$ requires $e\equiv2\pmod3$,
or $d\equiv4\pmod6$.

Finally suppose the distinguished rational map extends to
$h:N\to C$. Let $q:S\to C$ be the ruling morphism.
Then $hr=q$, since they agree generically and $C$ is separated.
The blowup and projective-line calculations give
$q_*\OO_S=\OO_C$ and $R^iq_*\OO_S=0$ for $i>0$.
The two-hyperplane argument, now for the projective map $h$
with one-dimensional fibres, gives $R^ih_*\OO_N=0$ for $i\geq2$.
The low-degree Leray sequence for $q=hr$ is
\[
 0\longrightarrow R^1h_*\OO_N
 \longrightarrow R^1q_*\OO_S
 \longrightarrow h_*\mathcal T
 \longrightarrow R^2h_*\OO_N.
\]
Hence $R^1h_*\OO_N=0$ and $h_*\mathcal T=0$.
Pushforward of a finite-support coherent sheaf is faithful here:
its scheme-theoretic support is finite over $k$, hence over $C$,
and pushforward is restriction of scalars on affine opens.
Consequently $\mathcal T=0$.
Also $h_*\OO_N=q_*\OO_S=\OO_C$, so $\chi(\OO_N)=1$.
This contradicts the congruence when $p\ne3$.
In characteristic three it requires the stated congruence on $d$.
\end{proof}

For $p=3$, the same calculation gives $\epsilon\geq2$ when
$d\equiv0\pmod6$ and $\epsilon\geq1$ when $d\equiv2\pmod6$.
The remaining residue class is a necessary numerical possibility;
no example is asserted.

Whenever a distinguished ruling $N\dashrightarrow C$ extends,
the last cohomological argument works for any genus and gives
$\chi(\OO_N)=1-g(C)$.
In the smallest equality case,
\eqref{eq:normalized-smallest-equality-euler} consequently gives
\[
 g(C)\equiv1\pmod p\quad(p\ne3),\qquad
 \frac{Y^2}{2}\equiv g(C)-1\pmod3\quad(p=3).
\]
The extension over a positive-genus base follows from
Proposition~\ref{prop:positive-genus-picard-fibration}.
For a genus-zero base the preceding proposition instead requires
indeterminacy when $p\ne3$.
For a very ample divisor, the canonical subsystem in
Corollary~\ref{lem:divided-differential-subsystem} supplies a further
positivity constraint, and Theorem~\ref{thm:strict-picard-all}
excludes the smallest equality in every characteristic.
The merely ample and higher-conductor questions remain at the
scope described in the final section.

\subsection{Rational normalization in characteristic two}

For a very ample curve in characteristic two, a positive-genus ruling is incompatible with the Euler characteristic and intersection form. The following proposition reduces the smallest equality to a rational normalized cover; the canonical argument and Theorem~\ref{thm:strict-picard-all} below exclude that case.

\begin{proposition}\label{prop:characteristic-two-rational-normalization}
Let $X$ be a smooth projective integral surface over an algebraically
closed field $k$ of characteristic two. Let $Y\subset X$ be a smooth
integral very ample divisor, and let $M\ne\OO_X$ have order two with
$M|_Y\simeq\OO_Y$. Suppose that
\[
 2g(Y)-2=2Y^2.
\]
Let $\pi:N\to X$ be the normalization of the $\mu_2$-torsor defined by a
trivialization of $M^2$. Then $N$ is birationally rational. In particular,
$X$ admits a dominant purely inseparable rational map of degree two
from a rational surface. The surface $X$ itself is not birationally
ruled.
\end{proposition}

\begin{proof}
Corollary~\ref{cor:very-ample-equality-almost-ruled} and
Proposition~\ref{prop:distinguished-ruled-cover} give a presentation
$k(N)=k(C)(t)$, with $C$ a smooth projective integral curve, and show
that $X$ is not birationally ruled. We prove that $g(C)=0$.

Suppose instead that $g=g(C)\geq1$, and put $d=Y^2>0$.
Adjunction gives
\[
 K_X\cdot Y=d.
\]
Proposition~\ref{prop:positive-genus-picard-fibration} supplies the
morphisms
\[
 f:X\longrightarrow B=C^{(2)},\qquad h:N\longrightarrow C,
 \qquad f\pi=F_{C/k}h,
\]
and gives
\[
 \chi(\OO_N)=1-g.
\]
For a general integral fibre $F$ of $f$, put $b=Y\cdot F$.
The same proposition gives
\[
 F^2=0,\qquad K_X\cdot F=b-2,\qquad
 2p_a(F)=b,\qquad b\geq4.
\]
In particular $b$ is even. The fibration and fibre identities here
require only ampleness once the ruled cover is given; very ampleness
is used for the last bound.

We first compute the Euler characteristic of the finite flat
rank-two algebra $\pi_*\OO_N$. Its unit is nonzero on every fibre,
so the quotient by $\OO_X$ is an invertible sheaf $Q$.
The inclusion of the Kummer algebra induces $M^{-1}\to Q$.
At codimension-one points outside $Y$ this is an isomorphism.
At the generic point of $Y$, the conductor basis is $1,w$, with
$w=(z-v)/t$, where $t$ is a local equation of $Y$. Modulo the unit
summand, $[z]=t[w]$. The divisor of the map of line bundles is
therefore exactly $Y$, giving
\[
 0\longrightarrow\OO_X\longrightarrow\pi_*\OO_N
 \longrightarrow M^{-1}(Y)\longrightarrow0;
\]
this is also the first assertion of
Proposition~\ref{prop:k3-picard-paracanonical}.
Since $M$ is numerically trivial, Riemann--Roch
\cite[Tag~02UO]{Stacks} and
$K_X\cdot Y=Y^2$ yield
\[
 \chi(M^{-1}(Y))
 =\chi(\OO_X)+\tfrac12Y\cdot(Y-K_X)=\chi(\OO_X).
\]
Finite pushforward and additivity of Euler characteristic thus give
\begin{equation}\label{eq:characteristic-two-cover-euler}
 2\chi(\OO_X)=\chi(\OO_N)=1-g.
\end{equation}
No splitting of the rank-two algebra is required.

For the classical adjunction argument used below, one may first extend
$k$ to an uncountable algebraically closed field and perform the
whole argument there. All the geometric hypotheses and the genus
of $C$ persist. The line bundle $M$ remains nontrivial: a nonzero
section of a torsion line bundle would have an effective divisor
with zero ample intersection and hence would trivialize it;
thus $H^0(X,M)=0$, and this vanishing persists under field extension.
The divided differential and
Proposition~\ref{prop:equality-rational-curves} again exclude
$(-1)$-curves after extension, and
Proposition~\ref{prop:distinguished-ruled-cover} again excludes
birational ruledness. The ruled blowup model and its degree-two
map also base change. A contradiction over the enlarged field
therefore suffices; no descent of ruledness is needed.

Proposition~\ref{prop:equality-rational-curves} excludes every
$(-1)$-curve on $X$. The classical adjunction-termination theorem
therefore makes $K_X$ nef: a negative canonical intersection on a
smooth projective surface without $(-1)$-curves would make that
surface birationally ruled
\cite[\S1, Steps~I--VII, pp.~326--330]{MumfordEnriquesI}.

Here $K_X^2\geq0$ can be checked directly from nefness and
$K_X\cdot Y=d>0$. If $K_X^2=-q<0$, then for rational $t>d/q$
sufficiently close to $d/q$, the class $L=Y+tK_X$ satisfies
\[
 L^2>0,\qquad L\cdot Y>0,\qquad L\cdot K_X<0.
\]
After clearing denominators, take $L$ integral.
For large $n$, Serre duality gives
\[
 H^2(X,\OO_X(nL))=0,
\]
because $(K_X-nL)\cdot Y<0$. Riemann--Roch then gives a nonzero
effective divisor in $|nL|$, contradicting the negative
intersection $K_X\cdot nL$. Thus $K_X^2\geq0$.

Fix a prime $\ell\ne2$, and write $b_i(X)$ for the
$\ell$-adic Betti numbers. The ruled model furnished by
Proposition~\ref{prop:distinguished-ruled-cover} is a smooth surface
$S$ obtained by point blowups of $C\times\PP^1$, with a purely
inseparable morphism $S\to X$ of degree two. Ekedahl's
Betti-number lemma gives
\[
 b_1(X)=b_1(S)=2g
 \quad\text{\cite[Lemma~0.1.5, p.~103]{EkedahlCanonical}}.
\]
The lemma does not require $S\to X$ to be finite. Its proof uses
injectivity of pullback in cohomology and a reverse map after
a Frobenius twist and point blowups.

The classes of $Y$ and $F$ in
$H^2_{\mathrm{\acute et}}(X,\mathbf Q_\ell(1))$
have intersection matrix
\[
 \begin{pmatrix}d&b\\ b&0\end{pmatrix}.
\]
Its determinant is $-b^2\ne0$, so the two classes are independent
and $b_2(X)\geq2$. Poincar\'e duality and algebraic Gauss--Bonnet
\cite[Tag~0FH0]{Stacks} give
\[
 c_2(X)=2-2b_1(X)+b_2(X)
       =2-4g+b_2(X)\geq4(1-g).
\]
On the other hand, Noether's formula and
\eqref{eq:characteristic-two-cover-euler} give
\[
 c_2(X)=12\chi(\OO_X)-K_X^2\leq6(1-g).
\]
For $g>1$ these two inequalities contradict each other.

It remains to consider $g=1$. The same formulas force
\[
 \chi(\OO_X)=0,\qquad c_2(X)=0,\qquad
 K_X^2=0,\qquad b_2(X)=2.
\]
The independent classes of $Y$ and $F$ now form a basis of
$H^2_{\mathrm{\acute et}}(X,\mathbf Q_\ell(1))$.
Write $[K_X]=\alpha[Y]+\beta[F]$. Pairing with $F$ and $Y$
gives
\[
 \alpha b=b-2,\qquad \alpha d+\beta b=d,
 \qquad
 \alpha=1-\frac2b,\qquad \beta=\frac{2d}{b^2}.
\]
Compatibility of the cup product with intersection then yields
\begin{equation}\label{eq:characteristic-two-canonical-square}
 K_X^2=\alpha^2d+2\alpha\beta b
      =d\left(1-\frac4{b^2}\right)>0,
\end{equation}
because $b\geq4$. This is the final contradiction.

Hence $g(C)=0$. Over the algebraically closed field, $C\simeq\PP^1$,
so $k(N)$ is a rational function field in two variables.
The extension $k(N)/k(X)$ is purely inseparable of degree two;
a rational model of $N$ supplies the asserted rational domination
of $X$.
\end{proof}

\begin{remark}\label{rem:characteristic-two-ample-residual}
Suppose that $Y$ is merely ample, with the same characteristic-two
equality and torsion hypotheses, and additionally suppose that $X$
has a purely inseparable ruled rational cover of degree two.
If the resulting ruling curve $C$ has positive genus, the proof
above still gives
\[
 g(C)=1,\qquad
 \chi(\OO_X)=K_X^2=c_2(X)=0,\qquad b_1(X)=b_2(X)=2.
\]
The general-fibre identities give a positive even integer $b=Y\cdot F$.
Equation~\eqref{eq:characteristic-two-canonical-square}, before its
strict inequality, then forces $b=2$ and $p_a(F)=1$.
Moreover, $Y\to B=C^{(2)}$ is separable of degree two. Its degree is
$Y\cdot F=2$; if it were inseparable, it would be a relative
Frobenius followed by an isomorphism and would preserve genus.
That would give $g(Y)=1$, contrary to $g(Y)=Y^2+1\geq2$.
Thus the purely inseparable multisection obstruction in
Proposition~\ref{prop:quasielliptic-multisection} does not already
exclude this configuration.

The genus-zero branch for merely ample $Y$ satisfies the necessary
numerical conditions in Proposition~\ref{prop:rational-normalization-canonical}.
The present proposition concerns the function field of the normalized
cover and does not extend a chosen rational ruling map over
every point of that cover. It supplies no existence example for
the Picard kernel, and does not settle the higher-conductor
or strict-inequality cases.
\end{remark}

\subsection{Canonical sections on a rational normalized cover}
\label{subsec:rational-canonical-equality}

For a rational normalized cover, Serre duality bounds the contribution
of its singularities. The argument requires only ampleness of the divisor
and supplies the rational-cover cases in the strict genus theorem.

\begin{proposition}\label{prop:rational-normalization-canonical}
Let $X$ be a smooth projective integral surface over an algebraically
closed field $k$ of characteristic $p>0$. Let $Y\subset X$ be a
smooth integral ample divisor, and let $M\ne\OO_X$ have order $p$
with $M|_Y\simeq\OO_Y$. Suppose
\[
 2g(Y)-2=pY^2,
\]
and let $\pi:N\to X$ be the normalization of the $\mu_p$-torsor
defined by a trivialization of $M^{\otimes p}$.
Assume that $N$ is birationally rational, and put $d=Y^2$.
Choose a proper birational morphism $r:S\to N$, where $S$ is
obtained from $\PP^1\times\PP^1$ by finitely many point blowups,
and put
\[
 \mathcal T=R^1r_*\OO_S,\qquad
 \epsilon=\dim_k H^0(N,\mathcal T).
\]
Then $\mathcal T$ has finite support and
\begin{equation}\label{eq:rational-normalization-canonical-cohomology}
 \begin{gathered}
 H^1(N,\OO_N)=0,\qquad
 H^2(N,\OO_N)\simeq H^0(N,\mathcal T),\\
 \chi(\OO_N)=1+\epsilon,\qquad
 \epsilon\in\{0,1\}.
 \end{gathered}
\end{equation}
If $\epsilon=1$, then
\[
 \omega_N\simeq\OO_N,\qquad
 \bigl(\omega_X(-(p-1)Y)\bigr)^{\otimes p}\simeq\OO_X,\qquad
 K_X\equiv(p-1)Y.
\]
In particular:
\begin{enumerate}[label=\textup{(\roman*)}]
\item No such birationally rational normalization exists for $p\geq5$.
\item If $p=2$, then $\epsilon=1$ and
\[
 \begin{gathered}
 \chi(\OO_X)=1,\qquad K_X\equiv Y,\qquad
 1\leq d=K_X^2\leq9,\\
 b_1(X)=0,\qquad b_2(X)=10-d.
 \end{gathered}
\]
\end{enumerate}
Here $b_i$ denotes the dimension of $\ell$-adic cohomology for any
prime $\ell\ne p$.
\end{proposition}

\begin{proof}
The conductor identity of
Proposition~\ref{prop:normalized-mup-conductor} makes its integer
$a$ equal to one and its nonnegative defects zero.
That proposition and
Corollary~\ref{cor:normalized-equality-torsor} show that $N$ is an
integral normal projective Gorenstein surface, that $\pi$ is finite
flat and radicial of rank $p$, and that
\[
 D=\pi^*Y
\]
is a smooth integral ample Cartier divisor.
They also give an isomorphism of invertible sheaves
\[
 \omega_N\simeq
 \pi^*\omega_X\otimes\OO_N(-(p-1)D)
 \simeq\pi^*\bigl(\omega_X(-(p-1)Y)\bigr).
\]
This is the absolute dualizing sheaf of $N$.
The formula follows from finite duality and the conductor ideal,
as proved there; it is not a trace calculation in the inseparable
function-field extension.

Adjunction on $Y$ gives $K_X\cdot Y=(p-1)d$, and the projection
formula therefore gives
\[
 K_N\cdot D
 =p\bigl(K_X\cdot Y-(p-1)Y^2\bigr)=0.
\]
Both divisors in this intersection are Cartier.

We include the cohomology calculation for the stated ample
hypotheses. Since $N$ is birationally rational, its function field
is that of $\PP^1\times\PP^1$. Resolve the resulting rational map
to $X$ by point blowups \cite[Tag~0C5H]{Stacks}.
The resulting surface $S$ is smooth and projective.
The morphism $S\to X$ extends to $r:S\to N$: sections of the
normalization algebra are integral over the local rings of $S$,
and belong to them by normality. This is the extension argument
in Proposition~\ref{prop:distinguished-ruled-cover}.
The resulting map $r$ is proper and birational, and $\pi r$ is a
purely inseparable morphism of generic degree $p$.

For a point blowup of a smooth surface, the pushforward of the
structure sheaf is the structure sheaf and the positive higher
direct images vanish \cite[Lemma~54.3.4, Tag~0AGP]{Stacks}.
The two affine charts give this vanishing directly: on their
intersection, each Laurent monomial in the chart parameter belongs
to one of the two chart rings, so the degree-one \v{C}ech
differential is surjective. The analogous calculation for
$\PP^1_R$ and Leray consequently give
\[
 H^0(S,\OO_S)=k,\qquad H^1(S,\OO_S)=H^2(S,\OO_S)=0.
\]

Normality gives $r_*\OO_S=\OO_N$. Indeed proper direct-image
coherence \cite[III$_1$, Theorem~3.2.1]{EGA} makes the pushforward
a finite algebra on every affine open of $N$.
Its sections lie in the common function field and are integral
over the normal affine ring, so they belong to that ring.
Moreover $r$ is an isomorphism near each codimension-one point of
$N$ \cite[Tag~0BFP]{Stacks}.
Thus the coherent sheaf $\mathcal T$ has finite support.

For completeness, $R^ir_*\OO_S=0$ for $i\geq2$ can be checked
without assumptions on smoothness of the fibres.
The closed graph gives a projective embedding of $S$ over $N$,
and every fibre has dimension at most one.
At a closed point of $N$, choose a relative hyperplane containing
no component of its fibre, then another avoiding the finite
intersection of the first with that fibre.
The infinite field $k$ permits these choices.
Properness gives an affine neighbourhood on which the two
hyperplanes have empty common intersection with $S$.
Their complements and their intersection are affine.
The resulting two-open \v{C}ech complex has no cohomology in
degree at least two. These neighbourhoods give the asserted
higher-direct-image vanishing.

The low-degree Leray sequence now reads
\[
 \begin{split}
 0\longrightarrow H^1(N,\OO_N)
 &\longrightarrow H^1(S,\OO_S)
 \longrightarrow H^0(N,\mathcal T)\\
 &\longrightarrow H^2(N,\OO_N)
 \longrightarrow H^2(S,\OO_S).
 \end{split}
\]
The two displayed positive cohomology groups on $S$ vanish.
Since $N$ is proper and integral over $k$, one also has
$H^0(N,\OO_N)=k$. This proves the cohomology and Euler
characteristic assertions in
\eqref{eq:rational-normalization-canonical-cohomology}.

The Gorenstein surface $N$ is Cohen--Macaulay of pure dimension
two. Serre duality for such a proper scheme
\cite[Lemma~48.27.5, Tag~0FVZ, and Remark~48.27.6, Tag~0FW0]{Stacks}
gives
\[
 H^0(N,\omega_N)\simeq H^2(N,\OO_N)^\vee.
\]
If $\epsilon>0$, choose a nonzero section of $\omega_N$.
As $N$ is integral, its zero scheme is an effective Cartier
divisor unless empty. Every nonzero effective divisor has
positive intersection with the ample divisor $D$.
The equality $K_N\cdot D=0$ therefore makes the zero divisor empty.
There cannot be an isolated zero: a nonzero nonunit in a
noetherian normal local domain has a height-one zero by the
principal ideal theorem.
The section is consequently nowhere vanishing, so
$\omega_N\simeq\OO_N$. It follows that
$h^0(N,\omega_N)=h^0(N,\OO_N)=1$, and duality gives $\epsilon=1$.
This proves $\epsilon\in\{0,1\}$.

Suppose $\epsilon=1$, and put $L=\omega_X(-(p-1)Y)$.
The canonical formula gives $\pi^*L\simeq\OO_N$.
The sheaf $\mathcal E=\pi_*\OO_N$ is locally free of rank $p$.
The finite projection formula yields
\[
 \mathcal E\otimes L
 \simeq\pi_*(\pi^*L)
 \simeq\mathcal E.
\]
Taking determinants and cancelling the invertible sheaf
$\det\mathcal E$ gives $L^{\otimes p}\simeq\OO_X$.
This is the scalar-line case of the determinant norm
\cite[Tag~0BD2]{Stacks}, with no separability requirement.
Its degree on every integral curve is zero, whence
$K_X\equiv(p-1)Y$ and $K_X^2=(p-1)^2d$.

Corollary~\ref{cor:normalized-equality-euler} gives
\begin{equation}\label{eq:rational-normalization-canonical-euler}
 1+\epsilon
 =p\chi(\OO_X)-\frac{p(p-1)(p-2)}{12}d.
\end{equation}
For a prime $p\geq5$, the equality $2g(Y)-2=pd$ makes $d=2e$
even. The integer $(p-1)(p-2)$ is divisible by six.
Thus the correction term in
\eqref{eq:rational-normalization-canonical-euler} is
$p(p-1)(p-2)e/6$, an integer divisible by $p$.
It follows that $1+\epsilon\equiv0\pmod p$, contradicting
$\epsilon\in\{0,1\}$.
This proves \textup{(i)}.

For $p=2$, the correction term vanishes, so
$1+\epsilon=2\chi(\OO_X)$.
Hence $\epsilon=1$ and $\chi(\OO_X)=1$, and the preceding
determinant argument gives $K_X\equiv Y$.
The purely inseparable morphism $\pi r:S\to X$ has smooth
projective source and target. Ekedahl's comparison
\cite[Lemma~0.1.5, p.~103]{EkedahlCanonical} gives
$b_1(X)=b_1(S)=0$; finiteness of that morphism is not required
for this assertion.
The class of $Y$ in $H^2_{\mathrm{et}}(X,\mathbf Q_\ell(1))$
is nonzero, since its square has degree $d>0$.
Therefore $b_2(X)\geq1$.
Poincaré duality and Gauss--Bonnet on the smooth surface $X$
\cite[Section~45.7, especially Lemma~45.7.8, Tag~0FH0]{Stacks}
give $c_2(X)=2+b_2(X)\geq3$.
Noether's formula
\cite[Section~42.66, Tag~02UO]{Stacks} then gives
\[
 12=12\chi(\OO_X)=K_X^2+c_2(X)=d+2+b_2(X).
\]
Consequently $1\leq d\leq9$ and $b_2(X)=10-d$, as asserted
in \textup{(ii)}.
\end{proof}

\begin{corollary}\label{cor:characteristic-three-rational-residual}
In the setting of Proposition~\ref{prop:rational-normalization-canonical},
suppose $p=3$.
Then either
\[
 \epsilon=0,\qquad d\equiv4\pmod6,
\]
or
\[
 \begin{gathered}
 \epsilon=1,\qquad d=2,\qquad
 K_X\equiv2Y,\qquad \chi(\OO_X)=1,\\
 K_X^2=8,\qquad c_2(X)=4,\qquad b_2(X)=2.
 \end{gathered}
\]
In the first alternative, every singular point of $N$ is a rational
double point.
\end{corollary}

\begin{proof}
The smallest equality makes $d=2e$ even, and
\eqref{eq:rational-normalization-canonical-euler} becomes
\[
 1+\epsilon=3\chi(\OO_X)-e.
\]
Since $\epsilon\in\{0,1\}$, the case $\epsilon=0$ requires
$e\equiv2\pmod3$, or $d\equiv4\pmod6$.

In this case $\mathcal T$ has finite support and
$H^0(N,\mathcal T)=0$, so $\mathcal T=0$.
Indeed its annihilator defines a zero-dimensional finite-type
$k$-scheme, hence a finite $k$-scheme, and pushforward of its
module to $\Spec k$ is faithful.
Localizing at a singular closed point of $N$ gives a smooth
modification with $H^1(\OO)=0$.
Any other normal modification is dominated by point blowups of
that model, by Tag~0C5H. Those blowups preserve the vanishing,
and the low-degree Leray sequence injects the $H^1(\OO)$ of
the normal modification into that of the dominating smooth model.
Thus $N$ has rational singularities in the all-modifications
sense \cite[Tag~0B4N]{Stacks}.
Its singular local rings are two-dimensional normal Gorenstein
algebraic local rings. Excellence makes them Nagata with normal
completions \cite[IV, 7.8.3(ii), (iii), (v), (vi)]{EGA}.
The rational-Gorenstein criterion therefore makes these
singularities rational double points in every characteristic
\cite[Sections~54.9 and~54.12, Tags~0B4V and~0BGB]{Stacks}.

Suppose $\epsilon=1$.
The canonical and determinant conclusions give
$K_X\equiv2Y$ and $K_X^2=4d$, while the Euler formula gives
$\chi(\OO_X)=(d+4)/6$.
The smooth rational model mapping purely inseparably to $X$
again gives $b_1(X)=0$, and the ample class gives $b_2(X)\geq1$.
Noether's formula and Gauss--Bonnet therefore yield
\[
 3\leq c_2(X)
 =12\chi(\OO_X)-K_X^2
 =2(d+4)-4d
 =8-2d.
\]
The positive even integer $d$ must be two.
Substitution gives all the displayed invariants.

\end{proof}

\begin{remark}\label{rem:rational-canonical-scope}
These conclusions concern merely ample divisors and assume that the
normalization is birationally rational. They do not exclude a ruling
with positive-genus base. In characteristic two they give a necessary
canonical numerical configuration; in characteristic three they give
the two alternatives above. The strict genus theorem below excludes
both rational and positive-genus ruled possibilities when $Y$ is very
ample. It does not settle strict inequalities or higher conductor values.
When $\epsilon=1$, the nonzero higher-direct-image defect witnesses a
nonrational singularity; no rational-singularity hypothesis is imposed.
\end{remark}

\subsection{Separable multisections and the Picard variety}
\label{subsec:positive-genus-separable-picard}

For a general hyperplane section, the classical comparison of Picard
identity components gives a finite unipotent restriction kernel in positive
characteristic; see the account in
\cite[Remark~5.8, pp.~39--41]{KleimanPicard}.
In the equality setting with the additional ruledness hypothesis,
the distinguished normalized cover identifies
the Picard variety of the surface and excludes algebraically trivial
kernel classes.

\begin{lemma}\label{lem:positive-selfintersection-separable}
Let $V$ be an integral projective surface over an algebraically
closed field $k$ of characteristic $p>0$, smooth on a neighbourhood
of a smooth integral Cartier curve $E$ with $E^2>0$.
Let $q:V\to C$ be a morphism to a smooth projective curve of genus
at least one. If $E$ dominates $C$ and $q$ is smooth at the generic
point of $E$, then $q|_E:E\to C$ is separable.
\end{lemma}

\begin{proof}
The conormal sequence on the smooth neighbourhood of $E$ is
short exact \cite[Tag~06AA]{Stacks}. Its three terms are locally
free, so dualizing gives
\[
 0\longrightarrow T_E\longrightarrow T_V|_E
 \longrightarrow\OO_E(E)\longrightarrow0.
\]
If $q|_E$ were inseparable, its differential would be zero
\cite[Tag~0C1C]{Stacks}. The tangent map
$T_V|_E\to(q|_E)^*T_C$ would therefore factor through
\[
 \OO_E(E)\longrightarrow(q|_E)^*T_C.
\]
This homomorphism is nonzero at the generic point, since $q$ is
smooth there. A nonzero map of line bundles on a smooth projective
curve has an effective zero divisor, so
\[
 E^2\leq\deg(q|_E)\bigl(2-2g(C)\bigr)\leq0,
\]
a contradiction. Smoothness of $V$ away from $E$ is not used.
\end{proof}

\begin{proposition}\label{prop:positive-genus-picard-variety}
Let $X$ be a smooth projective integral surface over an algebraically
closed field $k$ of characteristic $p>0$, let $Y\subset X$ be a smooth
integral ample divisor, and let $M\ne\OO_X$ have order $p$ with
$M|_Y\simeq\OO_Y$. Suppose
\[
 2g(Y)-2=pY^2.
\]
Let $\pi:N\to X$ be the distinguished normalized cover, and put
$D=\pi^*Y$. Assume that $N$ is birationally ruled, with a fixed
presentation $k(N)=k(C)(t)$ whose smooth projective ruling curve
$C$ has positive genus. Use the morphisms
\[
 h:N\longrightarrow C,\qquad f:X\longrightarrow B=C^{(p)}
\]
of Proposition~\ref{prop:positive-genus-picard-fibration}.
Then the following assertions hold.
\begin{enumerate}[label=\textup{(\roman*)}]
\item The finite morphisms $h_D:D\to C$ and $f_Y:Y\to B$ are
separable. Writing $d=Y^2$, $g=g(C)$, and
$b=\deg h_D=\deg f_Y$, one has
\[
 pd-2b(g-1)=\deg R_{D/C}=\deg R_{Y/B}\geq0,
\]
where $R$ denotes the ramification divisor defined by the different.
\item Every morphism from $X$ to an abelian variety factors uniquely
through $f$. After choosing compatible $k$-points, the induced
homomorphism
\[
 \operatorname{Alb}(X)\longrightarrow\operatorname{Alb}(B)
\]
is an isomorphism.
\item Pullback gives an isomorphism of abelian varieties
\[
 f^*:\Pic^0_{B/k}
 \xrightarrow{\ \sim\ }(\Pic^0_{X/k})_{\mathrm{red}}.
\]
\item The image of $M$ in $\operatorname{NS}(X)$ has order
exactly $p$. In particular, $M$ is not algebraically equivalent
to zero.
\item Put $F=k(B)$ and $\Gamma=X\times_B\Spec F$. The restriction
$M|_\Gamma$ is nontrivial and has order $p$.
\end{enumerate}
\end{proposition}

\begin{proof}
\emph{Separability.}
The conductor equality and
Proposition~\ref{prop:positive-genus-picard-fibration} give that
$D$ is smooth and ample, $N$ is smooth on a neighbourhood of $D$,
and $D^2=pd>0$. They also give a nonempty open $C_0\subset C$
with
\[
 N_{C_0}\simeq C_0\times\PP^1.
\]
Since $D$ is ample, its intersection with a general ruling fibre
is positive, so it dominates $C$. Its generic point lies in this
product open, where $h$ is smooth.
Lemma~\ref{lem:positive-selfintersection-separable} proves that
$h_D$ is separable.

The finite flat radicial map $\pi_D:D\to Y$ has degree $p$.
Both curves are smooth, so it identifies with relative Frobenius
through an isomorphism $Y\simeq D^{(p)}$
\cite[Tag~0CCY]{Stacks}. The relation
$f\pi=F_{C/k}h$ identifies $f_Y$ with the Frobenius twist of
$h_D$. Equivalently, the $p$-power isomorphism of fields carries
$k(D)/k(C)$ to $k(D)^p/k(C)^p$.
Hence $f_Y$ too is separable, and both maps have degree $b$.
Riemann--Hurwitz, including the wild different multiplicities,
gives
\[
 2g(D)-2=b\bigl(2g(C)-2\bigr)+\deg R_{D/C};
\]
see \cite[Tag~0C1B, Lemma~53.12.2 and its preceding proof]{Stacks}.
Since $g(D)=g(Y)$, this gives the displayed formula in (i);
the Frobenius-twisted formula gives the same ramification degree
for $Y\to B$.

\emph{Maps to abelian varieties.}
We recall the constancy argument for a morphism
$\PP^1\to Q$, where $Q$ is an abelian variety.
Translation trivializes $\Omega^1_{Q/k}$, and its invariant
differentials generate it everywhere. Their pullbacks to
$\PP^1$ vanish, since $H^0(\PP^1,\Omega^1_{\PP^1/k})=0$.
Thus the differential of the morphism is zero.
In characteristic $p$ one must still remove the possibility
of an inseparable nonconstant map.
The kernel of
$d:k(t)\to\Omega^1_{k(t)/k}$ is $k(t^p)$, as follows from
the basis $1,t,\ldots,t^{p-1}$ over $k(t^p)$.
Consequently the generic map factors through relative Frobenius
of $\PP^1$. The resulting rational map from its smooth
Frobenius twist to $Q$ extends to a morphism, because $Q$ is
proper \cite[Tag~0BXZ]{Stacks}.
Repeating this factorization, a nonconstant rational function
on the reduced image curve would pull back to an element of
$k(t^{p^n})$ for every $n$. Its nonzero order at a zero or pole
would be divisible by every $p^n$, which is impossible.
Therefore every morphism $\PP^1\to Q$ is constant.

Put $A=k(C)$, $K=k(X)$, and $L=k(N)$.
Proposition~\ref{prop:distinguished-ruled-cover} gives
\[
 A\cap K=A^p,\qquad AK=L,
\]
and a proper birational morphism $r:S\to N$, where $S$ is
obtained from $C\times\PP^1$ by finitely many point blowups.
Write $q=hr:S\to C$. After shrinking $C_0$, both
$S_{C_0}$ and $N_{C_0}$ identify with $C_0\times\PP^1$.

Let $\varphi:X\to Q$ be any morphism to an abelian variety, and set
$\Phi=\varphi\pi r:S\to Q$.
On every closed ruling fibre over $C_0$, the map $\Phi$ is
constant by the preceding argument.
Choose the constant section $[1:0]$ of the product ruling.
Restriction to that section defines a morphism
$\beta_0:C_0\to Q$.
The maps $\Phi$ and $\beta_0q$ agree at every closed point of
the reduced finite-type scheme $S_{C_0}$; their closed equalizer
is therefore all of $S_{C_0}$.
The map $\beta_0$ extends uniquely to $\beta:C\to Q$ by
normality of $C$ and properness of $Q$
\cite[Tag~0BXZ]{Stacks}.
Equality on the dense product open then gives
\[
 \varphi\pi r=\beta q\quad\hbox{on all of }S.
\]

If $\beta$ is constant, so is $\varphi$.
Otherwise let $W$ be its reduced image in $Q$.
It is an integral projective curve, also the reduced image of
$\varphi$; it need not be smooth or normal.
The two embeddings of $k(W)$ into $K$ and into $A$ agree in
$L=k(S)$. Their image therefore lies in
$A\cap K=A^p$.
The latter field is the image of $k(B)$ under relative Frobenius.
Thus the inclusion $k(W)\hookrightarrow k(B)$ gives a rational
map $B\dashrightarrow W$.
Since $B$ is a normal curve and $W$ is proper, it extends to a
morphism. Composing with $W\hookrightarrow Q$ gives
$\gamma:B\to Q$, with $\beta=\gamma F_{C/k}$.
The identities
$f\pi r=F_{C/k}q$ and $\varphi\pi r=\beta q$ now imply
\[
 \varphi=\gamma f.
\]
Indeed the equality holds on the generic point of $X$, since
$K\hookrightarrow L$ is injective, and then everywhere by
separatedness. The factorization is unique because $f$ is
surjective.

\emph{Albanese and Picard varieties.}
Choose $x_0\in X(k)$ and put $b_0=f(x_0)$.
Let $a_X:X\to\operatorname{Alb}(X)$ and
$a_B:B\to\operatorname{Alb}(B)$ be the pointed Albanese maps.
Their universal property and Picard duality are
\cite[Expos\'e~236, Corollary~3.2 and Theorem~3.3(iii),
pp.~236--238]{FGAPicard}.
Functoriality gives a homomorphism
\[
 u:\operatorname{Alb}(X)\to\operatorname{Alb}(B),
 \qquad ua_X=a_Bf.
\]
The factorization just proved writes $a_X=\gamma f$ for a
pointed morphism $\gamma:B\to\operatorname{Alb}(X)$.
Universality of $a_B$ supplies a homomorphism
$v:\operatorname{Alb}(B)\to\operatorname{Alb}(X)$ with
$va_B=\gamma$.
Then $uv\,a_Bf=a_Bf$, and surjectivity of $f$ gives
$uv\,a_B=a_B$.
Uniqueness in the Albanese universal property gives $uv=1$.
Likewise
$vu\,a_X=v\,a_Bf=\gamma f=a_X$ gives $vu=1$.
This proves (ii).

For completeness, the Picard schemes of $X$ and $B$ exist
because they are smooth projective geometrically integral
over $k$
\cite[Expos\'e~232, Theorem~3.1, pp.~148--149;
proof in Sections~4--5, pp.~150--155]{FGAPicard}.
Their identity components are proper
\cite[Expos\'e~236, Theorem~2.1(ii), pp.~231--232]{FGAPicard}.
Over the perfect field $k$, the canonical abelian subscheme
of the identity component is its reduction.
For $B$, the whole identity component is smooth:
$H^2(B,\OO_B)=0$, so
\cite[Expos\'e~236, Proposition~2.10(ii), pp.~235--236]{FGAPicard}
applies.
The functorial duality of the Albanese and Picard varieties
therefore identifies the dual of $u$ with
\[
 f^*:\Pic^0_{B/k}\longrightarrow(\Pic^0_{X/k})_{\mathrm{red}}.
\]
It is an isomorphism, proving (iii).
This argument makes no assumption that $\Pic^0_{X/k}$ itself
is reduced.

\emph{The torsion class.}
First, $f^*:\Pic(B)\to\Pic(X)$ is injective on line-bundle
classes. Indeed $f_*\OO_X=\OO_B$ by the proof of
Proposition~\ref{prop:positive-genus-picard-fibration}, so the
projection formula gives
$P\simeq f_*f^*P\simeq\OO_B$ whenever $f^*P$ is trivial.

We also need the following consequence of separability.
Let $j:V\to B$ be a finite separable morphism of smooth
projective integral curves, and let $P\ne\OO_B$ have order $p$.
Write $P=\OO_B(E)$ and choose $u\in k(B)^*$ with
$\operatorname{div}_B(u)=pE$.
The element $u$ is not a $p$th power in $k(B)$:
otherwise $u=v^p$ would give
$p(E-\operatorname{div}v)=0$, and the free abelian divisor
group would imply $P\simeq\OO_B$.
Thus $T^p-u$ defines the nontrivial purely inseparable
degree-$p$ generic extension of the associated $\mu_p$-torsor.

If $j^*P$ were trivial, a rational function $w$ on $V$ would
satisfy $\operatorname{div}_V(w)=j^*E$.
Then $j^*u/w^p$ would have divisor zero, and hence would lie
in $k^*$. Since $k$ is algebraically closed, this constant is
a $p$th power. It follows that $u$ becomes a $p$th power in
$k(V)$.
A finite separable extension cannot contain a nontrivial
purely inseparable extension of its base field, so this is
impossible. Therefore $j^*P$ is nontrivial.

Suppose now that $M$ is algebraically trivial.
A $k$-point of $\Pic^0_{X/k}$ is a $k$-point of its reduction.
Part (iii) gives $M=f^*P$ for some $P\in\Pic^0(B)(k)$.
Injectivity of $f^*$ makes $P$ nontrivial of order $p$.
But $f_Y:Y\to B$ is finite separable by (i), whereas
$f_Y^*P=M|_Y$ is trivial. The preceding argument is a
contradiction.
Thus $M$ is not algebraically trivial, and its nonzero image in
$\operatorname{NS}(X)=\Pic(X)/\Pic^0(X)(k)$ is killed by $p$.
Its order is therefore exactly $p$, proving (iv).

\emph{The generic fibre.}
The curve $\Gamma$ is proper, integral and regular: its local rings
are localizations of those of the smooth surface $X$.
Moreover $f_*\OO_X=\OO_B$ gives
$H^0(\Gamma,\OO_\Gamma)=F$. This equality can be checked by
localizing a finite affine cover and its compatibility conditions.
Geometric smoothness of $\Gamma$ is not required.

Suppose $M|_\Gamma$ were trivial. A nowhere-vanishing section
on $\Gamma$ is a rational section $e$ of $M$ on $X$.
Its divisor $E$ has no horizontal component, because the
horizontal prime divisors of $X$ are the closed points of
$\Gamma$, at which the section is a unit.
Thus $E$ is a vertical Cartier divisor and $M\simeq\OO_X(E)$;
see \cite[Tags~0BE6 and 0BE8]{Stacks} for rational-section
divisors and their detection on a normal scheme.
Under a fixed trivialization $M^p\simeq\OO_X$, the section
$e^p$ is a rational function $u$ with
\[
 \operatorname{div}_X(u)=pE.
\]
On the proper regular curve $\Gamma$ this function has no
zero or pole, so it is an invertible regular function and
$u\in F^*$. The detection of units at codimension-one points
is \cite[Tag~031T, Lemma~10.157.6(2)]{Stacks}.

Since $Y$ dominates $B$, it is not a component of $E$.
Choose a trivialization of $M|_Y$ and write the restricted
rational section as $w\in k(Y)^*$.
The two trivializations of $M^p|_Y$ differ by an invertible
global function on $Y$, hence by a constant $c\in k^*$.
Therefore $u|_Y=cw^p$. The constant is a $p$th power, and
the finite extension $k(Y)/F$ is separable by (i).
It cannot contain a nontrivial purely inseparable extension
of $F$, so $u=v^p$ for some $v\in F$.
It follows that
\[
 p\bigl(E-\operatorname{div}_X(v)\bigr)=0.
\]
The Weil divisor group is free abelian on the prime divisors.
Hence $E$ is principal and $M$ is trivial, a contradiction.
This proves (v). The argument allows arbitrary vertical
component coefficients, including those of multiple fibres;
it never assumes that $E$ is pulled back from $B$.
\end{proof}

The proposition leaves a specific torsion obstruction in the
positive-genus branch: a kernel class must have nonzero order-$p$
image in the N\'eron--Severi group and a nontrivial order-$p$
restriction to the regular generic fibre. Its geometric
normalization being rational does not make the Picard group
of that regular curve over $k(B)$ torsion-free.
It does not assert that every torsion line bundle descends to the
ruling base. When $p\geq5$, the ample equality formula
\eqref{eq:normalized-smallest-equality-euler}, together with
$\chi(\OO_N)=1-g(C)$, gives $g(C)=1+pn$ for an integer $n\geq0$.
Part~(i) then gives $d\geq2bn$ without requiring very ampleness.
If $Y$ is very ample, the additional bound $b\geq p+2$ follows from
Proposition~\ref{prop:positive-genus-picard-fibration}.

\subsection{Adjoint systems and the strict genus inequality}
\label{subsec:canonical-subsystem-strictness}

A logarithmic form vanishing on a very ample curve controls much more
than its first nonzero coefficient. It gives a system of canonical
forms whose rational map has surface image. We first prove this
without a conductor equality. We then use the positive square to split an Euler extension at
conductor equality. The resulting structure theorem excludes the
first four conductors and gives the strict genus inequality.

\begin{proposition}\label{prop:residual-adjoint}
Let $X$ be a smooth projective integral surface over an algebraically
closed field $k$ of characteristic $p>0$, let $Y\subset X$ be a smooth
integral very ample divisor, and let $s$ be its canonical section.
Suppose that a nonzero closed global one-form $\omega$ vanishes along
$Y$. Let $m$ be its exact order along $Y$, and let $Z$ be the
divisorial zero part, with multiplicities, of
$\theta=\omega/s^m\in H^0(X,\Omega^1_{X/k}(-mY))$.
Then $m=pa$ for an integer $a\geq1$, $Y$ is not a component of $Z$, and
\[
 0\leq Z\cdot Y\leq 2g(Y)-2-mY^2.
\]
The divisor
\[
 A=K_X-(m-2)Y-Z
\]
is nef and big, with $A^2>0$.
More precisely, the linear system
\[
 W_K=\{d((u/s)\omega):u\in H^0(X,\OO_X(Y))\}
 \subset H^0(X,\omega_X)
\]
has exact fixed divisor $(m-2)Y+Z$ and two-dimensional image.
Its residual system has no fixed curve, and on a dense open of $Y$
its rational map is the given hyperplane embedding. The map
$u\mapsto d((u/s)\omega)$ has kernel $ks$. In particular,
\begin{align*}
 h^0(X,\Omega^1_{X/k})&\geq h^0(X,\OO_X(Y)),\\
 \operatorname{rank}\bigl[d:H^0(X,\Omega^1_{X/k})\to
 H^0(X,\omega_X)\bigr]&\geq h^0(X,\OO_X(Y))-1,\\
 p_g(X)&\geq h^0(X,\OO_X(Y))-1.
\end{align*}
Both $K_X-(m-2)Y$ and $K_X-(p-2)Y$ are big, and the canonical
map of $X$ has two-dimensional image.
\end{proposition}

\begin{proof}
The leading restriction of $\theta$ to $Y$ is nonzero. Its tangential
image is nonzero as well, because the kernel on global sections in
the twisted conormal sequence is
$H^0(Y,\OO_Y(-(m+1)Y))=0$. Locally write
$\omega=t^m(b\,dt+\eta)$, with $\eta$ nonzero tangentially modulo $t$.
The term of order $m-1$ in $d\omega$ is $m\,dt\wedge\eta$.
Thus $p\mid m$, and $m\geq p\geq2$.
The nonzero tangential section of $\omega_Y(-mY)$ gives
$2g(Y)-2\geq mY^2>0$, so $g(Y)\geq2$.

Regular local rings are factorial
\cite[Lemma~15.123.2, Tag~0AG0]{Stacks}.
The minimum valuations of the two coefficients of $\theta$ therefore
define an effective Cartier divisor $Z$. Dividing those coefficients
by their common prime factors gives a regular section
\[
 \theta'\in H^0(X,\Omega^1_{X/k}(-mY-Z)),\qquad
 \theta=s_Z\theta',
\]
with no divisorial zero. Exactness of $m$ excludes $Y$ from $Z$.
The conormal term after restriction to $Y$ has degree
$-(m+1)Y^2-Z\cdot Y<0$. Hence the tangential image
\[
 \sigma'\in H^0(Y,\omega_Y(-mY-Z))
\]
is nonzero. Its finite zero divisor has degree
$2g(Y)-2-mY^2-Z\cdot Y$, proving the first bound. Notice also that
\[
 A\cdot Y=Y^2+2g(Y)-2-mY^2-Z\cdot Y>0.
\]
Division by $s_Z$ need not preserve closedness; the argument below
uses closedness before that division.

\paragraph{The two-form system.}
Choose local frames $e_i$ of $\OO_X(Y)$ and write
$s=t_i e_i$, $u=u_i e_i$, and $\eta_i=t_i^{-m}\omega$.
The forms $\eta_i$ are regular and closed because $p\mid m$.
If $t_i=v_{ij}t_j$, then $u_i=v_{ij}u_j$ and
$\eta_i=v_{ij}^{-m}\eta_j$. The identity
\[
 t_i\,du_i-u_i\,dt_i
 =v_{ij}^2(t_j\,du_j-u_j\,dt_j)
\]
shows that
\begin{equation}\label{eq:residual-adjoint-local}
 \gamma_i(u)=(t_i\,du_i-u_i\,dt_i)\wedge\eta_i
\end{equation}
glues to a section of $K_X-(m-2)Y$.
Every such section is divisible by $Z$; denote the quotient by
$\gamma'(u)\in H^0(X,\OO_X(A))$ and its image space by $W$.

Along $Y$, equation~\eqref{eq:residual-adjoint-local} becomes
$-u_i\,dt_i\wedge\eta_i$. After dividing by $Z$, adjunction identifies
this restriction with $-u|_Y\,\sigma'$ in
$A|_Y=\omega_Y(-(m-1)Y-Z)$.
Multiplication by the nonzero section $\sigma'$ is injective on
sections of the integral curve. The exact sequence
\[
 0\longrightarrow\OO_X\xrightarrow{s}\OO_X(Y)
 \longrightarrow\OO_Y(Y)\longrightarrow0
\]
therefore gives $\ker\gamma'=ks$. Away from the zeros of $\sigma'$,
the rational map of $W$ is precisely the original embedding of $Y$
in the hyperplane defined by $s$. The restriction need not have its
complete linear series, and it may have finite base points.

On $X\setminus Y$, the functions $u/s$ are affine coordinates of the
very ample embedding and their differentials generate $\Omega^1_{X/k}$.
Up to invertible factors, the sections of $W$ are
$d(u/s)\wedge\theta'$. For a free rank-two module, the ideal of these
wedges is the ideal of the two coefficients of $\theta'$.
Thus no curve outside $Y$ is fixed. At the generic point of $Y$,
$\sigma'\ne0$, so $Y$ is not fixed either. This proves that $W$ has
no fixed curve and that the exact fixed divisor before division is $Z$.
It also proves nefness of $A$: for each integral curve $E$, choose a
member of $W$ not vanishing identically on $E$. Its effective zero
divisor has nonnegative intersection with $E$.

\paragraph{Why the system has surface image.}
Resolve its finite base locus by point blowups
\cite[Lemma~54.4.3, Tag~0C5H]{Stacks}, obtaining
\[
 \rho:R\longrightarrow X,\qquad
 \psi:R\longrightarrow\PP(W^\vee).
\]
The centres may lie on $Y$ or be infinitely near it.
The strict transform $Y'$ is isomorphic to $Y$: blowing up a smooth
curve at a point does not change that curve. For example, for local
parameters $t,x$ with $Y=(t=0)$, the chart $t=xv$ has strict transform
$v=0$, mapping isomorphically to $Y$.
The morphism $\psi|_{Y'}$ agrees with the original embedding on a dense
open. Since the target is separated, it agrees everywhere.

Suppose the reduced image $B=\psi(R)$ is a curve. It equals the
embedded image of $Y'$, so $B\simeq Y$ and $Y'$ is a section of
$\psi:R\to B$. Every exceptional curve of $\rho$ is a projective line
\cite[Section~54.3, Tag~0AGP]{Stacks}, and maps constantly to $B$.
To include inseparable maps in this assertion, a nonconstant map
$\PP^1\to B$ with nonzero differential would give a nonzero map from
a line bundle of positive degree to $\omega_{\PP^1}$, of degree $-2$.
If its differential is zero, its function-field image lies in
$k(t^p)$. Factoring through this rational function field and extending
the rational map by properness divides the degree by $p$.
Iteration eventually gives a nonzero differential and the same
contradiction. Here $g(B)=g(Y)\geq2$ is essential.

Keep the full Cartier pullback
\[
 \rho^*Y=Y'+\sum_i r_iE_i,\qquad r_i\geq0.
\]
All the exceptional components $E_i$ are vertical over $B$.
Set $F=k(B)$ and $\Gamma=R\times_B\Spec F$.
This is a proper integral regular curve: its local rings are
localizations of regular local rings of $R$.
Smoothness over the possibly imperfect field $F$ has not been assumed.
The section $Y'$ gives a single reduced $F$-rational point on $\Gamma$;
its equation is a uniformizer in the corresponding discrete valuation
ring. Therefore the globally generated line bundle
$\OO_R(\rho^*Y)|_\Gamma$ has degree one.

A nonzero section of this line bundle has a zero divisor consisting
of one reduced $F$-rational point. Global generation gives a second
section avoiding that point. The two sections define a finite morphism
$\Gamma\to\PP^1_F$ of degree one, by the degree formula on proper curves
\cite[Lemma~37.44.1, Tag~02LS; Lemma~33.44.11, Tag~0AYQ]{Stacks}.
Finite birationality onto a normal target makes this an isomorphism
\cite[Lemma~29.55.8, Tag~0AB1]{Stacks}.
Thus $\Gamma\simeq\PP^1_F$.

The restricted canonical bundle $\omega_R|_\Gamma$ has degree $-2$.
Indeed, differential transitivity gives
\[
 \Omega^1_{F/k}\otimes_F\OO_\Gamma\longrightarrow
 \Omega^1_{R/k}|_\Gamma\longrightarrow\Omega^1_{\Gamma/F}
 \longrightarrow0.
\]
The middle and final terms have ranks two and one and are locally
free. Their kernel is a line bundle, onto which the first, also a
line bundle, surjects. It is therefore isomorphic to that kernel.
Taking determinants proves the degree assertion.

For a point blowup, the local formula
$dx\wedge d(xv)=x\,dx\wedge dv$ gives an effective exceptional
canonical correction. Consequently
\[
 \rho^*\OO_X(A)\hookrightarrow
 \omega_R(-(m-2)\rho^*Y-\rho^*Z).
\]
A nonzero section in $W$ stays nonzero on the generic fibre, which
contains the generic point of $R$. The restriction of $\rho^*Z$ is
effective, including any horizontal components, so the target has
degree
\[
 -2-(m-2)-\deg(\rho^*Z|_\Gamma)<0.
\]
This contradicts the existence of that section. Hence the image of
$\psi$ has dimension two.

\paragraph{The square and the canonical system.}
Let $P=\psi^*\OO_B(1)$. Section ratios give an injection
$P\hookrightarrow\rho^*\OO_X(A)$ with effective exceptional zero
divisor $E_{\rm b}$; hence $\rho^*A=P+E_{\rm b}$.
Pullback is orthogonal to exceptional curves. Their intersection form
is negative definite, as the total transforms of successive exceptional
curves have diagonal matrix with entries $-1$.
The proper intersection-degree formula gives
\[
 P^2=[k(R):k(B)]\deg_{\OO_B(1)}B>0,
 \qquad A^2=P^2-E_{\rm b}^2>0
\]
\cite[Section~54.3, Tag~0AGP; Lemma~33.45.7, Tag~0BEL]{Stacks}.
No separability is needed. Products of the sections in $W$ have the
quadratic Hilbert growth of the homogeneous coordinate ring of $B$,
so $A$ is big.

For every $u\in H^0(X,\OO_X(Y))$, the one-form
$\beta(u)=(u/s)\omega$ is regular, since $m\geq2$.
The map $u\mapsto\beta(u)$ is injective. Closedness and $p\mid m$ give
\[
 d\beta(u)=s^{m-2}\gamma(u)=s^{m-2}s_Z\gamma'(u).
\]
Its kernel is $ks$, and multiplication by the fixed section does not
change the rational map. This proves the fixed-divisor, image and rank
assertions. Multiplication by $s^{m-p}s_Z$ places $W$ inside
$H^0(X,\omega_X(-(p-2)Y))$ without changing its rational map.
Adding effective divisors to $A$ proves the two bigness claims;
the canonical system contains $W_K$ and therefore has surface image
as well.
\end{proof}

The bigness consequence also has an argument that does not use
closedness or a distinguished divisor once a negatively twisted
cotangent section is known. It clarifies the role of the embedding.

\begin{lemma}\label{lem:negative-cotangent-adjoint}
Let $S$ be a smooth projective integral surface over an algebraically
closed field, let $H$ be very ample, let $a>0$ be an integer,
and let $Z$ be an effective Cartier divisor.
If $0\ne\alpha\in H^0(S,\Omega^1_S(-aH-Z))$, then
$|K_S+(2-a)H-Z|$ has a linear subsystem with two-dimensional image.
In particular $K_S+(2-a)H-Z$ is big. If $\alpha$ has no
divisorial zero, this divisor is also nef and has positive square.
\end{lemma}

\begin{proof}
Choose coordinate sections $s_0,\ldots,s_N$ for the embedding by $H$.
The expressions $\tau_{ij}=s_i\,ds_j-s_j\,ds_i$ glue to sections of
$\Omega^1_S(2H)$: the terms differentiating a change of frame cancel.
They generate this bundle, since on $s_i\ne0$ the $\tau_{ij}$ are,
in the induced trivialization, the differentials of the affine
coordinate functions $s_j/s_i$.
For the nonzero section $\alpha$, the wedges
$\alpha\wedge\tau_{ij}$ therefore give a rational map, regular where
$\alpha$ is nonzero, defined by a subsystem of $K_S+(2-a)H-Z$.

At such a point $x$, choose a nonzero tangent vector $v$ spanning the
kernel of $\alpha_x$. Up to one common scalar the coordinates of this
map are
\[
 s_i(x)\,ds_j(v)-s_j(x)\,ds_i(v).
\]
These are the Pl\"ucker coordinates of the projective line through
$x$ with tangent direction $v$. Thus each fibre of the rational map
lies in its represented line. The image cannot be a point, since
that would put a dense open of $S$ in one line.

If the image were a curve $B$ in the Grassmannian of lines, a general
fibre would be one-dimensional. Its closure would be the entire
represented line $C\subset S$. The universal family of lines is a
projective-line bundle over the Grassmannian; its restriction to a
dense smooth open of $B$ is flat and contained in $S\times B$.
The containment extends from the general fibres because the condition
that a line lie in $S$ is closed. These fibres are Cartier divisors on
the smooth surface, and the family is a relative Cartier divisor.
Indeed, flatness identifies the fibre of its ideal with the ideal of
the fibre; lifting a local generator and applying Nakayama's lemma
makes that ideal principal. Its generator is a nonzerodivisor in the
regular integral scheme $S\times B$ after shrinking to the smooth
base open. They form an algebraic family, so distinct members
$C,C'$ are numerically equivalent and
$C^2=C\cdot C'\geq0$.
For clarity, numerical equivalence can also be checked by restricting
the relative Cartier divisor to any fixed curve: the degree of its
line bundle is constant in this connected family.

Choose a general member $C$ that is not a component of $Z$ and meets
the open set where $\alpha\ne0$. There are only finitely many
components of $Z$, so this choice is possible.
Put $D=aH+Z$. Then $C\simeq\PP^1$, $H\cdot C=1$ and
$D\cdot C=a+Z\cdot C\geq a>0$. Its twisted conormal sequence is
\[
 0\longrightarrow\OO_C(-C-D)\longrightarrow
 \Omega^1_S(-D)|_C\longrightarrow\omega_C(-D)
 \longrightarrow0.
\]
The end terms have degrees $-C^2-D\cdot C<0$ and $-2-D\cdot C<0$.
Thus the middle term has no global section. This contradicts the
restriction of $\alpha$, which is nonzero at points of the original
fibre. The image therefore has dimension two. Its coordinate ring
embeds in the section ring of the adjoint and has quadratic Hilbert
growth, proving bigness. The argument uses no separability assumption.
If $\alpha$ has no divisorial zero, generation by the $\tau_{ij}$
shows that the wedge system has no fixed curve. The nefness and
exceptional-intersection argument in the preceding proof then gives
nefness and positive square. No closedness of $\alpha$ is required.
\end{proof}

\subsection{The structure of a conductor equality}
\label{subsec:conductor-equality-structure}

We now prove Theorem~\ref{thm:conductor-equality-structure}. The
positive square just obtained turns a determinant divisor into a
connected curve. Its constant functions split an Euler extension.
This supplies an ambient point through which all the tangent planes
pass, and makes the quotient and its possible singularity explicit.

The exclusion of $a=6$ will be completed in
Proposition~\ref{prop:conductor-six-excluded}; that argument uses
only the structural conclusions and $a\geq5$ proved here.

\begin{proof}[Proof of Theorem~\ref{thm:conductor-equality-structure}]
Put $H=\OO_X(Y)$, $m=pa$ and $d=Y^2$. Let $s_0$ define $Y$.
The equality conclusions of
Propositions~\ref{prop:normalized-mup-conductor}
and~\ref{prop:normalized-mup-singularity-defect} give
\[
 \beta=\omega/s_0^m\in H^0(X,\Omega^1_{X/k}\otimes H^{-m})
\]
whose tangential restriction to $Y$ is nowhere zero. Its zero scheme
$Z$ is finite and disjoint from $Y$: any divisorial zero would meet
the ample curve $Y$. Adjunction gives $K_X\cdot Y=(m-1)d$.
The annihilator of $\beta$ is the saturated invertible subsheaf
\[
 \mathcal F=H^m\otimes\omega_X^{-1}\longrightarrow T_X.
\]
Locally its generator is the Koszul relation between the two
coefficients of $\beta$, which have no common prime divisor.
It is a subbundle outside $Z$, and its normal projection on $Y$
is an isomorphism $\mathcal F|_Y\simeq H|_Y$.

Set
\[
 A=H\otimes\mathcal F^{-1}=\omega_X\otimes H^{-(m-1)},
 \qquad B=A\otimes H.
\]
Thus $A\cdot Y=0$ and $A|_Y\simeq\OO_Y$.
Proposition~\ref{prop:residual-adjoint}, with no additional
divisorial zero, shows that $B$ is nef and big.

\paragraph{Splitting the Euler extension.}
Let $V=H^0(X,H)$ and embed $X$ in $\PP(V^*)$. Pull back the
restricted Euler sequence along
$\mathcal F\to T_{\PP(V^*)}|_X$. This gives a rank-two bundle
$P$ and a commutative diagram
\[
\begin{array}{ccccccccc}
0&\longrightarrow&\OO_X&\longrightarrow&P&\longrightarrow&\mathcal F&\longrightarrow&0\\
&&\Vert&&\downarrow&&\downarrow\\
0&\longrightarrow&\OO_X&\longrightarrow&V^*\otimes H&\longrightarrow&T_{\PP(V^*)}|_X&\longrightarrow&0.
\end{array}
\]
The bundle $P$ remains locally free at $Z$, since it is an extension
of line bundles. Put $E=P^*\otimes H$. There are maps
\[
 0\longrightarrow A\longrightarrow E\longrightarrow H
 \longrightarrow0,\qquad V\otimes\OO_X\longrightarrow E,
\]
and their composite is evaluation. For $s\in V$, denote its
section of $E$ by $e_s$. Locally its components record the value
of $s$ and differentiation in the $\mathcal F$ direction. At a
zero of $s$, the derivative component is independent of the frame.
The section $e_0=e_{s_0}$ is nowhere zero: off $Y$ its value
component is nonzero, and on $Y$ its normal derivative is nonzero.
Consequently
\begin{equation}\label{eq:conductor-euler-extension}
 0\longrightarrow\OO_X\xrightarrow{e_0}E
 \xrightarrow{e_0\wedge-}B\longrightarrow0
\end{equation}
is exact. Its determinant sections $t_s=e_0\wedge e_s$ restrict
on $Y$ to the hyperplane sections $s|_Y$, under the identification
$B|_Y\simeq H|_Y$ given by $e_0$.

Choose $s$ so that $s|_Y$ has $d$ distinct simple zeros. Such
hyperplanes form a nonempty open set, since $Y$ is a smooth
embedded curve. The divisor $C=\operatorname{div}(t_s)$ is
nonempty and does not contain $Y$. Every component meets $Y$;
simple intersection with $Y$ therefore forces every component to
have multiplicity one. A Cartier divisor on a smooth surface has
no embedded components, so $C$ is reduced.

It is also connected. If $C=C_1+C_2$ were a disjoint decomposition
into nonzero effective divisors, then
$C_i^2=B\cdot C_i\geq0$. A zero square would make $C_i$
orthogonal to the positive-square class $B$, and Hodge index
would make $C_i$ numerically zero, contrary to $Y\cdot C_i>0$.
Two positive-square classes cannot be orthogonal either.
Thus $H^0(C,\OO_C)=k$.

On $C$, the image of $e_s$ in the quotient in
\eqref{eq:conductor-euler-extension} vanishes. Hence
$e_s|_C=\lambda e_0|_C$ for a scalar $\lambda\in k$.
Dividing $e_s-\lambda e_0$ by the equation $t_s$ gives a map
$B\to E$ whose composite with $E\to B$ is the identity.
Thus $E\simeq\OO_X\oplus B$.
The surjection $E\to H$ has a nonzero component $B\to H$;
otherwise its image would be generated by $s_0$ and would vanish
on $Y$. This gives a nonzero section of $A^{-1}$.
Its effective zero divisor has intersection zero with the ample
curve $Y$, so it is empty. Therefore
\begin{equation}\label{eq:conductor-linear-canonical}
 A\simeq\OO_X,\qquad B\simeq H,\qquad
 \mathcal F\simeq H,\qquad K_X\sim(m-1)Y.
\end{equation}
This proves linear equivalence, with no injectivity assumption on
restriction of line bundles to $Y$.

\paragraph{The constant tangent point.}
In $E\simeq\OO_X\oplus H$, the component $H\to H$ of the
surjection $E\to H$ is a nonzero scalar. It splits that surjection.
Dualizing and twisting therefore splits the original Euler pullback
compatibly. The inclusion $H\simeq\mathcal F\to T_{\PP(V^*)}|_X$
lifts to a map $H\to V^*\otimes H$, given by a constant nonzero
vector $v\in V^*$. Write $v$ also for its projective point.
The induced tangent direction at $x\ne v$ is that of the line
$xv$. In particular every embedded tangent plane of $X$ contains
$v$. This is the usual notion of a strange embedding
\cite[\S2.1 and Proposition~2.5, author pp.~2--5]{FurukawaStrange}.

The point $v$ lies outside the hyperplane defining $Y$, since the
direction on $Y$ is transverse to that hyperplane. Choose
coordinates with $v=[1:0:\cdots:0]$ and $Y=(x_0=0)\cap X$.
The corresponding section of $T_X\otimes H^{-1}$ has no zero
except possibly $v$. If $v\in X$, its zero ideal is exactly the
maximal ideal of $v$: the ambient coefficients $x_i/x_0$ generate
that ideal, and the tangent inclusion is locally a direct summand.
Thus $Z$ is empty or one reduced point. Let $\epsilon$ be its
length.

At this point choose two ambient affine coordinates $u,w$ forming
regular parameters of $X$. The radial field satisfies
$D(u)=-u$, $D(w)=-w$. The first jet of its annihilating form is
a nonzero scalar multiple of $u\,dw-w\,du$.
Since $v\notin Y$, the same is true of $\omega$.
The derivative of a local unit contributes zero at a zero of the
form. Thus $d\omega(v)=2c\,du\wedge dw$ with $c\ne0$.
Closedness proves
\begin{equation}\label{eq:conductor-epsilon-parity}
 \epsilon=1\quad\Longrightarrow\quad p=2.
\end{equation}
The zero-length formula gives, in both cases,
\begin{equation}\label{eq:conductor-c2X}
 \epsilon=c_2(\Omega_X^1\otimes H^{-m})=c_2(X)+md,
 \qquad c_2(X)=\epsilon-pa\,d.
\end{equation}

\paragraph{The quotient and the normalized torsor.}
Multiplying the twisted field by $x_0$ gives the restriction of
$\delta=x_0\partial/\partial x_0$, with $\delta^p=\delta$.
The homogeneous ideal of $X$ is stable under this derivation.
The eigenspace projections are polynomials in $\delta$, since
$T^p-T$ has distinct roots, so the ideal is graded by weights
modulo $p$. Hence the ambient $\mu_p$ action restricts to $X$.
Its quotient $q:X\to S$ is finite radicial of degree $p$ and
normal. On invariant affine rings, normality follows by
intersection with the invariant function field. The quotient is
projective; equivalently, relative Frobenius factors through it
and makes $S$ finite over $X^{(p)}$.

Let $K=k(X)$ and write the original torsor generically as
$t^p=u$, with $u\notin K^p$. Then $k(N)=L=K(u^{1/p})$.
The field $\delta$ annihilates $du/u$, and hence $u$.
Extending $u$ to a $p$-basis of $K$ shows that
\[
 K^\delta=K^p(u)=L^p.
\]
This identifies the schemes as well as their function fields.
For an affine open $\Spec A\subset S$, let its inverse image in
$N$ be $\Spec C$. Both rings are normal, $C$ is finite over $A$,
and $\operatorname{Frac}A=(\operatorname{Frac}C)^p$.
For $c\in C$, the element $c^p$ is integral over $A$ and lies
in its fraction field, so it belongs to $A$.
Conversely, every $a\in A$ has a $p$th root in
$\operatorname{Frac}C$, integral over $C$, and therefore in $C$.
Thus $A=C^p$. Since $k$ is perfect, these identifications give
\begin{equation}\label{eq:conductor-frobenius-quotient}
 S\simeq N^{(p)},\qquad q\pi=F_{N/k}.
\end{equation}

Away from $v$, the saturated foliation is regular, so its quotient
is smooth. The local result used here is Seshadri's regular
quotient theorem \cite[Proposition~6 and Theorem~2,
pp.~6-11--6-16]{SeshadriCartier}. In rank one, a normalized
generator has $\partial(t)=1$ and $\partial^p=0$; invariant
coordinates give a quotient with regular parameters $t^p,z$.
At $Y$ one can also use the multiplicative coordinates
$\delta(t)=t$, $\delta(z)=0$.

If $v\in X$, then $p=2$ and its completed invariant ring is
\[
 k[[u^2,uw,w^2]]\simeq k[[a,b,c]]/(ac-b^2).
\]
Indeed, invariant monomials have even total degree, and are
products of $u^2,uw,w^2$, with the displayed sole relation.
Invariants commute with completion through the weight-zero
direct summand. This is an $A_1$ singularity. Blowing up $v$
in $X$, with exceptional curve $R$, resolves the quotient:
on the chart $w=uz$ the lifted field satisfies
$\delta(u)=-u$, $\delta(z)=0$, and the quotient chart has
coordinates $u^2,z$. The other chart is analogous.
Let $q':X'\to T'$ be this smooth quotient and $\rho:T'\to S$
the induced resolution. Its exceptional curve $E$ satisfies
$q'^*E=2R$ and $E^2=-2$.
It is the minimal resolution. It is crepant as well: since $S$
is Gorenstein, write $K_{T'}=\rho^*K_S+cE$; adjunction gives
$K_{T'}\cdot E=0$, whence $c=0$.
For $\epsilon=0$ set $X'=X$ and $T'=S$.
Equation~\eqref{eq:conductor-frobenius-quotient} transfers these
smoothness and singularity statements to $N$; its minimal
resolution $T$ has Frobenius twist $T'$.

\paragraph{Canonical class and Chern numbers.}
The normalized conductor formula gives
\[
 K_N=\pi^*(K_X-a(p-1)Y)\sim(a-1)\pi^*Y.
\]
For $a\geq2$ it is ample. The crepant resolution has nef and
big canonical divisor, so $T$ and $T'$ are minimal surfaces of
general type. Their canonical squares are
\[
 K_T^2=K_{T'}^2=K_N^2=p(a-1)^2d.
\]

We compute their second Chern numbers directly on smooth surfaces.
For a regular rank-one foliation $F\subset T_{X'}$, let $Q$
be its annihilator in $\Omega^1_{X'}$.
The invariant coordinates just used give exact sequences
\[
 0\to Q\to\Omega^1_{X'}\to F^{-1}\to0,
 \qquad
 0\to F^{-p}\to q'^*\Omega^1_{T'}\to Q\to0.
\]
Here $Q=\omega_{X'}\otimes F$; the $p$th power in the second
kernel follows from the transition of the local coordinate $t^p$.
Writing $f=c_1(F)$, determinants and second Chern classes give
\[
 q'^*K_{T'}=K_{X'}-(p-1)f,\qquad
 c_2(X')=-f\cdot(K_{X'}+f),\qquad c_2(T')=c_2(X').
\]
For $\epsilon=0$, $f=Y$. For $\epsilon=1$, the lifted
field has fixed divisor $b^*Y+R$, so $f=b^*Y+R$ and
$K_{X'}=b^*K_X+R$.
Using $R^2=-1$ and $b^*Y\cdot R=0$, one obtains
\begin{equation}\label{eq:conductor-structure-invariants}
\begin{gathered}
 K_T^2=p(a-1)^2d,\qquad c_2(T)=-pa\,d+2\epsilon,\\
 12\chi(\OO_T)=p(a^2-3a+1)d+2\epsilon.
\end{gathered}
\end{equation}
The last equality is Noether's formula.
As a separate check, Noether's formula on $X$ and
Corollary~\ref{cor:normalized-equality-euler} give
$12\chi(\OO_N)=p(a^2-3a+1)d+p\epsilon$;
equation~\eqref{eq:conductor-epsilon-parity} makes the two
expressions equal.

\paragraph{The first four conductors are impossible.}
If $a=1$, $K_T$ is numerically trivial.
The classification of smooth projective surfaces with numerically
trivial canonical divisor gives $c_2(T)\geq0$
\cite[\S2.1, Table~1, author p.~3]{MatsumotoMuN2023}.
The table has Betti pairs $(b_1,b_2)=(4,6),(0,22),(0,10),(2,2)$,
giving Euler numbers $0,24,12,0$.
For $\epsilon=0$, this contradicts $c_2(T)=-pd<0$.
For $\epsilon=1$, $p=2$ and $2-2d\geq0$, hence $d=1$.
But $H|_Y$ is very ample of degree one, so $Y$ is a line.
This contradicts $2g(Y)-2=pa\,d>0$.

For $a\geq2$, Gu--Sun--Zhou's inequality applies to the smooth
minimal surface $T$ of general type:
\[
 K_T^2\leq32\chi(\OO_T)
 \qquad\text{\cite[Theorem~1.5, p.~614;
 Corollary~3.4, pp.~623--624]{GuSunZhou}.}
\]
In particular $\chi(\OO_T)>0$.
If $a=2$, equation~\eqref{eq:conductor-structure-invariants}
gives $12\chi(\OO_T)=-pd+2\epsilon\leq0$, a contradiction.
If $a=3$ and $\epsilon=0$, it gives
$K_T^2=4pd>32pd/12=32\chi(\OO_T)$.
If $a=3$ and $\epsilon=1$, it gives
\[
 K_T^2=8d,\qquad \chi(\OO_T)=(d+1)/6.
\]
The inequality forces $d\leq2$, but neither positive integer
$d=1,2$ makes this Euler characteristic integral.
Thus $a\geq4$, proving ampleness and minimality.

We have established $a\geq4$. Suppose $a=4$, and let
$\rho:T\to N$ be the minimal resolution, with $\rho$ the identity
when $N$ is smooth. Its canonical identity and crepancy give
\[
 K_T\sim3A,\qquad A=\rho^*\pi^*Y,
\]
where $A$ is an integral Cartier divisor class. The invariant
formulas give
\[
 K_T^2=9pd,\qquad 12\chi(\OO_T)=5pd+2\epsilon.
\]
Since $\epsilon=1$ implies $p=2$, it follows that
\[
 12\bigl(K_T^2-18\chi(\OO_T)\bigr)
 =18pd-36\epsilon=18p(d-\epsilon)>0.
\]
For $\epsilon=0$, positivity follows from $d>0$.
For $\epsilon=1$, the value $d=1$ would make the very ample
curve $Y$ a line, contradicting $2g(Y)-2=pa\,d>0$.

The surface $T$ is smooth minimal of general type.
Gu--Sun--Zhou's theorem therefore gives a genus-two Albanese
fibration
\cite[Theorem~1.5, p.~614; Corollary~3.4 and its proof,
pp.~623--624]{GuSunZhou}.
For a fibre divisor $F$, one has $F^2=0$ and
$K_T\cdot F=2$ by adjunction.
The genus is the arithmetic genus in the flat-fibration
convention of that theorem, so this calculation does not require
a smooth fibre. It contradicts $K_T\sim3A$, since
$A\cdot F$ is an integer.
Therefore $a\geq5$.

\paragraph{The divisibility of the degree.}
If $v\notin X$, projection from $v$ is a morphism defined by
$x_1,\ldots,x_r$, with pullback of $\OO(1)$ equal to $H$.
Its coordinate ratios are invariant under $\mu_p$, so it factors
through $q:X\to S$. For the descended line bundle $L$ on $S$,
$q^*L=H$. Hence $d=pL^2$.
If $v\in X$, the base ideal of the projection is the maximal
ideal of that reduced point. On $X'$ the resolved projection
has line bundle $b^*H-R$ and factors through $q':X'\to T'$.
Its descended line bundle satisfies $q'^*L=b^*H-R$; thus
$d-1=pL^2$. This proves $p\mid(d-\epsilon)$.
No very ampleness of the descended line bundle is required.
\end{proof}

\begin{corollary}\label{cor:conductor-not-complete-intersection}
In Theorem~\ref{thm:conductor-equality-structure}, the embedding
of $X$ by $\OO_X(Y)$ is not a complete intersection.
\end{corollary}

\begin{proof}
The embedding is strange, as the proof shows, and is not linear
since $K_X\cdot Y=(pa-1)d>0$.
Furukawa's theorem states that a nonlinear smooth complete
intersection is strange exactly when it is an odd-dimensional
quadric in characteristic two
\cite[Theorem~1.1, author p.~2, proved in \S3]{FurukawaStrange}.
Our surface has even dimension, so it cannot be such a complete
intersection.
\end{proof}


\begin{corollary}\label{prop:higher-conductor-positive-genus-exclusion}
Let $X$ be a smooth projective integral surface over an algebraically
closed field of characteristic $p>0$, let $Y\subset X$ be a smooth
integral very ample divisor, and let $M\ne\OO_X$ have order $p$ with
$M|_Y\simeq\OO_Y$. Let $pa$ be the exact vanishing order of its
logarithmic form along $Y$, and let $N$ be the normalization of its
$\mu_p$-torsor. If $2g(Y)-2=paY^2$, then $N$ is not birationally ruled.
\end{corollary}

\begin{proof}
Theorem~\ref{thm:conductor-equality-structure} makes the minimal
resolution of $N$ a surface of general type. A birationally ruled
surface has Kodaira dimension $-\infty$, so it cannot occur.
\end{proof}

\begin{corollary}\label{lem:divided-differential-subsystem}
In Proposition~\ref{prop:residual-adjoint}, if $m=p$ and the tangential
restriction of $\omega/s^p$ to $\omega_Y(-pY)$ is nowhere vanishing,
then $Z=0$. The residual system is a subsystem of
$\mathcal L=\omega_X(-(p-2)Y)$, has no base point on $Y$, and restricts
to its given embedding. In particular $\mathcal L$ is nef and big and
\begin{equation}\label{eq:divided-subsystem-positive-square}
 \mathcal L^2>0.
\end{equation}
\end{corollary}

\begin{proof}
A nowhere-vanishing tangential section makes
$2g(Y)-2-pY^2=0$. Proposition~\ref{prop:residual-adjoint} gives
$Z\cdot Y=0$. Since $Y$ is ample and $Z$ effective, $Z=0$.
The restriction formula in its proof has a nowhere-vanishing common
factor, so there is no base point on $Y$. All the other conclusions
follow from that proposition.
\end{proof}

\begin{proof}[Proof of Theorem~\ref{thm:strict-picard-all}]
The nontrivial order-$p$ class gives the nonzero closed logarithmic
form of Proposition~\ref{prop:normalized-mup-conductor}. Its full
restriction to $Y$ vanishes, by the conormal sequence and negative
conormal degree. Proposition~\ref{prop:residual-adjoint} proves the
adjoint and canonical-image assertions.
The conductor inequality is $2g(Y)-2\geq paY^2\geq pY^2$.
If the last bound were an equality, then $a=1$ and equality would
hold with the exact conductor. This contradicts
Theorem~\ref{thm:conductor-equality-structure}, which requires
$a\geq5$. Hence $2g(Y)-2>pY^2$.
\end{proof}

\subsection{Conductor equalities in odd characteristic}
\label{subsec:odd-higher-conductor}

The canonical divisor in the structure theorem is divisible by the
conductor minus one as a Cartier class. Adjunction on an Albanese
fibre combines this divisibility with Tate's genus-change theorem
and the bounds of Gu, Sun and Zhou to leave at most three possible
fibre genera.

\begin{proposition}\label{prop:odd-higher-conductor-equality}
Let $X$ be a smooth projective integral surface over an algebraically
closed field $k$ of odd characteristic $p$, let $Y\subset X$ be a smooth
integral very ample curve, and let $M$ be a nontrivial line bundle of
order $p$ whose restriction to $Y$ is trivial. Let $\omega$ be its
logarithmic form, let $pa$ be its exact vanishing order along $Y$, and
let $\pi:N\to X$ be the normalization of its $\mu_p$-torsor. Suppose
that
\[
 2g(Y)-2=paY^2.
\]
Put $d=Y^2$. Then $a\geq5$, $N$ is smooth with ample
canonical bundle, and
\begin{equation}\label{eq:odd-conductor-invariants}
\begin{gathered}
 K_X\sim(pa-1)Y,\qquad K_N\sim(a-1)\pi^*Y,\qquad p\mid d,\\
 K_N^2=p(a-1)^2d,\qquad c_2(N)=-pa\,d,\\
 12\chi(\OO_N)=p(a^2-3a+1)d.
\end{gathered}
\end{equation}
Let $f:N\to C$ be its Albanese fibration, let $g$ be its
arithmetic fibre genus, and put
\[
 L=\pi^*Y,\qquad \ell=L\cdot F
\]
for a general fibre $F$. Exactly the following three rows can
survive the canonical divisibility, genus-change and slope bounds:
\begin{equation}\label{eq:odd-conductor-three-genera}
\begin{array}{c|c|c|c}
 \ell&\text{range of }a&g&\text{divisibility}\\ \hline
 1&a\geq5\text{ odd}&(a+1)/2&p-1\mid a+1\\
 2&a\geq7&a&p-1\mid2a\\
 3&a\geq21\text{ odd}&(3a-1)/2&p-1\mid3a-1
\end{array}
\end{equation}
The second and third rows require a nonhyperelliptic fibration.
In particular $a\ne6$, and $a=5$ forces $g=3$ and
$p\in\{3,7\}$. Every case satisfies
\begin{equation}\label{eq:odd-conductor-characteristic-bound}
 p\leq2a+1.
\end{equation}
The rows describe necessary numerical conditions and do not assert
the existence of equality cases.
\end{proposition}

\begin{proof}
Theorem~\ref{thm:conductor-equality-structure} gives the identities
\eqref{eq:odd-conductor-invariants}, the smoothness of $N$, and its
ample canonical bundle. Thus $N$ is minimal of general type,
$c_2(N)<0$ and $\chi(\OO_N)>0$.
The negative-Chern-number theorem
\cite[Theorems~2 and~7 and their proofs, pp.~264--266]{ShepherdBarronGeography}
gives an Albanese fibration $f:N\to C$, with $C$ smooth,
$f_*\OO_N=\OO_C$, arithmetic fibre genus $g\geq2$, and
integral singular rational geometric generic fibre.

The fibre need not be smooth. A general closed fibre $F$ is a
Cartier curve, with $F^2=0$, and adjunction gives
\[
 (a-1)L\cdot F=K_N\cdot F=2g-2.
\]
The dualizing degree formula applies to a Gorenstein curve, so
singular fibres cause no change here
\cite[Tags~0B4A, 0C19 and~0BS6]{Stacks}.
The morphism to the smooth curve $C$ is flat: its local rings
are torsion free over the base discrete valuation rings.
Consequently its arithmetic fibre genus equals its generic
arithmetic genus. The generic curve is regular by localization,
and is proper and geometrically integral over $k(C)$.
Its geometric normalization is rational. Tate's genus-change
theorem, in the form of
\cite[Corollary~2.3, author p.~4]{SchroerGenusChange}, gives
\[
 p-1\mid2g.
\]

Set $s=a-1\geq4$. Since $L$ is an ample Cartier divisor,
\begin{equation}\label{eq:odd-conductor-integer-fibre-degree}
 \ell=L\cdot F\in\mathbb Z_{>0},\qquad
 g=1+\frac{s\ell}{2}.
\end{equation}
In particular $g=2$ is impossible. The exact Chern ratio is
\[
 R(a)=\frac{K_N^2}{\chi(\OO_N)}
      =\frac{12s^2}{s^2-s-1}.
\]
For a hyperelliptic fibration, Gu--Sun--Zhou's bound is
\[
 R(a)\leq h(g):=
 \frac{(12g+8)(g-1)}{g^2-g-1}.
\]
For a nonhyperelliptic fibration it is $18$ at $g=3$,
$840/47$ at $g=4$, and
\[
 R(a)\leq n(g):=
 \frac{12(g-1)(3g^2-4g-4)}{g(3g^2-12g+15)}
 \qquad(g\geq5)
\]
\cite[Theorem~3.2 and its proof, pp.~622--623]{GuSunZhou}.
All denominators below are positive in their stated ranges.

We first record the monotonicity and comparison needed to make
the genus range finite:
\[
 h'(g)=-\frac{4(2g^2+2g+1)}{(g^2-g-1)^2}<0
 \qquad(g\geq2),
\]
\[
 n'(g)=-\frac{4\{5g^2((g-3)^2-2)+4(g-1)(3g-5)\}}
 {(g^3-4g^2+5g)^2}<0
 \qquad(g\geq5),
\]
\[
 n(g)-h(g)=
 \frac{4(g-1)(3g+2)(g^2-4g+2)}
 {g(g^2-4g+5)(g^2-g-1)}>0
 \qquad(g\geq5).
\]
In the derivative of $n$, both summands in braces are positive
for $g\geq5$.

Define
\[
 P(s,g)=(s^2-s-1)(5g^2-15g+4)
       -(s+1)g(3g^2-12g+15).
\]
Direct subtraction gives
\begin{equation}\label{eq:odd-conductor-slope-polynomial}
 n(g)-R(a)=
 \frac{12P(s,g)}{(s^2-s-1)g(3g^2-12g+15)}.
\end{equation}
If $\ell\geq4$, then $g\geq2s+1\geq9$.
Both slope bounds are at most $n(2s+1)$, whereas
\[
 P(s,2s+1)=-2s(2s^3+21s^2+2s-5)<0.
\]
Equation~\eqref{eq:odd-conductor-slope-polynomial} gives a
contradiction. Hence $\ell\leq3$.

If $\ell=1$, then $s$ is even and $g=1+s/2=(a+1)/2$.
For every even $s\geq4$,
\[
 h(1+s/2)-R(a)=
 \frac{4s(s^2-s-10)}{(s^2-s-1)(s^2+2s-4)}>0.
\]
The nonhyperelliptic bound also exceeds $h(g)$: for $g=3,4$
this follows from $18>88/5$ and $840/47>168/11$;
for $g\geq5$ it follows from the displayed comparison of
$n$ and $h$. Thus neither slope envelope further excludes
this row. Tate's condition becomes $p-1\mid a+1$.

For $\ell\geq2$ one has $g\geq s+1$, and
\[
 h(s+1)-R(a)=
 -\frac{4s(s^2+5s+5)}{(s^2-s-1)(s^2+s-1)}<0.
\]
Monotonicity of $h$ therefore excludes the hyperelliptic case
in the remaining rows.

If $\ell=2$, then $g=s+1\geq5$ and
\[
 P(s,s+1)=s(2s^3-10s^2-3s+5).
\]
The cubic factor is $-39$ at $s=4$ and $-10$ at $s=5$.
For $s\geq6$ it is
\[
 2s^2(s-5)-3s+5\geq2s^2-3s+5>0.
\]
The slope inequality consequently holds exactly for $a\geq7$,
and Tate's condition is $p-1\mid2a$.

If $\ell=3$, then $s$ is even, $g=1+3s/2\geq7$, and
\[
 P(s,1+3s/2)=\frac{3s}{8}(3s^3-59s^2-8s+20).
\]
For $4\leq s\leq19$ the cubic factor is at most
$-2s^2-8s+20<0$, because $3s-59\leq-2$.
For $s\geq20$ it is at least
$s^2-8s+20=(s-4)^2+4>0$.
Thus this row is permitted exactly for odd $a\geq21$,
and its Tate condition is $p-1\mid3a-1$.
This proves the complete numerical reduction
\eqref{eq:odd-conductor-three-genera}.

In the first row $p\leq a+2$, and in the second
$p\leq2a+1$. In the third, $p-1$ is a proper divisor
of $3a-1$: equality would give the composite number $p=3a$.
It follows that
\[
 p\leq\frac{3a+1}{2}<2a+1.
\]
This proves \eqref{eq:odd-conductor-characteristic-bound}.
When $a=5$ only the first row remains, giving $g=3$ and
$p-1\mid6$, hence $p=3$ or $7$. None of the rows allows $a=6$.
\end{proof}

\begin{corollary}\label{cor:odd-conductor-numerical-threshold}
In Proposition~\ref{prop:odd-higher-conductor-equality}, the following
are necessary lower bounds on the conductor:
\[
 a\geq A_{\mathrm{num}}(p):=
 \begin{cases}
 5,&p=3\text{ or }7,\\
 7,&p=5,\\
 9,&p=11,\\
 11,&p=13,\\
 (p-1)/2,&p\geq17.
 \end{cases}
\]
Each right-hand side is the first integer for which some genus
and fibration type satisfy the numerical tests in the proposition.
For $p\geq17$, equality in this numerical lower bound forces
$g=a=(p-1)/2$ and a nonhyperelliptic Albanese fibration.
For example,
\[
\begin{array}{c|rrrrrrrrrr}
 p&3&5&7&11&13&17&19&23&29&31\\ \hline
 A_{\mathrm{num}}(p)&5&7&5&9&11&8&9&11&14&15\\
 g\text{ at this value}&3&4&3&5&6&8&9&11&14&15
\end{array}
\]
These thresholds do not assert the existence of geometric equality
cases.
\end{corollary}

\begin{proof}
Put $q=(p-1)/2$. In the first row of
\eqref{eq:odd-conductor-three-genera}, $a$ is odd and
$a\equiv-1\pmod{p-1}$. In the second, $a$ is a multiple of
$q$ and $a\geq7$. The third starts at $a=21$.
These conditions give the first five cases in the statement
directly.

If $p\geq17$, then $q\geq8$. The first row has
$a\geq p-2=2q-1\geq q$, and the second has $a\geq q$.
If the third had $a<q$, the positive integer
$(3a-1)/(2q)$ would be less than $3/2$, hence equal to one.
This would imply $p=3a$, impossible for the prime $p$.
Thus every row requires $a\geq q$.
The pair $a=g=q$ satisfies the second row and its slope bound,
as proved above. The first row cannot have $a=q$ for $q\geq8$,
and the third would require $(3q-1)/(2q)$ to be an integer,
which it is not. This proves all assertions.
\end{proof}

For any fixed conductor, the complete characteristic test is the
union of the three divisibility conditions in
\eqref{eq:odd-conductor-three-genera}, with their stated ranges.
It can have gaps above the first permitted conductor. The third
row is impossible when $p\equiv1\pmod3$.
All assertions in this subsection concern equality with the exact
conductor; they impose no replacement inequality on cases with
strict excess in $2g(Y)-2\geq paY^2$.

\subsection{The sixth conductor is impossible}

We combine the integral canonical class with the same slope bound
to exclude conductor six, including the singular characteristic-two case.

\begin{proposition}\label{prop:conductor-six-excluded}
Under the hypotheses of Theorem~\ref{thm:conductor-equality-structure},
the conductor $a=6$ is impossible in every characteristic.
\end{proposition}

\begin{proof}
Suppose $a=6$, and retain $d=Y^2$ and $\epsilon\in\{0,1\}$
from the structure theorem. On the minimal smooth normalization
resolution $T$, the Cartier canonical identity and Chern numbers are
\[
 K_T=5L,\qquad K_T^2=25pd,\qquad
 c_2(T)=-6pd+2\epsilon,\qquad
 12\chi(\OO_T)=19pd+2\epsilon.
\]
The surface $T$ is minimal of general type, its second Chern number
is negative, and its Euler characteristic is positive.
The negative-Chern-number theorem supplies an Albanese fibration
of arithmetic genus $g\geq2$
\cite[Theorems~2 and~7, pp.~264--266]{ShepherdBarronGeography}.
Adjunction on a fibre gives $5\mid2g-2$, so $g\equiv1\pmod5$
and $g\geq6$. This uses arithmetic genus and does not require a
smooth general fibre.

For $g\geq6$, Gu--Sun--Zhou's nonhyperelliptic bound $n(g)$
dominates the hyperelliptic bound $h(g)$ and decreases with $g$.
Indeed,
\[
 n(g)-h(g)=
 \frac{4(g-1)(3g+2)(g^2-4g+2)}
 {g(g^2-4g+5)(g^2-g-1)}>0,
\]
\[
 n'(g)=-\frac{4\{5g^2((g-3)^2-2)+4(g-1)(3g-5)\}}
 {(g^3-4g^2+5g)^2}<0.
\]
Thus their theorem gives
\[
 \frac{K_T^2}{\chi(\OO_T)}\leq n(6)=\frac{800}{51}
 \qquad\text{\cite[Theorem~3.2, pp.~622--623]{GuSunZhou}.}
\]

If $\epsilon=0$, the ratio is $300/19$, and
$300/19-800/51=100/969>0$, a contradiction.
If $\epsilon=1$, then $p=2$, and the ratio is
$300d/(19d+1)$. The same inequality gives
\[
 15300d\leq15200d+800,
\]
so $d\leq8$.
On the original smooth surface, $K_X=11Y$ and
$c_2(X)=-12d+1$. The latter follows from the structure theorem's
quotient calculation: the blowup of the single internal center
increases $c_2$ by one and has the same $c_2$ as $T$.
Noether's formula on $X$ therefore yields
\[
 12\chi(\OO_X)=109d+1.
\]
Hence $d\equiv11\pmod{12}$, so $d\geq11$, contradicting
$d\leq8$. This excludes the internal case as well.
\end{proof}

\subsection{Positive self-intersection and a quantitative divisibility obstruction}
\label{subsec:positive-square-picard}

We prove the quantitative bound of Theorem~\ref{thm:positive-excess-picard}
using only positive
self-intersection of the distinguished curve. The ambient surface may
be normal and singular away from that curve. The argument keeps track
of the arithmetic genus of its inverse image under a normalized
$p$-root cover and of the loss of self-intersection needed to make that
inverse image smooth. This permits a finite sequence of root transfers
when the genus excess is sufficiently small.

\begin{lemma}\label{lem:positive-square-logarithmic}
Let $S$ be a normal integral projective surface over an algebraically
closed field $k$ of characteristic $p>0$. Let $C\subset S$ be a smooth
integral effective Cartier divisor contained in the smooth locus, with
$C^2>0$. If $B\ne\OO_S$ has order $p$ and $B|_C\simeq\OO_C$, its
logarithmic form $\theta$ is nonzero and has exact order $m\geq1$ along
$C$. Moreover $p\mid m$, and its first tangential coefficient is a
nonzero section of
\[
 \Omega_C^1\otimes\OO_C(-mC).
\]
In particular, $2g(C)-2\geq mC^2$. At equality that section is nowhere
zero.
\end{lemma}

\begin{proof}
Choose local frames of $B$ and a trivialization of $B^p$.
Let $u_i$ be the resulting units. Their logarithmic differentials
$du_i/u_i$ agree, since the ratios $u_i/u_j$ are $p$th powers.
They define a global K\"ahler form and, in particular, a rational
form $\theta$ that is regular on the smooth locus and closed there.
If this rational form were zero, each $u_i$ would have a $p$th root
in $k(S)$: over a perfect ground field, the kernel of the field
differential is $k(S)^p$. Each root and its inverse are integral over
the normal local rings, hence are regular units. Uniqueness of the
roots makes them compatible with the transition functions of $B$,
trivializing $B$. This is a contradiction.

The restriction of $\theta$ to $\Omega_C^1$ is zero. Indeed, after
trivializing $B|_C$, its $p$th-power trivialization is a global unit
on the proper integral curve, hence a constant. In the conormal
sequence its remaining component would be a section of
$\OO_C(-C)$, which has negative degree. It too is zero.
Thus the exact order $m$ is positive and finite.

Division by the $m$th power of a local equation of $C$ gives a
regular section of $\Omega_S^1(-mC)$ on a neighborhood of $C$.
This assertion holds at every point of $C$: in a regular local
ring, divisibility by a power of the equation of a smooth prime
divisor can be checked at its generic point. The resulting leading
restriction is nonzero. In
\[
0\longrightarrow\OO_C(-(m+1)C)
\longrightarrow\Omega_S^1(-mC)|_C
\longrightarrow\Omega_C^1(-mC)\longrightarrow0,
\]
the first term has negative degree. Therefore the tangential image
is nonzero. Locally write
\[
 \theta=h^m(A\,dh+B_1\,dx),\qquad C=(h=0),
\]
with $B_1$ nonzero at the generic point of $C$.
The coefficient of order $m-1$ in $d\theta$ is
$mB_1h^{m-1}dh\wedge dx$. Closedness gives $p\mid m$.
The nonzero tangential section has a line bundle of degree
$2g(C)-2-mC^2$, proving the inequality. If its degree is zero,
the section has no zeros.
\end{proof}

We will also use the conormal argument without assuming that the
one-form is logarithmic or closed. If a regular one-form has exact
order $j\geq0$ along a smooth proper curve of positive
self-intersection, its nonzero leading restriction has nonzero
tangential image. Otherwise it would give a nonzero section of
$\OO_C(-(j+1)C)$, a line bundle of negative degree.

\begin{lemma}[Resolution loss]\label{lem:positive-square-resolution-loss}
Let $N$ be a normal integral projective surface over an algebraically
closed field, and let $D$ be a reduced integral effective Cartier
divisor. Let $\overline D$ be its normalization, and put
\[
 \Delta_D=p_a(D)-g(\overline D).
\]
There is a projective resolution $f:T\to N$ for which the strict
transform $C_T$ of $D$ is smooth and
\[
 0\leq D^2-C_T^2\leq4\Delta_D.
\]
Here $C_T\simeq\overline D$. The assertion concerns a suitable
resolution, not every resolution obtained by further blowups.
\end{lemma}

\begin{proof}
Choose a projective resolution of $N$, an isomorphism over its
regular locus; such a resolution exists for surfaces of finite type
over a field \cite[Tag~0BGP]{Stacks}. Contract an exceptional smooth
rational curve of self-intersection $-1$ whenever one remains.
The contraction is projective, its target is smooth, and the map to
$N$ factors through it
\cite[Tags~0C5J, 0C2M and~0C2N]{Stacks}.
Each contraction removes an exceptional component, so the process
terminates. Denote the resulting resolution by $f_0:T_0\to N$.

For every exceptional integral curve $E_i$ on $T_0$, the Hodge index
theorem gives $E_i^2<0$: its intersection with the pullback of an
ample divisor on $N$ is zero, whereas that pullback has positive
square. Adjunction gives
\[
 K_{T_0}\cdot E_i=2p_a(E_i)-2-E_i^2\geq0.
\]
Indeed, the only possible negative value would require
$p_a(E_i)=0$ and $E_i^2=-1$. An integral proper curve of arithmetic
genus zero over an algebraically closed field is smooth rational,
and such an exceptional curve was contracted. Consequently
$K_{T_0}\cdot E\geq0$ for every effective exceptional divisor $E$.

Write
\[
 f_0^*D=D_0+E,
\]
where $D_0$ is the strict transform and $E$ is effective exceptional.
Since a pulled-back Cartier divisor has zero intersection with every
exceptional curve, $E^2=-D_0\cdot E$. Therefore
\begin{equation}\label{eq:positive-square-minimal-loss}
 \lambda_0:=D^2-D_0^2=D_0\cdot E\geq0.
\end{equation}

We check that $p_a(f_0^*D)=p_a(D)$, including the contribution of
nonrational surface singularities. Normality gives
$f_{0*}\OO_{T_0}=\OO_N$. The sheaf
$R^1f_{0*}\OO_{T_0}$ is supported on finitely many points, and
$R^if_{0*}\OO_{T_0}=0$ for $i>1$, since the fibres have dimension
at most one. Projection formula gives
\[
 R^1f_{0*}\OO_{T_0}(-f_0^*D)
   \simeq R^1f_{0*}\OO_{T_0}\otimes\OO_N(-D).
\]
Tensoring a finite-support sheaf by an invertible sheaf preserves
its length and hence its Euler characteristic. The two Leray
Euler-characteristic formulas consequently give
\begin{align*}
 \chi(\OO_{f_0^*D})
 &=\chi(\OO_{T_0})-\chi(\OO_{T_0}(-f_0^*D))\\
 &=\chi(\OO_N)-\chi(\OO_N(-D))
 =\chi(\OO_D).
\end{align*}
This proves the asserted arithmetic-genus equality, also when
$f_0^*D$ is nonreduced.

Set $\Delta_0=p_a(D_0)-g(\overline D)\geq0$.
Adjunction on the smooth surface $T_0$, applied to $D_0$ and to
$f_0^*D$, now yields
\begin{equation}\label{eq:positive-square-genus-loss}
 2(\Delta_D-\Delta_0)
    =\lambda_0+K_{T_0}\cdot E\geq\lambda_0.
\end{equation}
Blow up a singular point of the strict transform of $D_0$ whenever
one remains. For one such blowup $\sigma:U'\to U$, let $D_U$ be
the current strict transform, let $\mu\geq2$ be its multiplicity
at the center, let $F$ be the exceptional curve, and let $D_{U'}$
be the new strict transform. Then
\[
 D_{U'}=\sigma^*D_U-\mu F,\qquad
 K_{U'}=\sigma^*K_U+F,\qquad F^2=-1.
\]
Thus $D_{U'}^2=D_U^2-\mu^2$, and adjunction gives
$p_a(D_U)-p_a(D_{U'})=\mu(\mu-1)/2$.
The normalization and its genus are unchanged.
Thus the nonnegative integer $p_a-g(\overline D)$ strictly decreases,
and the process ends with a smooth strict transform $C_T$.
The total arithmetic-genus drop in this second stage is $\Delta_0$.
At each blowup,
\[
 \mu^2\leq2\mu(\mu-1)
       =4\left(\frac{\mu(\mu-1)}2\right),
\]
so the additional self-intersection loss is at most $4\Delta_0$.
Combining this with \eqref{eq:positive-square-genus-loss} gives
\[
 0\leq D^2-C_T^2
 \leq2(\Delta_D-\Delta_0)+4\Delta_0
 \leq4\Delta_D.
\]
The smooth strict transform is the normalization of $D$, as required.
\end{proof}

\begin{lemma}[Genus of the normalized root curve]
\label{lem:positive-square-root-genus}
Let $S$ be a smooth integral projective surface over an algebraically
closed field $k$ of characteristic $p>0$, and let $C\subset S$ be a
smooth integral curve with $d=C^2>0$. Let $B\ne\OO_S$ have order $p$
and be trivial on $C$. Write $m=pa$ for the exact order of its
logarithmic form, and put
\[
 G=2g(C)-2,\qquad \delta=G-md.
\]
Let $\pi:N\to S$ be the normalization of its $\mu_p$-torsor,
and set $D=\pi^*C$ scheme-theoretically. Then $D$ is a reduced
integral effective Cartier divisor,
\[
 D^2=pd,\qquad g(\overline D)=g(C),
\]
and there is an integer $\ell\geq0$ such that
\begin{equation}\label{eq:positive-square-conductor-genus}
 (p-1)\delta=\ell+2\Delta_D,\qquad
 \Delta_D=p_a(D)-g(C).
\end{equation}
In particular $\Delta_D\leq(p-1)\delta/2$.
\end{lemma}

\begin{proof}
Let $\tau:W\to S$ be the original torsor and $n:N\to W$ its
normalization. A local equation of $W$ is $z^p=u$, with $u$ a unit.
The normality argument in
Lemma~\ref{lem:positive-square-logarithmic} shows that $u$ is not a
$p$th power in $k(S)$. Thus $W$ is integral: its locally free
algebra embeds in the degree-$p$ purely inseparable function field.
Normalization is finite for varieties over a field
\cite[Tag~0BXR]{Stacks}, so $N$ is normal, integral and projective.
The finite map $\pi$ is radicial. Indeed, on a normal affine base
the $p$th power of an integral element of $k(N)$ belongs to $k(S)$
and is integral over the base ring, hence belongs to that ring.
This implies uniqueness of the prime above each prime and purely
inseparable residue extensions.

We first compute at the generic point of $C$, keeping its imperfect
residue field. On a sufficiently small smooth neighborhood choose
\'etale coordinates $h,x$, with $C=(h=0)$. If its ring is $R_0$,
the monomials $h^ix^j$, $0\leq i,j<p$, form a basis of $R_0$
over $R_0^p$, by the relative Frobenius isomorphism for \'etale
morphisms \cite[Tag~0EBS]{Stacks}. Write
\[
 u=\sum_{0\leq i,j<p}r_{ij}^{\,p}h^ix^j.
\]
Since $du$ is divisible by $h^{pa}$, comparison of the $dh$
coefficient detects every $i>0$, and comparison of the $dx$
coefficient detects every $j>0$. The factor $h^{pa}$ lies in
$R_0^p$, so freeness gives
$h^{pa}\mid r_{ij}^{\,p}$ for $(i,j)\ne(0,0)$.
Normality along the prime $(h)$ then gives $h^a\mid r_{ij}$.
Consequently
\begin{equation}\label{eq:positive-square-local-root}
 u=f^p+h^{pa}v,\qquad du=h^{pa}dv.
\end{equation}
The function $f$ is a unit along $C$. The tangential leading
coefficient of $du/u$ is $(v_x/u)\,dx$, so
$d\overline v\ne0$ in $\Omega^1_{k(C)/k}$.

Put $R=\OO_{S,\eta_C}$ and $F=k(C)$. The field $F$ is not being
assumed perfect. Since $k$ is perfect and $F/k$ has transcendence
degree one, $[F:F^p]=p$ and the kernel of
$d:F\to\Omega^1_{F/k}$ is $F^p$. Hence $\overline v\notin F^p$.
In the extension field set $t=(z-f)/h^a$. Then
\[
 t^p=v,\qquad z=f+h^at.
\]
The ring
\[
 \mathcal B=R[t]/(t^p-v)
\]
is a DVR with uniformizer $h$: it is finite free of rank $p$ over
$R$, and its quotient by $h$ is the field
$F(\overline v^{1/p})$. It is therefore the integral closure of
$R$ in the extension field. Its residue field is exactly
$F^{1/p}$, since both extensions of $F$ have degree $p$.
In particular $D$ is generically reduced.

The underlying inverse image of $C$ is irreducible because $\pi$
is radicial. A normal surface is Cohen--Macaulay, so its effective
Cartier divisor $D$ satisfies $S_1$. A generically reduced
$S_1$ scheme is reduced. Thus $D$ is reduced and integral.
Its normalization has function field $F^{1/p}$ and is a
Frobenius twist of $C$; base change by the inverse Frobenius
automorphism of the perfect field $k$ preserves its genus.
Thus $g(\overline D)=g(C)$. The finite projection formula gives
$D^2=pC^2$.

We now compute the arithmetic-genus defect. Multiplication in the
graded algebra of $W$, followed by projection to degree zero, is a
perfect pairing. It identifies its finite dualizing module with
the algebra itself, so
\[
 \omega_W\simeq\tau^*\omega_S.
\]
This pairing is not the field trace. Let
\[
 \mathfrak c
   =\mathcal Hom_{\OO_W}(n_*\OO_N,\OO_W)
\]
be the conductor, viewed as an ideal of $\OO_N$. Finite duality
and the preceding identification give
\[
 \omega_N\simeq\pi^*\omega_S\otimes\mathfrak c.
\]
Because $N$ is a normal surface, $\omega_N$ is a rank-one
reflexive sheaf. The same is true of $\mathfrak c$.

At the generic point of $D$, the original torsor order
$\mathcal A=R[z]$ and its normalization have the respective
$R$-bases
\[
 \mathcal A=\bigoplus_{i=0}^{p-1}Rh^{ai}t^i,
 \qquad
 \mathcal B=\bigoplus_{i=0}^{p-1}Rt^i.
\]
The first equality follows by using the basis
$1,z-f,\ldots,(z-f)^{p-1}$ of $R[z]$.
Since $\mathcal B$ is a DVR with uniformizer $h$, its conductor
into $\mathcal A$ is a power of $h$. The displayed bases show that
$h^b\mathcal B\subset\mathcal A$ if and only if
$b\geq a(p-1)$. Thus
\[
 \mathfrak c_{\eta_D}=h^{a(p-1)}\mathcal B.
\]

Define the rank-one reflexive sheaf
\[
 J=\mathfrak c\otimes\OO_N(a(p-1)D).
\]
At the generic point of $D$ it is the unit ideal; at every other
codimension-one point it is an ideal because $\mathfrak c$ is.
Reflexivity on the normal surface therefore gives an inclusion
$J\subset\OO_N$. Put $Q=\OO_N/J$.
At every closed point, both $J$ and $\OO_N$ have depth two.
The depth lemma shows that $Q$ has depth at least one wherever it
is nonzero. Hence $Q$ has no associated point of codimension two.
It has no associated component equal to $D$, since $J$ is the unit
ideal generically on $D$. A local equation of $D$ is consequently a
nonzerodivisor on $Q$. Tensoring
$0\to J\to\OO_N\to Q\to0$ with $\OO_D$ gives an injection
\[
 J|_D\longrightarrow\OO_D.
\]
Its cokernel is supported at finitely many points. Denote its
length by $\ell\geq0$.

Adjunction for a Cartier divisor on a Cohen--Macaulay surface gives
\[
 \omega_D\simeq H\otimes J|_D,\qquad
 H=\pi^*\!\left(\omega_S((1-a(p-1))C)\right)|_D.
\]
Here $H$ is invertible, although $\omega_D$ need not be.
Duality on the integral projective curve $D$ gives
$\chi(\omega_D)=-\chi(\OO_D)$.
The injection just proved, tensored with $H$, has cokernel of
length $\ell$. Riemann--Roch for the invertible sheaf $H$ therefore
yields
\begin{align*}
 2p_a(D)-2
 &=\deg H-\ell\\
 &=p\bigl(K_S\cdot C+(1-a(p-1))d\bigr)-\ell\\
 &=pG-m(p-1)d-\ell.
\end{align*}
Since $2p_a(D)-2=G+2\Delta_D$, rearrangement gives
\eqref{eq:positive-square-conductor-genus}.
No ampleness assumption entered this calculation.
\end{proof}

\begin{lemma}[Root transfer after resolution]
\label{lem:positive-square-resolved-transfer}
Under the hypotheses and notation of
Lemma~\ref{lem:positive-square-root-genus}, suppose
$B\simeq A^{\otimes p}$ for a line bundle $A$ on $S$.
There is a projective resolution $f:T\to N$ on which the strict
transform $C_T$ of $D$ is smooth and
\begin{equation}\label{eq:positive-square-transfer-square}
 g(C_T)=g(C),\qquad
 d_T:=C_T^2\geq pd-2(p-1)\delta.
\end{equation}
If $d_T>0$, the bundle $B_T=(\pi f)^*A$ is nontrivial of order $p$,
is trivial on $C_T$, and its logarithmic form has exact order
$a=m/p$. In particular $p\mid a$, and
\begin{equation}\label{eq:positive-square-excess-recurrence}
 0\leq\delta_T:=G-a d_T
       \leq\bigl(1+2(p-1)a\bigr)\delta.
\end{equation}
\end{lemma}

\begin{proof}
Apply Lemma~\ref{lem:positive-square-resolution-loss} to $D$.
Lemma~\ref{lem:positive-square-root-genus} gives
\[
 d_T\geq D^2-4\Delta_D
       \geq pd-2(p-1)\delta
\]
and identifies the normalization genus with $g(C)$.
This proves \eqref{eq:positive-square-transfer-square}.

Use frames of $A$, with transitions $g_{ij}$, and their $p$th
powers as frames of $B$. The torsor coordinates satisfy
$z_i=g_{ij}^p z_j$. They are regular units on the original torsor
and remain units after normalization and resolution. Thus
\[
 \theta_T=\frac{dz_i}{z_i}
\]
is a global regular closed one-form on $T$. It is the logarithmic
form of $B_T$, whose $p$th power is trivialized by the $z_i$.

At the generic point of $C_T$, the equations in
\eqref{eq:positive-square-local-root} give
$t^p=v$ and $z=f+h^at$. The derivative $v_x$ is a unit there.
The normalization is smooth on this open set, the resolution is an
isomorphism at the generic point of $D$, and $dh,dt$ are a
cotangent basis. The derivation dual to $dt$ kills the original
base ring: it kills $h$, and differentiation of $t^p=v$ then makes
it kill $x$; \'etaleness deals with the remaining functions.
In particular it kills $f$, even when $f$ is not constant.
The coefficient of $dt$ in $\theta_T$ is therefore exactly
\begin{equation}\label{eq:positive-square-transfer-coefficient}
 \frac{h^a}{z}.
\end{equation}
Hence $\theta_T$ is nonzero and its order $j$ along $C_T$ satisfies
$0\leq j\leq a$.

Assume $d_T>0$. If $j<a$, the leading restriction of $\theta_T$
would have zero tangential image: generically the coefficient of
$dt$ has order $a$, while $dh$ is conormal. This would give a
nonzero global section of
$\OO_{C_T}(-(j+1)C_T)$, contradicting its negative degree.
The section is global because divisibility by $h^j$ can be checked
at the generic point in the regular local rings along $C_T$.
Thus $j=a$.

A trivial line bundle on a proper integral variety has zero
logarithmic form for every $p$th-power trivialization, because
global units are constants. Consequently $B_T$ is nontrivial.
Its $p$th power is trivial, so its order is exactly $p$.
Since $a\geq1$, its logarithmic form restricts to zero on $C_T$.
The logarithmic-form map on $p$-torsion line bundles on a smooth
proper curve is injective: a defining unit of zero differential is
a $p$th power in the function field, normality makes the unique
roots local units, and those roots glue. Therefore
$B_T|_{C_T}\simeq\OO_{C_T}$.
Lemma~\ref{lem:positive-square-logarithmic} now gives $p\mid a$
and $\delta_T\geq0$. Finally
\[
 \delta_T
 =G-a d_T
 \leq G-apd+2(p-1)a\delta
 =\bigl(1+2(p-1)a\bigr)\delta,
\]
as asserted.
\end{proof}

The original equality transfer remains a useful special case:
normalization alone already gives the required smooth curve.

\begin{lemma}\label{lem:positive-square-root-transfer}
Under the hypotheses of Lemma~\ref{lem:positive-square-logarithmic},
suppose the tangential leading coefficient is nowhere zero and
$B\simeq A^{\otimes p}$ for an invertible sheaf $A$ on $S$.
Write $m=pq$. Let $\pi:S'\to S$ be the normalization of the
$\mu_p$-torsor associated with $B$.
Then $S'$ is a normal integral projective surface and is smooth
near the smooth integral Cartier divisor $C'=\pi^*C$.
The bundle $B'=\pi^*A$ is nontrivial of order $p$, is trivial on
$C'$, and its logarithmic form $\theta'$ has exact order $q$,
with nowhere-zero tangential leading coefficient. Furthermore,
\[
 (C')^2=pC^2>0,\qquad p\mid q.
\]
\end{lemma}

\begin{proof}
The construction and integrality of the normalized torsor in
Lemma~\ref{lem:positive-square-root-genus} use only normality of
the ambient surface; the local calculation uses only its smooth
neighborhood of $C$. Thus they apply here even when $S$ is singular
away from $C$. At every point of $C$, not just generically, the
nowhere-zero leading coefficient makes $v_x$ a unit in
\eqref{eq:positive-square-local-root}. The finite algebra
\[
 R_0[t]/(t^p-v),\qquad z=f+h^qt,
\]
is then smooth near $h=0$, and its divisor $h=0$ is smooth and
reduced. Its function field is the required extension. After
shrinking around that divisor it is normal, hence is the full
integral closure of $R_0$. These charts identify it with $S'$
near the scheme-theoretic inverse image $C'$. Radiciality makes
$C'$ integral. Its map to $C$ has degree $p$, and its normal
bundle is the pullback of the normal bundle of $C$; hence
$(C')^2=pC^2$.

The proof of Lemma~\ref{lem:positive-square-resolved-transfer}
now applies directly on the smooth neighborhood of $C'$.
The form $dz/z$ has exact order $q$, is nonzero, and gives the
stated nontrivial order-$p$ bundle trivial on $C'$; closedness
gives $p\mid q$. No resolution away from this neighborhood is
needed for these conclusions. Finally
\eqref{eq:positive-square-transfer-coefficient}, with $a=q$,
holds in each of the smooth charts along $C'$. After division by
$h^q$ its tangential coefficient is $1/z$, a unit. It is therefore
nowhere zero, proving the final assertion.
\end{proof}

\begin{proof}[Proof of Theorem~\ref{thm:positive-excess-picard}]
Lemma~\ref{lem:positive-square-logarithmic} gives $m>0$, $p\mid m$
and $\delta=G-md\geq0$, where $G=2g(Y)-2$ and $d=Y^2$.
In particular $r=v_p(m)\geq1$.

First suppose that $M=P^{\otimes p^r}$ for a line bundle
$P$ on $X$. Resolve $X$ by a projective resolution that is an
isomorphism near $Y$. The curve, its self-intersection, its genus,
and the logarithmic order do not change. Pullback on Picard groups
is injective: if $f^*L$ is trivial, normality gives
$f_*\OO=\OO_X$ and projection formula gives $L\simeq\OO_X$.
Thus the pulled-back $M$ is still nontrivial, and the chosen root
$P$ pulls back as well. We may therefore assume $X$ smooth.

Set
\[
 m_i=\frac{m}{p^i},\qquad
 Q_0=1,\qquad
 Q_i=\prod_{j=1}^{i}\bigl(1+2(p-1)m_j\bigr)
 \quad(1\leq i\leq r).
\]
Suppose for contradiction that $Q_r\delta<G$.
We construct successive data $X_i,Y_i,M_i$ with
$X_0=X$, $Y_0=Y$ and $M_0=M$. Write
\[
 d_i=Y_i^2,\qquad \delta_i=G-m_i d_i.
\]
At stage $i$, the bundle $M_i$ is the pullback of
$P^{\otimes p^{r-i}}$ through the composite map to $X$.
For $i<r$ it therefore has the $p$th root given by the pullback
of $P^{\otimes p^{r-i-1}}$.

Inductively assume $d_i>0$, that the logarithmic order is $m_i$,
and that $0\leq\delta_i\leq Q_i\delta$.
Apply Lemma~\ref{lem:positive-square-resolved-transfer} using that
chosen root. Its resolution and smooth strict transform satisfy
\begin{align*}
 m_{i+1}d_{i+1}
 &\geq G-\bigl(1+2(p-1)m_{i+1}\bigr)\delta_i\\
 &\geq G-Q_{i+1}\delta\\
 &\geq G-Q_r\delta>0.
\end{align*}
Thus the next distinguished curve still has positive
self-intersection. The lemma gives a nontrivial order-$p$ bundle
$M_{i+1}$ trivial on it, with exact logarithmic order $m_{i+1}$,
the same genus, and
\[
 0\leq\delta_{i+1}
 \leq\bigl(1+2(p-1)m_{i+1}\bigr)\delta_i
 \leq Q_{i+1}\delta.
\]
This proves the induction. At the final stage,
Lemma~\ref{lem:positive-square-logarithmic} forces $p\mid m_r$.
But $m_r=m/p^r$ is prime to $p$ by the definition of $r$.
The contradiction proves
\[
 Q_r(m)\delta\geq G
\]
whenever $M$ is a $p^r$th tensor power.

It remains to lift divisibility in the N\'eron--Severi group.
For a normal integral projective variety over $k$, the reduced
identity component of the Picard scheme is an abelian variety,
and its $k$-points are the algebraically trivial line-bundle classes
\cite[Theorems~2.5 and~5.4, Proposition~5.10 and
Remark~5.25, pp.~18--20, 37--38, 41--42 and~51]{KleimanPicard}.
Multiplication by $p^r$ is surjective on these points
\cite[Tag~03RP, parts~(3)--(4)]{Stacks}.
If $[M]=p^r[L_0]$ in $\operatorname{NS}(X)$, then
$M\otimes L_0^{-p^r}$ is algebraically trivial. Choose an
algebraically trivial line bundle $L_1$ with
\[
 L_1^{\otimes p^r}\simeq M\otimes L_0^{-p^r}.
\]
Then $M\simeq(L_0\otimes L_1)^{\otimes p^r}$, so the inequality
just proved applies. Since $G=md+\delta$, it is equivalent to
\[
 \bigl(Q_r(m)-1\bigr)\delta\geq md.
\]
An algebraically trivial $M$ has zero N\'eron--Severi class and
satisfies the divisibility hypothesis. This proves the product bound.
The strictness assertion for $r=1$ and odd $p$ is proved in
Theorem~\ref{thm:odd-positive-excess-strict} below, using this weak
bound and its equality cases.
\end{proof}

\begin{corollary}\label{thm:positive-square-picard}
Let $X$ be a normal integral projective surface over an algebraically
closed field $k$ of characteristic $p>0$, and let $Y\subset X$ be a
smooth integral effective Cartier divisor contained in the smooth
locus, with $Y^2>0$. Let $M\ne\OO_X$ have order $p$ and satisfy
$M|_Y\simeq\OO_Y$. If $m$ is the exact order along $Y$ of its
logarithmic form, then
\[
 m\geq1,\qquad p\mid m,\qquad 2g(Y)-2\geq mY^2.
\]
If equality holds and $r=v_p(m)$, then
\[
 [M]\notin p^r\operatorname{NS}(X).
\]
In particular, every algebraically trivial such $M$ satisfies
\[
 2g(Y)-2>mY^2.
\]
\end{corollary}

\begin{proof}
The first assertions are part of
Theorem~\ref{thm:positive-excess-picard}; its quantitative inequality
makes $\delta=0$ impossible when
$[M]\in p^r\operatorname{NS}(X)$, since $mY^2>0$.
\end{proof}

When $v_p(m)=1$, write $m=pb$ with $p\nmid b$.
Then $Q_1(m)-1=2(p-1)b$, and
Theorem~\ref{thm:positive-excess-picard} gives the particularly simple
necessary condition
\begin{equation}\label{eq:positive-square-first-valuation}
 \delta\geq\frac{p}{2(p-1)}Y^2
\end{equation}
for $[M]\in p\operatorname{NS}(X)$, in particular for algebraically
trivial $M$. In characteristic two this says $\delta\geq Y^2$.
The argument uses only the finite root chosen on the original
surface; it requires no Picard divisibility on an intermediate
surface. It imposes a quantitative restriction on positive excess,
but does not decide the existence of a kernel class beyond that
bound.


\subsection{Strictness of the first-valuation bound in odd characteristic}
\label{subsec:odd-positive-excess-strict}

For odd characteristic the first-valuation estimate
\eqref{eq:positive-square-first-valuation} is strict.
The proof keeps the equality cases in both the genus identity and the
resolution-loss estimate. These equality cases would force cusps of
type $(2,p)$ on the normalized root curve. A regular logarithmic form
on the normalized surface then has a conormal leading coefficient
that vanishes above each cusp, whereas the zero self-intersection
of the smooth strict transform requires that coefficient to be
nowhere zero.

\begin{theorem}\label{thm:odd-positive-excess-strict}
Let $X$ be a normal integral projective surface over an algebraically
closed field $k$ of odd characteristic $p$, and let $C\subset X$ be
a smooth integral effective Cartier divisor contained in the smooth
locus, with $d=C^2>0$.
Let $M\ne\OO_X$ have order $p$ and satisfy
$M|_C\simeq\OO_C$.
Suppose that its logarithmic form has exact order
$m=pa$, where $p\nmid a$.
If $[M]\in p\operatorname{NS}(X)$, then
\[
 \delta:=2g(C)-2-pa\,d>\frac{p\,d}{2(p-1)}.
\]
Equivalently,
\[
 \delta\geq
 \left\lfloor\frac{p\,d}{2(p-1)}\right\rfloor+1.
\]
In particular the strict inequality holds for every algebraically
trivial such $M$.
\end{theorem}

\begin{proof}
Theorem~\ref{thm:positive-excess-picard} gives the weak inequality.
We shall exclude
\begin{equation}\label{eq:odd-endpoint-assumption}
 2(p-1)\delta=pd.
\end{equation}

\emph{Step 1: the root cover and equality in the losses.}
We first choose the root on the original surface.
By the Picard-variety argument in the proof of
Theorem~\ref{thm:positive-excess-picard}, the reduced identity
component of the Picard scheme is an abelian variety, and
multiplication by $p$ is surjective on its $k$-points.
Choose $L_0$ representing a N\'eron--Severi root of $M$.
Then $M\otimes L_0^{-p}$ is algebraically trivial, so it has
an algebraically trivial $p$th root $L_1$.
Thus
\[
 A:=L_0\otimes L_1,\qquad A^{\otimes p}\simeq M.
\]
Resolve $X$ projectively by a morphism $\rho:S\to X$ that is an
isomorphism near $C$.
The curve, its self-intersection, its genus and the logarithmic
order are unchanged.
Normality gives $\rho_*\OO_S=\OO_X$, and projection formula makes
pullback on Picard groups injective.
We retain the notation $C,M,A$ for their pullbacks to the smooth
surface $S$.

Let $\pi:N\to S$ be the normalization of the $\mu_p$-torsor of $M$,
and put $D=\pi^*C$ scheme-theoretically.
Lemma~\ref{lem:positive-square-root-genus} says that $N$ is normal,
integral and projective, that $D$ is a reduced integral Cartier
divisor, and that
\[
 D^2=pd,\qquad g(\overline D)=g(C),\qquad
 (p-1)\delta=\ell+2\Delta,
\]
where
\[
 \Delta=p_a(D)-g(C),\qquad \ell\geq0.
\]
The normalization map $\phi:\overline D\to C$ is the degree-$p$
purely inseparable map with function field $k(C)^{1/p}$.

Use the resolution $f:T\to N$ constructed in the proof of
Lemma~\ref{lem:positive-square-resolution-loss}, and write $F$
for the smooth strict transform of $D$.
By \eqref{eq:odd-endpoint-assumption},
\[
 F^2\geq D^2-4\Delta
       \geq pd-2(p-1)\delta=0.
\]
If $F^2>0$, Lemma~\ref{lem:positive-square-resolved-transfer},
applied to the chosen root $A$, gives exact logarithmic order $a$
and forces $p\mid a$.
Therefore
\begin{equation}\label{eq:odd-endpoint-zero-square}
 F^2=0.
\end{equation}
It follows that equality holds throughout
\[
 pd=D^2-F^2\leq4\Delta\leq2(p-1)\delta=pd.
\]
In particular,
\begin{equation}\label{eq:odd-endpoint-genus}
 \Delta=\frac{pd}{4}>0,\qquad \ell=0.
\end{equation}

We retain the equality information in the resolution construction.
Its first stage is the resolution $f_0:T_0\to N$ with no exceptional
smooth rational curve of square $-1$.
Write
\[
 f_0^*D=D_0+E,\qquad
 \lambda_0=D^2-D_0^2=D_0\cdot E=-E^2,\qquad
 \Delta_0=p_a(D_0)-g(C).
\]
The proof of Lemma~\ref{lem:positive-square-resolution-loss} gives
\[
 2(\Delta-\Delta_0)
   =\lambda_0+K_{T_0}\cdot E\geq\lambda_0\geq0.
\]
In its second stage, let $\mu_i\geq2$ be the multiplicities
at the successive singular points of the strict curve.
Then
\[
 \sum_i\frac{\mu_i(\mu_i-1)}2=\Delta_0,\qquad
 D^2-F^2=\lambda_0+\sum_i\mu_i^2.
\]
Consequently
\begin{align*}
 pd
 &=\lambda_0+\sum_i\mu_i^2\\
 &\leq2(\Delta-\Delta_0)+4\Delta_0\\
 &=4\Delta-2(\Delta-\Delta_0)\\
 &\leq4\Delta=pd.
\end{align*}
Equality of the endpoints forces
\begin{equation}\label{eq:odd-endpoint-resolution-equalities}
 \Delta=\Delta_0,\qquad \lambda_0=0,\qquad
 \mu_i=2\ \text{for every }i.
\end{equation}

Since $E^2=0$, we have $E=0$.
Indeed, a nonzero effective exceptional divisor is orthogonal to
the pullback of an ample divisor on $N$ and has negative square
by the Hodge index theorem.
It cannot be numerically zero, since its intersection with an
ample divisor on $T_0$ is positive.
Thus $D$ avoids the singular locus of $N$:
if a singular point belonged to $D$, a local equation of $D$
would vanish along every exceptional curve above it and give
a positive coefficient in $E$.
A resolution of a singular normal surface point has an
exceptional curve; otherwise proper quasi-finiteness and
normality would make the resolution an isomorphism there.
We conclude that $N$ is smooth on a neighborhood of $D$.
This conclusion uses $\lambda_0=0$, rather than merely $\ell=0$.

\emph{Step 2: the singularities of the root curve.}
Let $\beta$ be the first tangential coefficient of the logarithmic
form of $M$, and let $Z$ be its zero divisor on $C$.
It is a nonzero section of
$\omega_C\otimes\OO_C(-paC)$, so
\[
 \deg Z=\delta.
\]
Where $\beta$ does not vanish, the local normalization chart
in Lemma~\ref{lem:positive-square-root-genus} has $v_x$ invertible.
The chart and its inverse-image curve are smooth there.
Since $\pi$ is radicial,
\begin{equation}\label{eq:odd-endpoint-count}
 \#\operatorname{Sing}(D)
 \leq\#\operatorname{Supp}(Z)\leq\deg Z.
\end{equation}

Let $P$ be a singular point of $D$.
By \eqref{eq:odd-endpoint-resolution-equalities}, its multiplicity
in the smooth surface $N$ is two.
The normalization has a unique point above $P$, because
$\phi:\overline D\to C$ is radicial.
Completion of the finite normalization injection embeds
$\widehat{\OO}_{D,P}$ in $k[[s]]$, the completed ring of that
normalization point.
The normalization defect is unchanged by completion
\cite[Tags~0C1R and~0C3V]{Stacks}.
Thus the completed plane curve has one branch.

We spell out the multiplicity-two calculation, using $p\ne2$.
The multiplicity of a plane branch is the minimum valuation of
a general linear combination of its two plane parameters.
For this valuation formula, compare the quotient by such a
parameter with the quotient in the finite normalization:
the finite-length normalization defect contributes equal
kernel and cokernel lengths and cancels.
There is therefore an element $x$ of valuation two.
Taking the square root of its unit factor in $k[[s]]$ permits a
normalization parameter with $x=s^2$.
Since every element of the square of the maximal ideal has
valuation at least four, $x$ is part of a minimal pair of
plane parameters.
Subtract the even part of the second parameter's power series
in $s$, and divide its remaining odd part by its unit power
series in $s^2$.
These are invertible changes of complete plane coordinates.
They give a second parameter $y=s^q$ for an odd integer $q\geq3$.
An odd term must occur: otherwise the completed branch's
fraction field would be contained in $k((s^2))$.
Finite birational normalization remains birational after this
one-branch completion, since a common nonzero denominator
for its finite module remains a common denominator after
completion.
We have proved
\[
 \widehat{\OO}_{D,P}=k[[s^2,s^q]],\qquad
 D:\ y^2-x^q=0.
\]

A parameter on $C$ pulls back under $\phi$ with valuation $p$.
Thus the semigroup of $k[[s^2,s^q]]$ contains the odd integer $p$.
Its least odd value is $q$, so $q\leq p$.
The missing monomials $s,s^3,\ldots,s^{q-2}$ show that the
local normalization defect is
\[
 \Delta_P=\frac{q-1}{2}\leq\frac{p-1}{2}.
\]
The normalization exact sequence gives
$\Delta=\sum_{P\in\operatorname{Sing}(D)}\Delta_P$
\cite[Tag~0C1R]{Stacks}.
Using \eqref{eq:odd-endpoint-genus} and
\eqref{eq:odd-endpoint-count}, we obtain
\[
 \frac{p-1}{2}\deg Z
 =\Delta
 \leq\frac{p-1}{2}\#\operatorname{Sing}(D)
 \leq\frac{p-1}{2}\deg Z.
\]
Every inequality is an equality.
Hence $Z$ is reduced, there is exactly one singular point of
$D$ over each point of $Z$, and every such point has the formal
equation
\begin{equation}\label{eq:odd-endpoint-cusp}
 y^2=x^p,\qquad x=s^2,\quad y=s^p.
\end{equation}
There is at least one such cusp, since $\Delta>0$.

\emph{Step 3: a degree-zero conormal leading coefficient.}
Use frames of $A$ with transitions $g_{ij}$ and their $p$th
powers as frames of $M$.
The torsor coordinates satisfy $z_i=g_{ij}^p z_j$ and remain
regular units on $N$.
Consequently
\[
 \theta_N=\frac{dz_i}{z_i}
\]
is a global regular K\"ahler one-form on $N$, closed on its
smooth neighborhood of $D$.
Its pullback $\theta_T=f^*\theta_N$ is regular and closed on $T$.
At the generic point of $D$, the calculation in
Lemma~\ref{lem:positive-square-resolved-transfer} gives
\[
 z=b+h^a t,\qquad t^p=v,
\]
with $dh,dt$ a cotangent basis.
The derivation dual to $dt$ kills the old base ring, including
the possibly nonconstant $b$.
Thus the $dt$-coefficient of $\theta_T$ is exactly $h^a/z$.
The form is nonzero, and its exact order $j$ along $F$ satisfies
$0\leq j\leq a$.
If $j=a$, its tangential leading coefficient is nonzero.
The coefficient of order $a-1$ in its exterior derivative
then forces $p\mid a$, a contradiction.
Therefore $j<a$, and its leading restriction is conormal.

Divisibility by a power of the local equation of the smooth
curve $F$ is checked at its generic point.
The nonzero leading restriction is thus a global section
\begin{equation}\label{eq:odd-endpoint-leading-section}
 0\ne\sigma\in
 H^0\bigl(F,\OO_F(-(j+1)F)\bigr).
\end{equation}
Its tangential image vanishes generically, hence everywhere.
By \eqref{eq:odd-endpoint-zero-square}, this line bundle has
degree zero.
The section $\sigma$ must therefore be nowhere zero.

\emph{Step 4: the zero forced at each cusp.}
Fix a cusp and use the complete smooth surface coordinates
in \eqref{eq:odd-endpoint-cusp}.
Put $h_0=y^2-x^p$.
The order of $\theta_N$ along $D$ is $j$, since the resolution
is an isomorphism at the generic point of $D$.
Its coefficients are divisible by $h_0^j$ in the regular
local ring: divisibility by this power of a prime Cartier
equation is detected at the generic point, even though
the Cartier curve is singular.
After completion we may write
\[
 \theta_N=h_0^j\bigl(A(x,y)\,dx+B(x,y)\,dy\bigr)
\]
with $A,B$ regular power series.
The expression in parentheses has zero tangential image
on the normalization of $D$.
Since $x=s^2$ and $y=s^p$, that image is
\[
 2sA(s^2,s^p)\,ds.
\]
Here $2s\ne0$, so $A(s^2,s^p)=0$.
Thus $A=h_0A_1$ and
\begin{equation}\label{eq:odd-endpoint-local-form}
 \theta_N
   =h_0^j\bigl(h_0A_1(x,y)\,dx+B(x,y)\,dy\bigr).
\end{equation}

Write $r=(p-1)/2$.
The successive curve blowups at this cusp have equations
$t_i^2=x^{p-2i}$.
There are exactly $r$ of them before the strict curve is smooth.
The final chart is
\[
 y=x^r t,\qquad F:\ t^2=x.
\]
Blowups at other cusps do not change this chart.
Set $H=x-t^2$, a local equation of $F$.
Then
\[
 x=t^2+H,\qquad y=x^r t,\qquad h_0=-x^{p-1}H.
\]
Along $F$,
\[
 dx|_F=2t\,dt+dH
\]
and
\begin{align*}
 dy|_F
 &=r t^{2r-1}\,dH+(2r+1)t^{2r}\,dt\\
 &=r t^{p-2}\,dH,
\end{align*}
because $2r+1=p=0$ in $k$.
The scalar $r$ is nonzero in $k$.
Pull \eqref{eq:odd-endpoint-local-form} to this chart, divide
by $H^j$, and restrict to $F$.
The term with $h_0A_1\,dx$ still has a factor $H$ and vanishes.
The remaining conormal coefficient is
\[
 (-1)^j r\,t^{\,2j(p-1)+p-2}B(t^2,t^p).
\]
It vanishes at $t=0$, since $B$ is regular and
\[
 2j(p-1)+p-2\geq p-2>0.
\]
Multiplying a local equation by a unit does not change
this vanishing.
Faithful flatness of completion shows that it is a zero
of the original section
\eqref{eq:odd-endpoint-leading-section}.
This contradicts its nowhere-vanishing property.
Therefore \eqref{eq:odd-endpoint-assumption} is impossible,
and the weak inequality is strict.
\end{proof}

The proof uses a regular form on the normalized surface and its
behavior at the resolved cusps.
It requires no inference of semiampleness or of a fibration
from the zero self-intersection of the strict transform.

\subsection{The equality boundary in characteristic two}
\label{subsec:positive-excess-boundary}

We now prove that the first positive-excess bound can be attained
in characteristic two only under a further restriction on the curve.
The proof examines equality in the resolution estimate. The resolved
curve has square zero, so a conormal coefficient need not vanish.
Its descent, together with a leading two-form on the original
surface, instead detects a nontrivial two-torsion line bundle on
the curve.

\begin{proof}[Proof of the discrepancy and nonvanishing assertions in Corollary~\ref{cor:positive-excess-characteristic-two-boundary}]
Write $C=Y$, $N=\OO_C(C)$ and $g=g(C)$. Thus
\[
 2g-2=(2a+1)d,\qquad
 \mathcal T=\omega_C\otimes N^{-(2a+1)},\qquad
 \deg\mathcal T=0.
\]
The lifting argument in the proof of
Theorem~\ref{thm:positive-excess-picard} gives a line bundle $L$
with $L^{\otimes2}\simeq M$. It has exact order four.
Resolve $X$ away from $C$, as in that proof. All the displayed
curve data remain unchanged, so we may assume that $X$ is smooth.

\smallskip\noindent
\emph{A finite model near the curve.}
Use frames of $L$ and their squares as frames of $M$. If $g_{ij}$
are the transitions of these frames of $M$, then $dg_{ij}=0$.
Let $f_i$ be the units defining the chosen trivialization of
$M^{\otimes2}$, and let $t_i=0$ define $C$. The local
$p$-basis calculation in
Lemma~\ref{lem:positive-square-root-genus} gives
\begin{equation}\label{eq:boundary-square-approximation}
 f_i=b_i^2+t_i^{2a}A_i,
\end{equation}
where $b_i$ is a unit after shrinking around $C$.
Since $f_i=g_{ij}^2f_j$, we have
$b_i-g_{ij}b_j\in I_C^a$.
Consequently the finite flat double covers
\[
 y_i^2=A_i,\qquad z_i=b_i+t_i^ay_i
\]
glue by
\[
 y_i=g_{ij}(t_j/t_i)^a y_j+
       \frac{g_{ij}b_j-b_i}{t_i^a}.
\]
Denote the resulting surface near $C$ by $V$ and its curve over $C$
by $B_V$. The tangential leading coefficient of the logarithmic
form of $M$ is nonzero, so $d(A_i|_C)\ne0$ generically.
Thus $B_V$ is generically a field extension of degree two.
Its algebra is locally free over the integral curve $C$, hence
torsion-free; generic integrality makes $B_V$ integral.

The submodule of constants in that rank-two algebra has quotient
with transition functions $g_{ij}(t_j/t_i)^a$. Since $M|_C$
is trivial, this gives the exact sequence of $\OO_C$-modules
\[
 0\longrightarrow\OO_C\longrightarrow
 (\pi_V)_*\OO_{B_V}\longrightarrow N^a\longrightarrow0.
\]
Riemann--Roch on $C$ therefore gives
\begin{equation}\label{eq:boundary-finite-genus}
 p_a(B_V)-g=g-1-ad=\frac d2.
\end{equation}
In particular $d$ is even.

Let $\pi:W\to X$ be the normalization of the original torsor.
Near $C$ it is also the normalization of $V$.
By Lemma~\ref{lem:positive-square-root-genus},
$B=\pi^*C$ is reduced, integral and Cartier, $B^2=2d$,
and its normalization has genus $g$.
Apply the resolution procedure in
Lemma~\ref{lem:positive-square-resolution-loss}: first take a
minimal ambient resolution $S_0\to W$, and then blow up singular
points of the strict curve until its transform $D$ on a smooth
projective surface $S$ is smooth. The map $q:D\to C$ is finite,
purely inseparable of degree two, and is relative Frobenius
\cite[Tag~0CCZ]{Stacks}. The preceding genus bound gives
\begin{equation}\label{eq:boundary-nonnegative-square}
 D^2\geq2d-2\delta=0.
\end{equation}

\smallskip\noindent
\emph{The conormal leading coefficient.}
The pullback $L'$ of $L$ to $S$ has square trivialized by the
units $z_i$. Its logarithmic form is the global regular form
$\eta=dz_i/z_i$. At the generic point of $C$, omit indices.
The differentials $dt,dA$ form a basis over $k(X)$, so write
\[
 db=R\,dt+S_1\,dA,\qquad R,S_1\in k(X).
\]
On the cover $A=y^2$, and $a$ is odd. Hence
\begin{equation}\label{eq:boundary-root-differential}
 \eta=\frac{(R+t^{a-1}y)\,dt+t^a\,dy}{z}.
\end{equation}
Here $t$ is a uniformizer along $D$, and $dt,dy$ form a
cotangent basis at its generic point. In particular $\eta\ne0$.
A trivial line bundle on a proper integral variety has zero
logarithmic differential for a square trivialization, because its
global units are constants. Thus $L'$ has exact order two.
The image of $\eta$ in $\Omega_D^1$ is zero: $dt$ is conormal
and the $dy$ coefficient is divisible by $t^a$.
The injectivity of the logarithmic-form map on torsion line bundles
of a smooth proper curve, established in the proof of
Lemma~\ref{lem:positive-square-resolved-transfer}, shows that
$L'|_D$ is trivial. None of these assertions requires $D^2>0$.

The residue $y_0$ lies in $k(D)=k(C)^{1/2}$ and not in $k(C)$.
Every normalized leading residue of $R$ lies in $k(C)$.
Therefore the two terms of $R+t^{a-1}y$ cannot cancel when
their orders coincide. If $h=\operatorname{ord}_D(\eta)$, then
\begin{equation}\label{eq:boundary-conormal-order}
 0\leq h=\min\{v_C(R),a-1\}\leq a-1,
\end{equation}
where the valuation of zero is infinity. Regularity gives the
lower bound. The leading coefficient is purely conormal, and
defines a nonzero global section
\begin{equation}\label{eq:boundary-conormal-section}
 \sigma\in H^0\bigl(D,N_{D/S}^{-(h+1)}\bigr).
\end{equation}
Here division by the $h$th power of a local equation is regular
along the entire smooth curve: in the regular local rings, this
divisibility can be checked at its generic point.
If $D^2>0$, this line bundle would have negative degree.
Combining with \eqref{eq:boundary-nonnegative-square} gives
$D^2=0$. The section $\sigma$ is therefore nowhere zero.

\smallskip\noindent
\emph{Equality in the genus and resolution losses.}
Let $B_0$ be the strict curve on $S_0$, and put
\[
 e=p_a(B_V)-p_a(B),\quad
 u=p_a(B)-p_a(B_0),\quad
 v=p_a(B_0)-g.
\]
The first number is the length of the quotient for the finite
birational map $B\to B_V$, and the other two are the
nonnegative genus losses in the resolution proof.
By \eqref{eq:boundary-finite-genus}, $e+u+v=d/2$.
Writing the ambient pullback as $B_0+F$,
\eqref{eq:positive-square-genus-loss} gives
$B^2-B_0^2=-F^2\leq2u$.
If the later blowup multiplicities are $\mu_j\geq2$, then
\[
 v=\sum_j\frac{\mu_j(\mu_j-1)}2,
 \qquad B_0^2-D^2=\sum_j\mu_j^2\leq4v.
\]
It follows that
\[
 0=D^2\geq2d-2u-4v=4e+2u.
\]
Thus $e=u=0$, $F^2=0$, and every $\mu_j$ equals two.

The surface $W$ is smooth along $B$. Indeed, if a singular
point lay on $B$, its local Cartier equation would have positive
order on every exceptional component over that point. Properness
forces $B_0$ to meet the exceptional fibre, so $B_0\cdot F>0$.
But $(B_0+F)\cdot F=0$ would give $F^2=-B_0\cdot F<0$,
a contradiction.
Since $e=0$, the finite birational map $B\to B_V$ is an
isomorphism. To see this directly, the quotient of the structure
sheaves has finite length; both integral proper curves have global
functions $k$, so equality of arithmetic genera makes that length
zero \cite[Tag~0CE4]{Stacks}.
The normalization $W\to V$ is then an isomorphism near $B_V$.
For a finite normalization extension $A_0\subset A_1$ near the
curve, the isomorphism modulo its equation $t$ gives
$A_1=A_0+tA_1$. Nakayama applied to $A_1/A_0$ proves the
claim. In particular $V$ is smooth along $B_V$.

\smallskip\noindent
\emph{The normal line after the cusps are resolved.}
Let $\alpha$ be the tangential leading coefficient of the
logarithmic form of $M$, and put $Z=\operatorname{div}(\alpha)$.
Then
\begin{equation}\label{eq:boundary-leading-divisor}
 \OO_C(Z)=\omega_C\otimes N^{-2a},\qquad \deg Z=d>0.
\end{equation}
At a point $P$ of $C$, the completed curve equation is
$y^2=A(x,0)$. The multiplicity of $Z$ is the order of
$d(A(x,0))$. In characteristic two that order is even, say
$2\ell$. The even powers of $A(x,0)$ form a square, while
the odd part is $x^{2\ell+1}$ times a square unit. Subtracting
the first square from $y$ and dividing by the unit gives
\[
 w^2=x^{2\ell+1}.
\]
For $\ell=0$ this is smooth. For $\ell\geq1$ its normalization
is $x=s^2$, $w=s^{2\ell+1}$, and its genus defect is $\ell$:
the missing exponents of $k[[s^2,s^{2\ell+1}]]$ are
$1,3,\ldots,2\ell-1$.

Because $V$ is smooth along the curve, at a cusp $x,w$ are
surface parameters. Exactly $\ell$ point blowups resolve the
curve. At each step its multiplicity is two and the exceptional
equation restricts to $s^2$ on the normalization. Thus
$D_{\rm new}=b^*D_{\rm old}-2E$ subtracts $4P'$ from the
normal line at each step, where $P'$ lies above $P$.
The point-blowup normal bundle is as in
\cite[Tag~0AGQ]{Stacks}; the coefficient here also follows
directly by substituting $w=xw_1$ successively.
The total correction is $4\ell P'$, exactly the coefficient
of $q^*Z$, since $q$ has ramification index two. Therefore
\begin{equation}\label{eq:boundary-normal-line}
 \begin{aligned}
 N_{D/S}
 &\simeq q^*N\otimes\OO_D(-q^*Z)\\
 &\simeq q^*\bigl(N\otimes\OO_C(-Z)\bigr)
 \simeq q^*(\mathcal T^{-1}).
 \end{aligned}
\end{equation}
Away from $Z$ this is the natural identification before any
blowup correction.

\smallskip\noindent
\emph{A cusp forces conormal order zero.}
The divisor $Z$ has positive degree, so choose one of the cusps above.
The form $\eta=dz/z$ is regular on $V$ before the cusp resolution.
At a cusp with complete equation $f=w^2-x^{2\ell+1}$, write
\[
 \eta=f^h\bigl(A(x,w)\,dx+B(x,w)\,dw\bigr).
\]
The coefficients $A,B$ are regular power series: divisibility by
the power of the prime Cartier equation can be checked at its
generic point in the regular local ring. On the branch
$x=s^2$, $w=s^{2\ell+1}$, the purely conormal leading condition
gives $s^{2\ell}B(s^2,s^{2\ell+1})=0$. Thus $B=fB_1$.
The composite chart of the $\ell$ cusp blowups is
\[
 w=x^\ell v,\qquad H=x-v^2,\qquad
 f=-x^{2\ell}H,\qquad dx=dH.
\]
After division by $H^h$ and restriction to the strict curve,
the term involving $fB_1$ vanishes. The conormal coefficient is
\[
 v^{4\ell h}A(v^2,v^{2\ell+1}).
\]
It would vanish at $v=0$ if $h>0$, contradicting the nowhere-zero
section $\sigma$. Hence $h=0$. Completion preserves the branch
normalization and finite-length defect
\cite[Tags~0C1R and~0C3V]{Stacks}; faithful flatness detects this zero
on the original curve.

\smallskip\noindent
\emph{For $a\geq3$ the discrepancy line is trivial.}
Suppose $a\geq3$. Then $h=0<a-1$. On the original surface put
$\theta_i=d\log b_i$. The overlap relation and $dg_{ij}=0$
give
\[
 \theta_i-\theta_j\in I_C^{a-1}\Omega_X^1.
\]
Write $\omega=d\log M$. Since
$\omega\in I_C^{2a}\Omega_X^1$, the two-forms
$\theta_i\wedge\omega$ agree modulo
$I_C^{3a-1}\Omega_X^2$.
At the generic point, \eqref{eq:boundary-square-approximation}
gives
\[
 \theta\wedge\omega
   =\frac{t^{2a}R}{bf}\,dt\wedge dA.
\]
Its exact order is $2a+h<3a-1$. Its leading coefficient
therefore glues to a nonzero global section of
\[
 \omega_X|_C\otimes N^{-(2a+h)}
 \simeq\omega_C\otimes N^{-(2a+h+1)}
 =\mathcal T\otimes N^{-h}.
\]
This line has degree $-hd$. Hence $h=0$, and its nonzero
section trivializes $\mathcal T$. Since $h<a-1$ and $a$ is
odd, this case requires $a\geq3$.

\smallskip\noindent
\emph{For $a=1$ the discrepancy line has exact order two.}
Suppose $a=1$, so $h=0=a-1$. Equations
\eqref{eq:boundary-conormal-section} and
\eqref{eq:boundary-normal-line} show that $\sigma$ trivializes
$q^*(\mathcal T^a)$. The norm along the finite flat
degree-two curve map gives $\mathcal T^{2a}\simeq\OO_C$.
We must also show that $\mathcal T^a$ is nontrivial.

In the conormal frame arising from $t$, the scalar of $\sigma$
at the generic point is
\[
 \gamma=\frac{c+y_0}{b_0},\qquad
 c\in k(C),\quad b_0\in k(C)^*,\quad y_0\notin k(C).
\]
Here $c$ is the residue of $t^{1-a}R$, or zero if that
element has positive valuation. Thus $\gamma\notin k(C)$.
By the natural compatibility of
\eqref{eq:boundary-normal-line} away from $Z$, replacing the
frame by one pulled back from a rational frame of $\mathcal T^a$
multiplies $\gamma$ only by an element of $k(C)^*$.
So $\sigma$ does not descend as a rational section on $C$.
If $\mathcal T^a$ were trivial, choose a nowhere-zero section
$\tau$ on $C$. The ratio $\sigma/q^*\tau$ would be a global
regular function on the integral projective curve $D$, hence a
constant in $k$. This contradicts non-descent. Consequently
$\mathcal T^a$ has exact order two.
Since $a=1$, the line $\mathcal T$ itself has exact order two.

\smallskip\noindent
\emph{Two-torsion on the curve in the remaining case.}
When $\mathcal T^a$ is nontrivial it already supplies the required
two-torsion. Suppose instead that $h=0$ and $\mathcal T$ is
trivial. The functions $b_i|_C$ trivialize $(L|_C)^{\otimes2}$;
therefore $\beta_i=d\log(b_i|_C)$ glue to its logarithmic form
$\beta$. The zero divisor $Z$ is nonempty by
\eqref{eq:boundary-leading-divisor}. At a point of $Z$, choose
coordinates $x,t$. There $A_x=0$. Smoothness of
$y^2=A(x,t)$ forces $A_t\ne0$. The nowhere-zero leading
coefficient of $\theta\wedge\omega$ is, up to a unit,
\[
 b_xA_t-b_tA_x.
\]
At this point it equals $b_xA_t$, so $b_x\ne0$. Thus
$\beta$ is nonzero, and $L|_C$ is a nontrivial line bundle
of exact order two. This proves $\Pic(C)[2]\ne0$ in both
cases. Combining with Theorem~\ref{thm:positive-excess-picard}
gives $\delta>d$ whenever $\Pic(C)[2]=0$.
The stated integer bounds follow because $\delta$ is even.
\end{proof}

The smallest case is $m=2$: equality would require the specific
line $\omega_C\otimes\OO_C(-3C)$ to have exact order two.
For $m=6$, the discrepancy line must be trivial. Neither calculation constructs a surface with these
properties. The rank-two refinement at $m=2$ is proved in
Proposition~\ref{prop:smallest-endpoint-rank}; its strictness consequence
completes Corollary~\ref{cor:positive-excess-characteristic-two-boundary}.
The actual conormal trivialization above, together with the
positive-characteristic moving lemma, yields the fibration in
Proposition~\ref{prop:boundary-resolved-fibration}.

\subsection{Two-torsion and fibrations at the characteristic-two boundary}
\label{subsec:smallest-endpoint-geometry}

At logarithmic order two, the boundary forces two independent actual
two-torsion classes on the distinguished curve.  The same local geometry
also gives a fibration on the resolved covering surface at every odd
half-order.  The latter conclusion uses the triviality of its actual
normal line and a positive-characteristic moving-divisor criterion.

\begin{proposition}[Two-torsion at logarithmic order two]
\label{prop:smallest-endpoint-rank}
Let $k$ be algebraically closed of characteristic two, let $X$ be a
normal integral projective surface, and let $C\subset X$ be a smooth
integral effective Cartier divisor in the smooth locus, with $d=C^2>0$.
Let $M\ne\OO_X$ have order two and satisfy $M|_C\simeq\OO_C$ and
$[M]\in2\operatorname{NS}(X)$.  Suppose its logarithmic form $\omega$
has exact order two along $C$ and
\[
 2g(C)-2=3d.
\]
Put $N=\OO_C(C)$ and $\mathcal T=\omega_C\otimes N^{-3}$, and let
$Z$ be the zero divisor of the tangential leading coefficient
\[
 \alpha\in H^0(C,\omega_C\otimes N^{-2})
\]
of $\omega$.  Then $\deg Z=d$, every multiplicity of $Z$ is even,
and the logarithmic form of $\mathcal T$ satisfies
\begin{equation}\label{eq:smallest-endpoint-log-divisor}
 \operatorname{div}(\omega_{\mathcal T})=3Z.
\end{equation}
For every square root $L^{\otimes2}\simeq M$, the line bundle
$U=L|_C$ has exact order two and differs from $\mathcal T$.
Its logarithmic form $\beta_U$ is nonzero at every point of
$\operatorname{Supp}Z$.  In particular,
\[
 \dim_{\F_2}\Pic(C)[2](k)\geq2,
\]
where the group consists of actual line-bundle classes.
\end{proposition}

\begin{proof}
The lifting argument in Theorem~\ref{thm:positive-excess-picard}
supplies a square root $L$, and every such root has exact order four.
Resolve $X$ away from $C$ and retain the notation for the pullbacks.
The curve, its normal line and the logarithmic order are unchanged.

We use the geometric construction in the proof of
Corollary~\ref{cor:positive-excess-characteristic-two-boundary}.
Let $\pi:W\to X$ be the normalized $M$-torsor, let $B=\pi^*C$,
and let $D\subset S$ be its smooth strict transform on the suitable
resolution constructed there.  Write $q:D\to C$ for the induced
degree-two purely inseparable map.  Near $B$, the surface $W$ is smooth
and has the local description
\begin{equation}\label{eq:smallest-endpoint-local-model}
 \begin{aligned}
 u&=b^2+t^2A,& z&=b+ty,& y^2&=A,\\
 \omega&=t^2\frac{dA}{u},& \eta&=\frac{dz}{z}.
 \end{aligned}
\end{equation}
Here $t=0$ defines $C$, and $b,z$ are units.  The form $\eta$ is
regular on $W$ and pulls back regularly to $S$.  The calculations
\eqref{eq:boundary-conormal-order}--\eqref{eq:boundary-conormal-section}
at $a=1$ and the normal-line identification
\eqref{eq:boundary-normal-line} give
\begin{equation}\label{eq:smallest-endpoint-conormal}
 D^2=0,\qquad N_{D/S}^{-1}\simeq q^*\mathcal T,
 \qquad
 0\ne\sigma\in H^0(D,N_{D/S}^{-1}),
\end{equation}
where $\sigma$ is the conormal component of the restriction of $\eta$
as a section of $\Omega_S^1|_D$, and is nowhere zero.
These calculations apply to each choice of $L$.

The section $\alpha$ is nonzero by
Lemma~\ref{lem:positive-square-logarithmic}.  Locally its differential
representative is $d(A|_C)/b_0^2$, where $b_0=b|_C$.
Equation~\eqref{eq:boundary-leading-divisor} gives $\deg Z=d$.
The local characteristic-two derivative calculation in that proof
shows that its multiplicities are even.  Above a point $P$ of
multiplicity $2\ell>0$, the curve $B$ has complete equation
$w^2=x^{2\ell+1}$.  In particular $Z$ is nonempty and each such point
corresponds to a cusp of $B$.

\smallskip\noindent
\emph{The logarithmic form of the restricted root.}
Let $\widetilde P\in B$ lie above $P\in\operatorname{Supp}Z$.
Since $B$ is singular at $\widetilde P$ in the smooth surface $W$,
its local equation $t$ has zero cotangent vector there.
The cotangent vector $\eta(\widetilde P)$ is nonzero: otherwise its
pullback to the point of $D$ above it would vanish, contrary to
\eqref{eq:smallest-endpoint-conormal} and the injective smooth
conormal sequence on $D$.
Choose local parameters $x,t$ on $X$.  At $\widetilde P$,
\[
 dz=db+y\,dt+t\,dy=db=b_x\,dx.
\]
It follows that $b_x(P)\ne0$.  Also $A_x(P)=0$, and smoothness of
$y^2=A(x,t)$ implies $A_t(P)\ne0$.

If local frames of $L$ differ by $g_{ij}$, then
$b_i|_C=g_{ij}^2b_j|_C$.  These units give the square trivialization
of $U=L|_C$, and its logarithmic form is
\[
 \beta_U=d\log(b_i|_C).
\]
The preceding calculation proves that $\beta_U(P)\ne0$ at every
point of $\operatorname{Supp}Z$.  Thus $U$ is nontrivial, since the
logarithmic form of a trivial line bundle on a proper integral curve
is zero.  Its square is trivial, so it has exact order two.

\smallskip\noindent
\emph{The logarithmic form of the discrepancy line.}
At the generic point of $C$, write
\[
 db=R\,dt+S_1\,dA,\qquad
 R=b_t-b_x\frac{A_t}{A_x}.
\]
The relation $dA=0$ on $W$ gives
\begin{equation}\label{eq:smallest-endpoint-root-form}
 \eta=\frac{(R+y)\,dt+t\,dy}{b+ty}.
\end{equation}
Set $A_0=A|_C$ and let $R_0$ denote the rational restriction of $R$
to $C$.  If $P$ has multiplicity $2\ell$ in $Z$, the functions
$b_x|_C$ and $A_t|_C$ are units at $P$, whereas $A_x|_C$ has order
$2\ell$.  The numerator $b_tA_x-b_xA_t$ is therefore a unit, so
\[
 \operatorname{ord}_P R_0=-2\ell.
\]

We use the conormal-map identification in
\eqref{eq:boundary-normal-line}.  If $e$ is the conormal frame on
$C$ defined by $t$, the map
$q^*N^{-1}\to N_{D/S}^{-1}$ has zero divisor $q^*Z$.
Consequently the frame corresponding to $e/x^{2\ell}$ in
$N^{-1}(Z)$ is the image of $q^*e/q^*(x^{2\ell})$.
The section $\alpha$ identifies $N^{-1}(Z)$ with $\mathcal T$.
Relative to the pullback of this frame,
\eqref{eq:smallest-endpoint-root-form} gives the unit coefficient
\[
 F=q^*(x^{2\ell})\frac{R_0+y_0}{b_0},
 \qquad y_0^2=q^*A_0.
\]
Any unit introduced by the blowup chart occurs both in the conormal
map and in the chosen frame, and cancels in this coefficient.
Its square is the pullback of
\begin{equation}\label{eq:smallest-endpoint-discrepancy-unit}
 V_0=x^{4\ell}\frac{R_0^2+A_0}{b_0^2}.
\end{equation}
More generally, in any local frame of $\mathcal T$, the square of
the coefficient of $\sigma$ belongs to $q^*k(C)$, because
$k(D)^2=k(C)$.  It descends to a regular unit: its pullback is a
unit, and $q$ is finite and surjective between smooth curves.
The transition rule is that of a section of $\mathcal T^{\otimes2}$.
Thus $\sigma^2$ descends to a nowhere-zero section of this line.
The resulting square trivialization has local value $V_0^{-1}$;
in characteristic two, $d\log(V_0^{-1})=d\log V_0$.

Differentiating \eqref{eq:smallest-endpoint-discrepancy-unit} in
characteristic two gives
\begin{equation}\label{eq:smallest-endpoint-discrepancy-log}
 \omega_{\mathcal T}=d\log V_0
      =\frac{dA_0}{R_0^2+A_0}.
\end{equation}
The numerator has order $2\ell$ at $P$, and the denominator has
pole order $4\ell$, since $A_0$ is regular.  Hence
$\operatorname{ord}_P\omega_{\mathcal T}=6\ell$.
This differential is nonzero because $dA_0\ne0$ generically, and it
is regular everywhere by its construction from the torsion line.
The zeros just calculated have total degree
$3\deg Z=3d=2g(C)-2$.  They therefore exhaust its zero divisor,
proving \eqref{eq:smallest-endpoint-log-divisor}.

At every point of the nonempty divisor $Z$, the form
$\omega_{\mathcal T}$ vanishes and $\beta_U$ does not.
Any two square trivializations differ by a nonzero constant and
therefore define the same logarithmic form.  Thus $U\ne\mathcal T$.
The line $\mathcal T$ has exact order two by
Corollary~\ref{cor:positive-excess-characteristic-two-boundary}.
The two distinct nonzero classes consequently generate a subgroup
of order four, as asserted.
\end{proof}

\begin{corollary}\label{cor:smallest-endpoint-rank-strictness}
Under the surface, curve and torsion hypotheses of
Proposition~\ref{prop:smallest-endpoint-rank}, retain logarithmic
order two and $[M]\in2\operatorname{NS}(X)$, but do not assume
$2g(C)-2=3d$.  If $\dim_{\F_2}\Pic(C)[2](k)\leq1$, then
\[
 2g(C)-2-2d>d.
\]
\end{corollary}

\begin{proof}
Theorem~\ref{thm:positive-excess-picard} gives
$2g(C)-2-2d\geq d$.  Equality would contradict
Proposition~\ref{prop:smallest-endpoint-rank}.
\end{proof}

The moving-divisor input is the trivial-normal case of Totaro's
criterion \cite[Lemma~4.1, pp.~7--8]{TotaroMovingDivisors}.
The following proof also gives a multiple that is a power of the
characteristic.

\begin{lemma}[Trivial normal bundle in positive characteristic]
\label{lem:trivial-normal-moving-divisor}
Let $Y$ be a normal integral projective variety over an algebraically
closed field of characteristic $p>0$, and let $A$ be an effective
Cartier divisor with $\OO_A(A)\simeq\OO_A$.  There is an integer
$e\geq0$ such that $\OO_Y(p^eA)$ is globally generated.
\end{lemma}

\begin{proof}
Choose the indicated trivialization.  The exact sequence
\[
 0\longrightarrow\OO_Y\longrightarrow\OO_Y(A)
   \longrightarrow\OO_A\longrightarrow0
\]
gives an obstruction $c\in H^1(Y,\OO_Y)$ to lifting
$1\in H^0(A,\OO_A)$.  If $c=0$, use the identity cover below and
put $e=0$.  Otherwise, finite-dimensionality of $H^1(Y,\OO_Y)$
gives a relation among the iterates of its Frobenius-semilinear map,
\[
 \sum_{i=0}^e a_iF^ic=0,\qquad a_e=1.
\]
The highest nonzero coefficient can be normalized to one.
Thus $P(T)=\sum_{i=0}^e a_iT^{p^i}$ is a monic additive polynomial
of degree $p^e$ with $P(c)=0$.

When $c\ne0$, represent it by an additive Cech cocycle $c_{ij}$ on a
finite affine cover with affine intersections.  There are regular
functions $b_i$ such that $P(c_{ij})=b_j-b_i$.  The schemes
\[
 Y'_i=\Spec_{U_i}
       \bigl(\OO_{U_i}[T_i]/(P(T_i)-b_i)\bigr)
\]
glue by $T_i=T_j-c_{ij}$: additivity gives
$P(T_j-c_{ij})=b_j-P(c_{ij})=b_i$, and the cocycle identity gives
compatibility on triple intersections.  The monic equations define
a finite locally free faithfully flat cover
$\nu:Y'\to Y$ of degree $p^e$.  The pulled-back cocycle is a
coboundary, so $\nu^*c=0$.  In either case we have such a finite
locally free cover of degree $p^e$ killing the obstruction.

Flatness preserves the displayed exact sequence.  Hence
$\nu^*\OO_Y(A)$ has a section $s$ restricting to $1$ on $\nu^*A$.
Together with the canonical section $t$ of that effective divisor,
it has no common zero.  For each closed point $y\in Y$, choose
$\lambda\in k$ so that $s+\lambda t$ is nonzero at every point of
the finite fibre $Y'_y$.  There are only finitely many forbidden
scalars, since $s$ and $t$ have no common zero, and $k$ is infinite.
In a local trivialization, the coefficient of this section is a unit
in the Artinian fibre algebra; this remains true for a nonreduced
fibre.

The finite-flat norm, defined locally by the determinant of
multiplication by this coefficient, is therefore nonzero at $y$.
Its transition factor is the $p^e$th power of that of $\OO_Y(A)$,
so it is a global section of $\OO_Y(p^eA)$.
Such a section exists at every closed point.  The base locus is
therefore empty, proving global generation.
\end{proof}

\begin{proposition}[A fibration on the resolved cover]
\label{prop:boundary-resolved-fibration}
Under the hypotheses of
Corollary~\ref{cor:positive-excess-characteristic-two-boundary},
write $C=Y$ and $m=2a$, with $a\geq1$ odd.
Use the normalized torsor and the suitable resolution of its inverse
curve constructed in that proof, and let $D\subset S$ be the smooth
strict transform.  Then $N_{D/S}\simeq\OO_D$, and there are a
morphism with connected fibres
$g:S\to\PP^1$, a point $b\in\PP^1(k)$, and an integer $s\geq0$
such that
\begin{equation}\label{eq:boundary-fibration-fibre}
 g_*\OO_S=\OO_{\PP^1},\qquad g^*b=2^sD.
\end{equation}
This conclusion requires only $C^2>0$, not ampleness of $C$.
\end{proposition}

\begin{proof}
Equations~\eqref{eq:boundary-nonnegative-square}--\eqref{eq:boundary-conormal-section}
and their proof give $D^2=0$ and a nowhere-zero section
of $N_{D/S}^{-(h+1)}$, where $0\leq h\leq a-1$.
The cusp calculation in the same proof applies before its division
into the cases $a=1$ and $a\geq3$.  At a cusp of multiplicity
$2\ell$ in $Z$, the conormal coefficient after resolution is
\[
 v^{4\ell h}A(v^2,v^{2\ell+1}),
\]
with $A$ regular.  Since $\deg Z=d>0$, there is such a cusp;
the nowhere-zero property forces $h=0$.  The section therefore
trivializes $N_{D/S}^{-1}$ itself.

Lemma~\ref{lem:trivial-normal-moving-divisor} makes
$\OO_S(2^eD)$ globally generated for some $e\geq0$.
This line bundle is numerically nontrivial, because the nonzero
effective curve $D$ has positive intersection with an ample divisor
on $S$, and it has square zero.  Its associated morphism has
one-dimensional image: image dimension zero would make the line
trivial, whereas image dimension two would make its square positive.
Take its Stein factorization $g:S\to H$.  The base $H$ is a smooth
integral projective curve, $g_*\OO_S=\OO_H$, and
\[
 \OO_S(2^eD)\simeq g^*A_H
\]
for an ample line bundle $A_H$ on $H$.

By projection formula, the canonical section with divisor $2^eD$
descends to a section of $A_H$.  Its zero divisor is supported at one
point $b$, and the fibre over $b$ is supported on $D$.  Thus
\[
 g^*b=rD,\qquad 2^e=r\deg A_H.
\]
Both factors on the right are positive integers, so $r=2^s$ for some
$s\geq0$.

At a cusp $w^2=x^{2\ell+1}$ the final resolution chart is
\[
 w=x^\ell v,\qquad D:(x-v^2)=0.
\]
The last exceptional curve $R\simeq\PP^1$ has equation $x=0$ and
meets $D$ only at this point, with intersection length
$\operatorname{length}(k[[v]]/(v^2))=2$.
The resolution away from the inverse curve and the blowups at other
cusps leave $R$ unchanged.  The restriction of $g^*\OO_H(b)$ to $R$
has degree $2r>0$, so $R$ dominates $H$, with degree $2r$.
L\"uroth's theorem gives $H\simeq\PP^1$, also in characteristic two.
This proves \eqref{eq:boundary-fibration-fibre}.
\end{proof}

The morphism is defined on the covering resolution $S$.  Its generic
fibre may have positive genus, and its distinguished fibre may be
multiple.  These conclusions give necessary restrictions on the
boundary configuration and do not establish its existence.

\subsection{Curve pullbacks and Albanese restrictions at the boundary}
\label{subsec:boundary-albanese-restrictions}

The first positive-excess boundary admits further restrictions when the
torsion class comes from a curve, or when the surface has independent
global one-forms.  The first argument retains the cotangent bundle of
the base curve in the conormal calculation.  The second uses ampleness
on the original surface.  Neither argument constructs a fibration from
a curve of square zero.

\begin{proposition}[Curve-pullback excess bound]
\label{prop:boundary-curve-pullback-excess}
Let $k$ be algebraically closed of characteristic $p>0$, let $X$ be a
normal integral projective surface, and let $C\subset X$ be a smooth
integral effective Cartier divisor contained in the smooth locus, with
$d=C^2>0$.  Let $M\ne\OO_X$ have order $p$ and be trivial on $C$.
Suppose its logarithmic form $\omega$ has exact order $pa$ along $C$,
where $a\geq1$ and $p\nmid a$, and put
\[
 \delta=2g(C)-2-pad.
\]
Suppose that $f:X\to B$ is a nonconstant morphism to a smooth integral
projective curve of genus $b\geq1$, and that
$M\simeq f^*P$ for a nontrivial order-$p$ line bundle $P$ on $B$.
Then $C$ dominates $B$. If $n=\deg(f|_C)$, then $p\mid n$, and there is an integer
$0\leq j\leq a-1$ such that
\begin{equation}\label{eq:boundary-curve-pullback-bound}
 \delta\geq\frac{pd}{2(p-1)}+
       \frac{n(b-1)}{(p-1)(j+1)}
 \geq\frac{pd}{2(p-1)}+\frac{n(b-1)}{(p-1)a}.
\end{equation}
In particular, $b\geq2$ excludes the endpoint
$\delta=pd/(2(p-1))$. In characteristic two this reads
$\delta\geq d+n(b-1)/a$, and at $a=1$ it reads
$\delta\geq d+n(b-1)$.  This proposition does not require ampleness of
$C$ or a separate N\'eron--Severi divisibility hypothesis.
\end{proposition}

\begin{proof}
If $f(C)=b_0$ is a point, write the effective Cartier fibre as
$f^*b_0=rC+E$, where $r>0$ and $E$ is effective without $C$ as a
component.  Since $f|_C$ is constant,
\[
 0=(f^*b_0)\cdot C=rd+E\cdot C>0,
\]
a contradiction.  Here intersections with $E$ are nonnegative because
$C$ is contained in the smooth locus.  Thus $f|_C$ is a finite
nonconstant morphism.  Resolve $X$ away from $C$; this preserves $f$,
$M$, and all the curve data.  We may therefore assume $X$ smooth.
The finite map $f|_C$ is flat because both curves are smooth.
Taking the norm of $(f|_C)^*P\simeq\OO_C$ gives $P^n\simeq\OO_B$.
Since $P$ has exact order $p$, it follows that $p\mid n$.

Multiplication by $p$ on the Jacobian of $B$ is surjective on
$k$-points.  Choose $Q\in\operatorname{Pic}^0(B)$ with $Q^p\simeq P$,
and put $L=f^*Q$.  Then $Q$ and $L$ have exact order $p^2$ and
$L^p\simeq M$.  Use a $p$th-power trivialization of $P$ and its pullback
for $M$.  Changing a $p$th-power trivialization on a proper integral
variety only multiplies it by a constant, so this choice does not
change $\omega$ or its order.

Let $r:B'\to B$ be the normalization of the $P$-torsor.  This torsor
is integral: a $p$th root of its defining unit in $k(B)$ would, by
normality, give local unit roots trivializing $P$.  Consequently $r$
is finite and purely inseparable of degree $p$.  The normal projective
curve $B'$ is smooth and is a Frobenius twist of $B$, so $g(B')=b$.
The line bundle $Q'=r^*Q$ has $p$th power trivial.  It is nontrivial, since
otherwise the norm formula for the finite flat morphism $r$ would give
\[
 \OO_B\simeq\operatorname{Nm}_r(r^*Q)\simeq Q^p\simeq P.
\]
Thus its logarithmic form $\theta$ is a nonzero regular one-form on
$B'$, by the injectivity of the logarithmic-form map on a smooth
proper curve recalled in
Lemma~\ref{lem:positive-square-resolved-transfer}.

Let $W\to X$ be the normalization of the $M$-torsor.  Before
normalization this torsor is the base change of the $P$-torsor.
Its map to that torsor factors through $B'$ by normality: pullbacks
of functions integral over the target order are integral, hence
regular, on $W$.  Take the suitable resolution $S\to W$ in
Lemma~\ref{lem:positive-square-resolved-transfer}, and let $D$ be
the smooth strict transform of the inverse curve.  We obtain a
morphism $\varphi:S\to B'$ compatible with $f$ and $r$.
The root-genus and resolution estimates give
\begin{equation}\label{eq:boundary-curve-pullback-square}
 D^2\geq pd-2(p-1)\delta.
\end{equation}
Indeed, the inverse Cartier curve on $W$ has square $pd$ and
normalization genus $g(C)$; its genus defect $\Delta$ satisfies
$2\Delta\leq(p-1)\delta$, while the suitable resolution loses at most
$4\Delta$ in self-intersection.  These are the assertions of
Lemmas~\ref{lem:positive-square-root-genus} and
\ref{lem:positive-square-resolution-loss}.

The map $q:D\to C$ is purely inseparable of degree $p$.  Since $D$
dominates $B'$ and $r\varphi|_D=f|_Cq$, comparison of degrees gives
\[
 \deg(\varphi|_D)=n.
\]
The pullback $L'$ of $L$ to $S$ is $\varphi^*Q'$, and its logarithmic
form is the regular form $\eta=\varphi^*\theta$.

We need the conormal-order calculation without assuming
$D^2>0$ or $\delta=d$.  At the generic point of $C$, the local
root calculation in Lemma~\ref{lem:positive-square-root-genus}
writes
\[
 u=b_0^p+t^{pa}A,\qquad z=b_0+t^ay,\qquad y^p=A,
\]
where $t$ defines $C$, the functions $b_0,z$ are units, and
$d(A|_C)\ne0$.  On the normalized cover $t$ is a uniformizer of
$D$ and $dt,dy$ form a cotangent basis at its generic point.
Write $db_0=R\,dt+S_1\,dA$ on $X$.  Since $p\nmid a$,
\begin{equation}\label{eq:boundary-curve-pullback-differential}
 \eta=\frac{(R+a t^{a-1}y)\,dt+t^a\,dy}{z}.
\end{equation}
The residue $y_0$ belongs to $k(D)=k(C)^{1/p}$ but not to $k(C)$.
Every normalized leading residue of $R$ belongs to $k(C)$.
Thus $R$ and $a t^{a-1}y$ cannot cancel when their orders coincide,
and
\[
 0\leq j:=\operatorname{ord}_D\eta
     =\min\{v_C(R),a-1\}\leq a-1.
\]
The valuation of zero is infinity.  In particular $\eta\ne0$, and
its leading coefficient is purely conormal because $j<a$.

Consider now the entire differential map
\[
 d\varphi:\varphi^*\omega_{B'}\longrightarrow\Omega_S^1.
\]
At the generic point of $D$, the form $\theta$ is a unit multiple
of a generator of $\omega_{B'}$, since $D$ dominates $B'$.
Therefore this map has exact order $j$ along $D$.
Division by $D^j$ gives a regular map
$\varphi^*\omega_{B'}\otimes\OO_S(jD)\to\Omega_S^1$
near the entire curve $D$: divisibility by a power of the smooth
prime divisor is checked at its generic point in the regular local
rings.  Its restriction has zero tangential image generically,
hence everywhere on the integral curve $D$.  The conormal sequence
gives a nonzero map
\begin{equation}\label{eq:boundary-curve-pullback-conormal}
 (\varphi|_D)^*\omega_{B'}\otimes N_{D/S}^{j}
       \longrightarrow N_{D/S}^{-1}.
\end{equation}
Taking degrees yields
\[
 (j+1)D^2+n(2b-2)\leq0.
\]
Together with \eqref{eq:boundary-curve-pullback-square}, this gives
the first inequality in \eqref{eq:boundary-curve-pullback-bound}.
The second follows from $j+1\leq a$ and $b\geq1$.
\end{proof}

For example, in characteristic two with $a=1$ and $b=2$,
the degree $n$ is a positive even integer. The proposition gives
$\delta\geq d+n\geq d+2$. This is a necessary bound and does
not construct a surface realizing it.

\begin{corollary}[Albanese image a curve]
\label{cor:boundary-albanese-curve}
In characteristic two, under the hypotheses of
Proposition~\ref{prop:boundary-curve-pullback-excess} concerning
$k,X,C,M,a$, suppose instead that $X$ is smooth, $M$ is algebraically
trivial, and the Albanese image of $X$ is one-dimensional.  Let
$f:X\to B$ be the Stein factorization of the Albanese map, and put
$b=g(B)$ and $n=\deg(f|_C)$, retaining
$d=C^2$ and $\delta=2g(C)-2-2ad$.  Then $b\geq1$ and
\[
 \delta\geq d+\frac{n(b-1)}a.
\]
If $\delta=d$, then $b=1$ and $\operatorname{Alb}(X)$ is an elliptic
curve with a nonzero two-torsion point over $k$.
\end{corollary}

\begin{proof}
Write the pointed Albanese map as $\alpha=h\circ f:X\to B\to A$.
The normal projective curve $B$ is smooth.  Its map $h$ induces
$u:\operatorname{Jac}(B)\to A$, while the composite
$X\to B\to\operatorname{Jac}(B)$ induces
$v:A\to\operatorname{Jac}(B)$ by the Albanese universal property.
The identities on $X$ show $uv=1_A$.  Surjectivity of $f$, followed
by the Jacobian universal property, gives $vu=1_{\operatorname{Jac}(B)}$.
Hence $A\simeq\operatorname{Jac}(B)$ and $b\geq1$.

Picard--Albanese duality identifies $f^*$ with an isomorphism
\[
 \operatorname{Pic}^0(B)\simeq
       (\operatorname{Pic}^0_{X/k})_{\mathrm{red}}
\]
\cite[Expos\'e~236, Corollary~3.2 and Theorem~3.3(iii),
pp.~236--238]{FGAPicard}.
Every $k$-point of $\operatorname{Pic}^0_{X/k}$ lies in its reduction.
Thus $M=f^*P$ for a unique nontrivial order-two line bundle $P$ on
$B$.  Apply Proposition~\ref{prop:boundary-curve-pullback-excess}.
\end{proof}

\begin{proposition}[Ample wedge restriction]
\label{prop:boundary-ample-wedge}
Let $X$ be a smooth integral projective surface over an algebraically
closed field of characteristic two, let $C\subset X$ be a smooth
integral ample divisor, and put $d=C^2$.  Let $M\ne\OO_X$ have
order two, be trivial on $C$, and have logarithmic form $\omega$ of
exact order $2a$ along $C$, where $a\geq1$ is odd.  Suppose
\[
 2g(C)-2=(2a+1)d,
 \qquad
 \mathcal T=\omega_C\otimes\OO_C(C)^{-(2a+1)}.
\]
If a global regular one-form $\beta$ satisfies $\omega\wedge\beta\ne0$,
then
\begin{equation}\label{eq:boundary-ample-wedge-splitting}
 K_X\sim2aC,\qquad
 \Omega_X^1\simeq\OO_X(2aC)\oplus\OO_X,\qquad
 c_2(X)=0,\qquad\mathcal T\simeq\OO_C.
\end{equation}
Consequently, if $\mathcal T$ is nontrivial, all global regular
one-forms are rationally collinear with $\omega$, their pairwise
wedge products vanish, and the Albanese map has generic differential
rank at most one.
\end{proposition}

\begin{proof}
Adjunction gives $K_X\cdot C=2ad$.  If
$\omega\wedge\beta\ne0$, division by the $2a$th power of the
canonical section of $\OO_X(C)$ gives a nonzero section
\[
 \tau\in H^0(X,\omega_X(-2aC)).
\]
Its effective zero divisor has intersection zero with $C$.
Ampleness forces this divisor to vanish.  A nonzero section of a
line bundle on a smooth surface cannot have an isolated zero
without a divisorial zero, since locally a nonunit nonzero function
has a height-one zero.  Thus $\tau$ is nowhere zero and
$K_X\sim2aC$.

The divided form $\omega/s_C^{2a}$ and the form $\beta$ define
\[
 \OO_X(2aC)\oplus\OO_X\longrightarrow\Omega_X^1
\]
with determinant $\tau$.  The determinant is invertible everywhere,
so this map is an isomorphism and $c_2(X)=0$.  Restricting
$K_X\sim2aC$ to $C$ and applying adjunction gives
$\mathcal T\simeq\OO_C$.

If $\mathcal T$ is nontrivial, no such $\beta$ exists.  The nonzero
form $\omega$ spans all global one-forms over $k(X)$, so all their
pairwise wedges vanish.  Invariant global one-forms generate the
cotangent bundle of an abelian variety.  Their pullbacks to $X$
therefore show that the Albanese differential has generic rank at
most one.
\end{proof}

\begin{corollary}[The smallest algebraically trivial endpoint]
\label{cor:boundary-smallest-albanese-gate}
Let $X$ be a smooth integral projective surface over an algebraically
closed field of characteristic two, and let $C\subset X$ be a smooth
integral ample divisor.  Let $M\ne\OO_X$ have order two and be
trivial on $C$.  Suppose its logarithmic form has exact order two,
\[
 2g(C)-2=3C^2,
 \qquad [M]\in2\operatorname{NS}(X).
\]
Then all global regular one-forms are rationally collinear, and the
Albanese map has generic differential rank at most one.  In particular,
it is not generically separable onto a two-dimensional image.

If $M$ is algebraically trivial, its Albanese geometry must be one
of the following:
\begin{enumerate}
 \item the Albanese image is a curve and $\operatorname{Alb}(X)$
       is elliptic;
 \item the Albanese image is a surface and the Albanese map has
       generic differential rank at most one.
\end{enumerate}
These conclusions apply, in particular, when $C$ is very ample.
\end{corollary}

\begin{proof}
Corollary~\ref{cor:positive-excess-characteristic-two-boundary}
makes $\mathcal T=\omega_C\otimes\OO_C(C)^{-3}$ a line bundle of
exact order two.  Apply Proposition~\ref{prop:boundary-ample-wedge}.
A morphism generically separable onto a surface image has generic
differential rank two, so it is excluded.

If $M$ is algebraically trivial, the N\'eron--Severi divisibility
condition is automatic.  The Albanese variety is nonzero, because
the $k$-points of a zero-dimensional reduced neutral Picard scheme
are trivial.  A curve image is handled by
Corollary~\ref{cor:boundary-albanese-curve}; the only other possible
image dimension is two.
\end{proof}

\begin{remark}[Scope and comparison]
\label{rem:boundary-albanese-scope}
The curve-pullback bound uses an existing morphism $X\to B$.
The curve-pullback bound does not require a fibration arising
from the resolved curve of square zero. Proposition~\ref{prop:boundary-resolved-fibration}
does produce a separate fibration on a suitable covering resolution,
using actual triviality of its normal line and the moving lemma.
The elliptic-base case and the inseparable surface-image case in
Corollary~\ref{cor:boundary-smallest-albanese-gate} remain unresolved
by these arguments.  Existence of either configuration remains open.
The nonclosed one-forms from Proposition~\ref{prop:residual-adjoint}
are compatible with rational collinearity: $d(r\omega)$ can be
nonzero although all pairwise wedges vanish.

For comparison, Cotner--Zavyalov prove a Lefschetz theorem for finite
flat commutative coefficients under a Hodge-equivalence or strongly
ample hypothesis \cite[Theorem~3.3.5, pp.~1861--1862]{CotnerZavyalov}.
Applied to $\mu_2$ and a strongly ample curve on a surface, it gives
injectivity of restriction on order-two line bundles.  Strong
ampleness is an additional hypothesis; their Example~3.3.6 concerns
failure for $\alpha_p$ and does not provide an order-two line-bundle
example.  Our bound instead retains $\deg\omega_{B'}$ in
\eqref{eq:boundary-curve-pullback-conormal}.

The remaining inseparable Albanese case belongs to an established
foliation setting: Conti describes the canonical inseparable
degree-two factorization for a characteristic-two surface whose
Albanese image is a surface and whose Albanese map is inseparable
\cite[Lemma~2.22 and Remark~2.23, p.~13]{ContiPositiveSeveriComparison}.
That description does not itself exclude the endpoint here.
\end{remark}

\subsection{Global structure at the smallest endpoint}
\label{subsec:smallest-endpoint-global-structure}

At logarithmic order two, ampleness makes the local normalization
calculation global.  The resulting cotangent line and foliation do
not require an Albanese hypothesis.

\begin{theorem}[Global structure at the smallest endpoint]
\label{thm:smallest-endpoint-global-structure}
Let $k$ be algebraically closed of characteristic two, let $X$ be a
smooth integral projective surface, and let $C\subset X$ be a smooth
integral ample divisor.  Put $d=C^2$ and $N=\OO_C(C)$.  Let
$M\ne\OO_X$ have exact order two and satisfy $M|_C\simeq\OO_C$.
Suppose that $M$ has a square root $L$, that its logarithmic form
$\omega$ has exact order two along $C$, and that
\[
 2g(C)-2=3d.
\]
The square-root condition is equivalent to
$[M]\in2\operatorname{NS}(X)$.
Let $\tau:W\to X$ be the $\mu_2$-torsor of $M$,
let $n:V\to W$ be its normalization, put $\pi=\tau n$ and
$B=\pi^*C$, and write $\mathfrak c\subset\OO_V$ for the conductor.
Then the following assertions hold.
\begin{enumerate}
\item The integer $d$ is even and
      $\mathcal T=\omega_C\otimes N^{-3}$ has exact order two.
      The surface $V$ is smooth along $B$, and
      \begin{equation}\label{eq:smallest-endpoint-global-conductor}
       \mathfrak c=\OO_V(-B),\qquad
       \omega_V\simeq\pi^*\bigl(\omega_X(-C)\bigr).
      \end{equation}
      In particular, $V$ is Gorenstein.  The morphism $\pi$ is finite
      flat of degree two, and the sequence
      \begin{equation}\label{eq:smallest-endpoint-global-cover-sequence}
       0\longrightarrow\OO_X\longrightarrow\pi_*\OO_V
       \longrightarrow M^{-1}(C)\longrightarrow0
      \end{equation}
      does not split as a sequence of $\OO_X$-modules.  Moreover,
      \[
       \chi(\OO_V)=2\chi(\OO_X)-d/2.
      \]
\item If $s_C$ is the canonical section of $\OO_X(C)$, then
      $\theta=\omega/s_C^2$ is a global section of
      $\Omega_X^1(-2C)$ whose zero scheme $Z_0$ is finite and disjoint
      from $C$.  It gives a saturated exact sequence
      \begin{equation}\label{eq:smallest-endpoint-global-cotangent}
       0\longrightarrow\OO_X(2C)
       \xrightarrow{\ \theta\ }\Omega_X^1
       \longrightarrow I_{Z_0}\otimes\omega_X(-2C)
       \longrightarrow0,
      \end{equation}
      and
      \begin{equation}\label{eq:smallest-endpoint-global-chern}
       \operatorname{length}(Z_0)=c_2(X)\ge0.
      \end{equation}
\item All global one-forms are rationally collinear with $\omega$;
      more precisely,
      \begin{equation}\label{eq:smallest-endpoint-global-forms}
       H^0(X,\Omega_X^1)
       =\theta\,H^0\bigl(X,\OO_X(2C)\bigr).
      \end{equation}
      Their common annihilator is the saturated, $2$-closed
      rank-one foliation
      \begin{equation}\label{eq:smallest-endpoint-global-foliation}
       \mathcal F
       :=\ker\bigl(\omega:T_X\longrightarrow\OO_X\bigr)
       \simeq\omega_X^{-1}(2C).
      \end{equation}
      It satisfies
      \[
       \mathcal F\cdot C=0,\qquad
       \mathcal F|_C\simeq\mathcal T^{-1}
       =\mathcal T\ne\OO_C,\qquad H^0(X,T_X)=0.
      \]
\item Every smooth integral rational curve $E\subset X$ satisfies
      \begin{equation}\label{eq:smallest-endpoint-global-rational-curve}
       E^2\le-2C\cdot E.
      \end{equation}
      In particular, $X$ has no $(-1)$-curve.
\end{enumerate}
\end{theorem}

\begin{proof}
The lifting argument in the proof of
Theorem~\ref{thm:positive-excess-picard} shows that
$[M]\in2\operatorname{NS}(X)$ is equivalent to the existence of $L$:
the discrepancy from a chosen divisor-class root lies in
$\Pic^0(X)$, whose $k$-points are divisible by two.
The line $L$ has exact order four.  Adjunction gives
$K_X\cdot C=2d$, and the displayed genus equality makes $d$ even.
Corollary~\ref{cor:positive-excess-characteristic-two-boundary}
gives the assertion about $\mathcal T$.
The torsor $W$ is integral: a square root of its defining unit in
$k(X)$ would extend as local unit roots by normality and would
trivialize $M$.

\smallskip\noindent
\emph{The local endpoint data.}
We recall the equality consequences needed below from the proof of
that corollary.  In frames of $M$ obtained by squaring frames of $L$,
the finite model near $C$ is
\begin{equation}\label{eq:smallest-endpoint-global-local-model}
 u_i=b_i^2+t_i^2A_i,\qquad
 z_i=b_i+t_i y_i,\qquad y_i^2=A_i,
\end{equation}
where $t_i=0$ defines $C$ and $u_i,b_i$ are units.
Its inverse curve $B_V$ is integral, and
$p_a(B_V)-g(C)=d/2$.
Lemma~\ref{lem:positive-square-root-genus} gives a reduced integral
Cartier curve $B$ on the normalization, with
\[
 B^2=2d,\qquad
 g(\overline B)=g(C),\qquad
 d=\ell+2\Delta,\quad
 \Delta=p_a(B)-g(C),\quad\ell\ge0.
\]
The suitable resolution in
Lemma~\ref{lem:positive-square-resolution-loss} has smooth strict
transform $D$ satisfying $D^2\ge2d-4\Delta\ge0$.
The logarithmic form of the pulled-back root $L$ has a nonzero
leading conormal coefficient of order zero.  A positive $D^2$
would make this a section of a negative-degree line bundle.
Thus $D^2=0$, $\Delta=d/2$, and $\ell=0$.

For the equality in the local normalization and ambient resolution,
write the successive genus losses as
\[
 e=p_a(B_V)-p_a(B),\quad
 u=p_a(B)-p_a(B_0),\quad
 v=p_a(B_0)-g(C).
\]
They are nonnegative and sum to $d/2$.  The resolution calculation
gives
\[
 0=D^2\ge2d-2u-4v=4e+2u.
\]
Hence $e=u=0$.  The ambient exceptional pullback over $B$ has square
zero and is therefore zero, so $V$ is smooth along $B$.
The finite birational map $B\to B_V$ has zero length quotient
and is an isomorphism.  Nakayama applied to the normalization
extension modulo $t_i$ identifies $V$ with
\eqref{eq:smallest-endpoint-global-local-model} near $B$.
These arguments use the square root of $M$, not algebraic
triviality of $M$.

\smallskip\noindent
\emph{The global conductor and the cover algebra.}
The root-genus proof constructs the rank-one reflexive ideal
\[
 J=\mathfrak c\otimes\OO_V(B)\subset\OO_V.
\]
It is the unit ideal at the generic point of $B$.  Both $J$ and
$\OO_V$ have depth two at closed points of the normal surface $V$.
Thus $Q_0=\OO_V/J$ has no zero-dimensional associated component.
It also has no associated component $B$, so a local equation of
$B$ is a nonzerodivisor on $Q_0$.  Consequently
$J|_B\to\OO_B$ is injective with cokernel of length $\ell=0$.
Nakayama gives $Q_0=0$ on a neighborhood of $B$.

The divisor $B$ is ample, being the finite pullback of $C$.
Any remaining support of $Q_0$ would contain a projective curve,
and such a curve meets $B$.  This contradicts the preceding
vanishing.  Hence $J=\OO_V$ globally.
Finite duality in the root-genus proof now gives
\eqref{eq:smallest-endpoint-global-conductor}.
Since $V$ is Cohen--Macaulay and its dualizing sheaf is invertible,
it is Gorenstein.

Off $C$ the conductor is the unit ideal, so $V$ agrees with the
original torsor there.  Near $C$ it is the finite flat model
\eqref{eq:smallest-endpoint-global-local-model}.  Thus $\pi$ is
finite flat everywhere.  Its algebra has an invertible quotient
$Q$ by the constants.  The original nonconstant summand $M^{-1}$
maps to $Q$; in the displayed basis its generator maps to
$t_i[y_i]$.  This map has zero divisor exactly $C$ and is invertible
off $C$, so $Q\simeq M^{-1}(C)$.

If \eqref{eq:smallest-endpoint-global-cover-sequence} split, its
complementary line $Q$ would have linear changes of local generators.
The square of such a generator has zero component in $Q$, because
this is so in the generic purely inseparable degree-two field.
It would therefore define a map
\[
 Q^{\otimes2}\longrightarrow\OO_X.
\]
This map is a section of $Q^{-2}\simeq\OO_X(-2C)$, hence is zero.
A nonzero generic generator would have square zero in a field,
a contradiction.  This proves nonsplittability.  Finally, $M$ is
numerically trivial because it is torsion, and Riemann--Roch gives
\[
 \chi(M^{-1}(C))
 =\chi(\OO_X)+\frac{d-K_X\cdot C}{2}
 =\chi(\OO_X)-d/2.
\]
The Euler-characteristic formula follows from the sequence.

\smallskip\noindent
\emph{The saturated cotangent line.}
Divisibility of a section by a power of the smooth prime divisor
$C$ can be checked at its generic point in the regular local rings
of $X$.  Hence $\theta=\omega/s_C^2$ is global.
Near $C$, the local expression is $dA_i/u_i$.
Smoothness of $y_i^2=A_i$ over every point of $C$ says that $dA_i$
is nonzero there, so $\theta$ has no ambient zero on $C$.
Off $C$, the original hypersurface torsor is normal.
A divisorial zero of $\omega=du_i/u_i$ would make $z_i^2=u_i$
singular along its inverse divisor, contradicting regularity in
codimension one.  Thus $\theta$ has no divisorial zero anywhere,
and its zero scheme $Z_0$ is finite and disjoint from $C$.

If $f,g$ are its coefficients in a local cotangent frame, wedge
with $\theta$ has image the ideal $(f,g)$ times the determinant
line.  The absence of a common height-one factor gives its kernel
as the line generated by $\theta$.  This proves
\eqref{eq:smallest-endpoint-global-cotangent}; its quotient is
torsion-free, so the first line is saturated.
At a zero, $f,g$ form a regular sequence in the regular local
surface ring.  The zero-scheme formula and $K_X\cdot C=2d$ give
\[
 \operatorname{length}(Z_0)
 =c_2\bigl(\Omega_X^1(-2C)\bigr)
 =c_2(X)-2K_X\cdot C+4d=c_2(X).
\]

\smallskip\noindent
\emph{Global forms and vector fields.}
If $\omega\wedge\beta\ne0$ for a global one-form $\beta$, division
by $s_C^2$ gives a nonzero section of $H=\omega_X(-2C)$.
Here $H\cdot C=0$ and $H|_C=\mathcal T$.
Since the nontrivial degree-zero line $\mathcal T$ has no nonzero
section, its zero divisor contains $C$.  Such an effective divisor
has positive intersection with $C$: its $C$-component contributes
$d>0$, and the remaining components have nonnegative intersection.
This contradicts $H\cdot C=0$.
All global forms are therefore rationally collinear with $\omega$.
Their images in the torsion-free quotient of
\eqref{eq:smallest-endpoint-global-cotangent} vanish generically
and hence everywhere, proving
\eqref{eq:smallest-endpoint-global-forms}.

Dualizing the cotangent sequence and using
$\mathcal Hom(I_{Z_0},\OO_X)=\OO_X$ identifies its annihilator
with $\omega_X^{-1}(2C)$.
Multiplication by $s_C^2$ does not change the sheaf kernel, so
this is precisely the kernel in
\eqref{eq:smallest-endpoint-global-foliation}.
The inclusion into $T_X$ is a line subbundle away from $Z_0$,
although it may fail to be a subbundle at $Z_0$.
Locally $\omega=du/u$, and a derivation $D$ lies in its kernel
exactly when $D(u)=0$.  Commutators and restricted squares of such
derivations again annihilate $u$.  This proves the foliation and
$2$-closedness assertions.  Adjunction gives
\[
 \mathcal F\cdot C=0,\qquad
 \mathcal F|_C\simeq\omega_C^{-1}\otimes N^3
 =\mathcal T^{-1}.
\]

For a global vector field $v$, the function $\omega(v)$ is a
constant on the projective integral surface and vanishes on $C$.
Thus $v$ is a section of $\mathcal F$.
A nonzero section of this line would have an effective zero
divisor of intersection zero with the ample curve $C$.
It would consequently be nowhere zero and would trivialize
$\mathcal F$, contrary to its restriction to $C$.
Hence $H^0(X,T_X)=0$.

Finally, for a smooth rational curve $E$, the restriction
$\OO_E(2C)\to\Omega_X^1|_E$ is nonzero because $\theta$ has no
divisorial zero.  Its composite with
$\omega_E=\OO_{\PP^1}(-2)$ is zero by degree.
The conormal sequence gives a nonzero map
$\OO_E(2C)\to\OO_E(-E)$, proving
\eqref{eq:smallest-endpoint-global-rational-curve}.
Since $C\cdot E\ge1$, a $(-1)$-curve is impossible.
\end{proof}

\begin{remark}[Positive self-intersection and global forms]
\label{rem:smallest-endpoint-positive-square-forms}
In the hypotheses of
Theorem~\ref{thm:smallest-endpoint-global-structure}, replace
ampleness of $C$ by $C^2>0$, retaining smoothness of $X$.
Then rational collinearity of all global one-forms and
$H^0(X,T_X)=0$ still hold.
Indeed, the endpoint proof still gives the nontrivial line
$\mathcal T$ and a divided form $\theta$ with no ambient zero
along $C$.  The preceding argument for a section of
$\omega_X(-2C)$ used only $C^2>0$, so proves collinearity.

For a global vector field $v$, the same constant-contraction
argument gives $\omega(v)=0$.  Near $C$, its annihilator is a
line bundle with restriction $\mathcal T^{-1}=\mathcal T$.
If $v\ne0$ and its exact order along $C$ in this local line is $j$,
division by $C^j$ gives a nonzero first coefficient in
\[
 H^0\bigl(C,\mathcal T\otimes N^{-j}\bigr).
\]
For $j=0$ this contradicts nontriviality of $\mathcal T$;
for $j>0$ the degree is negative.
Thus $v=0$.  This argument does not establish the global conductor
or the exact global cotangent sequence without ampleness.
\end{remark}

\subsection{The elliptic pullback at logarithmic order two}
\label{subsec:elliptic-endpoint-geometry}

At the smallest characteristic-two boundary, the curve-pullback bound
leaves an elliptic base.  The two logarithmic forms on the distinguished
curve determine the exact inseparable degree of that multisection and
constrain the first normalized base change of the surface.

\begin{proposition}[Exact height and ramification]
\label{prop:elliptic-endpoint-height}
Let $X$ be a smooth integral projective surface over an algebraically
closed field $k$ of characteristic two, and let $C\subset X$ be a
smooth integral ample divisor with $d=C^2$.
Suppose $f:X\to E$ is a morphism to an ordinary elliptic curve and
$P\in\Pic^0(E)$ has exact order two, with
\[
 M=f^*P\ne\OO_X,\qquad M|_C\simeq\OO_C.
\]
Assume that the logarithmic form of $M$ has exact order two along $C$
and that $2g(C)-2=3d$.  Let $Z$ be the leading tangential divisor in
Proposition~\ref{prop:smallest-endpoint-rank}, and put
$n=\deg(f|_C)$.
The $P$-torsor $r:E_1\to E$ is a smooth integral degree-two purely
inseparable cover of genus one.  There is a factorization
\begin{equation}\label{eq:elliptic-endpoint-factorization}
 f|_C=r\circ h,\qquad h:C\longrightarrow E_1,
\end{equation}
where $h$ is separable of degree $n/2\geq2$.
Thus the inseparable degree of $f|_C$ is exactly two and $n\geq4$.
The ramification divisor $R_h$, with its different multiplicities,
satisfies
\begin{equation}\label{eq:elliptic-endpoint-ramification}
 \deg R_h=3d,\qquad
 \operatorname{Supp}R_h\cap\operatorname{Supp}Z=\varnothing.
\end{equation}
In particular $h$ is unramified at every point of
$\operatorname{Supp}Z$.
\end{proposition}

\begin{proof}
The nontriviality of $M$ makes $f$ nonconstant.
Proposition~\ref{prop:boundary-curve-pullback-excess} shows that $C$
dominates $E$ and that $n$ is even.  Multiplication by two on
$\Pic^0(E)(k)$ is surjective.  Choose $Q$ with $Q^{\otimes2}\simeq P$
and put $L=f^*Q$.  Then $L^{\otimes2}\simeq M$, so the hypotheses
of Proposition~\ref{prop:smallest-endpoint-rank} hold; in particular
the required N\'eron--Severi divisibility follows from this actual
square root.

The $P$-torsor is integral.  Indeed, a square root of one of its
defining units in $k(E)$ would extend to local unit roots by normality,
and the unique roots would trivialize $P$.  Its logarithmic form
$\omega_P$ is nonzero by the injectivity of the logarithmic-form map
on actual two-torsion of a smooth proper curve, as used in
Lemma~\ref{lem:positive-square-resolved-transfer}.
Since $E$ has genus one, $\omega_P$ is nowhere zero.
The local equation $z^2=u$ of the torsor is therefore smooth, since
$du/u=\omega_P$ has no zero.  Thus the torsor is already normal and
smooth.  Its function field is $k(E)^{1/2}$, so it is an inverse
Frobenius twist of $E$ and has genus one.  The map $r$ is finite,
flat and purely inseparable of degree two.

Set $Q'=r^*Q$.  Its square is trivial, since the torsor trivializes
$r^*P$.  It is nontrivial: otherwise the norm formula would give
\[
 \OO_E\simeq\operatorname{Nm}_r(r^*Q)
       \simeq Q^{\otimes2}\simeq P.
\]
Its logarithmic form $\theta$ is therefore a nonzero, nowhere-zero
regular differential on $E_1$.

The pullback of the $P$-torsor to $C$ is trivial.  Indeed,
$M|_C\simeq\OO_C$, and in a global frame the square trivialization
is a nonzero constant, which has a square root in $k$.
A section of this pulled-back torsor gives the morphism $h$ in
\eqref{eq:elliptic-endpoint-factorization}.
It is nonconstant, hence finite, and
\[
 U=L|_C=h^*Q',\qquad \beta_U=h^*\theta.
\]
By Proposition~\ref{prop:smallest-endpoint-rank}, the form $\beta_U$
is nonzero and is nonvanishing on $\operatorname{Supp}Z$.
A finite map between smooth curves is separable exactly when its
function-field differential is nonzero.  Thus $h$ is separable.
Degree comparison gives $\deg h=n/2$.
The integer $d=\deg Z$ is positive and even, so
$g(C)=1+3d/2\geq4$.
A finite degree-one map between smooth projective curves is an
isomorphism; consequently $h$ cannot have degree one.
The factorization has a separable part of degree $n/2$ and a purely
inseparable part of degree two, proving the asserted exact height.

For the separable map $h$, the differential map
$h^*\omega_{E_1}\to\omega_C$ has zero divisor $R_h$, with the
different multiplicities also in the presence of wild ramification.
Since $\theta$ trivializes $\omega_{E_1}$, this says
\[
 R_h=\operatorname{div}(h^*\theta)
    =\operatorname{div}(\beta_U).
\]
Its degree is $2g(C)-2=3d$, and its support avoids
$\operatorname{Supp}Z$.  This proves
\eqref{eq:elliptic-endpoint-ramification}.
\end{proof}

\begin{proposition}[Smoothness after the first normalized base change]
\label{prop:elliptic-endpoint-smoothing}
Under the hypotheses of Proposition~\ref{prop:elliptic-endpoint-height},
let $W$ be the normalization of $X\times_E E_1$, and write
\[
 \pi:W\longrightarrow X,\qquad
 \varphi:W\longrightarrow E_1,\qquad B=\pi^*C.
\]
Then $\varphi$ is smooth at every point of $B$.
Its non-smooth locus is finite, and its generic fibre is smooth.
If $f_*\OO_X=\OO_E$, its generic fibre is also geometrically integral.
\end{proposition}

\begin{proof}
The base change $X\times_E E_1$ is the $M$-torsor, so $W$ is its
normalized torsor.  Lemma~\ref{lem:positive-square-root-genus}
makes $B$ a reduced integral effective Cartier divisor, and the proof
of Corollary~\ref{cor:positive-excess-characteristic-two-boundary}
shows that $W$ is smooth on a neighbourhood of $B$.
Use $Q,L,Q'$ and $\theta$ from the preceding proof.
The pullback of $L$ to $W$ is $\varphi^*Q'$, and its logarithmic
form is
\[
 \eta=\varphi^*\theta.
\]
On the suitable resolution $S\to W$, its restriction to the smooth
strict transform $D$ is the nowhere-zero conormal section
\eqref{eq:smallest-endpoint-conormal}.

Every point $w\in B$ has a point of $D$ above it.
If the cotangent vector $\eta(w)$ vanished, its pullback at that
point would vanish in $\Omega_S^1|_D$, contradicting the nonzero
conormal section.  Thus $\eta$ is nonzero at every point of $B$.
The surface $W$ is smooth there, and $\theta$ generates the
cotangent line of the smooth curve $E_1$.
The differential criterion for smoothness therefore shows that
$\varphi$ is smooth along $B$.

The non-smooth locus of a finite-type morphism is closed.
The divisor $B$ is ample, since it is the finite pullback of the
ample divisor $C$.  A curve in the non-smooth locus would have to
meet $B$: the restriction of an ample line bundle to a complete
curve has positive degree.  This contradicts smoothness along $B$.
The non-smooth locus is therefore finite.  Its image in $E_1$ is
finite as well, so the generic fibre of $\varphi$ is smooth.

Finally assume $f_*\OO_X=\OO_E$, and put $F=k(E)$ and $K=k(X)$.
The Stein condition says that $F$ is algebraically closed in $K$.
Also $f^*\omega_P=\omega\ne0$, so the map
$\Omega^1_{F/k}\to\Omega^1_{K/k}$ is nonzero.
A separating parameter of $F/k$ consequently has nonzero
differential in $K$ and can be extended to a separating
transcendence basis of $K/k$.  It follows that $K/F$ is separably
generated.  Thus $K/F$ is a regular field extension and remains
geometrically integral after extension to $k(E_1)$.
Normalization does not change this generic function field.
The smooth generic fibre of $\varphi$ is therefore geometrically
integral under the Stein hypothesis.
\end{proof}

The smooth-generic-fibre assertion uses ampleness of $C$.
The separate fibration on the covering resolution given by
Proposition~\ref{prop:boundary-resolved-fibration} already follows
from positive self-intersection.  The simultaneous occurrence of
these structures remains a necessary condition on an elliptic
endpoint configuration.

\subsection{The inseparable Albanese surface branch}
\label{subsec:smallest-endpoint-albanese-geometry}

We now impose the Albanese hypotheses.  Throughout this subsection,
$k$ is algebraically closed of characteristic two, $X$ is a smooth
integral projective surface, and $C\subset X$ is a smooth integral
ample divisor.  Let $M\ne\OO_X$ be algebraically trivial of exact
order two, with $M|_C\simeq\OO_C$.  Assume that its logarithmic form
$\omega$ has exact order two along $C$ and that
\[
 d=C^2,\qquad 2g(C)-2=3d.
\]
Choose a base point on $C$, let
$\alpha:X\to A=\operatorname{Alb}(X)$ be the pointed Albanese map,
and suppose that $\dim\alpha(X)=2$.  Put $q_A=\dim A\ge2$.
Since an algebraically trivial line bundle has a square root,
Theorem~\ref{thm:smallest-endpoint-global-structure} applies.
We retain its notation
\[
 N=\OO_C(C),\qquad
 \mathcal T=\omega_C\otimes N^{-3},\qquad
 \pi:V\longrightarrow X,\qquad
 \theta=\omega/s_C^2,\qquad
 \mathcal F=\omega_X^{-1}(2C).
\]
In particular, $\mathcal T$ has exact order two.

\begin{proposition}[The canonical quotient and its Albanese variety]
\label{prop:smallest-endpoint-albanese-quotient}
The Albanese differential has generic rank exactly one, and
\[
 T_{X/A}=\mathcal F.
\]
Let $q:X\to Y$ be the quotient by this foliation and write
$\alpha=hq$.  Then $Y$ is normal, $q$ is finite and purely inseparable
of degree two, and
\begin{equation}\label{eq:smallest-endpoint-albanese-quotient}
 Y\simeq V^{(2)},\qquad q\pi=F_{V/k}.
\end{equation}
Here $F_{V/k}:V\to V^{(2)}$ denotes relative Frobenius.
The map $h:Y\to A$ is the Albanese map of $Y$.
Consequently
\begin{equation}\label{eq:smallest-endpoint-albanese-frobenius}
 \operatorname{Alb}(V)^{(2)}\simeq A,
\end{equation}
and, under this identification, the homomorphism induced by
$\pi$ is the relative Frobenius of $\operatorname{Alb}(V)$.
There is no assertion that $h$ is generically separable.
\end{proposition}

\begin{proof}
Picard--Albanese duality identifies $\alpha^*$ on $\Pic^0(k)$
with the algebraically trivial line bundles on $X$
\cite[Expos\'e~236, Corollary~3.2 and Theorem~3.3(iii),
pp.~236--238]{FGAPicard}.
Thus $M=\alpha^*P$ for a unique nontrivial order-two line bundle
$P$ on $A$.  Its logarithmic form $\omega_P$ is a nonzero invariant
one-form and satisfies $\alpha^*\omega_P=\omega$.
Nonzeroness follows from logarithmic injectivity: a zero derivative
would give local unit square roots and trivialize $P$.
The global collinearity in
Theorem~\ref{thm:smallest-endpoint-global-structure} gives the upper
bound one for the rank of $d\alpha$, while $\omega$ gives the lower
bound.  Invariant forms generate $\Omega_A^1$, so their common
annihilator is exactly $\mathcal F=T_{X/A}$.

This is the canonical degree-two foliation associated with an
inseparable Albanese map with surface image
\cite[Lemma~2.22 and Remark~2.23, p.~13]{ContiPositiveSeveriComparison}.
We verify the normal quotient and its identification explicitly.
Put $K=k(X)$, and write a rational torsor equation as $z^2=u$.
Since $u\notin K^2$, it can be completed to a $2$-basis $u,v$ of $K$.
The constants of a nonzero derivation annihilating $u$ are
$K^2(u)$: using the basis $1,u,v,uv$ over $K^2$, this is the
kernel of differentiation with respect to $v$.
On the other hand,
\[
 k(V)^2=K^2(u).
\]

Let $R$ be a normal affine coordinate ring in $X$, and let
$\widetilde R$ be its integral closure in $k(V)$.
Every $b\in\widetilde R$ satisfies $b^2\in R$, because $b^2$
belongs to $K$ and is integral over $R$.
Frobenius on $V$ therefore factors through a map
$q':X\to V^{(2)}$.
This map is finite: locally $R$ is finite over $R^2$, hence over
the intermediate Frobenius image ring $\widetilde R^2$.
It has degree two by the field calculation above.
Moreover,
\[
 R\cap k(V)^2=\widetilde R^2.
\]
Indeed, if $r=b^2\in R$ with $b\in k(V)$, then $b$ is integral over
$R$ and belongs to $\widetilde R$.
Thus the structure sheaf of $V^{(2)}$ is precisely the invariant
subsheaf for the foliation.  These identifications use the usual
Frobenius-twist convention over the perfect field $k$.
They identify $q'$ with the normal quotient $q$.
Equivalently, both normal quotients are the integral closures of
the Frobenius base ring in the same intermediate field.
This proves \eqref{eq:smallest-endpoint-albanese-quotient}.
Since the foliation annihilates the pullbacks of local functions
on $A$, the invariant-sheaf description also gives the factor $h$.

For any pointed map $\psi:Y\to A_0$ to an abelian variety,
universality of $\alpha$ gives a unique homomorphism $v:A\to A_0$
such that
\[
 \psi q=v\alpha=vhq.
\]
Equality descends through $q$: it is dominant and $Y$ is reduced,
so its pullback on functions is injective, and the target is
separated.  Hence $\psi=vh$, proving the Albanese universal
property for $h$.
Formation of the Albanese commutes with base change in the cited
FGA theorem.  Applied to $Y\simeq V^{(2)}$, this gives
\eqref{eq:smallest-endpoint-albanese-frobenius}.
The map induced by $q$ is the identity under
$\operatorname{Alb}(Y)=A$, and $q\pi=F_{V/k}$.
The universal property of relative Frobenius now identifies the
map induced by $\pi$ with $F_{\operatorname{Alb}(V)/k}$.
This calculation does not remove any further inseparability
from $h$.
\end{proof}

\begin{proposition}[Invariant forms on the distinguished curve]
\label{prop:smallest-endpoint-albanese-forms}
Let
\[
 \mathcal W=\alpha^*H^0(A,\Omega_A^1)\subset H^0(X,\Omega_X^1).
\]
Then $\dim\mathcal W=q_A$, and tangential restriction to $C$ gives
\begin{equation}\label{eq:smallest-endpoint-albanese-restriction}
 \ker\bigl(\mathcal W\longrightarrow H^0(C,\omega_C)\bigr)
   =k\omega,\qquad
 \dim\operatorname{im}\bigl(\mathcal W\longrightarrow
 H^0(C,\omega_C)\bigr)=q_A-1.
\end{equation}
In particular, $C$ maps finitely and separably to the normalization
of $\alpha(C)$.
\end{proposition}

\begin{proof}
Pullback of invariant one-forms by the Albanese map is injective
for a smooth projective variety in positive characteristic
\cite[Lemma~1.3 and its proof, p.~195]{MehtaSrinivasAlbanese}.
Thus $\dim\mathcal W=q_A$.
All these forms are closed, and $\omega\in\mathcal W$.

Let $0\ne\beta\in\mathcal W$ have zero tangential restriction.
Its full restriction to $\Omega_X^1|_C$ lies in the conormal line
$N^{-1}$, which has negative degree, and is therefore zero.
Let $j\ge1$ be the exact vanishing order of $\beta$ along $C$.
After division by $C^j$, its leading coefficient has nonzero
tangential image: otherwise it would be a nonzero section of
$N^{-(j+1)}$, again of negative degree.
In local parameters $t,x$, write
\[
 \beta=t^j(a\,dt+b\,dx)
\]
with $b$ nonzero at the generic point of $C$.
The coefficient of $t^{j-1}dt\wedge dx$ in $d\beta$ is $jb$.
Closedness forces $j$ to be even, hence $j\ge2$.

By \eqref{eq:smallest-endpoint-global-forms},
$\beta=\theta s$ for a section $s$ of $\OO_X(2C)$.
The divided form $\theta$ has exact order zero along $C$,
so $s$ is divisible by $s_C^2$.
The quotient is a global regular function on $X$, hence a constant.
Thus $\beta\in k\omega$; the reverse inclusion is immediate.
This proves \eqref{eq:smallest-endpoint-albanese-restriction}.
Its image has dimension $q_A-1>0$, so $\alpha|_C$ has nonzero
differential.  The resulting nonconstant map to the normalized
image curve is finite and separable.
This assertion concerns the curve map, not the surface map.
\end{proof}

\begin{proposition}[The curve divisor and Albanese dimension]
\label{prop:smallest-endpoint-albanese-dimension}
Let $R$ be the zero divisor of the tangential component of
$\theta|_C$, a section of $\omega_C\otimes N^{-2}$.
There is an effective divisor $E_C$ such that
\begin{equation}\label{eq:smallest-endpoint-albanese-curve-lines}
 R=2E_C,\qquad \deg E_C=d/2,\qquad
 N\simeq\OO_C(2E_C)\otimes\mathcal T,\qquad
 \omega_C\simeq\OO_C(6E_C).
\end{equation}
Let $R_\alpha$ be the common zero divisor of the nonzero image
space in \eqref{eq:smallest-endpoint-albanese-restriction}.
Then
\begin{equation}\label{eq:smallest-endpoint-albanese-dimension}
 d\le\deg R_\alpha\le3d-2q_A+4,\qquad
 2\le q_A\le d+2.
\end{equation}
\end{proposition}

\begin{proof}
The tangential component is nonzero.  Otherwise the nonzero leading
restriction of $\theta$ would lie in $N^{-3}$, which has negative
degree.
The local model
\eqref{eq:smallest-endpoint-global-local-model} gives the
tangential component as
\[
 \frac{d(A_i|_C)}{(b_i|_C)^2}.
\]
In a completed local parameter on $C$, the derivative of a power
series in characteristic two has only even powers.  Its zero order
is even, and the denominator is a square unit.
Consequently $R=2E_C$.
Its degree is $2g(C)-2-2d=d$, and
\[
 \OO_C(R)=\omega_C\otimes N^{-2}=N\otimes\mathcal T.
\]
Since $\mathcal T^2\simeq\OO_C$, these identities give
\eqref{eq:smallest-endpoint-albanese-curve-lines}.

Equation~\eqref{eq:smallest-endpoint-global-forms}, restricted
tangentially to $C$, shows that $R_\alpha\ge R$.
Put $H=\omega_C(-R_\alpha)$, and write
\[
 r=h^0(C,H),\qquad s=h^0(C,\omega_C\otimes H^{-1}).
\]
We have $r\ge q_A-1$ and
$s=h^0(C,\OO_C(R_\alpha))\ge1$.
We give the Clifford estimate needed here directly.
At a smooth point, choose bases of these two spaces adapted to
vanishing order, with strictly increasing orders
$a_1<\cdots<a_r$ and $b_1<\cdots<b_s$.
The products with orders
\[
 a_1+b_1<\cdots<a_r+b_1<a_r+b_2<\cdots<a_r+b_s
\]
are $r+s-1$ independent canonical sections.
Hence $r+s-1\le g(C)$.
Combining this with Riemann--Roch,
$r-s=\deg H+1-g(C)$, gives
$2r\le\deg H+2$.
Therefore
\[
 2q_A-2\le3d-\deg R_\alpha+2.
\]
Together with $\deg R_\alpha\ge d$, this proves
\eqref{eq:smallest-endpoint-albanese-dimension}.
\end{proof}

\begin{remark}
\label{rem:smallest-endpoint-albanese-zero-divisor-scope}
The divisor $R_\alpha$ records zeros of the differential of the map
to $A$.  It need not be the ramification divisor of the separable
map from $C$ to its normalized image: the map from that normalized
curve into $A$ can itself have zeros in its differential.
No identification of these two divisors is used.
\end{remark}

\begin{proposition}[The normalized image curve]
\label{prop:smallest-endpoint-albanese-image}
Let $B_0$ be the normalization of $\alpha(C)$, write
$f:C\to B_0$ for the finite separable map, and
$j:B_0\to A$ for the normalization map followed by inclusion.
For the line bundle $P$ in the proof of
Proposition~\ref{prop:smallest-endpoint-albanese-quotient},
\begin{equation}\label{eq:smallest-endpoint-albanese-image-torsion}
 j^*P\simeq\OO_{B_0}.
\end{equation}
The curve $j(B_0)$ generates $A$.  In particular, with $b=g(B_0)$,
one has $b\ge q_A\ge2$, and
\[
 \operatorname{Jac}(B_0)\longrightarrow A
\]
is surjective and inseparable.  Its pullback on invariant
cotangent spaces has one-dimensional kernel.
\end{proposition}

\begin{proof}
If $j^*P$ were nontrivial, its $\mu_2$-torsor over the smooth curve
$B_0$ would give a nontrivial purely inseparable quadratic extension
of $k(B_0)$.
Its pullback line to $C$ is trivial because $M|_C$ is trivial.
The pulled-back square trivialization differs from a square frame
by a constant; every constant in $k^*$ has a square root.
Thus the pulled-back torsor is trivial, placing its defining square
root in $k(C)$.  This contradicts separability of
$k(C)/k(B_0)$ and proves
\eqref{eq:smallest-endpoint-albanese-image-torsion}.

Let $A_0$ be the abelian subvariety generated by $j(B_0)$.
Our base point on $C$ makes this curve pass through the origin.
The composite $X\to A/A_0$ contracts $C$.
If it were nonconstant, choose a hyperplane section of a projective
embedding of its image which avoids the image of $C$ and meets
that positive-dimensional image.  Its pullback is a nonzero
effective divisor disjoint from $C$, contrary to ampleness.
Hence $\alpha(X)\subset A_0$, and the universal property of
$\alpha$ forces $A_0=A$.
The Jacobian homomorphism is therefore surjective, giving
$b\ge q_A\ge2$.

Pullback of regular one-forms by the separable curve map $f$ is
injective.  Thus the kernel of $j^*$ on invariant one-forms is
the kernel after pulling back further to $C$.
By \eqref{eq:smallest-endpoint-albanese-restriction} and injectivity
of $\alpha^*$, it is precisely $k\omega_P$.
Under the usual identification of invariant differentials on
$\operatorname{Jac}(B_0)$ with $H^0(B_0,\omega_{B_0})$,
this is the cotangent kernel of the Jacobian homomorphism.
Its differential has rank $q_A-1$, so the surjective homomorphism
is inseparable.
\end{proof}

\begin{remark}
\label{rem:smallest-endpoint-albanese-normalized-image-scope}
The normalized image curve $B_0$ is not thereby the base of a
morphism from $X$.  In particular, the curve-pullback excess
inequality cannot be applied to $B_0$ without an additional
construction of such a morphism and a pullback description of $M$.
\end{remark}

\begin{proposition}[The image divisor on an abelian surface]
\label{prop:smallest-endpoint-albanese-surface-defect}
Suppose $q_A=2$ and put $D_A=\alpha(C)\subset A$ with its reduced
integral structure.  Let $\rho:A'\to A$ be the $\mu_2$-isogeny
associated with $P$.
The map $j:B_0\to A$ lifts to $j':B_0\to A'$.
If $D'$ is its integral image and
$\epsilon=p_a(D')-b$, then $\epsilon\ge0$ and
\begin{equation}\label{eq:smallest-endpoint-albanese-surface-defect}
 \rho^*D_A=2D',\qquad
 p_a(D_A)-b=b-1+2\epsilon,\qquad
 D_A^2=4(b-1+\epsilon).
\end{equation}
Thus $D_A$ is singular, its normalization defect is at least $b-1$,
and $D_A^2$ is a positive multiple of four.
Equality in the defect bound holds exactly when $D'$ is smooth.
These are assertions about the image divisor, not a formula for
$C^2$.
\end{proposition}

\begin{proof}
The $\mu_2$-cover of the nontrivial line $P$ is integral.
It is smooth because its local defining equation has differential
the nonzero invariant form $\omega_P$, which has no zeros.
With a point over the origin it is an abelian variety and
$\rho$ is an isogeny; this is the abelian-cover construction in
the proof of \cite[Lemma~1.3, p.~195]{MehtaSrinivasAlbanese}.
Equation~\eqref{eq:smallest-endpoint-albanese-image-torsion},
with the same square-scalar adjustment as above, gives the lift
$j'$.

Since $j$ is birational onto $D_A$, the intermediate fields satisfy
\[
 k(D_A)\subset k(D')\subset k(B_0)=k(D_A).
\]
Thus $D'\to D_A$ is birational, and $B_0$ normalizes both curves.
The radicial degree-two isogeny has a unique divisor above $D_A$.
Its residue-field degree there is one, so the ramification index
is two and $\rho^*D_A=2D'$.
Finite intersection projection gives
\[
 2D_A^2=(\rho^*D_A)^2=4(D')^2.
\]
Adjunction on the two smooth abelian surfaces gives
\[
 p_a(D_A)=1+\frac{D_A^2}{2},\qquad
 p_a(D')=1+\frac{(D')^2}{2}.
\]
Substituting $p_a(D')=b+\epsilon$ yields
\eqref{eq:smallest-endpoint-albanese-surface-defect}.
The normalization defect $\epsilon$ is nonnegative.
It vanishes exactly when the integral curve $D'$ is normal,
equivalently smooth over the perfect field $k$.
Since $b\ge2$, the asserted strict singularity of $D_A$ follows.
\end{proof}

\begin{remark}[The image-curve pattern]
\label{rem:smallest-endpoint-albanese-image-pattern}
The image-divisor conditions alone are compatible with examples.
Take ordinary elliptic curves $E_1,E_2$ and the degree-two
$\mu_2$-isogeny
\[
 \rho=F_{E_1/k}\times1:
 A'=E_1\times E_2\longrightarrow E_1^{(2)}\times E_2=A.
\]
Choose a smooth sufficiently ample integral curve $D'\subset A'$
whose projection to $E_2$ has nonzero differential.
Such a choice follows from Bertini in a sufficiently ample linear
system separating first-order jets: impose at one smooth point
a tangent direction not killed by the second projection, and
choose a general member smooth elsewhere.
The restriction of $\rho$ to $D'$ is radicial and has nonzero
differential.  It is therefore birational onto its image $D_A$.
Its normalization kills the line $P$ defining $\rho$, and the
defect equality in
\eqref{eq:smallest-endpoint-albanese-surface-defect} holds with
$\epsilon=0$.

This constructs only the image-curve behavior.
Taking $X=A'$ would make $\rho^*P$ trivial, and $\rho$ would not be
the Albanese map of $A'$.
It is not a construction of a surface satisfying the endpoint
hypotheses.  The existence or exclusion of the full inseparable
Albanese endpoint configuration remains open.
\end{remark}

\subsection{Numerical restrictions at the smallest endpoint}
\label{subsec:smallest-endpoint-numerical}

The nontrivial discrepancy line at logarithmic order two restricts
the space of cubic sections of the normal bundle.  The hyperplane-section
estimate then gives a lower bound for the degree of a very ample
endpoint.  A second restriction arises when the Albanese variety is a
surface and its Albanese morphism has degree two: the normalized torsor
is then birational to an abelian surface.

Throughout this subsection, let $k$ be algebraically closed of
characteristic two, let $X$ be a smooth integral projective surface,
and let $C\subset X$ be a smooth integral ample divisor.  Put
\[
 d=C^2,\qquad N=\OO_C(C).
\]
Let $M\ne\OO_X$ have exact order two, be trivial on $C$, and admit a
square root; equivalently, $[M]\in2\operatorname{NS}(X)$.
Assume that its logarithmic form has exact order two along $C$ and
that
\[
 2g(C)-2=3d.
\]
We use the notation of
Theorem~\ref{thm:smallest-endpoint-global-structure}:
$\pi:V\to X$ is the normalized torsor, $B=\pi^*C$, and
$\mathcal T=\omega_C\otimes N^{-3}$.
In particular, $d$ is even, $\mathcal T$ has exact order two, and
$K_X\cdot C=2d$.

\begin{proposition}[The degree of a very ample endpoint]
\label{prop:smallest-endpoint-very-ample-degree}
If $C$ is very ample, then
\[
 d\ge12,\qquad
 h^0(C,N)\le\left\lfloor\frac d4\right\rfloor+1,\qquad
 h^0(X,\OO_X(C))\le\left\lfloor\frac d4\right\rfloor+2.
\]
More generally, an embedding of $C$ by a subspace of $H^0(C,N)$
as a nondegenerate curve in $\PP^n$, with $n\ge2$, satisfies
$d\ge4n$.
\end{proposition}

\begin{proof}
Since $\mathcal T$ is a nontrivial degree-zero line bundle, Serre
duality and Riemann--Roch give
\begin{equation}\label{eq:smallest-endpoint-cubic-dimension}
 h^1(C,N^3)=h^0(C,\mathcal T)=0,\qquad
 h^0(C,N^3)=3d+1-g(C)=\frac{3d}{2}.
\end{equation}

We apply the hyperplane-section argument in the proof of
Lemma~\ref{lem:small-degree-curve-bound} before its final
small-degree specialization.  That argument applies to every
$n\ge2$.  Indeed, the twisted tangent--normal sequence shows that a
smooth embedded curve of genus at least two is nonstrange: a common
tangent point would give a nonzero section of $T_C\otimes N^{-1}$,
whose degree is negative.  Kadets's sectional-monodromy theorem
holds for $n\ge2$
\cite[Proposition~2.8 and proof, pp.~4--5]{Kadets}.
Thus the geometric generic hyperplane section $\Gamma$ consists
of $d$ points in linear general position in $\PP^{n-1}$.
The spanning and interpolation arguments in the cited proof give
\[
 d\ge n,\qquad
 h_\Gamma(j)\ge\min\{d,j(n-1)+1\}\quad(j\ge1).
\]
If $h_C$ denotes the homogeneous-coordinate Hilbert function,
multiplication by a hyperplane equation is injective in the
coordinate ring of $C$, and its quotient surjects onto the
coordinate ring of $\Gamma$.  Consequently
\begin{equation}\label{eq:smallest-endpoint-cubic-bound}
 \begin{split}
 \frac{3d}{2}
 &=h^0(C,N^3)\ge h_C(3)\\
 &\ge 1+n+\min\{d,2n-1\}+\min\{d,3n-2\}.
 \end{split}
\end{equation}
The inequality uses neither projective normality nor uniform
position in higher degree.

If $d<2n-1$, the last expression is $1+n+2d$, which is greater
than $3d/2$.  If $2n-1\le d<3n-2$, it is $3n+d$, again greater
than $3d/2$.  Hence $d\ge3n-2$, and
\eqref{eq:smallest-endpoint-cubic-bound} gives
$3d/2\ge6n-2$.  Since $d$ is even, this implies $d\ge4n$.

The restriction $N$ of the very ample bundle $\OO_X(C)$ is very
ample.  Its complete embedding cannot have image in $\PP^1$,
because $g(C)\ge2$.  If its image were a plane curve, adjunction
would give $d(d-3)=3d$, hence $d=6$ and $\omega_C=N^3$.
This contradicts $\mathcal T\ne\OO_C$.
Thus $n=h^0(C,N)-1\ge3$, proving the degree and normal-section
bounds.  Finally the kernel of
\[
 H^0(X,\OO_X(C))\longrightarrow H^0(C,N)
\]
is the one-dimensional space generated by the defining section
of $C$.  This proves the bound on sections on $X$.
\end{proof}

\begin{proposition}[An Albanese morphism of degree two]
\label{prop:smallest-endpoint-albanese-degree-two}
Suppose in addition that $M$ is algebraically trivial,
$A=\operatorname{Alb}(X)$ has dimension two, and the Albanese
morphism $\alpha:X\to A$ is generically finite of degree two.
Then $V$ is birational to the abelian surface
$A_V=\operatorname{Alb}(V)$, the canonical divisor $K_X$ is nef
and big, and
\begin{equation}\label{eq:smallest-endpoint-albanese-chern-interval}
 d=4\chi(\OO_X),\qquad
 2d\le K_X^2<3d,\qquad
 0<c_2(X)=3d-K_X^2\le d.
\end{equation}
In particular, $4\mid d$ and $X$ is a minimal surface of general type.
\end{proposition}

\begin{proof}
Let $\beta:V\to A_V$ be the Albanese morphism.
Proposition~\ref{prop:smallest-endpoint-albanese-quotient}
identifies $A_V^{(2)}$ with $A$ and gives
\[
 \alpha\pi=F_{A_V/k}\beta.
\]
Relative Frobenius of the abelian surface has degree four, whereas
$\pi$ has degree two.  Thus
\[
 2\deg\alpha=4\deg\beta.
\]
It follows that $\deg\beta=1$, so $\beta$ is proper and birational.

We first compute $\chi(\OO_V)$ without assuming that $V$ is smooth.
Choose a projective resolution $r:S\to V$ and put $b=\beta r$.
The proper birational morphism $b$ between smooth surfaces is a
composition of blowups at closed points
\cite[Lemma~54.17.1, Tag~0C5Q]{Stacks}.
A point blowup has direct image $\OO$ and vanishing higher direct
images of $\OO$.  For the first higher direct image, this follows
on an affine neighborhood with regular parameters $x,y$ from the
two affine blowup charts: every Laurent expression in $y/x$ on
their overlap is a sum of expressions on the two charts.
The resulting first \v{C}ech cohomology vanishes; higher direct
images above degree one vanish by fibre dimension.
These statements persist under composition.
Normality gives $r_*\OO_S=\OO_V$ and
$\beta_*\OO_V=\OO_{A_V}$.
The low-degree Leray sequence gives
\[
 R^1\beta_*\OO_V\hookrightarrow R^1b_*\OO_S=0.
\]
The higher direct images of $\beta$ vanish above degree one.
Consequently
\[
 \chi(\OO_V)=\chi(\OO_{A_V})=0.
\]
The Euler-characteristic formula in
Theorem~\ref{thm:smallest-endpoint-global-structure} now gives
$d=4\chi(\OO_X)$.

Pull back a nowhere-zero two-form on $A_V$ to the regular locus
of $V$.  It extends over the omitted codimension-two set because
$\omega_V$ is reflexive.  Its zero divisor $E_V$ is effective
and supported on the exceptional locus of $\beta$.
It is nonzero, since
\[
 E_V\cdot B=K_V\cdot B
 =2(K_X-C)\cdot C=2d>0.
\]
Here $V$ is Gorenstein and
$\omega_V\simeq\pi^*(\omega_X(-C))$ by
\eqref{eq:smallest-endpoint-global-conductor}.
On $S$, the Cartier pullback $r^*E_V$ is a nonzero effective
$b$-exceptional divisor.  It is orthogonal to $b^*H$ for an ample
divisor $H$ on $A_V$, and $(b^*H)^2>0$.
The Hodge index theorem gives $(r^*E_V)^2<0$: equality would make
this effective nonzero divisor numerically zero, contradicting
its positive intersection with an ample divisor on $S$.
Therefore
\[
 0>K_V^2=2(K_X-C)^2=2K_X^2-6d.
\]

The finite-flat norm of the nonzero canonical section on $V$
is a nonzero section of $\omega_X^{\otimes2}(-2C)$.
Indeed, the determinant of multiplication defines the norm of
a section, and $\operatorname{Nm}_\pi(\pi^*L)=L^{\otimes2}$
for every line bundle $L$ on $X$.
Thus an effective divisor $E_X$ satisfies
\begin{equation}\label{eq:smallest-endpoint-albanese-canonical-norm}
 2K_X\sim2C+E_X.
\end{equation}
Since $C$ is ample, $K_X$ is big.

To prove nefness, suppose an integral curve $E$ has $K_X\cdot E<0$.
Equation~\eqref{eq:smallest-endpoint-albanese-canonical-norm},
ampleness of $C$, and nonnegative intersection of distinct effective
curves force $E$ to be a component of $E_X$ with $E^2<0$.
Adjunction gives
\[
 2p_a(E)-2=K_X\cdot E+E^2.
\]
Both terms on the right are negative integers, so $p_a(E)\ge0$
forces $K_X\cdot E=E^2=-1$ and $p_a(E)=0$.
Then $E$ is a smooth rational $(-1)$-curve, contrary to
Theorem~\ref{thm:smallest-endpoint-global-structure}.
Thus $K_X$ is nef.  Intersecting
\eqref{eq:smallest-endpoint-albanese-canonical-norm} with $K_X$
now gives $2K_X^2\ge2K_X\cdot C=4d$.
Finally Noether's formula, obtained from
Grothendieck--Riemann--Roch in every characteristic
\cite[Tag~02UO]{Stacks}, gives
\[
 c_2(X)=12\chi(\OO_X)-K_X^2=3d-K_X^2.
\]
This proves \eqref{eq:smallest-endpoint-albanese-chern-interval}.
Nefness and bigness give the general-type assertion, and the absence
of $(-1)$-curves gives minimality.
\end{proof}

\begin{corollary}[Very ample degree in the Albanese degree-two case]
\label{cor:smallest-endpoint-albanese-very-ample-degree}
Under Proposition~\ref{prop:smallest-endpoint-albanese-degree-two},
suppose that $C$ is very ample and let
$r=h^0(X,\OO_X(C))-1$ be the dimension of its complete embedding.
Then
\[
 r\ge5,\qquad 4\mid d,\qquad d\ge16,\qquad
 6\le h^0(X,\OO_X(C))\le d/4+2.
\]
If $d=16$, then
\[
 h^0(X,\OO_X(C))=6,\qquad h^0(C,N)=5.
\]
\end{corollary}

\begin{proof}
An embedded surface has $r\ge2$.  For $r=2$ its image is $\PP^2$,
contrary to $K_X\cdot C=2d>0$.
For $r=3$ it is a degree-$d$ hypersurface, so
$K_X=(d-4)C$.  The same intersection identity gives $d=6$
and $K_X=2C$, contrary to $\mathcal T\ne\OO_C$.

Suppose $r=4$ and write $E=N_{X/\PP^4}$.
The tangent--normal and Euler sequences, followed by the
self-intersection formula, give
\[
 \deg c_2(E)=10d+5K_X\cdot C+K_X^2-c_2(X)=d^2.
\]
These are the identities used in the proof of
Proposition~\ref{prop:very-ample-numerical-obstruction};
the self-intersection statement is valid in every characteristic
\cite[Tags~0FV9 and~0FEZ]{Stacks}.
Using $K_X\cdot C=2d$ and
\eqref{eq:smallest-endpoint-albanese-chern-interval} gives
\begin{equation}\label{eq:smallest-endpoint-p4-double-point}
 d^2=17d+2K_X^2.
\end{equation}
The interval $2d\le K_X^2<3d$ yields $21\le d<23$.
There is no multiple of four in this interval, so $r\ne4$.

Thus $r\ge5$.  Proposition~\ref{prop:smallest-endpoint-very-ample-degree}
gives $r+1\le d/4+2$, hence $d\ge4(r-1)\ge16$.
At $d=16$ the upper and lower section bounds give
$h^0(X,\OO_X(C))=6$.
Restriction to $C$ then gives $h^0(C,N)\ge5$, whereas the same
proposition gives $h^0(C,N)\le5$.
\end{proof}

\begin{remark}[Scope of the numerical bounds]
\label{rem:smallest-endpoint-numerical-scope}
The very ample degree bound has no Albanese hypothesis.
The Chern interval requires both $\dim\operatorname{Alb}(X)=2$
and Albanese degree two.  A surface image in an Albanese variety
of larger dimension, or a larger degree over an abelian surface,
does not give the birational morphism used in its proof.
For a merely ample endpoint the cubic-section bound does not apply.
The degree-sixteen conclusion is a necessary condition and makes
no existence assertion.
\end{remark}

\subsection{Canonical and image restrictions in Albanese degree two}
\label{subsec:smallest-endpoint-strict-canonical}

The finite Stein factor and the Albanese image of the distinguished
curve impose related restrictions on the smallest endpoint.
The canonical bound is strict.  More precisely, the normalization
defect of the lifted image determines the resolution defect of the
Stein factor, and two exceptional cycles give a further inequality.
All these assertions require only ampleness of the curve.

Let $k$ be algebraically closed of characteristic two, let $X$ be a
smooth integral projective surface, and let $C\subset X$ be a smooth
integral ample divisor.  Put $d=C^2>0$.  Let $M\ne\OO_X$ be
algebraically trivial of exact order two, with $M|_C\simeq\OO_C$.
Suppose its logarithmic form has exact order two along $C$ and
\[
 2g(C)-2=3d.
\]
Assume that $A=\operatorname{Alb}(X)$ has dimension two and that
the Albanese morphism $\alpha:X\to A$ is generically finite of
degree two.  Proposition~\ref{prop:smallest-endpoint-albanese-degree-two}
gives
\begin{equation}\label{eq:strict-canonical-prior-numerics}
 \begin{gathered}
 d=4\chi(\OO_X),\qquad K_X\cdot C=2d,\qquad
 K_X\ \text{nef and big},\\
 2d\le K_X^2<3d,\qquad c_2(X)=3d-K_X^2>0.
 \end{gathered}
\end{equation}
Let $\pi:V\to X$ be the normalized torsor and let
$\beta:V\to A_V=\operatorname{Alb}(V)$ be its proper
birational Albanese morphism.
The global conductor and Albanese quotient results give
\begin{equation}\label{eq:strict-canonical-cover-data}
 \omega_V\simeq\pi^*(\omega_X(-C)),\qquad
 \alpha\pi=F_{A_V/k}\beta,\qquad A_V^{(2)}\simeq A.
\end{equation}
Here $V$ is normal and Gorenstein, and $\pi$ is finite flat of
degree two.

Retain the image notation of
Proposition~\ref{prop:smallest-endpoint-albanese-surface-defect}.
Thus $D_A=\alpha(C)$ has its reduced integral structure, $B_0$ is
its normalization, and $b=g(B_0)$.
Write $M=\alpha^*P$ and let $\rho:A'\to A$ be the degree-two
$\mu_2$-isogeny defined by $P$.
The map $B_0\to A$ lifts to $A'$; let $D'$ be its integral image and put
\begin{equation}\label{eq:strict-canonical-image-defect-definition}
 \epsilon=p_a(D')-b\ge0.
\end{equation}
This is the defect of the image in $A'$, not of an image in $A_V$.
No upper bound on $\epsilon$ is imposed.

\begin{theorem}[The canonical interval, Stein defect and image divisor]
\label{thm:smallest-endpoint-strict-canonical}
Under these hypotheses,
\begin{equation}\label{eq:strict-canonical-interval}
 2d<K_X^2<3d,\qquad 1\le c_2(X)\le d-1.
\end{equation}
Write the Stein factorization of $\alpha$ as
\[
 X\xrightarrow{\,r\,}X_0\xrightarrow{\,a_0\,}A.
\]
Then $X_0$ is normal and Gorenstein, $a_0$ is a finite flat
double cover, and $r$ is proper birational.
The quotient $Q=a_{0*}\OO_{X_0}/\OO_A$ is invertible.
Set
\[
 \mathcal L=Q^{-1},\qquad
 t=\operatorname{length}(R^1r_*\OO_X),\qquad
 D=K_X-r^*K_{X_0},
\]
using compatible canonical divisors.
The sheaf defining $t$ has finite support.
There is a nonzero effective exceptional divisor $Z$ such that
\begin{equation}\label{eq:strict-canonical-discrepancy}
 D=-Z,\qquad D^2<0,\qquad p_a(Z)=1.
\end{equation}

The map $C\to B_0$ is an isomorphism, and
\begin{equation}\label{eq:strict-canonical-image-class}
 D_A\equiv2\mathcal L
\end{equation}
numerically on $A$.
Let $E_X$ be the norm divisor constructed in the proof; it is
effective, has full exceptional support, and satisfies
$2K_X\sim2C+E_X$.
Then the actual divisor equality is
\begin{equation}\label{eq:strict-canonical-image-pullback}
 \alpha^*D_A=2C+E_X+2Z.
\end{equation}
The two defects satisfy
\begin{equation}\label{eq:strict-canonical-stein-defect}
 K_X^2=d+4t+Z^2,\qquad 2t=d+\epsilon.
\end{equation}
In particular $\epsilon$ is even.
Moreover,
\begin{equation}\label{eq:strict-canonical-double-cover-line}
 \omega_{X_0}\simeq a_0^*\mathcal L,\qquad
 \mathcal L^2=2\bigl(\chi(\OO_X)+t\bigr)
             =\frac{3d}{2}+\epsilon,
\end{equation}
and
\begin{equation}\label{eq:strict-canonical-image-cycle}
 C\cdot Z=d+2\epsilon,\qquad
 Z^2=-c_2(X)-2\epsilon.
\end{equation}

The exceptional intersection form gives
\begin{equation}\label{eq:strict-canonical-defect-tradeoff}
 4c_2(X)t\ge d^2,\qquad
 t\ge\max\left\{\frac d2,
          \left\lceil\frac{d^2}{4c_2(X)}\right\rceil\right\}.
\end{equation}
Equivalently, $2c_2(X)(d+\epsilon)\ge d^2$.
Equality in $4c_2(X)t\ge d^2$ holds precisely when
\begin{equation}\label{eq:strict-canonical-tradeoff-equality}
 (d-c_2(X))E_X=2c_2(X)Z
\end{equation}
as real exceptional divisors.
The conditions $\epsilon=0$, $t=d/2$, and smoothness of $D'$
are equivalent; in that case $c_2(X)\ge d/2$.
In particular $X_0$ has a nonrational singularity and $\alpha$
is not finite.
\end{theorem}

The canonical formula for a normal finite flat double cover is
recalled in \cite[Section~2.2, equations~(2.9)--(2.15),
pp.~8--9]{ContiPositiveSeveriComparison}.
The algebra calculation below makes its nonsplit inseparable case
explicit. The numerical range here lies above
$8\chi(\OO_X)$; the characteristic-two Severi comparison in
\cite[Theorem~8.4, p.~51]{ContiPositiveSeveriComparison} concerns
$K_X^2<(9/2)\chi(\OO_X)$ with an additional condition.

\begin{proof}
\par\smallskip\noindent\textbf{The full exceptional support of the norm divisor.}
Pull back a nowhere-zero two-form on $A_V$ by $\beta$.  On the
regular locus of $V$ this is a nonzero regular canonical section;
reflexivity of $\omega_V$ extends it over the omitted
codimension-two set.  Let $E_V$ be its effective zero divisor.

Its support is exactly the set of $\beta$-exceptional curves.
A proper birational map to a normal surface is an isomorphism
in codimension one away from its contracted curves, so there is
no zero elsewhere.  Conversely, at the generic point of a
contracted curve $F$, the surface $V$ is regular.  Two local
parameters at $\beta(F)$ restrict to constants on $F$, so their
pulled-back differentials lie in its conormal line.  Their wedge
vanishes.  Thus every contracted curve occurs in $E_V$ with
positive coefficient.

The finite-flat norm of this section, using
\eqref{eq:strict-canonical-cover-data}, gives
\begin{equation}\label{eq:strict-canonical-full-support-norm}
 2K_X\sim2C+E_X,\qquad E_X\ge0.
\end{equation}
The support of $E_X$ is the image of the support of $E_V$.
Indeed, the determinant of multiplication by a section is a unit
exactly when that section is a unit in the finite locally free
algebra.  This observation does not require a separable norm.
For a curve $F\subset V$, finiteness of $\pi$ makes
$E=\pi(F)$ a curve.  The Frobenius relation and finiteness of
Frobenius give
\[
 \beta(F)\text{ is a point}
 \quad\Longleftrightarrow\quad
 \alpha(E)\text{ is a point}.
\]
Every $\alpha$-exceptional curve has a curve above it under $\pi$.
Consequently $E_X$ contains every $\alpha$-exceptional curve,
each with positive coefficient, and has no other component.

\par\smallskip\noindent\textbf{The exceptional intersection form.}
Choose an ample divisor $H$ on $A$ and put $P=\alpha^*H$.
Then
\[
 P^2=2H^2>0,\qquad P\cdot E=0
\]
for every exceptional curve $E$.  Hodge index gives negative
square for every nonzero effective exceptional divisor.
It also gives negative definiteness for signed exceptional
divisors, as follows.  If a nonzero such divisor $U$ were
numerically zero, write $U=U_+-U_-$ with disjoint nonnegative
coefficient supports.  When $U_+\ne0$, one would have
\[
 U_+^2=U_+\cdot U_-\ge0,
\]
contrary to Hodge index applied to $U_+$.
When $U_+=0$, the divisor $-U$ is effective, nonzero and
numerically trivial, contrary to its positive intersection
with an ample divisor on $X$.  Thus a nonzero signed
exceptional divisor is never numerically zero.  Hodge index
on the orthogonal complement of $P$ proves negative definiteness.

\par\smallskip\noindent\textbf{The finite double cover.}
Stein factorization gives $r_*\OO_X=\OO_{X_0}$ and a finite
map $a_0:X_0\to A$
\cite[Theorem~37.53.4, Tag~03GX]{Stacks}.
Since $X$ is normal, $X_0$ is the normalization of $A$ in $k(X)$.
Thus $X_0$ is normal, $a_0$ has generic degree two, and $r$ is
birational.  Normal surfaces are Cohen--Macaulay.  Miracle
flatness therefore makes $a_0$ flat over the regular surface
$A$ \cite[Lemma~10.128.1, Tag~00R4]{Stacks}.
Put $\mathcal B=a_{0*}\OO_{X_0}$, a locally free algebra of rank two.

The constants form a line subbundle: over a local ring of $A$,
the element $1$ is nonzero modulo the maximal ideal and extends
to a basis of $\mathcal B$.  Hence
\begin{equation}\label{eq:strict-canonical-cover-algebra}
 0\longrightarrow\OO_A\longrightarrow\mathcal B
 \longrightarrow Q\longrightarrow0
\end{equation}
has invertible quotient.  We do not need this sequence to split.
Set $\mathcal L=Q^{-1}$.

For the canonical formula, choose a local basis $1,z$ with
$z^2=az+b$ and let $\lambda:\mathcal B\to\OO_A$ take the
coefficient of $z$.  The pairing $(u,v)\mapsto\lambda(uv)$
has matrix
\[
 \begin{pmatrix}0&1\\1&a\end{pmatrix},
\]
whose determinant is a unit in every characteristic.
Under a change $z_i=c_{ij}z_j+b_{ij}$, its functional changes
by $\lambda_i=c_{ij}^{-1}\lambda_j$.  The pairings therefore
glue to an isomorphism of $\mathcal B$-modules
\[
 \mathcal Hom_{\OO_A}(\mathcal B,\OO_A)
 \simeq\mathcal B\otimes Q^{-1}.
\]
Finite duality identifies the left side with the pushforward
of the relative dualizing sheaf
\cite[Sections~49.5 and~49.15, Tags~0FKW and~0DWM]{Stacks}.
Since $\omega_A\simeq\OO_A$, this proves that $X_0$ is
Gorenstein and
\begin{equation}\label{eq:strict-canonical-finite-identities}
 \omega_{X_0}\simeq a_0^*\mathcal L,\qquad
 \chi(\OO_{X_0})=\frac{\mathcal L^2}{2},\qquad
 K_{X_0}^2=2\mathcal L^2=4\chi(\OO_{X_0}).
\end{equation}
The Euler identity follows from
\eqref{eq:strict-canonical-cover-algebra} and Riemann--Roch on
the abelian surface, where $K_A=0$ and $\chi(\OO_A)=0$.
The intersection identity follows from finite projection.
Thus these formulas apply also to a nonsplit inseparable
double-cover algebra.

\par\smallskip\noindent\textbf{Cohomology of exceptional $(-2)$-cycles.}
We need the following vanishing: if every exceptional component
of $r$ is a smooth rational $(-2)$-curve, then
\begin{equation}\label{eq:strict-canonical-cycle-vanishing}
 R^1r_*\OO_X=0.
\end{equation}
To prove it, let $G>0$ be any effective exceptional cycle.
Negative definiteness gives $G^2<0$, so some component $E$
appearing in $G$ satisfies $G\cdot E<0$.
The Cartier-divisor sequence
\[
 0\longrightarrow\OO_E(-(G-E))
 \longrightarrow\OO_G\longrightarrow\OO_{G-E}
 \longrightarrow0
\]
is exact, including when $E$ has multiplicity in $G$.
The degree of its first term is
\[
 -(G-E)\cdot E=-G\cdot E-2\ge-1.
\]
Its first cohomology vanishes on $E\simeq\PP^1$.
Induction on the sum of the coefficients of $G$ proves
$H^1(G,\OO_G)=0$ for every effective exceptional cycle.

The map $r$ is an isomorphism away from finitely many points
and has connected fibres.  A nontrivial fibre has support
equal to the union of its exceptional curves: connectedness
excludes an isolated point component outside that union.
Over the local ring at such a point, let $I$ be the extended
maximal ideal and let $J=\OO_X(-G_{\mathrm{red}})$, where
$G_{\mathrm{red}}$ is the reduced exceptional fibre.
These ideals have the same radical, so their powers are cofinal
by Noetherianity and quasi-compactness.
For any $n$, choose $m,N$ with
\[
 I^N\subset J^m\subset I^n.
\]
Then the map
\[
 H^1(\OO_X/I^N)\longrightarrow H^1(\OO_X/I^n)
\]
factors through
$H^1(\OO_X/J^m)=H^1(\OO_{mG_{\mathrm{red}}})=0$.
Its inverse limit is therefore zero.
Formal functions identifies that limit with the completed
stalk of $R^1r_*\OO_X$
\cite[Theorem~30.20.5, Tag~02OC]{Stacks}.
The stalk is a finite module supported at the closed point;
its zero completion makes it zero.  This proves
\eqref{eq:strict-canonical-cycle-vanishing}.
No normal-crossing hypothesis on the exceptional divisor was used.

\par\smallskip\noindent\textbf{Strictness of the canonical bound.}
Suppose that $K_X^2=2d$.  Intersecting
\eqref{eq:strict-canonical-full-support-norm} with $K_X$ gives
\[
 K_X\cdot E_X=2K_X^2-2K_X\cdot C=0.
\]
By nefness and the full exceptional support of $E_X$,
$K_X\cdot E=0$ for every $\alpha$-exceptional curve.
These are exactly the exceptional curves of $r$, since
$a_0$ is finite.  Each has negative square, and adjunction gives
\[
 2p_a(E)-2=K_X\cdot E+E^2=E^2<0.
\]
Thus $p_a(E)=0$ and $E^2=-2$.
The normalization genus and normalization defect of an integral
proper curve are nonnegative and sum to its arithmetic genus.
Hence each $E$ is smooth rational, and
\eqref{eq:strict-canonical-cycle-vanishing} applies.

The divisor $K_{X_0}$ is Cartier by
\eqref{eq:strict-canonical-finite-identities}.
For compatible canonical divisors,
$D=K_X-r^*K_{X_0}$ is exceptional, and
$D\cdot E=K_X\cdot E=0$ for every exceptional $E$.
Negative definiteness forces $D=0$.
Thus $K_X^2=K_{X_0}^2$.
Normality gives $r_*\OO_X=\OO_{X_0}$, and higher direct
images above degree one vanish by fibre dimension.
Together with \eqref{eq:strict-canonical-cycle-vanishing},
this gives $\chi(\OO_X)=\chi(\OO_{X_0})$.
Equation~\eqref{eq:strict-canonical-finite-identities} now yields
\[
 K_X^2=4\chi(\OO_X)=d,
\]
contrary to $K_X^2=2d$ and $d>0$.
The prior interval \eqref{eq:strict-canonical-prior-numerics}
therefore gives \eqref{eq:strict-canonical-interval}.
In particular, some exceptional curve has positive intersection
with $K_X$; otherwise the same argument would give $K_X^2=d$.

\par\smallskip\noindent\textbf{The discrepancy and singularity defect.}
Retain $D=K_X-r^*K_{X_0}$ without the equality assumption.
It is exceptional and satisfies
\[
 D\cdot E=K_X\cdot E\ge0
\]
for every exceptional curve.
If its positive part $D_+$ were nonzero, write
$D=D_+-D_-$ with disjoint supports.  Then
\[
 0\le D_+\cdot D=D_+^2-D_+\cdot D_-<0,
\]
using negative definiteness and the nonnegative intersection
of distinct effective curves.  This contradiction proves $D\le0$.
If $D=0$, every exceptional curve again has zero canonical
intersection.  The preceding cycle argument would give
$R^1r_*\OO_X=0$ and $K_X^2=d$, contradicting
\eqref{eq:strict-canonical-prior-numerics}.
Thus $D=-Z\ne0$, with $Z$ effective and exceptional, and $D^2<0$.

The sheaf $R^1r_*\OO_X$ has finite support.  Leray, with the
degree-one direct image contributing with negative sign, gives
\[
 \chi(\OO_X)=\chi(\OO_{X_0})-t.
\]
The Cartier pullback of $K_{X_0}$ is orthogonal to $D$.
Using \eqref{eq:strict-canonical-finite-identities} and
$4\chi(\OO_X)=d$, we obtain
\[
 \begin{split}
 K_X^2
 &=K_{X_0}^2+D^2\\
 &=4\chi(\OO_X)+4t+D^2
 =d+4t+D^2.
 \end{split}
\]
Since $D^2<0$ and $K_X^2>2d$, one has
$t>(K_X^2-d)/4>0$.
Thus $R^1r_*\OO_X\ne0$, so $X_0$ has a nonrational
singularity and $r$ is not an isomorphism.
Consequently $\alpha$ is not finite.
The same Euler calculation gives
\[
 \mathcal L^2
 =2\chi(\OO_{X_0})
 =2\bigl(\chi(\OO_X)+t\bigr).
\]
Since $Z=-D$ is exceptional, $K_X\cdot Z=-Z^2$, and adjunction
for the effective Cartier cycle gives
\[
 p_a(Z)=1+\frac{Z^2+K_X\cdot Z}{2}=1.
\]

\par\smallskip\noindent\textbf{The degree of the distinguished curve.}
By Proposition~\ref{prop:smallest-endpoint-albanese-quotient},
the Albanese differential has generic rank one.
The extension $k(X)/k(A)$ has degree two and is not separable,
so it is purely inseparable of exponent one.
On a normal affine ring $R$ of $A$, let $S$ be its integral
closure in $k(X)$.
Every $s\in S$ has $s^2\in k(A)$ integral over $R$, hence
$s^2\in R$.
There is therefore a unique prime of $S$ over each prime of $R$:
if two primes have the same contraction, membership of $s$ in
one puts $s^2$ in that contraction and then puts $s$ in the other.
The residue-field extensions are purely inseparable as well.

By Proposition~\ref{prop:smallest-endpoint-albanese-forms},
the map $C\to B_0$ is finite and separable; in particular $C$
is not contracted.
The proper birational map $r$ is an isomorphism at the generic
point of $C_0=r(C)$, because $X_0$ is normal.
Thus $k(C)=k(C_0)$, and the preceding description of the finite
Stein cover makes $k(C)/k(D_A)$ purely inseparable.
It is also separable, and hence has degree one.
A finite birational map between the smooth curves $C$ and $B_0$
is an isomorphism.  In particular,
\begin{equation}\label{eq:strict-canonical-curve-degree}
 b=g(C),\qquad \alpha_*[C]=[D_A].
\end{equation}

We also need the coefficient of $C$ in $\alpha^*D_A$.
Localize $A$ at the generic point of $D_A$.
Its local ring $R$ is a discrete valuation ring, and its integral
closure $S$ in $k(X)$ is a local normal one-dimensional ring,
hence a discrete valuation ring.
The finite flat Stein cover makes $S$ free of rank two over $R$.
Its residue-field degree is one, by
\eqref{eq:strict-canonical-curve-degree}.
If a uniformizer of $R$ has valuation $e$ in $S$, the filtration
by powers of a uniformizer of $S$ gives
\[
 2=\dim_{\kappa(R)}S/\mathfrak m_RS
   =e[\kappa(S):\kappa(R)]=e.
\]
This argument applies to the purely inseparable cover directly.
The same valuation occurs along $C$ since $r$ is an isomorphism
at its generic point.
The unique nonexceptional curve over $D_A$ is $C$, so
\begin{equation}\label{eq:strict-canonical-image-initial-pullback}
 \alpha^*D_A=2C+E_b
\end{equation}
for some effective exceptional divisor $E_b$.

\par\smallskip\noindent\textbf{The image class and exceptional correction.}
For any Cartier divisor $H$ on $A$, exceptional divisors have
zero intersection with $\alpha^*H$.
The canonical identity and proper intersection projection give
\[
 K_X\cdot\alpha^*H
 =\alpha^*\mathcal L\cdot\alpha^*H
 =2\mathcal L\cdot H.
\]
Here projection holds for a proper morphism, including one
with contracted curves
\cite[Lemmas~42.26.3--42.26.4, Tags~02ST and~02SU]{Stacks}.
Intersecting \eqref{eq:strict-canonical-full-support-norm}
with $\alpha^*H$ and using
\eqref{eq:strict-canonical-curve-degree} therefore gives
\[
 4\mathcal L\cdot H
 =2K_X\cdot\alpha^*H
 =2C\cdot\alpha^*H
 =2D_A\cdot H.
\]
Every integral curve on the smooth surface $A$ is Cartier.
Thus $D_A\equiv2\mathcal L$, proving
\eqref{eq:strict-canonical-image-class}.

Now \eqref{eq:strict-canonical-image-initial-pullback}, the norm
relation, and $K_X\equiv\alpha^*\mathcal L-Z$ give
\[
 E_b-E_X-2Z\equiv0.
\]
This is a signed exceptional divisor.
Negative definiteness forces it to vanish as a divisor.
Consequently $E_b=E_X+2Z$, proving
\eqref{eq:strict-canonical-image-pullback}.
Only numerical equivalence on $A$ has been used; the pullback
formula on $X$ is an equality of divisors.

\par\smallskip\noindent\textbf{The Stein and image defects.}
The image formula in
Proposition~\ref{prop:smallest-endpoint-albanese-surface-defect},
together with $b=g(C)$, yields
\[
 D_A^2=4(b-1+\epsilon)=6d+4\epsilon.
\]
The numerical image class and the finite-cover Euler identity give
\[
 D_A^2=4\mathcal L^2
       =8\bigl(\chi(\OO_X)+t\bigr)=2d+8t.
\]
Comparing them proves $2t=d+\epsilon$.
Since $4\mid d$, the integer $\epsilon$ is even.
Its nonnegativity gives $t\ge d/2$, and substitution gives
\[
 \mathcal L^2=\frac{3d}{2}+\epsilon.
\]
This proves the defect and line identities in the statement.

The norm relation gives $C\cdot E_X=2d$.
Intersecting the exact pullback formula with $C$ and applying
\eqref{eq:strict-canonical-curve-degree}, we obtain
\[
 6d+4\epsilon=D_A^2
 =2d+C\cdot E_X+2C\cdot Z=4d+2C\cdot Z.
\]
Thus $C\cdot Z=d+2\epsilon$.
Also
\[
 Z^2=K_X^2-d-4t
    =K_X^2-3d-2\epsilon=-c_2(X)-2\epsilon,
\]
proving \eqref{eq:strict-canonical-image-cycle}.
The equivalence of $\epsilon=0$ with smoothness of $D'$ was
proved in Proposition~\ref{prop:smallest-endpoint-albanese-surface-defect};
the defect identity makes it equivalent to $t=d/2$.

\par\smallskip\noindent\textbf{The exceptional-cycle inequality.}
Put $c=c_2(X)$ and consider the effective real exceptional divisors
\[
 U=\frac12E_X,\qquad W=Z+\frac12E_X.
\]
Their numerical classes are $U\equiv K_X-C$ and
$W\equiv\alpha^*\mathcal L-C$.
Hence
\[
 U^2=(K_X-C)^2=K_X^2-3d=-c.
\]
The image class and the curve degree give
\[
 \alpha^*\mathcal L\cdot C
 =\mathcal L\cdot D_A
 =2\mathcal L^2=(\alpha^*\mathcal L)^2.
\]
Consequently
\[
 W^2=d-2\mathcal L^2=-4t.
\]
Since $U$ is exceptional,
\[
 U\cdot W=-U\cdot C=-(K_X\cdot C-d)=-d.
\]
The negative of the intersection form is a positive-definite inner
product on the real exceptional space.
Cauchy--Schwarz gives $d^2\le4ct$.
Combining this with $t\ge d/2$ and integrality of $t$ proves
\eqref{eq:strict-canonical-defect-tradeoff}.
Replacing $2t$ by $d+\epsilon$ gives its equivalent form.

Equality holds precisely when $W=(d/c)U$, since the inner
products of $U$ with itself and with $W$ are $c$ and $d$.
This is equivalent to $(d-c)E_X=2cZ$, proving
\eqref{eq:strict-canonical-tradeoff-equality}.
Finally, when $\epsilon=0$, the inequality becomes
$2cd\ge d^2$, and hence $c\ge d/2$.
\end{proof}

\begin{remark}[Degree sixteen and scope]
\label{rem:smallest-endpoint-strict-degree-sixteen}
If $C$ is very ample and $d=16$, then
Corollary~\ref{cor:smallest-endpoint-albanese-very-ample-degree}
gives
\[
 h^0(X,\OO_X(C))=6,\qquad h^0(C,\OO_C(C))=5.
\]
Theorem~\ref{thm:smallest-endpoint-strict-canonical} requires
\[
 33\le K_X^2\le47,\qquad 1\le c_2(X)\le15,\qquad
 t\ge\max\left\{8,\left\lceil\frac{64}{c_2(X)}\right\rceil\right\},
\]
and
\[
 \epsilon=2t-16,\qquad \mathcal L^2=24+\epsilon.
\]
For example, $c_2(X)=1,2,4,8$ require respectively
$t\ge64,32,16,8$.
If $D'$ is smooth, then $t=8$ and $8\le c_2(X)\le15$.
These are necessary conditions, with no existence assertion.

All conclusions of the theorem require only ampleness of $C$;
very ampleness is used here solely for the degree and section counts.
The surface $X_0$ is the finite Stein factor of $\alpha$,
not the foliation quotient $V^{(2)}$.
The cycle $Z$ may be reducible or nonreduced; $p_a(Z)=1$
does not assert that an exceptional component is a smooth elliptic
curve.
The numerical identity $D_A\equiv2\mathcal L$ does not assert
linear equivalence.
The normalization defect $\epsilon$ is not bounded above, and
the exceptional-cycle inequality does not exclude the entire
degree-two branch.
The image-pattern examples in
Remark~\ref{rem:smallest-endpoint-albanese-image-pattern}
do not satisfy all the endpoint surface hypotheses.
\end{remark}

\subsection{Formal depth and descent of the torsion class}
\label{subsec:conductor-picard-descent}

The exact order of the logarithmic form also measures the order to
which the torsion line bundle is trivial on the infinitesimal
neighborhoods of the curve. At conductor equality, this gives a
precise obstruction to lifting the radial action to the line bundle.
A $p$th tensor power supplies a lift; at an external center
its descent again decreases the conductor.

\begin{proposition}\label{prop:torsion-formal-depth}
Let $X$ be a normal projective integral surface over an algebraically
closed field $k$ of characteristic $p>0$, let $Y\subset X$ be a
smooth integral effective Cartier divisor contained in the smooth
locus, with $Y^2>0$, and let $M\ne\OO_X$ have order $p$
with $M|_Y\simeq\OO_Y$.
Let $pa$ be the exact order along $Y$ of its logarithmic form.
Then
\[
 M|_{aY}\simeq\OO_{aY},
 \qquad M|_{(a+1)Y}\not\simeq\OO_{(a+1)Y}.
\]
In particular, $a$ is the largest integer $n$ for which $M|_{nY}$
is trivial.
\end{proposition}

\begin{proof}
Put $I=\OO_X(-Y)$. For every $n\geq1$ one has
\begin{equation}\label{eq:thickening-constant-functions}
 H^0(nY,\OO_{nY})=k.
\end{equation}
Indeed, the successive kernels in the restriction to smaller
thickenings are
$I^j/I^{j+1}\simeq\OO_Y(-jY)$ for $j\geq1$.
Their degrees are negative, so they have no nonzero sections.
Starting with $H^0(Y,\OO_Y)=k$ proves the assertion; constants
provide the surviving sections. Consequently
\begin{equation}\label{eq:thickening-cohomology-injection}
 H^1(X,I^n)\longrightarrow H^1(X,\OO_X)
 \quad\text{is injective}
\end{equation}
for every $n\geq1$.

We first establish a local approximation. Choose a sufficiently
small smooth affine neighborhood of a point of $Y$ with
\'etale coordinates $t,z$, where $Y=(t=0)$, and denote its ring
by $R$. The monomials $t^iz^j$, $0\leq i,j<p$, form a basis of
$R$ over $R^p$. This is the polynomial-ring calculation transported
along the \'etale map by relative Frobenius
\cite[Tag~0EBS]{Stacks}. If $u\in R^*$ and
$du\in t^{pa}\Omega^1_{R/k}$, write
\[
 u=\sum_{0\leq i,j<p}r_{ij}^{\,p}t^iz^j.
\]
In the $dt$ coefficient of $du$, the monomials with $i>0$ are
distinct members of the $p$-basis; the $dz$ coefficient treats
those with $j>0$. Since $t^{pa}\in R^p$, freeness implies that
$r_{ij}^{\,p}$ is divisible by $t^{pa}$ whenever $(i,j)\ne(0,0)$.
Thus $r_{ij}$ is divisible by $t^a$, and
\begin{equation}\label{eq:local-torsion-pth-approximation}
 u=v^p+t^{pa}b.
\end{equation}
Here $v=r_{00}$ is a unit after shrinking, because $u$ is a unit
along $Y$. This calculation takes place in the algebraic local
rings.

Choose local frames $e_i$ of $M$ and a trivialization of $M^p$,
and let $u_i$ be the corresponding units. Then
$\omega=du_i/u_i$. Replacing the frames using
\eqref{eq:local-torsion-pth-approximation} gives, near $Y$,
\[
 u_i=1+t_i^{pa}b_i.
\]
If $e_i=g_{ij}e_j$, then
$g_{ij}^p=u_i/u_j\equiv1\pmod{I^{pa}}$.
Hence $(g_{ij}-1)^p\in I^{pa}$, and regularity of the parameter
defining $Y$ gives $g_{ij}-1\in I^a$.
The frames therefore glue on $aY$, proving the first assertion.

Suppose that $M$ were trivial on $(a+1)Y$. Lift a trivializing
frame locally, so that $g_{ij}\equiv1\pmod{I^{a+1}}$.
The units $u_i$ then agree modulo $I^{p(a+1)}$, and glue to a
global unit on $p(a+1)Y$. By
\eqref{eq:thickening-constant-functions} this unit is a scalar.
Rescaling the frames by a $p$th root of that scalar gives
\[
 u_i=1+t_i^{p(a+1)}b_i.
\]
Differentiation produces no term of order $p(a+1)-1$, so
$\omega$ has order at least $p(a+1)$ along $Y$. This contradicts
its exact order $pa$.
\end{proof}

\begin{proposition}\label{prop:radial-linearization-obstruction}
Assume the hypotheses and equality of
Theorem~\ref{thm:conductor-equality-structure}.
Let $\delta$ be the radial derivation induced by
$x_0\partial/\partial x_0$ in the homogeneous coordinates used there, and
write $I=\OO_X(-Y)$.
For transition functions $g_{ij}$ of $M$, put
\[
 o_\delta(M)=
 \bigl[\delta(g_{ij})/g_{ij}\bigr]\in H^1(X,\OO_X).
\]
Then $M$ admits a $\mu_p$-linearization for the radial action if
and only if $o_\delta(M)=0$.
Moreover,
\begin{equation}\label{eq:radial-leading-obstruction}
 p\nmid a\quad\Longrightarrow\quad o_\delta(M)\ne0.
\end{equation}
Thus a linearization, and in particular descent of $M$ to the
radial quotient, requires $p\mid a$.
\end{proposition}

\begin{proof}
A lift of $\delta$ to $M$ is a $k$-linear operator $\nabla$ with
$\nabla(fe)=\delta(f)e+f\nabla(e)$.
In the frame $e_i$, write $\nabla(e_i)=\alpha_i e_i$.
Compatibility is precisely
\[
 \alpha_i-\alpha_j=\delta(g_{ij})/g_{ij}.
\]
Thus the displayed class is independent of the frames and is the
obstruction to an ordinary lift.

If a lift exists, $\nabla^p-\nabla$ is $\OO_X$-linear:
its commutator with multiplication by $f$ is multiplication by
$\delta^p(f)-\delta(f)=0$.
It is therefore multiplication by a scalar $c\in k$.
Choose $\lambda\in k$ such that $\lambda^p-\lambda=-c$.
The adjusted operator satisfies
\[
 (\nabla+\lambda)^p=\nabla+\lambda.
\]
The eigenspace projections for the distinct roots of $T^p-T$
are polynomials in this operator. They commute with restriction
to open sets, and the Leibniz rule makes the resulting character
grading compatible with the grading on $\OO_X$ induced by
$\delta$. This is a $\mu_p$-linearization.
Conversely, differentiating a linearization gives such a lift,
with the restricted identity already satisfied.
This proves the equivalence.

By Proposition~\ref{prop:torsion-formal-depth}, the first
nontrivial infinitesimal class of $M$ is
\[
 0\ne\xi\in H^1(Y,I^a/I^{a+1}).
\]
Indeed, the square-zero unit sequence for
$aY\subset(a+1)Y$ identifies the kernel of restriction on Picard
groups with this cohomology group: the global units on both
thickenings are $k^*$ and restriction is surjective.
Choose frames compatible with the trivialization on $aY$.
Their transitions satisfy $g_{ij}-1\in I^a$ and represent $\xi$
modulo $I^{a+1}$.

The derivation vanishes on $Y$ and has weight one on $I/I^2$.
It therefore has weight $a$ on $I^a/I^{a+1}$, so
\[
 \frac{\delta(g_{ij})}{g_{ij}}
 \equiv a(g_{ij}-1)\pmod{I^{a+1}}.
\]
The obstruction is represented by a cocycle in $I^a$, whose image
in $H^1(Y,I^a/I^{a+1})$ is $a\xi$.
When $p\nmid a$, it is nonzero. Its image in $H^1(X,\OO_X)$
remains nonzero by
\eqref{eq:thickening-cohomology-injection}.
This proves \eqref{eq:radial-leading-obstruction}.
A bundle pulled back from the quotient has its canonical
linearization, giving the final assertion.
\end{proof}



\begin{proposition}\label{prop:external-conductor-descent}
Assume the hypotheses and equality of
Theorem~\ref{thm:conductor-equality-structure}, and suppose that
its radial center $v$ lies outside $X$. Let
$q:X\to S=X/\mu_p$ be the radial quotient and let $Y_S=q(Y)$.
Then $S$ is smooth projective, $q|_Y:Y\to Y_S$ is an isomorphism,
and $Y_S$ is a smooth very ample divisor satisfying
\begin{equation}\label{eq:external-quotient-polarization}
 q^*Y_S=pY,\qquad Y_S^2=pY^2.
\end{equation}
The homomorphism $q^*:\Pic(S)\to\Pic(X)$ is injective.

If $M$ admits a $\mu_p$-linearization, it descends to a nontrivial
line bundle $M_S$ of exact order $p$, trivial on $Y_S$.
The exact order of its logarithmic form along $Y_S$ is $a$.
Thus $p\mid a$ and
\[
 2g(Y_S)-2=p(a/p)Y_S^2,\qquad a/p\geq5.
\]
\end{proposition}

\begin{proof}
Smoothness and projectivity of $S$ follow from the quotient
calculation in
Theorem~\ref{thm:conductor-equality-structure}.
To identify its polarization, let $R$ be the homogeneous
coordinate ring of $X$, with $\mu_p$ scaling $x_0$ and fixing
$x_1,\ldots,x_r$. The invariant ring is a quotient of
\[
 k[z,x_1,\ldots,x_r],\qquad z=x_0^p,\quad
 \deg z=p,\quad\deg x_i=1.
\]
Indeed, an equivariant surjection is surjective on each character
summand. The $p$th Veronese of this weighted polynomial ring
is generated in degree one: a monomial of weight $np$ is a
product of $n$ monomials of weight $p$. The same holds for its
quotient. Hence $S=\operatorname{Proj}(R^{\mu_p})$ is embedded
by $\OO_S(p)$.

There is an isomorphism
\[
 R^{\mu_p}/(x_0^p)\simeq R/(x_0).
\]
Surjectivity follows from the invariant generators $x_i$, $i>0$.
An invariant element divisible by $x_0$ is $x_0$ times an element
of weight $p-1$. Lift the latter to a polynomial of that weight;
each of its monomials is divisible by $x_0^{p-1}$ with invariant
quotient. This proves the kernel assertion.
Thus $z$ cuts out $Y_S$ scheme-theoretically, identifies it with
$Y$, and gives $\OO_S(Y_S)=\OO_S(p)$.
Its pullback is the divisor of $x_0^p$, proving the first
identity in \eqref{eq:external-quotient-polarization}.
The second follows by intersection and $\deg q=p$.
This argument does not require projective normality of $X$.

We next prove injectivity on Picard groups. If $q^*P$ is trivial,
its canonical linearization acts on its one-dimensional space
of global sections by a character. A nonzero section is nowhere
zero. Evaluate it at a $k$-point of the fixed divisor $Y$.
The action on that fibre of a pullback bundle is trivial, so the
character is trivial. The invariant section descends, since
\[
 (q_*q^*P)^{\mu_p}=P
\]
by the projection formula and
$(q_*\OO_X)^{\mu_p}=\OO_S$.
Surjectivity of $q$ makes the descended section nowhere zero.
Therefore $P$ is trivial.

Suppose now that $M$ is linearized. Its action on the trivial
bundle $M|_Y$ is one character, since $Y$ is proper integral.
Twist the linearization by the inverse character so that its
fibre action on $Y$ is trivial.
On an invariant affine neighborhood of a point of $Y$, a
weight-zero section of $M$ lifts a nonzero vector in the fibre:
take the weight-zero part of any lift.
After shrinking, it is an invariant frame.
Away from $Y$, the function $x_0/x_i$ is an invertible function
of weight one on suitable charts. A homogeneous local frame
exists by decomposing a frame into its character components
and choosing one nonzero at the point. Multiplication by a
power of $x_0/x_i$ makes its weight zero.
The nonvanishing opens are invariant and descend to opens of
$S$, since $q$ is a universal homeomorphism.
Thus invariant frames glue to a line bundle
\[
 M_S=(q_*M)^{\mu_p},\qquad q^*M_S\simeq M.
\]
Injectivity of $q^*$ makes $M_S$ nontrivial of exact order $p$,
and $Y\simeq Y_S$ gives $M_S|_{Y_S}\simeq\OO_{Y_S}$.

Let $\omega_S$ be its logarithmic form. Its $p$th-power
trivialization pulls back to one differing from that of $M$
by a scalar, so
\begin{equation}\label{eq:descended-logarithmic-form}
 q^*\omega_S=\omega.
\end{equation}
In particular $\omega_S\ne0$.
Let $n$ be its exact order along $Y_S$. The restriction to
$\Omega^1_{Y_S}$ is zero, since $M_S|_{Y_S}$ is trivial.
The remaining conormal section has negative degree, so $n\geq1$.
The leading tangential coefficient is nonzero: otherwise the
leading section would lie in
\[
 H^0\bigl(Y_S,\OO_{Y_S}(-(n+1)Y_S)\bigr)=0.
\]
Write locally $\omega_S=s^n(A\,ds+B\,dz)$ with $B$ nonzero
modulo $s$. Closedness has leading coefficient
$n s^{n-1}B\,ds\wedge dz$, so $p\mid n$.

At a general closed point of $Y$, the completed quotient
coordinates are $s=t^p$, $z=z$. Pullback kills $ds$, and the
nonzero tangential term has exact order $pn$. Thus
\[
 pa=pn,\qquad n=a.
\]
The use of completed local rings at a closed point makes no
perfectness assumption on $k(Y)$.
Consequently the new conductor integer is $a/p$, and
\eqref{eq:external-quotient-polarization} gives
\[
 2g(Y_S)-2=pa\,Y^2=aY_S^2=p(a/p)Y_S^2.
\]
Theorem~\ref{thm:conductor-equality-structure}, applied to the new
smooth datum, gives $a/p\geq5$.
\end{proof}

\begin{lemma}\label{lem:radial-power-descent}
Under the hypotheses and equality of
Theorem~\ref{thm:conductor-equality-structure}, suppose
$M\simeq L^{\otimes p}$ for an invertible sheaf $L$ on $X$.
Then $M$ admits a radial $\mu_p$-linearization, $p\mid a$,
and $M=q^*P$ for an invertible sheaf $P$ on the radial quotient
$q:X\to S$. If the center is external, $P$ is the nontrivial
order-$p$ descended bundle of
Proposition~\ref{prop:external-conductor-descent}.
\end{lemma}

\begin{proof}
Let $g_{ij}$ be the transition functions of $L$.
The transitions $g_{ij}^p$ of $M$ are invariant under the
$\mu_p$-coaction. Indeed, writing its coordinate algebra as
$k[\epsilon]/(\epsilon^p)$, the difference between the coaction
of a function and that function lies in the ideal $(\epsilon)$,
whose elements have $p$th power zero.
Every open is stable under the height-one action.
Giving the local tensor-power frames weight zero therefore
glues a linearization of $M$.
Proposition~\ref{prop:radial-linearization-obstruction} gives
$p\mid a$.

For the invertible descent, on every invariant affine chart
one has $R^p\subset R^{\mu_p}\subset R$.
Since $k$ is perfect, relative Frobenius factors as
\[
 X\xrightarrow{q}S\xrightarrow{g}X^{(p)},
 \qquad gq=F_{X/k}.
\]
The ring inclusions glue; moreover $R^{\mu_p}$ is finite over
$R^p$ as a submodule of the finite $R^p$-module $R$.
Set $P=g^*(L^{(p)})$. Then
\[
 q^*P=F_{X/k}^*L^{(p)}\simeq L^{\otimes p}\simeq M.
\]
For an external center, injectivity of $q^*$ in
Proposition~\ref{prop:external-conductor-descent} identifies
$P$ with that proposition's descended bundle.
\end{proof}

The radial descent uses a very ample embedding, whereas the
normalization argument in Section~\ref{subsec:positive-square-picard}
uses only positive self-intersection. The finite-root obstruction
therefore applies also to ample curves and to curves on normal
surfaces that are smooth in their neighborhood. For a general
N\'eron--Severi torsion class the necessary roots may not exist;
the obstruction in Corollary~\ref{thm:positive-square-picard}
states exactly which divisibility is excluded at equality.

\subsection{Cartier descent at an internal radial center}
\label{subsec:internal-cartier-descent}

The finite-root obstruction in
Corollary~\ref{thm:positive-square-picard} depends on the $p$-adic
valuation of the exact logarithmic order. In the singular case of
Theorem~\ref{thm:conductor-equality-structure}, the radial quotient
gives a stronger obstruction: the torsion bundle cannot descend as
an invertible sheaf, and its N\'eron--Severi class is not even a
double. The argument takes place on the smooth resolution of that
quotient. Its polarization contracts the exceptional curve, so the
reducedness step in the Euler-extension argument requires an exact
local multiplicity calculation.

\begin{proposition}\label{prop:internal-cartier-descent-obstruction}
Let $X$ be a smooth projective integral surface over an algebraically
closed field $k$ of characteristic two. Let $Y\subset X$ be a smooth
integral very ample divisor, and let $M\ne\OO_X$ have order two with
$M|_Y\simeq\OO_Y$. Let $2a$ be the exact order along $Y$ of its
logarithmic form, and suppose that
\[
 2g(Y)-2=2aY^2.
\]
Assume that the radial center constructed in the proof of
Theorem~\ref{thm:conductor-equality-structure} lies on $X$, and let
$q:X\to S=X/\mu_2$ be the radial quotient. There is no invertible
sheaf $P$ on $S$ such that $q^*P\simeq M$. Consequently
\[
 M\notin 2\Pic(X),\qquad
 [M]\notin 2\operatorname{NS}(X).
\]
Here $2\Pic(X)$ denotes the subgroup of tensor squares. Equivalently,
the conclusion applies to the case in which the distinguished
normalized torsor has its single $A_1$ singularity.
\end{proposition}

\begin{proof}
\emph{The quotient and its resolution.}
Put $H=\OO_X(Y)$ and $d=Y^2$. We first specify the structural
conclusions that will be used. The proof of
Theorem~\ref{thm:conductor-equality-structure} gives
\[
 K_X\sim(2a-1)Y
\]
and homogeneous coordinates $x_0,\ldots,x_r$ for the embedding by
$H$, with
\[
 v=[1:0:\cdots:0]\in X,
 \qquad Y=X\cap(x_0=0).
\]
The radial action scales $x_0$ and fixes the other homogeneous
coordinates. The divided logarithmic form has no divisorial zero,
and its zero scheme is the single reduced point $v$.

The normal projective quotient $S$ has one singular point
$o=q(v)$, of type $A_1$. Two of the ambient affine coordinates
$u=x_i/x_0$ and $w=x_j/x_0$ are regular parameters at $v$. They
have weight one for the $\mu_2$-action, and the completed invariant
ring is
\begin{equation}\label{eq:internal-cartier-invariant-ring}
 \widehat{\OO}_{S,o}
 =k[[A,B,C]]/(AC-B^2),
 \qquad A=u^2,\quad B=uw,\quad C=w^2.
\end{equation}
Let $b:R=\operatorname{Bl}_vX\to X$ have exceptional curve $F$.
The lifted action has a finite purely inseparable quotient
$q':R\to T$ of degree two, with $T$ smooth and projective. It fits
into
\[
 \rho q'=qb,\qquad \rho:T\longrightarrow S,
\]
where $\rho$ is the minimal resolution, its only exceptional curve
$E$ has square $-2$, and $q'^*E=2F$. On the chart $w=uz$ the
quotient coordinates are $A=u^2,z$, so
\[
 B=Az,\qquad C=Az^2,\qquad E=(A=0).
\]
The saturated lifted foliation on $R$ has line bundle
$\OO_R(b^*Y+F)$. The regular-foliation determinant formula used in
the structure theorem therefore gives
\begin{equation}\label{eq:internal-cartier-canonical-pullback}
 q'^*K_T=K_R-(b^*Y+F)\sim(2a-2)b^*Y.
\end{equation}

\emph{The polarization and Picard pullback.}
The invariant-ring polarization also works with an internal center.
Let $R_X$ be the homogeneous coordinate ring of the given embedding.
Its invariant ring is a graded quotient of
\[
 k[z,x_1,\ldots,x_r],\qquad z=x_0^2,
 \qquad \deg z=2,\quad\deg x_i=1.
\]
Indeed the homogeneous quotient map is equivariant, and its
restriction to each character summand is surjective. The second
Veronese of this weighted polynomial ring is generated in degree
one: each monomial of weighted degree $2n$ is a product of $n$
monomials of weighted degree two. The same holds for its quotient.
Thus $S=\operatorname{Proj}(R_X^{\mu_2})$ has a very ample line
bundle $J=\OO_S(2)$.

There is an isomorphism
\[
 R_X^{\mu_2}/(x_0^2)\simeq R_X/(x_0).
\]
Surjectivity follows from the invariant generators $x_i$ with
$i>0$. If an invariant element is $x_0g$, then $g$ has weight
one, since $x_0$ is a nonzerodivisor. A weight-one polynomial
lifting $g$ is divisible by $x_0$ with invariant quotient. This
proves the kernel assertion. The section $z$ therefore cuts out
a smooth Cartier curve $Y_S\simeq Y$, and
\begin{equation}\label{eq:internal-cartier-polarization}
 J=\OO_S(Y_S),\qquad q^*Y_S=2Y,\qquad Y_S^2=2d.
\end{equation}
The point $o$ is outside $Y_S$. In particular this divisor is
Cartier at $o$, where its section is a unit. Projective normality
of $X$ is not required.

We will use injectivity of Picard pullback for both $q$ and $q'$.
Here is the relevant fixed-point argument, which does not require
the quotient morphism to be flat. Suppose a projective integral
$k$-scheme $W$ with $H^0(W,\OO_W)=k$ has a finite $\mu_p$-quotient
$f:W\to Q$ and a scheme-theoretically fixed $k$-point. If $f^*L$
is trivial, its canonical linearization acts on its one-dimensional
space of global sections by a character. A nonzero section is
nowhere zero. Its value at the fixed point is nonzero, and the
canonical pullback action on that fibre is trivial. The character
is therefore trivial. This invariant section descends, because
\[
 (f_*f^*L)^{\mu_p}
 =L\otimes(f_*\OO_W)^{\mu_p}=L.
\]
Its descent is nowhere zero since $f$ is surjective. Hence $L$
is trivial, proving that $f^*$ on Picard groups is injective.

Every point of $Y$ is scheme-theoretically fixed, and the blowup
does not change a neighborhood of $Y$. The argument applies to
$q$ and $q'$. Put
\[
 D=\rho^*Y_S,\qquad J_T=\OO_T(D)=\rho^*J.
\]
The curve $D$ is smooth, isomorphic to $Y$, and disjoint from $E$.
Equations~\eqref{eq:internal-cartier-canonical-pullback} and
\eqref{eq:internal-cartier-polarization}, together with injectivity
for $q'$, give
\begin{equation}\label{eq:internal-cartier-canonical-class}
 D^2=2d,\qquad q'^*D=2b^*Y,
 \qquad K_T\sim(a-1)D.
\end{equation}
Thus $J_T$ is globally generated, nef and big. We use the pulled-back
space
\[
 V_0=H^0(S,J)\subset H^0(T,J_T)
\]
throughout. It defines the morphism
\[
 f:T\xrightarrow{\rho}S\hookrightarrow\PP(V_0^*).
\]
The only integral curve of intersection zero with
$D$ is $E$: every other curve maps onto a curve on $S$, and its
intersection is positive by the projection formula and ampleness
of $J$.

\emph{The descended logarithmic form.}
Suppose, for a contradiction, that an invertible $P$ on $S$
satisfies $q^*P\simeq M$. Injectivity of $q^*$ gives
\[
 P\ne\OO_S,\qquad P^2\simeq\OO_S,
 \qquad P|_{Y_S}\simeq\OO_{Y_S}.
\]
Its square trivialization gives a global K\"ahler form
$\alpha_S=d\log P$ on $S$, defined by logarithmic differentials
of local units. This definition is valid on the singular surface.
After pullback its square trivialization differs from the chosen
one for $M$ by a scalar, whose differential is zero. Therefore
\begin{equation}\label{eq:internal-cartier-form-pullback}
 q^*\alpha_S=\omega.
\end{equation}
The form $\alpha=\rho^*\alpha_S$ is regular and closed on $T$,
and is nonzero since $q'^*\alpha=b^*\omega$.

Its exact order on $D$ is $a$. To check this, its tangential
restriction is zero because $P|_{Y_S}$ is trivial. The possible
remaining conormal section lies in $\OO_D(-D)$, of negative
degree, and is also zero. Let $n\geq1$ be the exact order. The
leading tangential coefficient is nonzero: a leading section
with zero tangential image would lie in
\[
 H^0(D,\OO_D(-(n+1)D))=0.
\]
At a general closed point of $Y$, the quotient coordinates are
$s=t^2,z=z$. Pullback kills $ds$, while the nonzero leading
tangential coefficient has exact order $2n$. Equation
\eqref{eq:internal-cartier-form-pullback} gives $2n=2a$, hence
$n=a$. This uses completed local rings at a closed point and
makes no perfectness assumption on $k(Y)$.

Let $s_D$ be the canonical section of $J_T$. The form
\[
 \beta=\alpha/s_D^a\in H^0(T,\Omega_T^1(-aD))
\]
is defined everywhere on $T$. Along $D$, divisibility of each
coefficient at its generic point implies divisibility by that
power of its prime local equation in every regular local ring.
Its tangential restriction is nowhere zero on $D$: it is a nonzero
section of a line bundle
of degree
\[
 2g(D)-2-aD^2=0.
\]
Closedness also forces $2\mid a$, by the leading tangential
coefficient calculation. The argument below requires no further
numerical restriction on $a$.

There is no divisorial zero of $\beta$ outside $E$. There is
none on $D$ by the preceding restriction. A zero divisor outside
$D\cup E$ would pull back under the finite map $q'$ to a zero
divisor of $b^*\omega$ away from $b^*Y\cup F$. Such a divisor
would give a divisorial zero of $\omega$ outside $Y$, contrary
to the structural input. The next calculation treats $E$.

\emph{The exceptional multiplicity.}
Choose a local frame of $P$ at $o$. In the completed invariant
ring~\eqref{eq:internal-cartier-invariant-ring}, its square
trivialization is a unit
\begin{equation}\label{eq:internal-cartier-unit-jet}
 U=c+\lambda A+\mu B+\nu C+U_2,
 \qquad c\ne0,\quad U_2\in(A,B,C)^2.
\end{equation}
We have $\mu\ne0$. Indeed, after pulling back to $k[[u,w]]$,
$dA=dC=0$ and
\[
 dB=u\,dw+w\,du.
\]
The radial annihilator and its reduced zero at $v$ make the
first jet of $\omega$ a nonzero scalar multiple of this form.
Indeed the radial derivation has coefficient pair $(u,w)$.
Its annihilating forms are multiples of $u\,dw+w\,du$ by the
Koszul relation, and the reduced zero forces the multiplier to
be a unit.
Terms in $(A,B,C)^2$ have pullback order at least four, and their
differentials cannot contribute to that first jet. Equation
\eqref{eq:internal-cartier-form-pullback} therefore forces
$\mu/c\ne0$.

On the resolution chart $B=Az$, $C=Az^2$, put
\[
 P_2(z)=\lambda+\mu z+\nu z^2.
\]
This polynomial is nonzero, and its derivative is $\mu$. Thus
\begin{equation}\label{eq:internal-cartier-exceptional-form}
 \alpha=
 \frac{(P_2(z)+O(A))\,dA+(A\mu+O(A^2))\,dz}{U}.
\end{equation}
The expansion can be taken in the completed discrete valuation
ring at the generic point of $E$. Substitution sends
$(A,B,C)^2$ into $(A^2)$, so its $z$-derivatives also lie in
$(A^2)$. Completion is faithfully flat, and the displayed orders
are the actual divisorial orders. In particular $\beta$, which
differs from $\alpha$ by a unit near $E$, is nonzero at the
generic point of $E$. It has no divisorial zero anywhere.

Its saturated annihilator is consequently an invertible subsheaf
\begin{equation}\label{eq:internal-cartier-annihilator}
 \mathcal F=\OO_T(aD)\otimes\omega_T^{-1}
 \simeq J_T\longrightarrow T_T.
\end{equation}
This is the Koszul relation between the two local coefficients
of $\beta$, which have no common prime divisor. Near the generic
point of $E$, a generator is a unit multiple of
\begin{equation}\label{eq:internal-cartier-local-field}
 (A\mu+O(A^2))\partial_A
 -(P_2(z)+O(A))\partial_z.
\end{equation}
It is nonzero and tangent to $E$ at its generic point, and its
value on $A$ has order exactly one along $E$.

The composite
\[
 \mathcal F\longrightarrow T_T
 \longrightarrow f^*T_{\PP(V_0^*)}
\]
vanishes along $E$, because $f$ contracts $E$. Its vanishing
order there is exactly one. Indeed $A=u^2$ is the ratio of the
global sections $x_i^2$ and $x_0^2$ of $J$, pulled back to $T$.
Its derivative, computed in
\eqref{eq:internal-cartier-local-field}, has order exactly one.
Every other affine embedding coordinate has derivative divisible
by $A$, since the field is tangent to the contracted curve.
This proves the exact multiplicity assertion for the ambient
tangent map, with its saturated source line bundle.

\emph{An Euler extension on the resolution.}
Pull back the restricted Euler sequence along that composite.
It gives an extension
\[
 0\longrightarrow\OO_T\longrightarrow\mathcal P
 \longrightarrow\mathcal F\longrightarrow0
\]
and a map $\mathcal P\to V_0^*\otimes J_T$.
The middle term is locally free, including along $E$, because
the end terms are line bundles. This does not require the map
to the ambient tangent bundle to be a subbundle at $E$.
Set $\mathcal E=\mathcal P^*\otimes J_T$. By
\eqref{eq:internal-cartier-annihilator} there are maps
\[
 0\longrightarrow\OO_T\longrightarrow\mathcal E
 \xrightarrow{\mathrm{ev}}J_T\longrightarrow0,
 \qquad V_0\otimes\OO_T\longrightarrow\mathcal E,
\]
whose composite is evaluation of sections of $J_T$.
Write $e_s$ for the section induced by $s\in V_0$ and
$e_0=e_{s_D}$. Locally its two components record value and
differentiation in the $\mathcal F$ direction.

The section $e_0$ is nowhere zero. Away from $D$ its value
component is nonzero. On $D$ the normal projection of
$\mathcal F$ is an isomorphism onto $\OO_D(D)$, since the
tangential restriction of $\beta$ is nowhere zero. Its
derivative component is therefore nonzero there. It follows that
\begin{equation}\label{eq:internal-cartier-determinant-extension}
 0\longrightarrow\OO_T\xrightarrow{e_0}\mathcal E
 \xrightarrow{e_0\wedge-}J_T\longrightarrow0
\end{equation}
is exact. The determinant sections $t_s=e_0\wedge e_s$ restrict
on $D$, after one nowhere-zero common factor, to $s|_D$.
On the chart $s_D\ne0$, their orders are those of the derivatives
of $s/s_D$ in the $\mathcal F$ direction.

Choose $s$ with two properties: $s|_D$ has distinct simple zeros,
and $t_s$ has multiplicity exactly one along $E$.
The first is a nonempty open condition, since $J$ embeds the
smooth curve $Y_S$. Explicitly, hyperplanes containing a tangent
line form the image of an incidence variety of smaller dimension
than the dual projective space. The second is a nonempty open
linear condition: all $t_s$ vanish along $E$, and the coordinate
$A$ just exhibited has derivative not vanishing modulo $A^2$.
Both conditions can therefore be satisfied over the infinite
field $k$.

The divisor $C_s=\operatorname{div}(t_s)$ is nonempty and does
not contain $D$. Every component other than $E$ meets $D$, and
the simple zeros of $t_s|_D$ force each such component to have
multiplicity one. The multiplicity of $E$ is one by construction.
A Cartier divisor on the smooth surface $T$ has no embedded
associated components, so $C_s$ is reduced.

It is also connected, although $J_T$ is not ample. If
$C_s=C_1+C_2$ were a disjoint decomposition into nonzero effective
divisors, then
\[
 C_i^2=J_T\cdot C_i\geq0.
\]
A zero square would make $C_i$ orthogonal to the positive-square
class $J_T$. Hodge index would then make $C_i$ numerically zero,
impossible for a nonzero effective divisor because it has positive
intersection with an ample class on $T$. Thus both squares would
be positive; but two positive-square classes cannot be orthogonal
by Hodge index. This contradiction proves connectedness.

The proper reduced connected curve $C_s$ has
$H^0(C_s,\OO_{C_s})=k$. Restricting
\eqref{eq:internal-cartier-determinant-extension} to $C_s$ gives
\[
 e_s|_{C_s}=\gamma e_0|_{C_s}
\]
for a scalar $\gamma\in k$. Dividing $e_s-\gamma e_0$ by
$t_s$ gives a homomorphism $J_T\to\mathcal E$ whose composite
with $e_0\wedge-$ is the identity. Consequently
\[
 \mathcal E\simeq\OO_T\oplus J_T.
\]
For the original evaluation quotient
$\mathrm{ev}:\mathcal E\to J_T$, the component from the first
summand is $s_D$, and the component from $J_T$ is a scalar.
That scalar is nonzero, since otherwise the quotient would
vanish on $D$. Thus the original evaluation quotient splits
as well.

Dualizing and twisting now splits the original Euler pullback
compatibly. The map
\[
 \mathcal F\simeq J_T\longrightarrow f^*T_{\PP(V_0^*)}
\]
lifts to $V_0^*\otimes J_T$. Since $H^0(T,\OO_T)=k$, this
lift is a constant vector $v_1\in V_0^*$. It is nonzero:
away from $E$, the map $f$ is an immersion and the foliation
is generically nonzero.

\emph{The contradiction at the contracted curve.}
The ambient tangent map vanishes along $E$, and $f(E)=o$.
A constant Euler vector has zero projective tangent value at
a point precisely when its projective point is that point.
Therefore $[v_1]=o$. This point is outside the hyperplane
$Y_S$, so the $s_D$-coordinate of $v_1$ is nonzero.

On the affine chart $s_D\ne0$, the induced projective field is
a nonzero scalar multiple of the radial field centered at $o$.
In particular, the affine section ratios $A,B,C$, all
zero at $o$, satisfy
\[
 A\longmapsto c_1A,\qquad
 B\longmapsto c_1B,\qquad
 C\longmapsto c_1C,
 \qquad c_1\ne0.
\]
On the resolution chart $B=Az$, its lift on $A\ne0$ thus sends
$z=B/A$ to zero and $A$ to $c_1A$. These identities hold on the
whole chart by regularity. This lift vanishes along $E$.

On the other hand, the generator of the source line
$\mathcal F$ is, up to a unit, the field in
\eqref{eq:internal-cartier-local-field}. Its $\partial_z$
coefficient is nonzero at the generic point of $E$, since
$P_2(z)\ne0$. The isomorphism $\mathcal F\simeq J_T$ changes
local generators only by units and cannot remove this discrepancy.
The contradiction proves that no invertible $P$ on $S$ pulls
back to $M$.

\emph{Tensor squares and the N\'eron--Severi class.}
Suppose $M\simeq L^{\otimes2}$ on $X$. On invariant affine
charts one has $R^2\subset R^{\mu_2}\subset R$, so relative
Frobenius factors as
\[
 X\xrightarrow{q}S\xrightarrow{g}X^{(2)},
 \qquad gq=F_{X/k}.
\]
The invertible sheaf $P=g^*(L^{(2)})$ is defined also at $o$,
and
\[
 q^*P=F_{X/k}^*L^{(2)}\simeq L^{\otimes2}\simeq M,
\]
contrary to what was just proved. Hence $M$ is not a tensor square.

Finally, if $[M]=2[L_0]$ in $\operatorname{NS}(X)$, then
$M\otimes L_0^{-2}$ is algebraically trivial. The reduced identity
component of the Picard scheme is an abelian variety, and
multiplication by two is surjective on its $k$-points
\cite[Theorems~2.5 and~5.4, Proposition~5.10 and Remark~5.25]{KleimanPicard}
\cite[Tag~03RP, parts~(3)--(4)]{Stacks}.
Choose a square root $L_1$ of that algebraically trivial bundle.
Then $M\simeq(L_0\otimes L_1)^{\otimes2}$, again a contradiction.
This proves the N\'eron--Severi assertion.
\end{proof}

For an internal center, this excludes every double in the
N\'eron--Severi group, independently of $v_2(2a)$. When that
valuation is greater than one, it strengthens the finite-height
obstruction of Corollary~\ref{thm:positive-square-picard}.
The proposition concerns descent as an invertible sheaf on the
singular quotient. A linearization can have different stabilizer
characters on the fixed divisor and at the isolated fixed point;
the proposition does not identify every linearization with such
a descent. Nondivisible N\'eron--Severi classes and positive genus
excess remain outside its conclusion.


\section{Uniform cohomology for tame coefficient systems}\label{sec:tame-coefficients}

We prove the coefficient estimate by bounding poles on a resolution. The problem is that tame inertia can have arbitrarily large order. Passing to a root stack makes the sheaf unramified; a character count will remove the dependence on the chosen root order.

Keep the notation of Theorem~\ref{thm:local-coefficients}. Write $B_{ij}=C_i\cdot C_j$, $g_i=g(C_i)$, and let $\nu_i$ be the number of crossing points on $C_i$. The matrix $B$ is negative definite \cite[Lemma 14.1, p.~224]{Lipman} and has nonnegative off-diagonal entries. Choose positive integers $c_i$ such that
\[
 Z=\sum_i c_iC_i,\qquad J=\OO_X(-Z),\qquad
 d_i=\deg(J|_{C_i})=-\sum_jc_jB_{ij}>0.
\]
Such integers exist because $(-B)^{-1}$ is entrywise nonnegative, as proved in Section~\ref{sec:abelian-inertia}; its product with the all-ones vector is strictly positive. Clear denominators. Put
\begin{equation}\label{eq:resolution-constants}
 L_i=\sum_jc_j|B_{ij}|,\qquad
 K_i=\max\{0,-B_{ii},2g_i-2+\nu_i\},\qquad
 b_i=\frac{2(p-1)K_i}{p}.
\end{equation}
We will use an integer $m_0\geq2$ larger than every $(b_i+L_i)/d_i$.

\subsection{Removing tame inertia}

Choose $N$ prime to $p$, divisible by all inertia orders of the monodromy of $M$, and form the stack $\mathfrak X$ obtained by taking $N$th roots of each $C_i$. Write $\pi:\mathfrak X\to X$ and $\mathfrak C_i$ for the reduced divisors above $C_i$. At a crossing with parameters $x,y$, its chart is
\[
 \left[\Spec R[u,v]/(u^N-x,v^N-y)\,/\,\mu_N^2\right].
\]
Away from crossings there is one root variable. These regular charts show that $\mathfrak X$ is a regular tame Deligne--Mumford stack, proper over $A$, with coarse space $X$ and $N\mathfrak C_i=\pi^*C_i$. This is the usual root construction \cite[\S\S2.3--2.4, author pp.~7--10]{Cadman}.

The sheaf $M$ extends to a locally constant sheaf on $\mathfrak X$. The root charts kill its tame codimension-one inertia by Abhyankar's lemma \cite[Expos\'e XIII, Proposition 5.2, original pp.~428--429]{SGA1}. Normalize each regular chart in the corresponding finite cover. Purity of the branch locus makes the normalization \'etale, and uniqueness of normalization gives the descent data on overlaps \cite[Tags 0BRM and 0BMB]{Stacks}. This argument places no prime-to-$p$ restriction on the entire monodromy group.

Set $V=M\otimes_{\F_p}\OO_{\mathfrak X}$ and write $r=\operatorname{rank}M$. It is a vector bundle with $F^*V\simeq V$ for absolute Frobenius, and is trivialized by a finite \'etale cover. The divisor $\mathfrak C_i$ is a $\mu_N$-gerbe over the curve obtained from $C_i$ by taking roots at its crossing points. Normalize degree on it so that pullback from $C_i$ preserves ordinary degree. A line bundle has its $N$th power descending to $C_i$; divide that degree by $N$. In particular
\begin{equation}\label{eq:root-degrees}
 \deg\bigl(\OO(\mathfrak C_j)|_{\mathfrak C_i}\bigr)=B_{ij}/N,
 \qquad \deg(J|_{\mathfrak C_i})=d_i.
\end{equation}
The degree of a vector bundle means the degree of its determinant. Frobenius multiplies it by $p$.

The restriction of $V$ to each component is semistable of degree zero. Indeed, on a finite \'etale trivializing cover the determinant of a subbundle maps into an exterior power of a trivial bundle, so has nonpositive degree. Degrees under a finite representable cover scale by a positive factor; their signs are preserved. The determinant of $V$ becomes trivial there, so has degree zero. This proof works even if the degree of the cover is divisible by $p$. The related finite-bundle statement and the root-stack interpretation of parabolic degree are treated by Borne \cite[Theorems 4--5 and Proposition 6, author pp.~12,17,22]{Borne}.

We record a cohomological consequence of this degree convention. Decompose a bundle on $\mathfrak C_i$ into characters of its generic $\mu_N$. Nontrivial characters have zero coherent cohomology, since taking invariants along that gerbe is exact. A character-zero summand $W$ descends to the rooted curve. Its coarse pushforward $E$ is a vector bundle of the same rank on $C_i$, and
\begin{equation}\label{eq:coarse-degree}
 0\leq\deg W-\deg E\leq\nu_i\operatorname{rank}W.
\end{equation}
To verify this locally, lift a character basis on a root chart by equivariant Nakayama. The adjunction map from the pullback of the invariant module has diagonal entries $z^{a_j}$, with $0\leq a_j<N$. The determinant loses $\sum_j a_j/N$ in degree. Summing at the crossing points proves \eqref{eq:coarse-degree}.

If $W$ is semistable of negative slope $t$, it has no sections, so Riemann--Roch gives
\begin{equation}\label{eq:curve-bound}
 h^1(\mathfrak C_i,W)\leq
 \operatorname{rank}W\,(g_i-1-t+\nu_i).
\end{equation}
In slope zero, descend the character-zero summand $W$ to the rooted
curve and choose a point away from its crossing points. Evaluation
there is injective: a section vanishing there would yield, after
saturation, a positive-degree line subbundle on that rooted curve. Thus in that case
\begin{equation}\label{eq:curve-zero}
 h^1(\mathfrak C_i,W)\leq
 \operatorname{rank}W\,(g_i+\nu_i).
\end{equation}
Every isotypic summand of $V$ is semistable of degree zero. Each is a direct summand and has nonpositive degree, and their degrees sum to zero. The preceding estimates therefore apply after twisting such a summand by any line bundle which cancels its generic character.

\subsection{Frobenius and a common pole cutoff}

Put $\mathcal B=F_*\OO_{\mathfrak X}/\OO_{\mathfrak X}$. This is locally free. A coefficient field makes $A$ a quotient of a power-series ring over perfect $k$, so $A$ and the finite-type chart rings are $F$-finite. Frobenius on the root stack is representable because it induces the automorphism $\zeta\mapsto\zeta^p$ on its stabilizers. It is finite flat on the regular charts, and its unit is a unimodular vector locally in $F_*\OO$.

The pullback $F^*\mathcal B$ identifies with the augmentation ideal in $F^*F_*\OO$. Its powers have locally free graded pieces the truncated symmetric powers
\[
 T^a(\Omega^1_{\mathfrak X/k}),\qquad 1\leq a\leq2(p-1).
\]
This is Sun's Frobenius filtration \cite[Theorem 3.7, author pp.~12--13]{Sun}. Its applicability here can be checked without a finite-type-over-$k$ assumption. At every closed point of the exceptional divisor, the completed regular chart is $k[[u,v]]$, and
\[
 k[[u,v]]\otimes_{k[[u^p,v^p]]}k[[u,v]]
       =k[[u,v]][\alpha,\beta]/(\alpha^p,\beta^p).
\]
The augmentation ideal is $(\alpha,\beta)$; its graded generators are $du,dv$. This parameter-independent calculation verifies the asserted finite-module filtration after faithfully flat completion, and hence near the exceptional divisor. A nonzero coherent error there would have a closed point in its support.

The conormal sequence on $\mathfrak C_i$ has line terms of degrees
\[
 -B_{ii}/N,\qquad 2g_i-2+\nu_i(1-1/N).
\]
For the second degree use $dx=Nz^{N-1}dz$ at each root point; the generic gerbe is \'etale and adds no differential direction. Each degree is at most $K_i$. The truncated symmetric powers have a filtration by monomials in this two-line flag, with both exponents at most $p-1$. Every line layer therefore has degree at most $2(p-1)K_i$. Tensoring with $V$ and using $F^*V\simeq V$ bounds the slope of every subbundle of $F^*(\mathcal B\otimes V)|_{\mathfrak C_i}$ by that same number. Pull back a subbundle by Frobenius and divide degrees by $p$. We obtain
\begin{equation}\label{eq:frobenius-slope}
 \mu(\text{any subbundle of }(\mathcal B\otimes V)|_{\mathfrak C_i})\leq b_i.
\end{equation}
The bound is independent of $N$ and $r$.

Let $\mathfrak D=\bigcup_i\mathfrak C_i$ be reduced, let $i:\mathfrak D\to\mathfrak X$, and let $j:U\to\mathfrak X$. Write $R^1\Psi M=i^*R^1j_*M$. Put $\mathfrak Z=\pi^*Z=\sum_iNc_i\mathfrak C_i$, whose ideal is again denoted $J$. Sheaves supported on a nilpotent thickening of $\mathfrak D$ are
viewed on its small \'etale site, which is invariant under such
thickenings. The inverse image $i^*$ is that of \'etale sheaves;
it preserves the infinitesimal layers of the pole quotients.
Tensor restriction of coherent modules to the reduced divisor
would discard those layers. The Artin--Schreier sequence gives
a pole filtration $U_m$ of $R^1\Psi M$, with
\begin{equation}\label{eq:pole-sequence}
 0\longrightarrow J^{-m}V/V\xrightarrow{F_V-1}J^{-pm}V/V
       \longrightarrow U_m\longrightarrow0.
\end{equation}
Here $F_V$ is the Frobenius-semilinear map given by the unit structure. We justify the exactness, including its bound on preimages. First $j_*M=M$: on a regular strict-local root chart the complement of the boundary is connected, and the extended local system is constant. On an affine chart $U$ is the complement of a principal divisor, so higher direct images of $V$ as a quasi-coherent sheaf vanish. Artin--Schreier computes $R^1\Psi M$ as the quotient of $j_*V/V$ by $F_V-1$. A positive pole of a vector, measured by its smallest component valuation in an \'etale trivialization, is multiplied by $p$ under $F_V$ and cannot cancel with its smaller pole under the identity. Thus a preimage of a vector with poles bounded by $pmZ$ has poles bounded by $mZ$. This proves \eqref{eq:pole-sequence}; it also proves that the increasing union is all of $R^1\Psi M$.

For $m\geq2$, the identity term lies in the preceding pole level because $m\leq p(m-1)$. Therefore the graded quotient is the cokernel of Frobenius between
\[
 J^{-m}V/J^{-m+1}V
       \quad\text{and}\quad J^{-pm}V/J^{-pm+p}V.
\]
Give the target its Frobenius scalar action, $a\cdot v=a^pv$. This is the coherent structure used for the following slope estimates; the underlying additive cokernel is unchanged. Since $F^{-1}\mathfrak Z=p\mathfrak Z$, projection formula and the locally free cokernel $\mathcal B$ identify it as
\begin{equation}\label{eq:pole-graded}
 U_m/U_{m-1}=
 (\mathcal B\otimes V)|_{\mathfrak Z}\otimes J^{-m}|_{\mathfrak Z}.
\end{equation}
This also follows by restricting $0\to\OO\to F_*\OO\to\mathcal B\to0$, which remains exact because $\mathcal B$ is locally free. The calculation is compatible on charts, so descends. It is the twisted root-stack form of Suzuki's calculation \cite[Proposition 3.1, author pp.~6--7]{Suzuki}.

Filter $\OO_{\mathfrak Z}$ by removing one component at a time. A line layer on $\mathfrak C_i$ has degree at most $L_i$: all coefficients of the removed cycle lie between zero and $Nc_j$, and \eqref{eq:root-degrees} applies. The corresponding layer of \eqref{eq:pole-graded} has maximal slope at most $b_i+L_i-md_i$, by \eqref{eq:frobenius-slope}. This is negative for $m\geq m_0$, so
\[
 H^0(\mathfrak D,U_m/U_{m-1})=0\qquad(m\geq m_0).
\]
Global sections commute with this filtered union on the quasi-compact exceptional stack; this can be checked on a finite affine atlas and its overlaps. It follows that
\begin{equation}\label{eq:cutoff}
 H^0(\mathfrak D,R^1\Psi M)=H^0(\mathfrak D,U_{m_0}).
\end{equation}

\subsection{Counting the low-pole contribution}

We now bound the fixed pole level without introducing a factor of $N$. We use the negative-cycle filtration underlying Suzuki's Proposition 2.1 \cite[author pp.~4--5]{Suzuki}. Its elementary argument applies to the normalized matrix $B/N$: if an effective nonzero cycle $E$ is supported on these curves, negative definiteness gives a component of its support with $E\cdot\mathfrak C_i<0$. Removing that component gives a negative-degree line layer of $\OO(E)$ modulo the preceding cycle. More generally, when the initial cycle has nonpositive intersection with every component, the same argument for the difference of cycles gives negative layers throughout. By induction this supplies filtrations of
\[
 J^{-pm}/\OO,\qquad J^{-pm}/J^{-m},\qquad J^{-m}/\OO
\]
by negative-degree line bundles on the components. Tensoring with the degree-zero semistable $V$ gives no global sections of these sheaves.

Consequently \eqref{eq:pole-sequence} injects $H^0(U_m)$ into the kernel of a Frobenius-semilinear map minus an injective linear map from $H^1(J^{-m}V/V)$ to $H^1(J^{-pm}V/V)$. If $f:W\to W'$ is Frobenius-semilinear and $g:W\to W'$ is injective and linear, this kernel has $\F_p$-dimension at most $\dim_kW$. Split $g$ linearly and reduce to fixed vectors of a semilinear endomorphism of $W$. Any minimal $k$-linear relation between fixed vectors, normalized to have one coefficient equal to one, can be subtracted from its Frobenius transform. Minimality forces every coefficient to lie in $\F_p$. Thus $\F_p$-independent fixed vectors are $k$-independent.

In a negative filtration of $J^{-m}/\OO$, a line layer $P$ on $\mathfrak C_i$ satisfies
\[
 -mL_i\leq\deg P<0.
\]
As the coefficient of $\mathfrak C_i$ runs through $mNc_i$ removals, the generic character of $P$ runs through every residue modulo $N$ exactly $mc_i$ times. Twists from other components have trivial character at its generic point. Decompose $V|_{\mathfrak C_i}=\bigoplus_\chi V_{i,\chi}$. Only the summand whose character cancels that of $P$ contributes to cohomology. The sum of its ranks over all these layers is exactly
\[
 mc_i\sum_\chi\operatorname{rank}V_{i,\chi}=mc_i r.
\]
This identity cancels the growing number of layers. Applying \eqref{eq:curve-bound} to each contributing summand, then using the long exact sequences of the filtration, gives
\begin{equation}\label{eq:low-pole-bound}
 h^1(J^{-m}V/V)\leq r m\sum_i c_i(g_i+\nu_i+mL_i+1).
\end{equation}
In particular \eqref{eq:cutoff} and this estimate with $m=m_0$ give a bound linear in $r$ for the nearby-cycle contribution.

The normalization sequence of the reduced exceptional stack has one node term for each crossing, with stabilizer $\mu_N^2$. Its invariant fibre has dimension at most $r$, and averaging kills its higher cohomology. On each normalization component only the generic-character-zero summand contributes. Equation~\eqref{eq:curve-zero} therefore gives
\begin{equation}\label{eq:exceptional-bound}
 h^1(\mathfrak D,V)\leq
 r\left(\sum_i(g_i+\nu_i)+\#\{\text{crossings}\}\right).
\end{equation}
On the finite-dimensional space $H^0(\mathfrak D,V)$, the map $F_V-1$ is surjective. Indeed its differential is minus the identity, so it is an \'etale homomorphism of additive algebraic groups with open image; an open subgroup of the connected additive group is the whole group. Artin--Schreier and the fixed-vector argument above bound $\dim_{\F_p}H^1(\mathfrak D,M)$ by \eqref{eq:exceptional-bound}.

For clarity, the nearby-cycle comparison used here follows from ordinary scheme proper base change, which applies to all torsion coefficients, including $\F_p$ \cite[Theorem 59.91.11, Tag 095S]{Stacks}. Coarse pushforward for the root stack is exact on coherent sheaves and on $\F_p$-sheaves. Locally it is invariants under a finite group of order prime to $p$, and averaging is exact. For the \'etale assertion, a strict henselian coarse neighborhood has a finite root cover which is a union of strictly henselian local schemes; their global-section functors are exact, and finite-group cohomology vanishes by averaging. This description commutes with passage to the closed fibre. Thus proper base change for $X$ over the strictly henselian $\Spec A$ gives the degree-one exact sequence
\[
 0\longrightarrow H^1(\mathfrak D,M)\longrightarrow H^1(U,M)
      \longrightarrow H^0(\mathfrak D,R^1\Psi M).
\]
Combining \eqref{eq:low-pole-bound} and \eqref{eq:exceptional-bound} proves Theorem~\ref{thm:local-coefficients}, for example with
\[
 C_X=m_0\sum_i c_i(g_i+\nu_i+m_0L_i+1)
          +\sum_i(g_i+\nu_i)+\#\{\text{crossings}\}.
\]

\section{Bounding conductors and abelian inertia}\label{sec:abelian-inertia}

We prove Theorem~\ref{thm:abelian-inertia} by turning the refined Swan conductor into an intersection inequality. Only the local integrality theorem is needed, in the published form of Kato's Theorem 7.1 restated by Bright and Newton \cite[Theorem 4.1, \S4.3, pp.~848--849]{BN}.

Fix a character $\chi\in H^1(U,\mathbf Q_p/\mathbf Z_p)$ and let $S=\sum_i s_i C_i$ be its logarithmic Swan divisor. At a closed point of the exceptional divisor the local ring of $X$ is excellent regular, its residue field is $k$, and its boundary equations form part of a regular parameter system. These are the hypotheses of the cited theorem: excellence and finite $p$-degree of the residue field suffice. They hold because $A$ is complete noetherian and $X$ is of finite type over $A$; $k$ is perfect. No finite-type-over-a-field globalization is required.

If $s_i>0$, the intrinsic generic refined conductor is nonzero. Kato's local integrality assertion at every closed point of $C_i$ extends it to a nonzero section
\begin{equation}\label{eq:rsw-section}
 0\ne\operatorname{rsw}_i(\chi)\in
 H^0\left(C_i,\Omega_X^1(\log D)|_{C_i}\otimes\OO_X(S)|_{C_i}\right).
\end{equation}
Indeed the local assertion bounds its poles along every other boundary equation by the corresponding Swan order. These are exactly the local lattices of the displayed bundle, and they agree on overlaps because the generic differential is intrinsic. The logarithmic differential bundle is locally free near $D$: this follows directly after completion to the regular two-variable power-series ring over perfect $k$.

The residue sequence on $C_i$ has the form
\[
 0\longrightarrow\Omega^1_{C_i/k}(\log\{\text{crossings}\})
 \longrightarrow\Omega_X^1(\log D)|_{C_i}
 \longrightarrow\OO_{C_i}\longrightarrow0.
\]
Its line terms have degrees $2g_i-2+\nu_i$ and zero. Hence every line subbundle of its middle term has degree at most $k_i=\max(0,2g_i-2+\nu_i)$. Saturating the map $\OO_{C_i}(-S)\to\Omega_X^1(\log D)|_{C_i}$ from \eqref{eq:rsw-section} gives $S\cdot C_i\geq-k_i$. If $s_i=0$, the same inequality holds because all off-diagonal intersections and all $s_j$ are nonnegative. In vector notation,
\[
 Bs\geq-k.
\]

For the matrix step put $C=-B$. It is positive definite with nonpositive off-diagonal entries, and $C^{-1}$ is entrywise nonnegative. To prove this, take $Cx=v\geq0$ and write $x=x_+-x_-$ with disjointly supported nonnegative parts. If $x_-\ne0$, then
\[
 0\leq x_-^Tv=x_-^TCx_+-x_-^TCx_-<0,
\]
a contradiction. Thus $x\geq0$. Applying the inverse to $Cs\leq k$ proves $s\leq(-B)^{-1}k=t$.

Choose $e$ with $p^e>\max_i t_i$. Multiplication of a character by $p$ divides its Swan bound by $p$, rounding down. Here is the characteristic-$p$ calculation. The conductor filtration is the image, under Artin--Schreier--Witt, of the Witt vectors $a=(a_{m-1},\ldots,a_0)$ satisfying $p^jv(a_j)\geq-n$ for all $j$ \cite[\S\S3.2.2--3.2.3, pp.~152--153]{KLS}. Witt restriction drops the last displayed coordinate, so every surviving weight is divided by $p$. It represents multiplication by $p$ under the inclusions into $\mathbf Q_p/\mathbf Z_p$. Therefore
\[
 s_i(p^e\chi)\leq\lfloor s_i(\chi)/p^e\rfloor=0.
\]
Conductor zero means that an integral Witt vector represents the class; its Artin--Schreier--Witt equation is \'etale over the discrete valuation ring. Thus $p^e\chi$ is unramified along $D$ and extends over $X$ by regular purity.

For a character of order $n$ prime to $p$, use the Kummer description as a line bundle on $U$ with a trivialization of its $n$th power. Extend the line bundle to the regular $X$. Its trivialization has divisor $\sum_i a_iC_i$; intersecting with every $C_j$ gives $Ba\equiv0\pmod n$. Multiplication by the adjugate of $B$ shows that $b=|\det B|$ kills all the residues $a_i$ modulo $n$, which are its tame inertia characters. Purity extends the resulting character to $X$.

The map $Q=\pi_1(U)^{\mathrm{ab}}\to\pi_1(X)^{\mathrm{ab}}$ is surjective: a connected finite \'etale cover of normal $X$ remains connected over its dense open $U$. Let $I$ be its kernel. The two character calculations imply $bp^eI=0$, since finite characters of the profinite abelian group $Q$ separate points. Suzuki's Theorem 1.2 makes $Q$ topologically finitely generated \cite[author p.~2]{Suzuki}. The bounded-exponent subgroup of a finitely generated profinite abelian group is finite. Thus $I$ is finite, proving Theorem~\ref{thm:abelian-inertia}.

\section{The maximal metabelian quotient}\label{sec:metabelian}

We deduce Theorem~\ref{thm:metabelian} from a uniform first-cohomology bound for modules with abelian monodromy. A single finite cover removes the abelian inertia for all such modules. We first give the tame generator bound used for coefficients of characteristic different from $p$.

\begin{lemma}\label{lem:tame-generators}
Let $\Delta$ be the dual graph of $D$. The full tame fundamental group relative to $(X,D)$ is generated by at most
\[
 d_X=\sum_i(2g_i+\nu_i+1)+b_1(\Delta)
\]
elements.
\end{lemma}
\begin{proof}
Every finite tame cover extends to a root stack $\mathfrak X$ of some prime-to-$p$ order $N$, as above. The reduced exceptional stack $\mathfrak D$ surjects on its fundamental group. To see this, a connected finite \'etale cover of $\mathfrak X$ has coarse space finite over $X$: on a root chart its invariant ring is a submodule of a finite module over a noetherian ring. This coarse space is proper over the henselian local base $A$. Its closed fibre is connected by proper henselian connectedness \cite[Tag 0A48]{Stacks}. Taking invariants commutes with reduction by averaging, so that fibre is the coarse space of the closed-fibre cover of $\mathfrak X$. The latter, and hence its reduction over $\mathfrak D$, is connected. This proves the asserted surjectivity.

On a normalization component $\mathfrak C_i$, the generic $\mu_N$ is central in the monodromy of every finite cover: its band acts by natural automorphisms of the fibre functor. Quotienting this cyclic image makes the cover descend to the rooted curve. Restriction to the complement of its crossing points gives a tame cover of the ordinary curve, and a connected cover stays connected there. SGA~1, Expos\'e XIII, Corollaire 2.12 bounds this full tame curve group by $2g_i+\nu_i$ generators \cite[original pp.~392--395; recomposed pp.~290--292]{SGA1}. Adding one generator for generic inertia gives $2g_i+\nu_i+1$ for the component group.

Finite \'etale covers of $\mathfrak D$ glue from its normalization with identifications at the two preimages of each node. On a local chart this is descent of finite $\mu_N^2$-sets on the strictly henselian nodal ring $k[[u,v]]/(uv)$ and its two strictly henselian branches. Choose a spanning tree of $\Delta$. Its identifications give a single fibre; the remaining edges each contribute one gluing element. The images of the component groups and these elements generate every connected finite monodromy group. Otherwise their orbits would decompose all the component covers compatibly with the gluing, contradicting connectedness. This proves the displayed bound for every finite quotient, independently of $N$. Compactness of the spaces of generating tuples proves the same bound for the profinite group.
\end{proof}

\begin{lemma}\label{lem:abelian-coefficients}
There is a constant $C_A$ such that, for every prime $\ell$ and every finite-dimensional $\F_\ell$-local system $M$ on $U$ with abelian monodromy,
\[
 \dim_{\F_\ell}H^1(U,M)\leq C_A\dim_{\F_\ell}M.
\]
\end{lemma}
\begin{proof}
The finite group $I$ of Theorem~\ref{thm:abelian-inertia} is separated from zero by a finite quotient of $Q=\pi_1(U)^{\mathrm{ab}}$. Thus choose an open subgroup $H\subset Q$ with $H\cap I=0$, and let $U'\to U$ be the associated fixed connected finite Galois cover, of group $G$. The integral closure $A'$ of $A$ in its function field is finite, complete, normal and local, with residue field $k$. Choose an SNC resolution $X'$ of $\Spec A'$ dominating the normalization of $X$ in that field. Exceptional inertia on $X'$ maps into $I$, because every \'etale cover of $X$ pulls back unramified there; it also maps into $H$, by the choice of $U'$. Its image in $Q$ is therefore zero.

Every abelian-monodromy $M$ factors through $Q$, so its pullback $M'$ to $U'$ is unramified along the boundary and extends to $X'$ by purity. When $\ell=p$, Theorem~\ref{thm:local-coefficients} on $X'$ bounds its first cohomology by a fixed multiple of its rank. When $\ell\ne p$, its first-cohomology classes factor through the full tame quotient. Indeed a continuous cocycle defines a finite extension of its monodromy by the additive $\F_\ell$-module. On wild inertia the original action is trivial and the cocycle is a homomorphism from a pro-$p$ group to an $\ell$-group, hence zero. Lemma~\ref{lem:tame-generators} now bounds the cocycle space by $d_{X'}\operatorname{rank}M'$: evaluation on generators is injective.

Hochschild--Serre for $U'/U$ injects the kernel of restriction in degree one into $H^1(G,H^0(U',M'))$. This has dimension at most $d(G)\operatorname{rank}M$, again by evaluation on a generating set. Adding this fixed bound to the two bounds on $U'$ proves the assertion, uniformly in $\ell$.
\end{proof}

\begin{proof}[Proof of Theorem~\ref{thm:metabelian}]
Let $d_0$ be a generator bound for $\pi_1(U)^{\mathrm{ab}}$, and choose an integer $D\geq\max\{d_0,C_A+1\}$. We show by induction on the order that every finite metabelian quotient $F$ of $\pi_1(U)$ is generated by $D$ elements. If $F$ is abelian this follows from $d_0$. Otherwise choose a minimal nonzero normal subgroup $M$ inside its abelian derived subgroup. Then $M$ is an elementary abelian $\ell$-group for some prime, simple as an $F$-module under conjugation. Its action factors through $F^{\mathrm{ab}}$.

Choose a generating $D$-tuple in $F/M$, padding with identities if necessary. Every lift generates a subgroup $K$ with $KM=F$. The intersection $K\cap M$ is normal in $F$: it is normalized by $K$, and by $M$ because $M$ is abelian. Simplicity gives $K\cap M=M$ or zero. In the first case $K=F$; in the second it is a complement to $M$. Thus failure of all lifts could occur only if $F$ splits as $M\rtimes H$ with $H=F/M$.

In the split case,
\[
 H^1(F,M)=H^1(H,M)\oplus\operatorname{End}_H(M).
\]
Indeed restricting a cocycle to $M$ gives an $H$-equivariant endomorphism, and each such endomorphism extends by $c(mh)=c(m)$; the remaining cocycle comes from $H$. The second summand is nonzero. Inflation from $F$ to $\pi_1(U)$ is injective, so Lemma~\ref{lem:abelian-coefficients} gives
\[
 \dim_{\F_\ell}H^1(H,M)\leq C_A\dim_{\F_\ell}M-1.
\]
Complements correspond to crossed homomorphisms $H\to M$. Their number is
\[
 |Z^1(H,M)|=|H^1(H,M)|\,|M|/|M^H|
       \leq |M|^{C_A+1}/\ell < |M|^D.
\]
There are $|M|^D$ lifts of the chosen tuple, and each complement contains exactly one such lift. A generating lift therefore exists. This proves the induction. Finally the closed sets of $D$-tuples generating each finite quotient have the finite intersection property in the compact space of $D$-tuples of the maximal metabelian quotient. Their intersection is nonempty and consists of topological generating tuples.
\end{proof}


\section{Further questions}

The construction leaves several precise questions. It would be useful to minimize the degree and the irregularity of a smooth surface with a very ample $L$ and $h^1(L^{-1})>0$, and to give such a surface over an explicit finite field. The independent roots give a large but completely controlled cover; they do not optimize either invariant. For polarization torsors, characteristic two requires a formulation using non-\'etale two-torsion and is outside Theorem~\ref{thm:polarization}.

For schematicity of the full Picard functor, the remaining question concerns proper flat families with geometrically normal or geometrically integral fibres that are not covered by the abelian-torsor corollary. Positive-dimensional and possibly nonflat Picard components obstruct the locally quasi-finite argument, and the torsor proof does not reduce an arbitrary family to an abelian one. Theorem~\ref{thm:surface-torsor-modification} also excludes a projective scheme modification of a nonschematic abelian-surface torsor when its source is flat with geometrically integral fibres. For higher-dimensional torsors the rank-one defect may have positive-dimensional fibres, so the finite splitting-cover argument does not extend directly. Theorem~\ref{thm:torsor-projective-bundle} excludes projective bundles of every positive rank over nonschematic abelian torsors on all bases. The remaining construction questions therefore concern higher-dimensional modifications or families without the torsor and projective-bundle structures used here.

For the local fundamental group, the unresolved extension is a uniform first-cohomology estimate for finite coefficient systems with arbitrary wild monodromy. The tame estimate and the fixed cover removing abelian inertia do not provide this: in further derived layers the coefficient action need no longer factor through the abelian quotient. The arguments here therefore stop at the stated metabelian conclusion.

For Picard restriction on geometric points, the remaining question is whether a smooth very ample curve can have a nontrivial kernel of the kind motivating Corollary~\ref{cor:smooth-torsor}; the weaker ample version is also open here. The tangent kernel in Theorem~\ref{thm:surface} does not supply a nontrivial torsion class. The conductor identity determines the normalized cover and gives a signed adjunction formula. Equality makes that cover Gorenstein and regular near its smooth ample inverse-image curve, and gives a global deformed-group torsor structure. At exact conductor equality with $Y$ very ample,
Theorem~\ref{thm:conductor-equality-structure} determines these
singularities: there are none in odd characteristic, and at most
one $A_1$ point in characteristic two.

Theorem~\ref{thm:strict-picard-all} excludes the smallest very-ample equality in every positive characteristic. Thus a possible nontrivial order-$p$ class must satisfy $2g(Y)-2>pY^2$. The remaining problem is to construct or exclude such a class with
this larger genus. The surface must also have a canonical map
with two-dimensional image and the nef and big residual adjoint
$K_X-(m-2)Y-Z$ of Proposition~\ref{prop:residual-adjoint}.
The vector space of sections defining its derivative subsystem has
dimension $h^0(X,\OO_X(Y))-1$.
These constraints do not settle the remaining cases.
They do not assert that $X$ is minimal: an exceptional curve can be
part of the additional zero divisor $Z$.

At every exact conductor equality $2g(Y)-2=paY^2$, the remaining
configuration has $a\geq5$ and $a\ne6$, $K_X\sim(pa-1)Y$ and an ample
canonical bundle on the distinguished normalization. Its minimal
resolution is of general type, so no ruled normalization remains.
There is a constant tangent point in the very ample embedding;
the embedding is not a complete intersection.
The degree satisfies $p\mid(Y^2-\epsilon)$, with $\epsilon=1$
allowed only for the single $A_1$ point in characteristic two.
In odd characteristic the three genus rows in
\eqref{eq:odd-conductor-three-genera} and the conductor thresholds
of Corollary~\ref{cor:odd-conductor-numerical-threshold} hold.
Corollary~\ref{thm:positive-square-picard} excludes every
algebraically trivial equality class, including characteristic two.
More precisely, if $r=v_p(pa)$, a remaining equality class must be
nonzero in $\operatorname{NS}(X)$ and cannot belong to
$p^r\operatorname{NS}(X)$. In the characteristic-two case with
one $A_1$ point, Proposition~\ref{prop:internal-cartier-descent-obstruction}
strengthens this to $[M]\notin2\operatorname{NS}(X)$.
It excludes invertible descent to the singular radial quotient;
a linearization with different stabilizer characters need not
provide such a descent. These smooth or one-node general-type
configurations with such N\'eron--Severi torsion remain to be
constructed or excluded. For positive genus excess
$\delta=2g(Y)-2-mY^2$, Theorem~\ref{thm:positive-excess-picard}
reduces the algebraically trivial case, and more generally the case
$[M]\in p^{v_p(m)}\operatorname{NS}(X)$, to
\[
 \delta\geq\frac{mY^2}{Q_{v_p(m)}(m)-1}.
\]
When $v_p(m)=1$ and $p$ is odd, the stronger strict bound is
$\delta>pd/(2(p-1))$.
At the characteristic-two boundary with $m=2a$, $a$ odd and
$\delta=d$, Corollary~\ref{cor:positive-excess-characteristic-two-boundary}
requires nontrivial actual two-torsion, and at $a=1$ requires
actual two-torsion rank at least two. The discrepancy line
$\mathcal T$ has exact order two if $a=1$ and is trivial if $a\geq3$.
At $a=1$ its logarithmic divisor is $3Z$, whereas the logarithmic
form of every restricted square root is nonzero on $Z$.
For all odd $a$ a suitable resolution of the normalized cover admits
a fibration over $\PP^1$ with a distinguished fibre $2^sD$.
These assertions neither construct an endpoint surface nor give
a fibration of $X$.
For a class pulled back from a curve of genus $b$, in any
characteristic $p>0$ the stronger bound is
$\delta\geq pd/(2(p-1))+n(b-1)/((p-1)a)$.
In characteristic two this excludes $b\geq2$ at the endpoint.
For a smooth surface and an algebraically trivial class at $a=1$
with the curve ample, the remaining Albanese possibilities are an
elliptic curve image or a surface image with generic differential
rank exactly one. The elliptic curve map from $Y$ has inseparable
degree exactly two and separable degree at least two, with its
ramification divisor disjoint from $Z$; the first normalized base
change has smooth generic fibre. In the surface-image case,
$2\leq\dim\operatorname{Alb}(X)\leq d+2$, and $Y$ maps separably
to the normalized image curve $B_0$. Its Jacobian maps surjectively
but inseparably to $\operatorname{Alb}(X)$, with differential corank one.
If the Albanese variety is a surface, the image curve has the
singularity defect and square computed above. No morphism
$X\to B_0$ is obtained. Both branches must also satisfy the nonsplit
conductor and saturated cotangent conclusions, with $c_2(X)\geq0$,
no global vector fields and no exceptional curves of the first kind.
When $Y$ is very ample, either branch must also satisfy
$d\ge12$ and
$h^0(X,\OO_X(Y))\le\lfloor d/4\rfloor+2$.
If the Albanese variety is a surface and the Albanese morphism
has degree two, the stronger restrictions are
\[
 d=4\chi(\OO_X),\qquad
 2d<K_X^2<3d,\qquad
 1\le c_2(X)=3d-K_X^2\le d-1.
\]
The map from $Y$ to the normalization of its Albanese image is
then an isomorphism.
Let $t$ be the resolution defect of the finite Stein factor and
let $\epsilon$ be the normalization defect of the lifted image
in the degree-two isogeny defined by $M$.
These defects must satisfy
\[
 2t=d+\epsilon,\qquad \epsilon\ge0\ \text{even},\qquad
 t\ge\max\left\{\frac d2,
       \left\lceil\frac{d^2}{4c_2(X)}\right\rceil\right\}.
\]
The canonical discrepancy on $X$ is the negative of a nonzero
effective exceptional divisor.
The image divisor has numerical class $2\mathcal L$, where
$\mathcal L$ is the canonical line of the finite double cover,
and its pullback has the exceptional correction determined in
\eqref{eq:strict-canonical-image-pullback}.
A smooth lifted image forces $t=d/2$ and $c_2(X)\ge d/2$.
The defects are otherwise not bounded above by these restrictions.
With $Y$ very ample as well, its complete embedding has dimension
at least five and $d$ is a multiple of four at least sixteen.
At $d=16$ one must have
\[
 \begin{gathered}
 h^0(X,\OO_X(Y))=6,\qquad h^0(Y,\OO_Y(Y))=5,\\
 t\ge\max\{8,\lceil64/c_2(X)\rceil\},\qquad
 \epsilon=2t-16.
 \end{gathered}
\]
These are the restrictions of
Proposition~\ref{prop:smallest-endpoint-very-ample-degree},
Proposition~\ref{prop:smallest-endpoint-albanese-degree-two},
Corollary~\ref{cor:smallest-endpoint-albanese-very-ample-degree}, and
Theorem~\ref{thm:smallest-endpoint-strict-canonical}.
Neither configuration is constructed or excluded in full.
The remaining questions also concern larger excess and the
classes without the specified N\'eron--Severi divisibility. The theorem
does not assert that either range contains an example.

For a merely ample curve in the smallest conductor equality $2g(Y)-2=pY^2$, Corollary~\ref{thm:positive-square-picard} forces a nonzero N\'eron--Severi class that is not a $p$th multiple. Subject to that restriction, the characteristic-two K3 alternative and the allowed ruled configurations remain open. In characteristic two, if the distinguished cover is rational, Proposition~\ref{prop:rational-normalization-canonical} gives $\chi(\OO_X)=1$, $K_X\equiv Y$ and $1\leq Y^2\leq9$, but the embedding obstruction needs very ampleness. If a degree-two ruled cover has positive-genus base, Remark~\ref{rem:characteristic-two-ample-residual} forces genus one, fibre degree two, the stated zero Chern invariants and a separable degree-two multisection. In characteristic three, the merely ample rational-cover case with a canonical section has the degree-two numerical package described in Corollary~\ref{cor:characteristic-three-rational-residual}. Higher conductors $a>1$, strict adjunction inequality and the general existence or exclusion of global torsion classes also remain unresolved.

\section*{Method and availability}

The mathematics and text were produced by GPT-6 Astra (ChatGPT Pro mode).
Separate same-model reconstructions from exact statements covered the
polarization criterion; the abelian-torsor, modification and projective-bundle
theorems; the very ample surface and its numerical invariants; the uniform
coefficient and metabelian theorems; and the main conductor, genus and adjoint
arguments. The normal-surface equality corollary with its finite
N\'eron--Severi divisibility obstruction also received a separate
statement-only reconstruction. The full quantitative theorem, including
the strict odd-characteristic first-valuation bound, received a complete
statement-only reconstruction without project-proof input. This proof
uses a finite flat cover for the genus calculation and explicit cusp
charts for strictness. A further same-model comparison checked the whole
reconstruction and its primary inputs.
The characteristic-two endpoint received a separate reconstruction;
the sharper endpoint alternatives and odd-characteristic strictness
were then reconstructed from their statements using that checking
model's own earlier argument as declared prior input. The latter
proof uses explicit blowup charts and a pole-order estimate.
The characteristic-two curve-pullback bound also received a fresh
statement-only reconstruction using successive blowups of the base.
The logarithmic-order-two rank and divisor statements, the resolved-cover
fibration in that case, and the elliptic conclusions received a further
statement-only reconstruction. Its degree-zero section establishes the
smooth local cover before the cusp resolution. The same attempt did not
complete the stronger fibration assertion for larger odd
half-orders; that failed extension is preserved separately.
A separate exact-statement reconstruction covers the full odd-half-order
fibration assertion, using a finite flat additive-polynomial cover
and its determinant norm to produce the pencil.
The common global conductor and cotangent theorem received a complete
separate reconstruction from its statement. A further reconstruction
of the five Albanese conclusions used that checking model's own
preceding global proof as declared prior input.
The first separate attempt at the numerical restrictions independently
excluded degree ten but did not complete the sharp general curve bound
or the Albanese calculation. A later complete reconstruction used its
own earlier argument and four explicitly stated proved dependencies.
Its curve argument uses the equality case of Castelnuovo's bound,
while the proof in the text uses cubic sections and sectional monodromy.
A separate exact-statement reconstruction also covers the strict
canonical bound and finite Stein-factor conclusions, conditional on
two explicitly stated packages of earlier endpoint results.
The numerical Stein and image identities received a separate
reconstruction conditional on four explicitly stated proved inputs.
A further reconstruction of the exact pullback formula and the
exceptional-cycle inequality, including its equality condition,
used that model's own preceding proof and the proved
negative-definiteness statement.
Neither reconstruction was supplied with the proof given here.
The arbitrary-base projective-bundle reconstruction
used its checking model's earlier normal-base argument. The conductor
reconstructions used only their declared earlier arguments as prior inputs.
Two unsuccessful all-base reconstructions remain documented in the records.
Three attempts to reconstruct the full internal Cartier assertion did not
close its even-conductor case; those failed attempts remain distinct
from the complete comparisons of its proof.

Complete targeted comparisons checked the proofs and the primary sources at
their applied hypotheses and locators. These comparisons include the exact
residual adjoint, normalization and singularity calculations, the conditional
ample-curve statements, formal depth, radial descent, Cartier descent
at the internal center, the three-genus
restriction, the exclusion of conductor six, and the characteristic-two
positive-excess boundary criterion, its sharper alternatives,
odd-characteristic strictness, and the curve-pullback and Albanese
restrictions. Further complete comparisons checked the rank-two and logarithmic-divisor refinement, the moving-divisor argument and rational base on the resolved cover, the elliptic height and smoothness conclusions, and the global conductor, cotangent and inseparable Albanese restrictions. The strict canonical boundary and the finite Stein-factor singularity defect also received complete proof and source comparisons. The Stein and image identities, exact exceptional pullback, and exceptional-cycle inequality and equality condition were also compared in full. The comparison of the exceptional-cycle calculation began with its statement and outline already known. The proof for normal surfaces and curves of positive self-intersection,
including the quantitative excess bound and its finite divisibility
condition, is written out in full. The records distinguish reconstruction from a statement alone from
checking an already known proof, and identify all prior mathematical inputs
and incidental exposure. The additional local-section identities were checked
by exact polynomial reduction and a separate binary-polynomial implementation.

The cohomology dimensions were recomputed by principal-parts calculations
and Riemann--Roch. Different methods checked the exact intersection and
elementary differential algebra. The programs \emph{Bundle curve intersection
verifier}, \emph{Ruled-cover field verifier}, \emph{Fibration numerical
verifier}, \emph{Equality Euler verifier}, \emph{Degree filtration operator
verifier}, \emph{Rational canonical verifier}, \emph{Projection and blowup
verifier}, \emph{Strict equality and source arithmetic verifier},
\emph{General adjoint arithmetic verifier}, \emph{Odd conductor arithmetic
verifier}, \emph{Conductor structure arithmetic verifier},
\emph{Joint conductor polynomial verifier}, \emph{Conductor descent verifier},
\emph{Internal Cartier jet verifier}, \emph{Manuscript arithmetic verifier},
\emph{Positive excess arithmetic verifier},
\emph{Finite-cover filtration verifier},
\emph{Boundary cusp and overlap verifier},
\emph{Endpoint cusp differential verifier},
\emph{Endpoint semigroup and degree verifier},
\emph{Smallest endpoint symbolic verifier}, and
\emph{Smallest endpoint bitset and semigroup verifier}
are available through the
workbench, with their inputs, original outputs and separate reruns.
These calculations support the algebra; the geometric proofs are in the text.

Earlier targeted checks were recorded by Claude Opus 5.5 at maximum effort;
they are historical checks, with their exposure boundaries preserved.
No human has checked the paper. The text and original programs are dedicated
to the public domain under CC0~1.0. Editable sources and detailed proof,
source and computation records are at \url{\workbenchurl}.
The \mbox{permanent} record address will be assigned at
\url{\recordurl}.

\begin{thebibliography}{99}
\small\setlength{\itemsep}{1pt}
\bibitem{Groth59} A.~Grothendieck, Montpellier archives, cote 59, sheet 13, Universit\'e de Montpellier. \url{https://grothendieck.umontpellier.fr/59.pdf}.
\bibitem{PR} B.~Poonen and E.~Rains, Self cup products and the theta characteristic torsor, \emph{Math. Res. Lett.} \textbf{18} (2011), 1305--1318. \S\S3.1--3.2 and 3.6, Proposition 3.6, diagram (3.7), and Proposition 3.12(f), pp.~1311--1315. \url{https://doi.org/10.4310/MRL.2011.v18.n6.a18}.
\bibitem{FP} R.~Fringuelli and R.~Pirisi, The Picard group of the universal abelian variety and the Franchetta conjecture for abelian varieties, \emph{Michigan Math. J.} \textbf{68} (2019), 651--671. Theorem C, author version p.~3. \url{https://arxiv.org/abs/1603.09190}; \url{https://doi.org/10.1307/mmj/1564106669}.
\bibitem{Raynaud} M.~Raynaud, Contre-exemple au ``vanishing theorem'' en caract\'eristique $p>0$, in \emph{C.~P.~Ramanujam---a tribute}, Tata Inst. Fund. Res. Studies in Math. \textbf{8}, Springer, 1978, 273--278. Official retypesetting, pp.~319--325: \url{https://mathweb.tifr.res.in/Documents/Publications/Studies/CP-Ramanujam-A-Tribute.pdf}.
\bibitem{HexagonSurface}
\emph{A very ample counterexample to Kodaira vanishing on a surface,
answering a question of Raynaud},
Hexagon:2611.00936v1, 23~November~2026.
Theorem~1 and Sections~3--6, pp.~2, 4--8.
\url{https://hexagonmath.org/2611.00936v1}.
\bibitem{LR} N.~Lauritzen and A.~P.~Rao, Elementary counterexamples to Kodaira vanishing in prime characteristic, \emph{Proc. Indian Acad. Sci. Math. Sci.} \textbf{107} (1997), 21--25. Author manuscript: \url{https://arxiv.org/abs/alg-geom/9604012}.
\bibitem{Zheng} X.~Zheng, Counterexamples of Kodaira vanishing for smooth surfaces of general type in positive characteristic, \emph{J. Pure Appl. Algebra} \textbf{221} (2017), 2431--2444. \url{https://arxiv.org/abs/1509.04993}.
\bibitem{Bhatt} B.~Bhatt, Derived splinters in positive characteristic, \emph{Compositio Math.} \textbf{148} (2012), 1757--1786. Proof of Proposition 7.2, p.~1771. \url{https://doi.org/10.1112/S0010437X12000309}.
\bibitem{BhattAnnihilation} B.~Bhatt, Annihilating the cohomology of group schemes, \emph{Algebra Number Theory} \textbf{6} (2012), 1561--1577. Example~6.1, Proposition~6.2 and Corollary~6.3, pp.~1574--1575. \url{https://doi.org/10.2140/ant.2012.6.1561}.
\bibitem{AjitLee} ChatGPT and Claude, in collaboration with R.~Ajit and S.~Lee, \emph{Counterexamples in positive characteristics}, preliminary draft, 5 November 2026, Theorem~A, p.~2. \url{https://hexagonmath.org/2611.00204v1}.
\bibitem{LangerS}
A.~Langer, On the $S$-fundamental group scheme,
\emph{Ann. Inst. Fourier} \textbf{61} (2011), 2077--2119.
Theorems~9.1 and~10.2 and their proofs, pp.~2105--2106 and~2109--2111.
\url{https://doi.org/10.5802/aif.2667}.
\bibitem{KweonNS}
H.~J.~Kweon, Bounds on the torsion subgroup schemes of
N\'eron--Severi group schemes,
\emph{Adv. Math.} \textbf{409}, part~B (2022), article~108687.
\url{https://doi.org/10.1016/j.aim.2022.108687}.
Locators refer to arXiv:2008.01908v2, 27~September~2021:
\S7, Theorem~7.2, pp.~17--18, and Question~7.7, p.~19.
\url{https://arxiv.org/abs/2008.01908v2}.
\bibitem{CotnerZavyalov}
S.~Cotner and B.~Zavyalov, Lefschetz theorems in flat cohomology
and applications, \emph{Compositio Math.} \textbf{160} (2024),
1850--1870. Definitions~2.2.1 and~2.3.5, Corollary~3.1.2, and
Theorem~3.4.4;
author version, pp.~4, 8--9 and~17.
\url{https://doi.org/10.1112/S0010437X24007231};
\url{https://bogdanzavyalov.com/refs/papers/Lefschetz.pdf}.
\bibitem{BiswasHolla}
I.~Biswas and Y.~I.~Holla, Comparison of fundamental group schemes
of a projective variety and an ample hypersurface,
\emph{J. Algebraic Geom.} \textbf{16} (2007), 547--597.
Theorem~1.1 and Section~2, author pp.~2--7.
\url{https://doi.org/10.1090/S1056-3911-07-00449-3};
\url{https://arxiv.org/abs/math/0603299v2}.
\bibitem{ZhangQuasiElliptic}
Y.~Zhang, The $d$-very ampleness of adjoint line bundles on
quasi-elliptic surfaces, \emph{J. Pure Appl. Algebra} \textbf{227}
(2023), article~107271. Theorem~3.1 (attributed to Langer),
author version arXiv:2103.04268v7, pp.~4--5.
\url{https://doi.org/10.1016/j.jpaa.2022.107271};
\url{https://arxiv.org/abs/2103.04268v7}.
\bibitem{Liedtke} C.~Liedtke, Algebraic surfaces in positive characteristic, in \emph{Birational Geometry, Rational Curves, and Arithmetic}, Springer, 2013, 229--292. Author version, Theorems~4.4--4.5, p.~21, Theorem~5.6, pp.~23--24, and \S7, pp.~26--30. \url{https://arxiv.org/abs/0912.4291}; \url{https://doi.org/10.1007/978-1-4614-6482-2_11}.
\bibitem{FMST} T.~Fassarella, W.~Mendson, J.~P.~dos Santos and F.~Touzet, \emph{Foliations with small singular set in arbitrary characteristic}, preprint, September 2026. Theorems~1.1--1.2, pp.~2--3, and Theorem~2.8, p.~13. \url{https://arxiv.org/abs/2607.18571}.
\bibitem{Kato} K.~Kato, Swan conductors for characters of degree one in the imperfect residue field case, in \emph{Algebraic K-theory and algebraic number theory}, Contemp. Math. \textbf{83}, AMS, 1989, 101--131. \url{https://doi.org/10.1090/conm/083/991978}.
\bibitem{BN} M.~Bright and R.~Newton, Evaluating the wild Brauer group, \emph{Invent. Math.} \textbf{234} (2023), 819--891. Section 4.3 and Theorem 4.1. \url{https://doi.org/10.1007/s00222-023-01210-8}.
\bibitem{Suzuki} T.~Suzuki, Finite generation of nilpotent quotients of fundamental groups of punctured spectra, \emph{J. Algebra} \textbf{633} (2023), 489--509. \url{https://doi.org/10.1016/j.jalgebra.2023.07.005}. Author version: \url{https://arxiv.org/abs/2207.01563}.
\bibitem{Sun} X.~Sun, Direct images of bundles under Frobenius morphism, \emph{Invent. Math.} \textbf{173} (2008), 427--447. \url{https://doi.org/10.1007/s00222-008-0125-y}. Author version: \url{https://arxiv.org/abs/math/0611360}.
\bibitem{Lipman} J.~Lipman, Rational singularities with applications to algebraic surfaces and unique factorization, \emph{Publ. Math. IHES} \textbf{36} (1969), 195--279. Lemma 14.1, p.~224. \url{https://www.numdam.org/item/PMIHES_1969__36__195_0.pdf}.
\bibitem{Cadman} C.~Cadman, Using stacks to impose tangency conditions on curves, \emph{Amer. J. Math.} \textbf{129} (2007), 405--427. \url{https://doi.org/10.1353/ajm.2007.0007}. Author version: \url{https://www.charlescadman.com/raw/pdf/stacks.pdf}.
\bibitem{Borne} N.~Borne, Fibr\'es paraboliques et champ des racines, \emph{Int. Math. Res. Not.} (2007), rnm049. \url{https://doi.org/10.1093/imrn/rnm049}. Author version: \url{https://arxiv.org/abs/math/0604458}.
\bibitem{KLS} K.~Kato, I.~Leal and T.~Saito, Refined Swan conductors mod $p$ of one-dimensional Galois representations, \emph{Nagoya Math. J.} \textbf{236} (2019), 134--182. \url{https://doi.org/10.1017/nmj.2019.13}.
\bibitem{SGA1} A.~Grothendieck and M.~Raynaud, \emph{Rev\^etements \'etales et groupe fondamental (SGA~1)}, Lecture Notes in Math. \textbf{224}, Springer, 1971. Expos\'e XIII, Corollaire 2.12. Recomposed edition: \url{https://arxiv.org/abs/math/0206203}.
\bibitem{Stacks} The Stacks Project Authors, \emph{The Stacks Project}. Exact tags are cited in the text. \url{https://stacks.math.columbia.edu}.
\bibitem{BLR}
S.~Bosch, W.~L\"utkebohmert and M.~Raynaud,
\emph{N\'eron Models}, Ergebnisse der Mathematik (3), vol.~21,
Springer, 1990. \S6.6, Theorem~2(c),(d), pp.~164--166.
\url{https://doi.org/10.1007/978-3-642-51438-8}.
\bibitem{RaynaudDescent}
M.~Raynaud, Un crit\`ere d'effectivit\'e de descente,
\emph{S\'eminaire d'alg\`ebre commutative}, vol.~2 (1967--1968),
expos\'e~5, 22~pp. Corollary~3.8 and its application, p.~5-15;
Theorem~4.2 and its proof, pp.~5-16--5-18.
\url{https://www.numdam.org/item/SAC_1967-1968__2__A5_0/}.
\bibitem{EGA}
A.~Grothendieck and J.~Dieudonn\'e,
\emph{\'El\'ements de g\'eom\'etrie alg\'ebrique}.
II, \emph{Publ. Math. IHES} \textbf{8} (1961),
Proposition~6.6.1 and Corollary~6.6.4;
III$_1$, \emph{Publ. Math. IHES} \textbf{11} (1961),
Theorem~3.2.1, pp.~116--117;
III$_2$, \emph{Publ. Math. IHES} \textbf{17} (1963),
\S\S7.7 and 7.9;
IV$_2$, \emph{Publ. Math. IHES} \textbf{24} (1965),
7.8.3, pp.~214--216;
IV$_3$, \emph{Publ. Math. IHES} \textbf{28} (1966),
\S\S8.8--8.10, Theorems~9.7.7, 11.2.6, 11.3.8, 11.3.10 and~13.1.3,
Proposition~15.5.9;
IV$_4$, \emph{Publ. Math. IHES} \textbf{32} (1967),
Proposition~17.7.8, Theorem~17.11.1, pp.~82--84,
Corollary~21.12.7 and Proposition~21.14.3.
\bibitem{LaurentSchroer}
B.~Laurent and S.~Schr\"oer, Para-abelian varieties and Albanese maps,
\emph{Bull. Braz. Math. Soc. (N.S.)} \textbf{55} (2024), article~4.
Corollary~4.6, author version pp.~14--15.
\url{https://arxiv.org/abs/2101.10829}.
\bibitem{Alexeev}
V.~Alexeev, Complete moduli in the presence of semiabelian group action,
\emph{Ann. of Math.} \textbf{155} (2002), 611--708.
Theorem~3.0.3, p.~664.
\url{https://arxiv.org/abs/math/9905103};
\url{https://doi.org/10.2307/3062130}.
\bibitem{ACMPrym}
J.~D.~Achter and S.~Casalaina-Martin,
\emph{Putting the P back in Prym}, author manuscript dated
18~December~2024. Standing base convention, p.~3;
full-Picard schematicity assertion, \S1.7.2, first paragraph, p.~13.
\url{https://www.math.colostate.edu/~achter/math/prym1224.pdf}.
\bibitem{Groth57} A.~Grothendieck, Montpellier archives, cote~57, sheet~314 (PDF p.~315), Universit\'e de Montpellier. \url{https://grothendieck.umontpellier.fr/57.pdf}.
\bibitem{CDL}
F.~Cossec, I.~Dolgachev and C.~Liedtke, \emph{Enriques Surfaces I},
author version dated 19 April 2024. Degree filtration (0.2.17), pp.~38--40; Theorem~4.1.6, pp.~364--365;
Corollary~4.1.12, p.~369; Proposition~4.1.14, pp.~369--370.
\url{https://sites.lsa.umich.edu/idolga/wp-content/uploads/sites/1334/2024/08/EnriquesOne.pdf}.
\bibitem{SuWang}
X.~Su and B.~Wang, \emph{Picard Groups of Spectral Varieties and Moduli of Higgs Sheaves}, preprint, arXiv:2409.10296v2, 17~November~2026.
Lemmas~3.22--3.23 and Remark~3.24, pp.~19--21.
\url{https://arxiv.org/abs/2409.10296v2}.
\bibitem{Workbench} Mathematical workbench for schemes and moduli, original programs and checking records. \url{\workbenchurl}.
\bibitem{MRT} A.~M\'ezard, M.~Romagny and D.~Tossici,
Sekiguchi--Suwa theory revisited, \emph{J. Th\'eor. Nombres Bordeaux}
\textbf{26} (2014), 163--200. \S2.2, p.~167, and Lemma~6.1, p.~188.
\url{https://doi.org/10.5802/jtnb.863}.

\bibitem{MilneAV} J.~S.~Milne, Abelian varieties, in \emph{Arithmetic Geometry}, Springer, 1986, 103--150. Author revision, Propositions~9.2 and~10.1 and Remark~9.3, pp.~14--17. \url{https://www.jmilne.org/math/xnotes/AVs.pdf}.
\bibitem{Zomervrucht} W.~Zomervrucht, \emph{Ineffective descent of genus one curves}, arXiv:1501.04304, 2015; author PDF dated 22 April 2018. Theorem~2.1 and \S3, pp.~2--3. \url{https://arxiv.org/abs/1501.04304}.
\bibitem{BrionProjective} M.~Brion, Homogeneous projective bundles over abelian varieties, \emph{Algebra Number Theory} \textbf{7} (2013), 2475--2510. Introduction and generalities, pp.~2475--2479. \url{https://doi.org/10.2140/ant.2013.7.2475}.
\bibitem{EkedahlCanonical} T.~Ekedahl, Canonical models of surfaces of general type in positive characteristic, \emph{Publ. Math. IHES} \textbf{67} (1988), 97--144. Lemma~0.1.5, p.~103; I, Theorem~2.4 and its proof, pp.~109--115; II, Theorem~1.6, p.~125, and Theorem~1.7, pp.~126--127. \url{https://doi.org/10.1007/BF02699128}.
\bibitem{MatsumotoActions} Y.~Matsumoto, $\mu_p$- and $\alpha_p$-actions on K3 surfaces in characteristic $p$, \emph{J. Algebraic Geom.} \textbf{32} (2023), 271--322. Author version, Proposition~2.10, p.~5; Theorem~3.3 and Table~2, pp.~10--11; Proposition~4.1, pp.~16--17. \url{https://doi.org/10.1090/jag/804}; \url{https://arxiv.org/abs/1812.03466v5}.
\bibitem{MatsumotoMuN} Y.~Matsumoto, $\mu_n$-actions on K3 surfaces in positive characteristic, author version, 23 March 2020. Proposition~4.8 and Table~3, pp.~13--15. \url{https://arxiv.org/abs/1710.07158v3}.
\bibitem{Kadets} B.~Kadets, Sectional monodromy groups of projective curves, arXiv:1809.07293v4, 14 November 2020. Proposition~2.8 and proof, pp.~4--5. \url{https://arxiv.org/abs/1809.07293v4}; \url{https://doi.org/10.1112/jlms.12375}.
\bibitem{MumfordEnriquesI}
D.~Mumford, Enriques' classification of surfaces in char $p$: I,
in \emph{Global Analysis (Papers in Honor of K.~Kodaira)},
D.~C.~Spencer and S.~Iyanaga, eds.,
University of Tokyo Press and Princeton University Press, 1969,
325--339. \S1, Steps~I--VII, pp.~326--330.
\url{https://www.dam.brown.edu/people/mumford/alg_geom/papers/1969b--EnrClass-I-NC.pdf}.

\bibitem{FGAPicard}
A.~Grothendieck, Technique de descente et th\'eor\`emes d'existence
en g\'eom\'etrie alg\'ebrique,
V.~Les sch\'emas de Picard: th\'eor\`emes d'existence,
and VI.~Les sch\'emas de Picard: propri\'et\'es g\'en\'erales,
\emph{S\'eminaire Bourbaki}, 14e ann\'ee, 1961/62,
expos\'es~232 and~236.
Expos\'e~232, Corollary~2.4, p.~147, Theorem~3.1, pp.~148--149,
and its proof,
Sections~4--5, pp.~150--155;
expos\'e~236, Theorem~2.1(ii), pp.~231--232,
Proposition~2.10(ii), pp.~235--236,
Proposition~3.1 and Corollary~3.2, p.~236, and
Theorem~3.3(iii), pp.~237--238.

\bibitem{KleimanPicard}
S.~L.~Kleiman, \emph{The Picard scheme}, author version,
1 April 2005, Theorem~2.5, pp.~18--20; Theorem~5.4, pp.~37--38;
Remark~5.8, pp.~39--41; Proposition~5.10, pp.~41--42;
and Remark~5.25, p.~51.
\url{https://arxiv.org/abs/math/0504020}.
Published in \emph{Fundamental Algebraic Geometry},
Mathematical Surveys and Monographs \textbf{123}, AMS, 2005,
235--321; corresponding Remark~9.5.8.
\bibitem{CataneseOguiso2020}
F.~Catanese and K.~Oguiso, The double point formula with isolated
singularities and canonical embeddings,
\emph{J. London Math. Soc.} \textbf{102} (2020), 1337--1356.
Section~5; author proof, pp.~12--18.
\url{https://doi.org/10.1112/jlms.12371};
\url{https://www.mathe8.uni-bayreuth.de/de/team/prof-fabrizio-catanese-old/pdf/168.pdf}.
\bibitem{BombieriMumfordIII}
E.~Bombieri and D.~Mumford, Enriques' classification of surfaces
in char.~$p$, III, \emph{Invent. Math.} \textbf{35} (1976),
197--232. Classification summary, pp.~197--198.
\url{https://doi.org/10.1007/BF01390138}.


\bibitem{PoonenSlavov}
B.~Poonen and K.~Slavov, The exceptional locus in the Bertini
irreducibility theorem for a morphism, \emph{Int. Math. Res. Not.}
\textbf{2022} (2022), 4503--4513. Theorem~1.8 and \S7, author
version pp.~2, 6.
\url{https://doi.org/10.1093/imrn/rnaa182};
\url{https://math.mit.edu/~poonen/papers/bertini_locus.pdf}.
\bibitem{GuSunZhou}
Y.~Gu, X.~Sun and M.~Zhou, Slope inequalities and a Miyaoka--Yau type
inequality, \emph{J. Eur. Math. Soc.} \textbf{25} (2023), 611--632.
Theorem~1.5, p.~614; Theorem~3.2, Remark~3.3 and Corollary~3.4,
pp.~622--624. \url{https://doi.org/10.4171/JEMS/1183}.
\bibitem{ShepherdBarronGeography}
N.~I.~Shepherd-Barron, Geography for surfaces of general type in
positive characteristic, \emph{Invent. Math.} \textbf{106} (1991),
263--274. Theorems~2 and~7 and their proofs, pp.~264--266.
\url{https://gdz.sub.uni-goettingen.de/id/PPN356556735_0106}.
\bibitem{SchroerGenusChange}
S.~Schr\"oer, On genus change in algebraic curves over imperfect
fields, \emph{Proc. Amer. Math. Soc.} \textbf{137} (2009), 1239--1243.
Theorems~2.1--2.2 and Corollary~2.3, author version pp.~3--4.
\url{https://doi.org/10.1090/S0002-9939-08-09712-8};
\url{https://www.math.uni-duesseldorf.de/~schroeer/publications_pdf/genus-change.pdf}.

\bibitem{MatsumotoMuN2023} Y.~Matsumoto,
On $\mu_n$-actions on K3 surfaces in positive characteristic,
\emph{Nagoya Math. J.} \textbf{249} (2023), 11--49.
Author version v4, \S\S2.1--2.3, Table~1, pp.~3--8.
\url{https://doi.org/10.1017/nmj.2022.20};
\url{https://arxiv.org/abs/1710.07158v4}.
\bibitem{SeshadriCartier} C.~S.~Seshadri,
L'op\'eration de Cartier. Applications,
\emph{S\'eminaire Claude Chevalley} \textbf{4} (1958--1959),
expos\'e~6, 1--26. Proposition~6, Lemma~3 and Theorem~2,
pp.~6-11--6-16.
\url{https://numdam.org/item/SCC_1958-1959__4__A6_0/}.
\bibitem{FurukawaStrange} K.~Furukawa,
Cohomological characterization of hyperquadrics of odd dimensions
in characteristic two, \emph{Math. Z.} \textbf{278} (2014), 119--130.
Theorem~1.1 and Propositions~2.5, 2.10; author version pp.~2--6,
with the proof of Theorem~1.1 in \S3.
\url{https://doi.org/10.1007/s00209-014-1308-4};
\url{https://arxiv.org/abs/1304.1634v3}.


\bibitem{LangPTorsion}
W.~E.~Lang, Remarks on $p$-torsion of algebraic surfaces,
\emph{Compositio Math.} \textbf{52} (1984), 197--202.
Proof of Lemma~3, pp.~198--199.
\url{https://www.numdam.org/item/CM_1984__52_2_197_0.pdf}.
\bibitem{FujinoSatoToric}
O.~Fujino and H.~Sato, \emph{Characteristic-free characterizations
of projective space for smooth toric varieties}, preprint, 2026.
Author version dated 10~September~2026, version~0.13:
Theorem~1.1, p.~1; proof in \S5, pp.~8--10; Remark~6.3, p.~12.
\url{https://arxiv.org/abs/2609.03282v2};
\url{https://www.math.kyoto-u.ac.jp/~fujino/characterization20260911.pdf}.


\bibitem{NagamachiTakamatsu}
I.~Nagamachi and T.~Takamatsu, \emph{On behavior of conductors,
Picard schemes, and Jacobian numbers of varieties over imperfect fields},
arXiv:2110.03917v3, 3~March~2022.
Introduction and Theorem~0.2, pp.~1--3.
\url{https://arxiv.org/abs/2110.03917v3}.


\bibitem{CastryckStrengTesta}
W.~Castryck, M.~Streng and D.~Testa,
Curves in characteristic $2$ with non-trivial $2$-torsion,
\emph{Adv. Math. Commun.} \textbf{8} (2014), 479--495.
Introduction and Theorem~2, author version pp.~1--3,
8 December 2020.
\url{https://arxiv.org/abs/1402.3241v4};
\url{https://doi.org/10.3934/amc.2014.8.479}.
\bibitem{EsnaultPartialConnection}
H.~Esnault, Partial connection for $p$-torsion line bundles in
characteristic $p>0$, in \emph{Inspired by S.~S.~Chern},
Nankai Tracts Math.~\textbf{11}, World Scientific, 2006, pp.~173--178.
Formula~(2.4), author p.~3.
\url{https://arxiv.org/abs/math/0608203}.
\bibitem{ContiPositiveSeveriComparison}
F.~C.~G.~Conti,
\emph{Surfaces close to the Severi lines in positive characteristic},
arXiv:2111.08622v1, 16~November~2021.
Section~2.2, equations~(2.9)--(2.15), pp.~8--9;
Lemma~2.22 and Remark~2.23, p.~13;
Theorem~8.4, p.~51.
\url{https://arxiv.org/abs/2111.08622v1}.

\bibitem{GounelasLiedtkeSymmetric}
F.~Gounelas and C.~Liedtke, \emph{Symmetric differentials on K3
surfaces}, arXiv:2608.17953v1, 18~August~2026.
Theorem~1.1, p.~2; used here only for the comparison with
regular tensors on the smooth resolution.
\url{https://arxiv.org/abs/2608.17953v1}.
\bibitem{TotaroMovingDivisors}
B.~Totaro, \emph{Moving codimension-one subvarieties over finite fields},
arXiv:0712.2049v1, 12~December~2007.
Lemma~4.1 and proof, author pp.~7--8.
\url{https://arxiv.org/abs/0712.2049v1};
\url{https://www.math.ucla.edu/~totaro/papers/public_html/moving.pdf}.
\bibitem{MehtaSrinivasAlbanese}
V.~B.~Mehta and V.~Srinivas,
Varieties in positive characteristic with trivial tangent bundle,
\emph{Compositio Math.} \textbf{64} (1987), 191--212.
Lemma~1.3 and proof, p.~195.
\url{https://www.numdam.org/item/CM_1987__64_2_191_0/}.
\end{thebibliography}
\end{document}
